\documentclass[12pt,b5paper,twoside]{article}

\usepackage{scrextend}
\usepackage[T1]{fontenc}
\usepackage{newpxtext}
\usepackage{amsmath,amssymb,mathtools}
\usepackage{newpxmath}
\usepackage{bm}
\usepackage[b5paper,left=0.7cm,right=0.7cm,top=0.91cm,bottom=0.91cm,
  includehead,includefoot,headheight=15pt,headsep=5pt,footskip=18pt]{geometry}
\usepackage{microtype}
\usepackage{enumitem}
\usepackage{needspace}
\usepackage{fancyhdr}
\usepackage{xcolor}
\usepackage{hyperref}

\definecolor{LinkBlue}{RGB}{0,75,135}

% Standard article has no 13pt option; set the true base size explicitly.
\normalfont
\changefontsizes{13pt}
\linespread{1.1}

\hypersetup{
  unicode=true,
  colorlinks=true,
  linkcolor=LinkBlue,
  citecolor=LinkBlue,
  urlcolor=LinkBlue,
  filecolor=LinkBlue,
  bookmarksnumbered=true,
  bookmarksopen=true,
  pdftitle={On Rings of Operators},
  pdfauthor={F. J. Murray and J. v. Neumann},
  pdfsubject={Annals of Mathematics 37 (1936), 116--229}
}

\setlength{\parindent}{1.25em}
\setlength{\parskip}{0pt}
\setlength{\emergencystretch}{3em}
\setlist{nosep,leftmargin=2.2em}
\allowdisplaybreaks[2]

\pagestyle{fancy}
\fancyhf{}
\fancyhead[LE]{\small\thepage\qquad F. J. MURRAY AND J. v. NEUMANN}
\fancyhead[RO]{\small ON RINGS OF OPERATORS\qquad\thepage}
\renewcommand{\headrulewidth}{0pt}
\fancypagestyle{firstpage}{%
  \fancyhf{}%
  \fancyfoot[C]{\raisebox{1.3mm}{\thepage}}%
  \renewcommand{\headrulewidth}{0pt}%
}

\makeatletter
\renewcommand{\@cite}[2]{(#1)\if@tempswa, #2\fi}
\renewcommand{\@biblabel}[1]{(#1)}
\newcommand{\solidotimes}{\mathbin{\mathpalette\solidotimes@\relax}}
\newcommand{\solidotimes@}[2]{%
  \begingroup
  \setbox\z@=\hbox{$\m@th#1\bm{\otimes}$}%
  \@tempdima=.04\wd\z@
  \@tempdimb=.04\ht\z@
  \rlap{\kern-\@tempdima\copy\z@}%
  \rlap{\kern\@tempdima\copy\z@}%
  \rlap{\raise\@tempdimb\copy\z@}%
  \rlap{\lower\@tempdimb\copy\z@}%
  \box\z@
  \endgroup}
\makeatother

\newcommand{\Hspace}{\mathfrak H}
\newcommand{\M}{\mathbf M}
\newcommand{\N}{\mathbf N}
\newcommand{\B}{\mathbf B}
\renewcommand{\P}{\mathbf P}
\newcommand{\Q}{\mathbf Q}
\newcommand{\R}{\mathrm R}
\newcommand{\clos}[1]{[\,#1\,]}
\newcommand{\setof}[2]{\{\,#1\,;\,#2\,\}}
\newcommand{\reldiv}[2]{\left[\frac{#1}{#2}\right]}

\newcommand{\PartHeading}[2]{%
  \par\bigskip
  \phantomsection\label{part:#1}%
  \addcontentsline{toc}{part}{Part #1: #2}%
  \begin{center}\large\scshape Part #1: #2\end{center}\nobreak
}
\newcommand{\PartHeadingTOC}[3]{%
  \par\bigskip
  \phantomsection\label{part:#1}%
  \addcontentsline{toc}{part}{Part #1: #3}%
  \begin{center}\large\scshape Part #1: #2\end{center}\nobreak
}
\newcommand{\ChapterHeading}[3]{%
  \par\medskip
  \phantomsection\label{chap:#1}%
  \addcontentsline{toc}{section}{Chapter #2: #3}%
  \begin{center}\scshape Chapter #2: #3\end{center}\nobreak
}
\newcommand{\ChapterHeadingTOC}[4]{%
  \par\medskip
  \phantomsection\label{chap:#1}%
  \addcontentsline{toc}{section}{Chapter #2: #4}%
  \begin{center}\scshape Chapter #2: #3\end{center}\nobreak
}
\newcommand{\SectionStart}[2]{%
  \par\medskip
  \phantomsection\label{sec:#1}%
  \addcontentsline{toc}{subsection}{#1. #2}%
  \noindent\textbf{#1.}\enspace
}
\newcommand{\SectionStartTOC}[3]{%
  \par\medskip
  \phantomsection\label{sec:#1}%
  \addcontentsline{toc}{subsection}{#1. #3}%
  \noindent\textbf{#1.}\enspace #2%
}
\newcommand{\IntroStart}[2]{%
  \par\medskip
  \phantomsection\label{intro:#1}%
  \addcontentsline{toc}{subsection}{#1. #2}%
  \noindent\textbf{#1.}\enspace
}
\newcommand{\Definition}[1]{%
  \par\smallskip\noindent\phantomsection\label{def:#1}\textbf{Definition #1.}\enspace
}
\newcommand{\Lemma}[1]{%
  \par\smallskip\noindent\phantomsection\label{lem:#1}\textbf{Lemma #1.}\enspace
}
\newcommand{\MainTheorem}[1]{%
  \par\smallskip\noindent\phantomsection\label{thm:#1}\textbf{Theorem #1.}\enspace
}
\newcommand{\Problem}[1]{%
  \par\smallskip\noindent\phantomsection\label{prob:#1}\textbf{Problem #1.}\enspace
}
\newcommand{\Proof}{\par\noindent\textit{Proof:}\enspace}
\newcommand{\secref}[1]{\hyperref[sec:#1]{\S#1}}
\newcommand{\introref}[1]{\hyperref[intro:#1]{\S#1 of the Introduction}}
\newcommand{\defref}[1]{\hyperref[def:#1]{Definition #1}}
\newcommand{\lemref}[1]{\hyperref[lem:#1]{Lemma #1}}
\newcommand{\thmref}[1]{\hyperref[thm:#1]{Theorem #1}}
\newcommand{\probref}[1]{\hyperref[prob:#1]{Problem #1}}

\begin{document}
\setcounter{page}{1}
\thispagestyle{firstpage}

\noindent\begin{minipage}{0.42\textwidth}
\small\bfseries Annals of Mathematics\\[-1pt]
Vol. 37, No. 1, January, 1936
\end{minipage}

\vspace{2.8em}
\renewcommand{\thefootnote}{\fnsymbol{footnote}}
\begin{center}
  {\Large\bfseries ON RINGS OF OPERATORS}\par
  \vspace{1.1em}
  {\large By \textsc{F. J. Murray}\footnotemark{} and \textsc{J. v. Neumann}}\par
  \vspace{0.8em}
  {(Received April 3, 1935)}\par
  \vspace{0.9em}
  {\large\itshape Introduction}
\end{center}
\footnotetext{National Research Fellow.}
\setcounter{footnote}{0}
\renewcommand{\thefootnote}{\arabic{footnote}}

\IntroStart{1}{Purpose and origin}
The problems discussed in this paper arose naturally in continuation of the
work begun in a paper of one of us (\cite{ref18}, chiefly parts I and II). Their
solution seems to be essential for the further advance of abstract operator
theory in Hilbert space under several aspects. First, the formal calculus with
operator-rings leads to them. Second, our attempts to generalise the theory of
unitary group-representations essentially beyond their classical frame have
always been blocked by the unsolved questions connected with these problems.
Third, various aspects of the quantum mechanical formalism suggest strongly
the elucidation of this subject. Fourth, the knowledge obtained in these
investigations gives an approach to a class of abstract algebras without a finite
basis, which seems to differ essentially from all types hitherto investigated.

The results which we shall obtain throw light on an entirely new side of
operator theory; they lead to a new notion of linear dimensionality (which,
under certain conditions, has a continuous range of numerical values); and
indicate a way out of the paradoxes of unbounded operator theory.

In the four following \hyperref[intro:2]{\S}\hyperref[intro:5]{\S}\ of this
Introduction we will give a somewhat more
detailed outline of the four aspects of our problems, which were enumerated
above.

\IntroStart{2}{Rings of operators. Types considered}
We consider a linear space $\Hspace$ with an inner product (that is a Hilbert
space or a finite dimensional Euclidean space, cf. (b) in \secref{1.1}), and in
it the ring $\B$ of all bounded operators (cf. (f) in \secref{1.1}). Subsets
$\M$ of $\B$ for which $A,B\in\M$ implies
$\alpha A,A^*,A+B,AB\in\M$ ($\alpha$ is any complex number, $A^*$ is the
adjoint of $A$, cf. the notation in \secref{1.1}), and which are closed in a
suitable topology of $\B$ are called rings. (All these notions are discussed in
detail in \cite{ref18}.) We will chiefly consider rings $\M$ which contain the
unit operator $1$: $1\in\M$. If $\M,\N$ are given rings, denote the smallest
ring containing $\M,\N$ by $R(\M,\N)$, and the greatest ring contained in
$\M$ and $\N$ by $\M\cdot\N$ (this, by the way, happens to be the set
theoretical intersection of $\M$ and $\N$).

The operations $R(\M,\N)$ and $\M\cdot\N$ obey both the commutative and
the associative law, therefore each of them could be analogised with addition
or with multiplication. Owing to
$R(\M,\M)=\M\cdot\M=\M$ the analogy is closer with the ``Boolean
algebras'' (as they occur in set theory or in logics), than with algebra proper.
As the analogue of the distributive law, which connects addition and
multiplication, does not hold in general, this analogy is not perfect. (The
distributive law would be, according to how we identify $R(\M,\N)$ and
$\M\cdot\N$ with addition and multiplication, either
\[
  R(\M\cdot\N,\M\cdot\mathbf P)=\M\cdot R(\N,\mathbf P)
\]
or
\[
  R(\M,\N)\cdot R(\M,\mathbf P)=R(\M,\N\cdot\mathbf P).
\]
None of these equations holds in general.) Thus it is somewhat arbitrary how
we carry out these identifications. We choose to make $R(\M,\N)$ correspond
to addition and $\M\cdot\N$ to multiplication.

In this notation the rings $\M$ (and similarly the rings with $1\in\M$) form
a lattice. (Cf. \cite{ref1}, \cite{ref9}, \cite{ref12a}.) We use the terminology of
\cite{ref9}, and therefore write in this \hyperref[intro:2]{\S}\
$\M\cup\N$ for $R(\M,\N)$
and $\M\cap\N$ for $\M\cdot\N$. One verifies immediately the lattice
postulates I--IV (cf. \cite[p.~422]{ref1}) for these operations.

The lattice $R$ of all rings with $1\in\M$ contains a smallest element: the
set $(\alpha1)$ of all operators of the form $\alpha1$; and a greatest element:
the set $\B$ of all bounded operators in $\Hspace$. We write in this
\hyperref[intro:2]{\S}, $0$
and $1$ for $(\alpha1)$ and $\B$.

Now $R$ has a property which brings it nearer to the Boolean algebras than
lattices usually are: there exists an operation $\M'$ in $R$ which dualises
addition and multiplication; that is, for which
\begin{align}
 (\M\cup\N)'&=\M'\cap\N', \tag{$D_1$}\label{eq:D1}\\
 (\M\cap\N)'&=\M'\cup\N'. \tag{$D_2$}\label{eq:D2}
\end{align}
To this end we define $\M'$ as the set of those $A$ which commute with all
$B\in\M$ (cf. \cite[p.~374]{ref18}); then $\M'\in R$ and
\begin{equation*}
 \M=\M''. \tag{$D_3$}\label{eq:D3}
\end{equation*}
(Cf. \cite[p.~397]{ref18}.) Now \eqref{eq:D1} is obvious; and
\eqref{eq:D2} follows from \eqref{eq:D1} by replacing $\M,\N$ by
$\M',\N'$, applying $'$ to the equation, and using \eqref{eq:D3}.

In the analogy with Boolean algebras (or logics) in which $\M\cup\N$
corresponds to the sum (resp. ``or'') and $\M\cap\N$ to the product (resp.
``and''), $\M'$ should correspond to the complement (resp. ``no'').

Another lattice of such a structure consists of all closed linear manifolds of
$\Hspace$. Denoting it by $M$ we define: If $\mathfrak M,\mathfrak N$ are two
closed linear manifolds in $\Hspace$, then $\mathfrak M\cup\mathfrak N$ is the
smallest closed linear manifold containing them (cf. (d) in \secref{1.1}, where
it is denoted by $[\mathfrak M,\mathfrak N]$); $\mathfrak M\cap\mathfrak N$
is the greatest closed linear manifold contained in them (their set theoretical
intersection, $\mathfrak M\cdot\mathfrak N$); $0$ is the smallest element of
$M$: the set $(0)$ consisting of $0$ alone; $1$ is the greatest element of $M$:
$\Hspace$. (We use these notations in this \hyperref[intro:2]{\S}\ only.)
Then we may define
$\mathfrak M'$ as the set of all elements of $\Hspace$ which are orthogonal to
all elements of $\mathfrak M$: $\Hspace-\mathfrak M$
(cf. \cite[p.~74]{ref16}). One verifies again the lattice postulates I--IV
(cf. \cite[p.~422]{ref1}) for $M$, while the distributive law (postulate VI,
eod., p.~453) does not hold in general. (Postulate V, eod., p.~445, holds,
thus $M$ is a $B$-lattice, cf. eod.) \eqref{eq:D1}--\eqref{eq:D3} too hold.

Following the analogy with the complement in Boolean algebras (and ``no'' in
logics) further, we are led to expect
\begin{align}
 \M\cup\M'&=1, \tag{$D_4$}\label{eq:D4}\\
 \M\cap\M'&=0. \tag{$D_5$}\label{eq:D5}
\end{align}
(In the lattice $M$ \eqref{eq:D4}, \eqref{eq:D5} hold. In fact there is a quite
intimate connection between $M$ and a certain type of logics: The
probability-logics of quantum mechanics. Cf. \cite[pp.~130--134]{ref20}.)
These equations, however, are not true in general in $R$. For instance, if
$\M$ is Abelian (cf. \cite[p.~374]{ref18}), then $\M\subset\M'$, and thus
$\M\cup\M'=\M'$, $\M\cap\M'=\M$. For this reason it is of interest to
find out, for which particular elements $\M$ of $R$ \eqref{eq:D4},
\eqref{eq:D5} hold. As $'$ dualises \eqref{eq:D4} and \eqref{eq:D5}
(apply $'$ and use \eqref{eq:D1}--\eqref{eq:D3}), they are equivalent. Thus it
suffices to discuss one of them. In our original notations they read:
\begin{align}
 R(\M,\M')&=\B, \tag{$\overline D_4$}\label{eq:D4bar}\\
 \M\cdot\M'&=(\alpha1). \tag{$\overline D_5$}\label{eq:D5bar}
\end{align}
The theory of these equations will be the chief subject of this paper
(cf. \secref{3.1}).

\IntroStart{3}{The finite dimensional case}
Let $\Hspace$ be the $n$-dimensional Euclidean space $\mathfrak E_n$,
$n=1,2,\ldots$. In this case it is easy to detail the meaning of
\eqref{eq:D5bar}.

Assume $\M\cdot\M'=(\alpha1)$. Let $\M^u$ be the set of all unitary
elements of $\M$, then $\M$ is the ring generated by $\M^u$
(cf. \cite[p.~392]{ref18}). Thus $\M'=(\M^u)'$; and as $\M^u$ is a group of
unitary matrices, $\M$ consists of all linear aggregates of elements of
$\M^u$. Now apply the theory of unitary group representations to $\M^u$.

As $\M^u$ consists of unitary matrices, it is completely reducible
(cf. \cite[pp.~11--12, footnote 1]{ref14}). Introduce a coordinate system in
$\Hspace=\mathfrak E_n$, in which $\M^u$ appears completely reduced. Let the
irreducible parts of $\M^u$ correspond to the sets of coordinates
\[
 p_1+\cdots+p_{s-1}+1,\ldots,p_1+\cdots+p_{s-1}+p_s
\]
for $s=1,\ldots,r$, where $r=1,2,\ldots$;
$p_1,\ldots,p_r=1,2,\ldots$; $p_1+\cdots+p_r=n$, and denote them by
$I_1,\ldots,I_r$. Denote the number of those which are equivalent to $I_1$ by
$q$ $(=1,2,\ldots,r)$, and arrange the $I_s$ so that these should be
$I_1,\ldots,I_q$. By a further coordinate transformation we can even make
$I_1,\ldots,I_q$ identical. Besides $p_1=\cdots=p_q=p$. Thus every element
of $\M^u$ looks like this: It consists of $r$ $(\geq q)$ matrices succeeding
each other along the main diagonal, the $q$ first ones being identical and of
degree $p$. The same is true for their linear aggregates, the elements of $\M$.

The matrices which occur in the $q$ first diagonal minors form an irreducible
system, inequivalent to those which occur in the $r-q$ other diagonal minors.
So the theorems of Burnside, Frobenius, and I. Schur
(cf. \cite{ref3}, \cite{ref8}; or \cite[p.~412]{ref14}) apply. Therefore these matrices
are perfectly arbitrary, and independent of the $r-q$ other ones. Thus we can
prescribe the $q$ first ones to be equal to the unit matrix, and the $r-q$ others
to vanish. So an element $E_0$ of $\M$ results, which commutes obviously with
all elements of $\M$; thus $E_0\in\M\cdot\M'$. As
$\M\cdot\M'=(\alpha1)$, we have $E_0=\alpha1$. This clearly necessitates
$\alpha=1$, and so the $r-q$ vanishing diagonal minors in $E_0$ cannot be
present. Therefore $r-q=0$, $q=r$, $n=p_1+\cdots+p_q=p\cdot q$.

If we write the general element $A$ of $\M$ as a matrix:
$A=\{a_{\mu,\nu}\}$, $\mu,\nu=1,\ldots,n$, then our discussions have
characterised $A$ as follows: Put
\[
 \mu=p(s-1)+\sigma,\qquad \nu=p(t-1)+\tau,
\]
where $s,t=1,\ldots,q$; $\sigma,\tau=1,\ldots,p$ (remember $n=pq$), then
\[
 a_{\mu,\nu}=
 \begin{cases}
  0,&s\ne t,\\
  \text{a function of $\sigma,\tau$ alone},&s=t.
 \end{cases}
\]
If we replace the index $\mu$ $(=1,\ldots,n)$ by the two indices $\sigma$
$(=1,\ldots,p)$ and $s$ $(=1,\ldots,q)$, and similarly $\nu$ by $\tau$ and
$t$, then we have
\[
 A=\{a_{\sigma,s,\tau,t}\}_{\substack{\sigma,\tau=1,\ldots,p\\
                                      s,t=1,\ldots,q}}
 \quad\text{with}\quad
 a_{\sigma,s,\tau,t}=\delta_{s,t}b_{\sigma,\tau}.
 \tag{S}\label{eq:S}
\]
Here $\delta_{s,t}$ is the Weierstrass--Kronecker symbol, and
$b_{\sigma,\tau}$ is arbitrary.

Conversely: If $\M$ consists of all
\[
 \{\delta_{s,t}b_{\sigma,\tau}\}_{\substack{\sigma,\tau=1,\ldots,p\\
                                             s,t=1,\ldots,q}},
\]
then it is a ring, $\M'$ consists of all
\[
 \{\delta_{\sigma,\tau}c_{s,t}\}_{\substack{\sigma,\tau=1,\ldots,p\\
                                             s,t=1,\ldots,q}},
\]
and so $\M\cdot\M'$ of all
$\{\delta_{\sigma,\tau}\delta_{s,t}\alpha\}=\alpha1$. Thus \eqref{eq:S}
gives the general solution of \eqref{eq:D5bar} for $\Hspace=\mathfrak E_n$.

In other words: The general solution of \eqref{eq:D5bar} obtains, by replacing
the one coordinate index $\mu=1,\ldots,n$ in $\Hspace=\mathfrak E_n$ by two
independent coordinate indices $\sigma=1,\ldots,p$ and $s=1,\ldots,q$
($n=pq$), and collecting in $\M$ all linear operators $A$ which operate on
$\sigma$ alone, while $\M'$ consists of all those which operate on $s$ alone.

This is the effect of the application of the classical
Burnside--Frobenius--I. Schur theory of unitary group representations. Thus the
solving of \eqref{eq:D5bar} connects directly with the fundamental facts about
unitary representations.

If the analogous situation held in Hilbert space, we would be led to expect that
if $\Hspace$ is Hilbert space, the solutions of \eqref{eq:D5bar} can be brought
in the following form: $\Hspace$ is isomorphic to the functional space of all
functions $f(x,y)$ in two independent variables $x,y$,
($\iint |f(x,y)|^2\,dx\,dy$ finite); the ranges of $x$ and of $y$ are continuous
or discrete, in the latter case $\int\cdots dx$ resp. $\int\cdots dy$ is meant to
indicate a sum. (Cf. \cite[pp.~69 and 108--111]{ref16}.) $\M$ consists of all
operators which operate on $x$ alone, and $\M'$ of all those which operate on
$y$ alone. (Cf. \secref{3.2}, where this is more precisely discussed.)

The essentially different character of Hilbert space, compared with the spaces
$\mathfrak E_n$, $n=1,2,\ldots$, is reflected in the fact that this is not so. We
will see that if $\Hspace$ is Hilbert space, \eqref{eq:D5bar} has further
solutions. (Cf. in particular \secref{8.6} and \secref{13.2}.) Their properties
seem to be of great importance for the general theory of operators.

The general importance of the solutions of \eqref{eq:D5bar} for operator
theory follows from this fact too, that their knowledge allows a characterisation
and classification of all rings of operators. This will be discussed in a
subsequent paper of the second author.

\IntroStart{4}{Relationship with quantum mechanics}
Another interpretation of \eqref{eq:D5bar} is suggested by quantum mechanics.
The operators of $\Hspace$ correspond there to all observable quantities which
occur in a mechanical system $\mathfrak S$. (Cf. \cite[pp.~55--60]{ref6}, and
\cite[p.~167]{ref20}. We restrict ourselves to bounded operators, which
correspond to those observables which have a bounded range. Thus $\B$
corresponds to the totality of these observables.) Now if $\mathfrak S$ can be
decomposed into two parts $\mathfrak S_1,\mathfrak S_2$ and if we denote the set
of the operators which correspond to observables situated entirely in
$\mathfrak S_1$ or in $\mathfrak S_2$ by $\M_1$ resp. $\M_2$, then we see:
\begin{enumerate}[label=(\arabic*)]
\item \phantomsection\label{intro4:condition1}$\M_1,\M_2$ are rings, and $1$
  (which corresponds to the ``constant''
  observable $1$) belongs to both $\M_1,\M_2$.
\item \phantomsection\label{intro4:condition2}If $A\in\M_1$, $B\in\M_2$ then
  the measurements of the observables of
  $A$ and $B$ do not interfere (being in different parts of $\mathfrak S$);
  therefore $A,B$ commute (cf. \cite[pp.~11--14 and 76]{ref6}, or
  \cite[pp.~117--121]{ref20}). Thus $\M_2\subset\M_1'$.
\item \phantomsection\label{intro4:condition3}As $\mathfrak S$ is the sum of
  $\mathfrak S_1,\mathfrak S_2$ therefore
  $R(\M_1,\M_2)=\B$.
\end{enumerate}
\hyperref[intro4:condition1]{(1)}--\hyperref[intro4:condition3]{(3)} describe
the problem of ``factorising'' $\B$ which is discussed in
more detail in \secref{3.1}; it leads to our old problem: As
$\M_1'\supset\M_2$ therefore
\[
 R(\M_1,\M_1')\supset R(\M_1,\M_2)=\B,
\]
so $R(\M_1,\M_1')=\B$, that is precisely \eqref{eq:D4bar}, which, as we know,
is equivalent to \eqref{eq:D5bar}. Conversely: If $\M$ fulfills
\eqref{eq:D4bar} (that is \eqref{eq:D5bar}), then $\M_1=\M$, $\M_2=\M'$
satisfy \hyperref[intro4:condition1]{(1)}--\hyperref[intro4:condition3]{(3)}.
(Cf. \secref{3.1} for more details.)

Thus our problem of solving \eqref{eq:D5bar} corresponds to the quantum
mechanical problem of dividing a system $\mathfrak S$ into two subsystems
$\mathfrak S_1,\mathfrak S_2$; and in particular the solutions $\M$ of
\eqref{eq:D5bar} correspond to the complete rings of all observables of
suitable quantum mechanical systems.

This interpretation of \eqref{eq:D5bar} suggests of course strongly the surmise
formulated at the end of \secref{2.2}: It should be possible to describe
$\Hspace$ as (isomorphic to) the space of all two variable functions $f(x,y)$,
($\iint|f(x,y)|^2\,dx\,dy$ finite), $\M$ operating on $x$ only, and $\M'$ on
$y$ only. In this case $\mathfrak S_1,\mathfrak S_2$ would be explicitly given:
$\mathfrak S_1$ being described by the coordinate $x$, and $\mathfrak S_2$ by
the coordinate $y$.

The fact that the surmise of \secref{2.2} is not true, is therefore the more
remarkable; particularly so because certain features of the ``exceptional'' rings
$\M$ seem to make them even better suited for quantum mechanical purposes
than the customary $\B$. We will now discuss these properties of $\M$.

\IntroStart{5}{\texorpdfstring{Properties of certain rings, in case
$\mathrm{II}_1$}{Properties of certain rings, in case II-1}}
The full system of solutions of \eqref{eq:D5bar} will be discussed in
\secref{8.3}--\secref{8.4}. While we refer to those sections for a complete
discussion, we would like to call attention here to the following: If the surmise
formulated at the end of \secref{2.2} holds, then $\M$ is obviously isomorphic
to the ring of all operators on $x$. That is, according to whether $x$ has a
finite or infinite range, $\M$ is isomorphic to the ring $\B$ of either a finite
dimensional Euclidean space $\mathfrak E_n$, $n=1,2,\ldots$, or of a Hilbert
space $\Hspace'$. In the systematic notation used in \secref{8.4} these
possibilities are labeled as cases $(\mathrm I_n)$, $n=1,2,\ldots$, respectively,
$(\mathrm I_\infty)$. (In the last case the range of $x$ may be either continuous
or discontinuous, but it is well known, that this does not affect the character
of the Hilbert space $\Hspace'$ at all.) As we mentioned, these cases do not
exhaust all possibilities for $\M$: further cases, labeled $(\mathrm{II}_1)$,
$(\mathrm{II}_\infty)$, $(\mathrm{III}_\infty)$ may occur. (All of them, except
$(\mathrm{III}_\infty)$ certainly exist; while the existence of
$(\mathrm{III}_\infty)$ is as yet undecided. Cf. \secref{13.3}.)

The cases $(\mathrm I_n)$, $n=1,2,\ldots$, are of course the simplest ones, and
they are the ones on which our notions as to what operators should look like
have been developed. This is true for general operator theory as well as for
quantum mechanics. It is therefore of particular importance to compare the
other cases $(\mathrm I_\infty)$--$(\mathrm{III}_\infty)$ under the following
aspects: Which one has the most properties in common with $(\mathrm I_n)$, the
ring of all matrices of $n$ rows and $n$ columns, and which one can be
considered as the limiting case of $(\mathrm I_n)$ for $n\to\infty$? The
investigations of operators in Hilbert space have always been carried on with
the idea that the $\B$ of Hilbert space, that is $(\mathrm I_\infty)$, is the
natural limiting case. We think however that our results indicate that there is
more point in assigning this r\^ole to $(\mathrm{II}_1)$. There are two chief
reasons for this assertion: The existence of a trace, and the behavior of
unbounded operators. Chaps. XV and XVI of this paper have been devoted to
the purpose of deriving these common features of $(\mathrm I_n)$ and
$(\mathrm{II}_1)$. In what follows we will make a few qualitative remarks on
the subject.

The cases $(\mathrm I_n)$ and $(\mathrm{II}_1)$ are characterised by this
property: It is possible to define for all Hermitian operators $A\in\M$ a
function $T(A)$ with finite real values, which has the formal properties of the
trace (cf. \secref{15.4}--\secref{15.5}), and there exists only one such function
$T(A)$. In the cases $(\mathrm I_n)$, where the operators $A\in\M$ correspond
to matrices $\{a_{\mu,\nu}\}_{\mu,\nu=1,\ldots,n}$, obviously
\[
 T(A)=\frac1n\sum_{1}^{n} a_{\mu,\mu}
\]
(the ``normalising factor'' $1/n$ is necessary, because we require $T(1)=1$),
and it is well known, that in case $(\mathrm I_\infty)$ (that is: for all bounded
infinite matrices $\{a_{\mu,\nu}\}_{\mu,\nu=1,2,\ldots}$) no such function
$T(A)$ can be defined. (The expression which is formed in analogy to
$(1/n)\sum_{1}^{n} a_{\mu,\mu}$ will not converge in general.) Now we
proved, as mentioned above, that the only further case where such a $T(A)$
exists and is unique, is $(\mathrm{II}_1)$.

Considering the immediate applicability of $T(A)$ to quantum mechanics (it
is the ``a priori'' expectation value of the observable $A$, which is correctly
normalised here, but cannot be in $(\mathrm I_\infty)$,
cf. \cite[pp.~165, 169]{ref20}), it is significant that its existence and
uniqueness is a characteristic of $(\mathrm{II}_1)$.

A more general problem which will be discussed by the second named author in
a forthcoming paper arises now from the following motive: It would be
desirable to characterise those abstract algebras, in which one and only one
function $T(A)$ with the properties of the trace (as discussed above) exists.
(They should be abstract, that is no connection with the operator rings should
be postulated.) The result is, that all algebras which meet these requirements
are isomorphic to an operator ring $\M$ in a (Euclidean or Hilbert) space
$\Hspace$, which solves \eqref{eq:D5bar} and belongs to cases $(\mathrm I_n)$,
$n=1,2,\ldots$, or $(\mathrm{II}_1)$.

Returning to case $(\mathrm{II}_1)$ let us point out, that the analogy with
$(\mathrm I_n)$ goes so far, that theorems of the following type hold: Every
system of linear equations (which can of course be written as one operatorial
equation $Af=0$ where $A\in\M$ is the ``coefficient matrix,'' and $f\in\Hspace$
represents the variables) has a rank. In case $(\mathrm I_n)$ this is a number
$0,1,\ldots,n-1,n$, but we replace it by its $n$-th part:
$0,1/n,\ldots,(n-1)/n,1$. Now in case $(\mathrm{II}_1)$ it is any number
$\geq0,\leq1$. Two adjoint systems of equations $Af=0$ and
$A^*f=0$ (cf. \secref{16.1}) have the same rank. (Observe the analogy with
$(\mathrm I_n)$, and the contrast with $(\mathrm I_\infty)$!) Similarly a
dimensionality for subspaces of $\Hspace$ (linear, closed sets) can be defined,
which has the values $0,1/n,\ldots,(n-1)/n,1$ (in the customary normalisation
$0,1,\ldots,n-1,n$) in case $(\mathrm I_n)$, and all values
$\geq0,\leq1$ in case $(\mathrm{II}_1)$.

With this notion of dimensionality even the ``minimax principle''
(cf. \cite[pp.~26--29]{ref5}) can be proved, which thus labels the proper values
of all Hermitian operators (in $\M$), even in continuous spectra! So it makes
sense to speak about the ``proper value No. $\alpha$ of the operator $A$ (from
below)'' for any $\alpha\geq0,\leq1$, even for irrational ones!
(Cf. \secref{15.2}.)

\IntroStart{6}{\texorpdfstring{Unbounded operators derived from a case
$\mathrm{II}_1$}{Unbounded operators derived from a case II-1}}
Let us now consider the unbounded operators. While $\M$ consists prima facie
of bounded operators only, there is a natural way to extend it, so as to include
unbounded ones too: An arbitrary (not necessarily bounded) operator $X$ is said
to belong to $\M$, if it is invariant under all unitary transformations
$U'\in\M'$; we denote the set of all those linear, closed operators $X$ with an
everywhere dense domain which belong to $\M$, by $U(\M)$. (Cf. \secref{4.2}
and \secref{16.4}; as to the notions of linearity and closure cf.
\cite[p.~70]{ref16}. The bounded elements of $U(\M)$ form precisely $\M$.)

While in general Hermitian operators possess a resolution of unity if and only
if they are hypermaximal (cf. \cite[p.~92]{ref16}, or
\cite[Theorem 9.3, p.~339]{ref15}) this is much simpler in cases
$(\mathrm I_n)$ and $(\mathrm{II}_1)$: Here every Hermitian operator
$A\in U(\M)$ is hypermaximal, and has therefore a (unique) resolution of unity.
(In case $(\mathrm I_n)$ this is due to the trivial reason, that every linear
operator is bounded in that case. But in case $(\mathrm{II}_1)$ unbounded
operators exist, in the same way as they do in $(\mathrm I_\infty)$!)

The operators in $U(\B)$ ($\Hspace$ a Hilbert space) behave very pathologically
from the point of view of operator algebra: Thus if $A,B\in U(\B)$, then $A+B$
and $AB$ cannot be formed in general, and the entire mechanism of replacing
operators by their extension leads to very paradoxic results.
(Cf. \cite[pp.~230--234]{ref17}, where these conditions are discussed in
detail.) This applies, of course, to every $U(\M)$, $\M$ in case
$(\mathrm I_\infty)$. In cases $(\mathrm I_n)$ and $(\mathrm{II}_1)$ again it is
not so: The algebra of $U(\M)$ works without any difficulties, and if $A$ is an
extension of $B$, $A,B\in U(\M)$, then $A=B$. (Cf. for details \secref{16.3}
and \secref{16.4}.)

\IntroStart{7}{Outline}
A detailed account of the content of this paper is given in the table of contents
which follows. The main results are summed up in the
\hyperref[thm:I]{Theorems I}--\hyperref[thm:XV]{XV}, while
the essential problems are formulated as Problems. It will be pointed out after
each problem to what extent it is solved and where the solutions are to be
found. The auxiliary considerations, which lead to the Theorems, are grouped
into Lemmas.

The notations are explained in \secref{1.1}, where the chief references too are
given. These, as well as all other quotations, refer to the bibliography, which
follows after the table of contents.

Various continuations of the investigations begun in this paper, which we
propose to undertake in several subsequent papers, are indicated throughout the
text. (Cf. also note at the end of this paper.)

Here however we would like to mention that very similar methods to those used
in this paper permit us to discuss the corresponding problems in the
non-separable analogues of Hilbert space (cf. \cite{ref11}, \cite{ref13}). The analogy
to Cantor's theory of alephs, which we stress in
\hyperref[chap:6]{Chapters VI} and \hyperref[chap:7]{VII}, becomes
then even more apparent. This subject, too, will be treated at a later occasion.
% --- End of parts/01-introduction.tex ---

\phantomsection
\pdfbookmark[1]{Contents}{bookmark:contents}
\renewcommand{\contentsname}{Contents}
\tableofcontents

\phantomsection
\pdfbookmark[1]{Bibliography}{bookmark:bibliography}
\renewcommand{\refname}{Bibliography}
\begin{thebibliography}{12a}
\bibitem[1]{ref1} G. Birkhoff: \emph{On Combinations of Sub-Algebras}. Proc. Cambridge Phil. Soc. Vol. 29, pp. 441--464, (1933).
\bibitem[2]{ref2} S. Bochner: \emph{Integration einer Funktion deren Werte elemente eines Vektorraumes sind}. Fund. Math. Vol. 20, pp. 262--276 (1933).
\bibitem[3]{ref3} W. Burnside: \emph{On the Condition of Reducibility of any Group of Linear Substitutions}. Proc. London Math. Soc. Vol. 3, pp. 430--440 (1905).
\bibitem[4]{ref4} C. Carath\'eodory: \emph{Reele Funktionen}. B. G. Teubner, Leipzig and Berlin, (1918 and 1927). (The quotations follow the 1918 edition.)
\bibitem[5]{ref5} R. Courant and D. Hilbert: \emph{Methoden der Mathematischen Physik}. O. J. Springer, Berlin, (1931).
\bibitem[6]{ref6} P. A. M. Dirac: \emph{The Principles of the Quantum Mechanics}. Clarendon Press. Oxford (1930).
\bibitem[7]{ref7} K. Friedrichs: \emph{Spektraltheorie halbbeschr\"ankter Operatoren}. Math. Ann. Vol. 109, pp. 465--487 (1934).
\bibitem[8]{ref8} G. Frobenius and I. Schur: \emph{\"Uber die \"Aquivalenz der Gruppen Linearen Substitutionen}. Berlin Ber. 1906, pp. 209--217.
\bibitem[9]{ref9} F. Klein: \emph{Abstrakte Verkn\"upfungen}. Math. Ann. Vol. 105, pp. 308--323, (1931).
\bibitem[10]{ref10} H. Lebesgue: \emph{Le\c{c}ons sur l'int\'egration et la recherche des fonctions primitives}. Gauthier, Paris, (1904 and 1928). (The quotations follow the 1904 edition.)
\bibitem[11]{ref11} H. L\"owig: \emph{Komplexe euklidische R\"aume von beliebiger endlicher oder transfiniter Dimensionszahl}. Acta Szeged, Vol. 7, pp. 1--33, (1934).
\bibitem[12]{ref12} F. J. Murray: \emph{A Theory for * operators analogous to the theory of reducibility for self-adjoint transformations in Hilbert Space}. Proc. Nat. Acad. Vol. 19, pp. 1057--1058, (1933).
\bibitem[12a]{ref12a} O. Ore: \emph{On the foundation of abstract algebra}. Annals of Math. Vol. 36, pp. 406--437, (1935).
\bibitem[13]{ref13} F. Rellich: \emph{Spektraltheorie in nichtseparablen R\"aumen}. Math. Ann. vol. 110, pp. 342--356, (1935).
\bibitem[13a]{ref13a} F. Riesz: \emph{Zur Theorie des Hilbertschen Raumes}. Acta Szeged, Vol. 7, pp. 34--38, (1934).
\bibitem[14]{ref14} I. Schur: \emph{Neue Begr\"undung der Theorie der Gruppencharaktere}. Sitzungber. Preuss. Akad. 1905, pp. 406--432.
\bibitem[15]{ref15} M. H. Stone: \emph{Linear Transformations in Hilbert Space}. Amer. Math. Soc. Colloquim Publications Vol. XV (1932).
\bibitem[16]{ref16} J. v. Neumann: \emph{Allgemeine Eigenwerttheorie Hermitescher Funktionaloperatoren}. Math. Ann. Vol. 102, pp. 49--131, (1929).
\bibitem[17]{ref17} J. v. Neumann: \emph{Zur Theorie der unbeschr\"ankten Matrizen}. Journal f\"ur Math. 161, pp. 208--236, (1929).
\bibitem[18]{ref18} J. v. Neumann: \emph{Zur Algebra der Funktionaloperatoren}. Math. Ann. Vol. 102, pp. 370--427, (1929).
\bibitem[19]{ref19} J. v. Neumann: \emph{\"Uber Funktionen von Funktionaloperatoren}. Annals of Math. Vol. 32, pp. 191--226, (1931).
\bibitem[20]{ref20} J. v. Neumann: \emph{Mathematische Begr\"undung der Quantenmechanik}. Springer, Berlin (1931).
\bibitem[21]{ref21} J. v. Neumann: \emph{\"Uber adjungierte Funktionaloperatoren}. Annals of Math. Vol. 33, pp. 294--310, (1932).
\bibitem[22]{ref22} J. v. Neumann: \emph{On a certain topology for rings of operators}. Annals of Math. Vol. 37 preceding immediately, (1936).
\end{thebibliography}

\Needspace{10\baselineskip}
\PartHeading{I}{Preparatory Considerations}

\ChapterHeading{1}{I}{Notations}

\SectionStart{1.1}{Notations}
In what follows we will have to assume that the reader is familiar with
the unitary--orthogonal--geometrical discussion of the elementary properties
of Hilbert space, as contained in the treatise \cite{ref16} (particularly
pp.~63--78) or in the remarkable exposition of the subject in
\cite{ref15} Chapter~I. Besides we will have to make frequent use of some
other papers of one of us, namely \cite{ref18}, \cite{ref21} and \cite{ref22},
which is a variant to a theorem of \cite{ref18}.

Certain notations will be used on this account subsequently, without further
references to their meaning. They are as follows:

\noindent (a) $x\in S$ means that $x$ is an element of the set $S$,
$S\subset T$ or $T\supset S$ means that $S$ is a subset of $T$, (including
$S=T$). The set of all elements $x$ with a certain property $\epsilon(x)$,
will be denoted by $(x;\epsilon(x))$; the set of all $f(x)$ with $x$ having
the property $\epsilon(x)$, by $(f(x);\epsilon(x))$; the set theoretical sum
of the sets $S,T,\ldots$, plus the elements $x_0,y_0,\ldots$ by
$(S,T,\ldots,x_0,y_0,\ldots)$. The set theoretical product (common part) of
the sets $S,T,\ldots$ is denoted by $S\cdot T\cdots$.

A symbol $\eta$, such that $x\mathbin{\eta}S$ has a meaning similar to that
of $x\in S$, will be defined in \secref{4.2}, \defref{4.2.1}.

\noindent (b) A linear space with a linear product, and which is separable and
complete, will be denoted by $\Hspace$. (We will use affixes and suffixes if
more than one such space occurs.) In other words: $\Hspace$ is a space in
which operations $\alpha\cdot f$, $f+g$, $(f,g)$ satisfying the conditions
A, B, C, E, of \cite{ref16} p.~64--66 are given. Condition D (loc. cit.) is
explicitly excepted. Thus $\Hspace$ is either a Hilbert space or a finite
dimensional Euclidean space accordingly as D is fulfilled or not. We only
consider $1,2,\ldots$, dimensional Euclidean spaces, excluding explicitly the
0-dimensional case where $\Hspace=(0)$.

\noindent (c) We denote complex or real numbers by
$\alpha,\epsilon,a,x,C,D$; integers by
$i,\allowbreak j,\allowbreak k,\allowbreak m,\allowbreak n,\allowbreak
p,\allowbreak q,\allowbreak s,\allowbreak t,\allowbreak u,\allowbreak v$ and $N$.
Elements of $\Hspace$ are denoted by $f,g,\varphi,\psi$, sometimes (in direct
products) by $\Phi,\Psi$.

\noindent (d) Arbitrary subsets of $\Hspace$ are denoted by $\mathfrak S$,
linear subsets by $\mathfrak A$, linear and closed subsets by
$\mathfrak E$, $\mathfrak F$, $\mathfrak M$, $\mathfrak N$,
$\mathfrak P$, $\mathfrak Q$.
As the latter are again spaces of the type of $\Hspace$, their symbols
sometimes replace $\Hspace$.

The smallest linear and the smallest closed linear set containing certain sets
or elements are denoted by $(\ldots)$ and $[\ldots]$ respectively (instead of
$(\ldots)$, cf. (a)). The set of all elements of $\mathfrak N$ which are
orthogonal to $\mathfrak M$ is denoted by $\mathfrak N-\mathfrak M$.

\noindent (e) Operators in an $\Hspace$ will be denoted by $A,B,X,Y$. For
certain special kinds of operators, the following letters will be used:
$E,F,G$ and $P$ for projection operators, $U,V$, and $W$ for unitary and
partially isometric operators (cf. \cite[p.~70]{ref16}).

\noindent (f) The ring of all bounded, everywhere defined, linear and closed
operators in $\Hspace$ is $\B$, its subrings are $\M,\N$, (cf.
\cite[p.~388]{ref18}).

Arbitrary subsets of $\B$ are $S$, the smallest ring containing $S$ is
$R(S)$, the ring of all $A$ which commute with $X$ and $X^*$ for every
$X\in S$ is $S'$.

\noindent (g) As discussed in \cite[pp.~378--388]{ref18}, two different
topologies can be used in $\Hspace$, and four in $\B$, ``Strong'' and
``Weak'' in $\Hspace$ and $\B$, ``uniform'' in $\B$ and finally there is the
``strongest'' topology in $\B$. If nothing in particular is said, we always
mean the ``strong'' topology in $\Hspace$; otherwise the topology to be used
will be specified. The properties of these topologies are not all of the
customary type and they are of some importance. Details are given in
\cite{ref18} and \cite{ref22} loc. cit.

\noindent (h) In part IV, a group $\mathfrak G$ and a space $S$, will be
considered, where $\mathfrak G$ consists of one-to-one mappings of $S$ on
itself. We will denote the elements of $\mathfrak G$ by $a,b$, and $c$; those
of $S$ by $x,y$; the unit, the inverse induced by $a\in\mathfrak G$ in $S$
are $1,a^{-1}$. These notations will however only be used in
\hyperref[part:IV]{Part IV}.

The results of operator theory which we will use most frequently (apart from
the general theory of Hilbert space and the spectral form of hypermaximal
operators) are these: Theorems 1, 5, 9 in \cite{ref18}, Theorems 5, 6, in
\cite{ref19}, and Theorem 7 in \cite{ref21}.

\ChapterHeading{2}{II}{Direct Products}

\SectionStart{2.1}{Construction of the direct product. Theorem I}
An operation which generates a new space $\Hspace$ with the help of $n$ given
ones, $\Hspace_1,\ldots,\Hspace_n$ is the \emph{direct multiplication}. As it
is one of the essential elements of the situations which we are going to
discuss, and as a general abstract treatment (valid for Hilbert spaces too and
not using special coordinate systems) of this notion does not occur in the
literature, we will give a systematic discussion.

\Definition{2.1.1}
Let $n=1,2,\ldots$ spaces $\Hspace_1,\ldots,\Hspace_n$ be given. Consider all
functionals $\Phi(f_1,\ldots,f_n)$, which are defined for all systems
\[
  f_i\in\Hspace_i,\quad i=1,\ldots,n,
\]
and have complex values, and which are conjugate linear in each $f_i$:
\begin{align*}
 \text{(i)}\quad
 \Phi(\ldots,\alpha f_i,\ldots)
   &=\overline\alpha\Phi(\ldots,f_i,\ldots),\\
 \text{(ii)}\quad
 \Phi(\ldots,f_i+g_i,\ldots)
   &=\Phi(\ldots,f_i,\ldots)+\Phi(\ldots,g_i,\ldots).
\end{align*}
Call their set $\prod_{i=1}^{n}\solidotimes\Hspace_i$.

\Definition{2.1.2}
If a system $f_i^0\in\Hspace_i$, $i=1,\ldots,n$, is given, form the functional
\[
 \Phi(f_1,\ldots,f_n)=\prod_{i=1}^{n}(f_i^0,f_i).
\]
Clearly $\Phi\in\prod_{i=1}^{n}\solidotimes\Hspace_i$. Define
\[
 \overline{\Phi}=\prod_{i=1}^{n}\otimes f_i^0
      =f_1^0\otimes\cdots\otimes f_n^0.
\]

\Definition{2.1.3}
Consider all finite linear aggregates of the form
\[
 \overline{\Phi}=\sum_{\nu=1}^{p}\prod_{i=1}^{n}\otimes f_{i,\nu}^0,
 \qquad p=0,1,\ldots,\quad f_{i,\nu}^0\in\Hspace_i.
\]
Call their set $\mathop{\prod\nolimits'}_{i=1}^{n}\otimes\Hspace_i$.
(Clearly
$\mathop{\prod\nolimits'}_{i=1}^{n}\otimes\Hspace_i
 \subset\prod_{i=1}^{n}\otimes\Hspace_i$.)

\Lemma{2.1.1}
{\itshape
If $\overline{\Phi},\overline{\Psi}\in
\mathop{\prod\nolimits'}_{i=1}^{n}\otimes\Hspace_i$, that is
\[
 \overline{\Phi}=\sum_{\nu=1}^{p}\prod_{i=1}^{n}\otimes f_{i,\nu}^0;
 \qquad
 \overline{\Psi}=\sum_{\mu=1}^{q}\prod_{i=1}^{n}\otimes g_{i,\mu}^0,
\]
then form
\[
 (\overline{\Phi},\overline{\Psi})=
 \sum_{\nu=1}^{p}\sum_{\mu=1}^{q}
              \prod_{i=1}^{n}(f_{i,\nu}^0,g_{i,\mu}^0).
\]
This expression depends on $\overline{\Phi},\overline{\Psi}$ only, but not on the particular
decomposition used.}

\Proof Owing to the linearity it suffices to consider the case when
$\overline{\Phi}=0$ or $\overline{\Psi}=0$; and owing to the symmetry in
$\overline{\Phi},\overline{\Psi}$, we may assume $\overline{\Psi}=0$.
We will show that each addend of the sum $\sum_{\nu=1}^{p}$ in
$(\overline{\Phi},\overline{\Psi})$
vanishes separately. Indeed:
\[
 \sum_{\mu=1}^{q}\prod_{i=1}^{n}(f_{i,\nu}^0,g_{i,\mu}^0)
 =\overline{\Psi}(f_{1,\nu}^0,\ldots,f_{n,\nu}^0)=0.
\]

\Lemma{2.1.2}
{\itshape
If $\overline{\Phi}\in\mathop{\prod\nolimits'}_{i=1}^{n}
\otimes\Hspace_i$, then
$(\overline{\Phi},\overline{\Phi})>0$ for $\overline{\Phi}\ne0$.}

\Proof Consider first any
$\overline{\Phi}\in\mathop{\prod\nolimits'}_{i=1}^{n}
\otimes\Hspace_i$. Then
\[
 (\overline{\Phi},\overline{\Phi})=\sum_{\mu,\nu=1}^{p}
              \prod_{i=1}^{n}(f_{i,\mu}^0,f_{i,\nu}^0).
\]
For any complex $x_1,\ldots,x_p$,
\[
 \sum_{\mu,\nu=1}^{p}(f_{i,\mu}^0,f_{i,\nu}^0)x_\mu\overline{x_\nu}
 =\left(\sum_{\mu=1}^{p}x_\mu f_{i,\mu}^0,
        \sum_{\nu=1}^{p}x_\nu f_{i,\nu}^0\right)
 =\left\lVert\sum_{\mu=1}^{p}x_\mu f_{i,\mu}^0\right\rVert^2\geq0;
\]
therefore each matrix
$((f_{i,\mu}^0,f_{i,\nu}^0))_{\mu,\nu=1,\ldots,p}$ ($p$ dimensional!) is
semidefinite. Thus it is the sum of (at most $p$) semidefinite matrices of
rank 1, that is of the form
$((\alpha_\mu^i\overline{\alpha_\nu^i}))_{\mu,\nu=1,\ldots,p}$. Thus
$(\overline{\Phi},\overline{\Phi})$ is a sum of terms
\begin{align*}
 \sum_{\mu,\nu=1}^{p}\prod_{i=1}^{n}
       \alpha_\mu^i\overline{\alpha_\nu^i}
 &=\sum_{\mu,\nu=1}^{p}\prod_{i=1}^{n}\alpha_\mu^i
       \cdot\overline{\prod_{i=1}^{n}\alpha_\nu^i}\\
 &=\sum_{\mu=1}^{p}\prod_{i=1}^{n}\alpha_\mu^i
       \cdot\overline{\sum_{\nu=1}^{p}\prod_{i=1}^{n}\alpha_\nu^i}
  =\left|\sum_{\mu=1}^{p}\prod_{i=1}^{n}\alpha_\mu^i\right|^2\geq0.
\end{align*}
So we have $(\overline{\Phi},\overline{\Phi})\geq0$.

From this, Schwarz's inequality
\[
 |(\overline{\Phi},\overline{\Psi})|
 \leq\sqrt{(\overline{\Phi},\overline{\Phi})
 \cdot(\overline{\Psi},\overline{\Psi})}
\]
follows literally as in \cite[p.~64]{ref16}. Thus
$(\overline{\Phi},\overline{\Phi})=0$ implies
$(\overline{\Phi},\overline{\Psi})=0$ for every
$\overline{\Psi}\in\mathop{\prod\nolimits'}_{i=1}^{n}\otimes\Hspace_i$.
Put $\overline{\Psi}=\prod_{i=1}^{n}\otimes f_i$, then
\[
 \overline{\Phi}(f_1,\ldots,f_n)
   =\left(\overline{\Phi},\prod_{i=1}^{n}\otimes f_i\right)=0,
\]
and so $\overline{\Phi}=0$. Thus $\overline{\Phi}\ne0$ implies
$(\overline{\Phi},\overline{\Phi})>0$.

\Lemma{2.1.3}
{\itshape
With the above definition of $(\overline{\Phi},\overline{\Psi})$,
$\mathop{\prod\nolimits'}_{i=1}^{n}\otimes\Hspace_i$ is a linear space with
an inner product, that is, it satisfies the conditions A, B, in
\cite[p.~64]{ref16}. Thus it can be metrized by defining
\[
 \operatorname{Distance}(\overline{\Phi},\overline{\Psi})
 =\lVert\overline{\Phi}-\overline{\Psi}\rVert,
 \qquad\text{where}\qquad
 \lVert\overline{\Phi}\rVert
 =\sqrt{(\overline{\Phi},\overline{\Phi})}.
\]}

\Proof This follows immediately from Lemmas
\hyperref[lem:2.1.1]{2.1.1} and \hyperref[lem:2.1.2]{2.1.2}, remembering
\cite[pp.~64--65]{ref16}.

\Lemma{2.1.4}
{\itshape
We have
\[
 \left\lVert\prod_{i=1}^{n}\otimes f_i\right\rVert
   =\prod_{i=1}^{n}\lVert f_i\rVert,
\]
and for every $\overline{\Phi}\in\mathop{\prod\nolimits'}_{i=1}^{n}
\otimes\Hspace_i$,
\[
 \overline{\Phi}(f_1,\ldots,f_n)
   =\left(\overline{\Phi},\prod_{i=1}^{n}\otimes f_i\right),
 \qquad
 |\overline{\Phi}(f_1,\ldots,f_n)|
   \leq\lVert\overline{\Phi}\rVert\cdot
   \prod_{i=1}^{n}\lVert f_i\rVert.
\]}

\Proof The first equation follows directly from the definition of
$\lVert\cdots\rVert$ in
$\mathop{\prod\nolimits'}_{i=1}^{n}\otimes\Hspace_i$; while the second is
obvious, and the third results from Schwarz's inequality:
\[
 |\overline{\Phi}(f_1,\ldots,f_n)|
 =\left|\left(\overline{\Phi},\prod_{i=1}^{n}\otimes f_i\right)\right|
 \leq\lVert\overline{\Phi}\rVert\cdot
      \left\lVert\prod_{i=1}^{n}\otimes f_i\right\rVert
 =\lVert\overline{\Phi}\rVert\cdot\prod_{i=1}^{n}\lVert f_i\rVert.
\]

\Definition{2.1.4}
Consider those functionals
$\Phi\in\prod_{i=1}^{n}\solidotimes\Hspace_i$, for which a sequence
$\overline{\Phi}_1,\overline{\Phi}_2,\ldots\in
\mathop{\prod\nolimits'}_{i=1}^{n}\otimes\Hspace_i$ exists, so that
\begin{align*}
 \text{(i)}\quad &
 \lim_{r\to\infty}\overline{\Phi}_r(f_1,\ldots,f_n)
 =\Phi(f_1,\ldots,f_n)
   \quad\text{for all }f_i\in\Hspace_i,\\
 \text{(ii)}\quad &
 \lim_{r,s\to\infty}\lVert\overline{\Phi}_r-
 \overline{\Phi}_s\rVert=0.
\end{align*}
We write for this briefly
$\Phi\sim(\overline{\Phi}_1,\overline{\Phi}_2,\ldots)$. Call their set
\[
 \prod_{i=1}^{n}\otimes\Hspace_i
   =\Hspace_1\otimes\cdots\otimes\Hspace_n,
\]
the \emph{direct product} of $\Hspace_1,\ldots,\Hspace_n$.

\Lemma{2.1.5}
{\itshape
Every sequence $\overline{\Phi}_1,\overline{\Phi}_2,\ldots$ from
$\mathop{\prod\nolimits'}_{i=1}^{n}\otimes\Hspace_i$ satisfying (ii) in
\defref{2.1.4} belongs to precisely one
$\Phi\in\prod_{i=1}^{n}\otimes\Hspace_i$ in the sense of (i), (ii) eod.

If $\Phi\sim(\overline{\Phi}_1,\overline{\Phi}_2,\ldots)$, then all other
sequences with $\Phi\sim(\overline{\Psi}_1,\overline{\Psi}_2,\ldots)$ are
characterised by
$\lim_{r\to\infty}\lVert\overline{\Phi}_r-\overline{\Psi}_r\rVert=0$.}

\Proof We have (by \lemref{2.1.4})
\[
 |\overline{\Phi}_r(f_1,\ldots,f_n)
 -\overline{\Phi}_s(f_1,\ldots,f_n)|
 \leq\lVert\overline{\Phi}_r-\overline{\Phi}_s\rVert
 \prod_{i=1}^{n}\lVert f_i\rVert.
\]
Thus
$\lim_{r,s\to\infty}|\overline{\Phi}_r(f_1,\ldots,f_n)-
\overline{\Phi}_s(f_1,\ldots,f_n)|=0$, so
$\lim_{r\to\infty}\overline{\Phi}_r(f_1,\ldots,f_n)$ exists. Call it
$\Phi(f_1,\ldots,f_n)$. Clearly
$\Phi\in\prod_{i=1}^{n}\solidotimes\Hspace_i$. Thus (i), (ii) hold; therefore
$\Phi\in\prod_{i=1}^{n}\otimes\Hspace_i$;
$\Phi\sim(\overline{\Phi}_1,\overline{\Phi}_2,\ldots)$.
$\Phi$ is unique by (i).

As to the second statement, we can assume, owing to the linearity, that
$\Phi=\overline{\Phi}_1=\overline{\Phi}_2=\cdots=0$. If
$\lim_{r\to\infty}\lVert\overline{\Psi}_r\rVert=0$ then we find, as above,
$0=\lim_{r\to\infty}\overline{\Psi}_r(f_1,\ldots,f_n)$. And
$\lVert\overline{\Psi}_r-\overline{\Psi}_s\rVert
\leq\lVert\overline{\Psi}_r\rVert+
\lVert\overline{\Psi}_s\rVert$, thus
$\lim_{r,s\to\infty}\lVert\overline{\Psi}_r-
\overline{\Psi}_s\rVert=0$. Thus
$0\sim(\overline{\Psi}_1,\overline{\Psi}_2,\ldots)$.

Assume conversely $0\sim(\overline{\Psi}_1,\overline{\Psi}_2,\ldots)$.
If we did not have
$\lim_{r\to\infty}\lVert\overline{\Psi}_r\rVert=0$, we would have
$\lVert\overline{\Psi}_r\rVert\geq a>0$ for a fixed $a>0$ and infinitely
many $r$'s.
Considering a suitable subsequence, we obtain it for all $r$'s. Now
$\lim_{r\to\infty}\overline{\Psi}_r(f_1,\ldots,f_n)=0$ implies
$\lim_{r\to\infty}(\overline{\Psi}_r,\overline{\Psi}^{0})=0$ for any fixed
$\overline{\Psi}^{0}\in
\mathop{\prod\nolimits'}_{i=1}^{n}\otimes\Hspace_i$. Put
$\overline{\Psi}^{0}=\overline{\Psi}_{r_0}$, where $r_0$ is so great that
$r,s\geq r_0$ imply
$\lVert\overline{\Psi}_r-\overline{\Psi}_s\rVert\leq a/2$.
Then we have for $r\geq r_0$
\begin{align*}
 |(\overline{\Psi}_r,\overline{\Psi}_{r_0})|
 &\geq |(\overline{\Psi}_{r_0},\overline{\Psi}_{r_0})|
       -|(\overline{\Psi}_{r_0}-\overline{\Psi}_r,
          \overline{\Psi}_{r_0})|\\
 &\geq\lVert\overline{\Psi}_{r_0}\rVert^2
       -\lVert\overline{\Psi}_{r_0}-\overline{\Psi}_r\rVert
        \cdot\lVert\overline{\Psi}_{r_0}\rVert\\
 &\geq\lVert\overline{\Psi}_{r_0}\rVert^2
       -\frac a2\lVert\overline{\Psi}_{r_0}\rVert
  =\left(\lVert\overline{\Psi}_{r_0}\rVert-\frac a2\right)
       \lVert\overline{\Psi}_{r_0}\rVert\\
 &\geq\left(a-\frac a2\right)a=\frac{a^2}{2},
\end{align*}
contradicting
$\lim_{r\to\infty}(\overline{\Psi}_r,\overline{\Psi}_{r_0})=0$.

\Lemma{2.1.6}
{\itshape
If $\Phi,\Psi\in\prod_{i=1}^{n}\otimes\Hspace_i$, that is
\[
 \Phi\sim(\overline{\Phi}_1,\overline{\Phi}_2,\ldots),\qquad
 \Psi\sim(\overline{\Psi}_1,\overline{\Psi}_2,\ldots),\qquad
 \overline{\Psi}_r,\overline{\Phi}_r\in
 \mathop{\prod\nolimits'}_{i=1}^{n}\otimes\Hspace_i,
\]
then $\lim_{r\to\infty}(\overline{\Phi}_r,\overline{\Psi}_r)$ exists.
Denote it by
$(\Phi,\Psi)$. This quantity depends on $\Phi,\Psi$ only, and not on the
particular representation used. If
$\overline{\Phi},\overline{\Psi}\in
\mathop{\prod\nolimits'}_{i=1}^{n}\otimes\Hspace_i$, it agrees
with the previous definition.}

\Proof
$\lim_{r,s\to\infty}\lVert\overline{\Phi}_r-\overline{\Phi}_s\rVert=0$,
$\lim_{r,s\to\infty}\lVert\overline{\Psi}_r-\overline{\Psi}_s\rVert=0$
imply by Schwarz's
inequality
$\lim_{r,s\to\infty}|(\overline{\Phi}_r,\overline{\Psi}_r)
-(\overline{\Phi}_s,\overline{\Psi}_s)|=0$, thus
$\lim_{r\to\infty}(\overline{\Phi}_r,\overline{\Psi}_r)$ exists.
If further
$\Phi\sim(\overline{\Phi}'_1,\overline{\Phi}'_2,\ldots)$,
$\Psi\sim(\overline{\Psi}'_1,\overline{\Psi}'_2,\ldots)$, then we have by
\lemref{2.1.5}
$\lim_{r\to\infty}\lVert\overline{\Phi}_r-\overline{\Phi}'_r\rVert=0$,
$\lim_{r\to\infty}\lVert\overline{\Psi}_r-\overline{\Psi}'_r\rVert=0$
and so
\[
 \lim_{r\to\infty}(\overline{\Phi}_r,\overline{\Psi}_r)
 =\lim_{r\to\infty}(\overline{\Phi}'_r,\overline{\Psi}'_r).
\]
If
$\overline{\Phi},\overline{\Psi}\in
\mathop{\prod\nolimits'}_{i=1}^{n}\otimes\Hspace_i$, then
$\Phi\sim(\overline{\Phi},\overline{\Phi},\ldots)$,
$\Psi\sim(\overline{\Psi},\overline{\Psi},\ldots)$ show that the
new $(\Phi,\Psi)$ agrees with the old one.

\Lemma{2.1.7}
{\itshape
If $\Phi\in\prod_{i=1}^{n}\otimes\Hspace_i$, then $(\Phi,\Phi)>0$ for
$\Phi\ne0$.}

\Proof $(\Phi,\Phi)\geq0$ follows by continuity from \lemref{2.1.2}. If
$(\Phi,\Phi)=0$, then for
$\Phi\sim(\overline{\Phi}_1,\overline{\Phi}_2,\ldots)$,
$\overline{\Phi}_r\in
\mathop{\prod\nolimits'}_{i=1}^{n}\otimes\Hspace_i$,
$\lim_{r\to\infty}(\overline{\Phi}_r,\overline{\Phi}_r)=0$,
\[
 \lim_{r\to\infty}\lVert\overline{\Phi}_r\rVert=0,
\]
and thus by \lemref{2.1.5} $\Phi=0$. So $\Phi\ne0$ implies
$(\Phi,\Phi)>0$.

\Lemma{2.1.8}
{\itshape
With the above definition of $(\Phi,\Psi)$,
$\prod_{i=1}^{n}\otimes\Hspace_i$ is a linear space with an inner product,
that is it satisfies the conditions A, B in \cite[p.~64]{ref16}. Thus it can
be metrized by defining.
\[
 \operatorname{Distance}(\Phi,\Psi)=\lVert\Phi-\Psi\rVert,
 \qquad\text{where}\qquad
 \lVert\Phi\rVert=(\Phi,\Phi)^{\frac12}.
\]}

\Proof This follows immediately from Lemmas
\hyperref[lem:2.1.6]{2.1.6} and \hyperref[lem:2.1.7]{2.1.7}, remembering
\cite{ref16} pp.~64--65.

\Lemma{2.1.9}
{\itshape
$\prod_{i=1}^{n}\otimes f_i$ is a continuous function of the
$f_i\in\Hspace_i$. (We use now the metric, strong, topology in all these
spaces.)}

\Proof Put $a=\max_{i=1,\ldots,n}\lVert f_i^0\rVert$, and assume that all
$\lVert f_i-f_i^0\rVert\leq\epsilon$ for an $\epsilon>0$. Then
\begin{align*}
 &\left\lVert\prod_{i=1}^{n}\otimes f_i
        -\prod_{i=1}^{n}\otimes f_i^0\right\rVert\\
 &\quad=\left\lVert\sum_{j=1}^{n}\left(
       \prod_{i=1}^{j}\otimes f_i\otimes
       \prod_{i=j+1}^{n}\otimes f_i^0
       -\prod_{i=1}^{j-1}\otimes f_i\otimes
       \prod_{i=j}^{n}\otimes f_i^0\right)\right\rVert\\
 &\quad=\left\lVert\sum_{j=1}^{n}
       \prod_{i=1}^{j-1}\otimes f_i\otimes(f_j-f_j^0)\otimes
       \prod_{i=j+1}^{n}\otimes f_i^0\right\rVert\\
 &\quad\leq\sum_{j=1}^{n}\left\lVert
       \prod_{i=1}^{j-1}\otimes f_i\otimes(f_j-f_j^0)\otimes
       \prod_{i=j+1}^{n}\otimes f_i^0\right\rVert\\
 &\quad\leq\sum_{j=1}^{n}\prod_{i=1}^{j-1}\lVert f_i\rVert\cdot
       \lVert f_j-f_j^0\rVert\cdot
       \prod_{i=j+1}^{n}\lVert f_i^0\rVert\\
 &\quad\leq\sum_{j=1}^{n}(a+\epsilon)^{j-1}\epsilon a^{n-j}
       =(a+\epsilon)^n-a^n.
\end{align*}

\Lemma{2.1.10}
{\itshape
\lemref{2.1.4} holds for
$\Phi\in\prod_{i=1}^{n}\otimes\Hspace_i$ too.}

\Proof Follows by continuity.

\Lemma{2.1.11}
{\itshape
If $\Phi\in\prod_{i=1}^{n}\otimes\Hspace_i$,
$\overline{\Phi}_1,\overline{\Phi}_2,\ldots\in
\mathop{\prod\nolimits'}_{i=1}^{n}\otimes\Hspace_i$ then
\[
 \Phi\sim(\overline{\Phi}_1,\overline{\Phi}_2,\ldots)
\]
is equivalent to
$\lim_{r\to\infty}\lVert\Phi-\overline{\Phi}_r\rVert=0$.}

\Proof If
$\lim_{r\to\infty}\lVert\Phi-\overline{\Phi}_r\rVert=0$, then
\lemref{2.1.10} gives \defref{2.1.4}, (i), and
$\lVert\overline{\Phi}_r-\overline{\Phi}_s\rVert
\leq\lVert\Phi-\overline{\Phi}_r\rVert+
\lVert\Phi-\overline{\Phi}_s\rVert$ gives \defref{2.1.4}, (ii). Thus
$\Phi\sim(\overline{\Phi}_1,\overline{\Phi}_2,\ldots)$. Conversely, if
$\Phi\sim(\overline{\Phi}_1,\overline{\Phi}_2,\ldots)$, then by definition,
$\lVert\Phi-\overline{\Phi}_s\rVert=
\lim_{r\to\infty}\lVert\overline{\Phi}_r-\overline{\Phi}_s\rVert$. Thus
$\lim_{r,s\to\infty}\lVert\overline{\Phi}_r-
\overline{\Phi}_s\rVert=0$ implies
$\lim_{s\to\infty}\lVert\Phi-\overline{\Phi}_s\rVert=0$.

\Lemma{2.1.12}
{\itshape
$\mathop{\prod\nolimits'}_{i=1}^{n}\otimes\Hspace_i$,
$\prod_{i=1}^{n}\otimes\Hspace_i$ are both separable spaces, that is they
satisfy condition D in \cite[p.~65]{ref16}.}

\Proof
$\mathop{\prod\nolimits'}_{i=1}^{n}\otimes\Hspace_i$ is dense in
$\prod_{i=1}^{n}\otimes\Hspace_i$ by \lemref{2.1.11}, so it suffices to show
that $\mathop{\prod\nolimits'}_{i=1}^{n}\otimes\Hspace_i$ is separable. By
\defref{2.1.3} the separability of the set of all
$\prod_{i=1}^{n}\otimes f_i$, $f_i\in\Hspace_i$, suffices; and this follows,
by \lemref{2.1.9} from the separability of the spaces
$\Hspace_1,\ldots,\Hspace_n$.

\Lemma{2.1.13}
{\itshape
$\prod_{i=1}^{n}\otimes\Hspace_i$ is topologically complete, that is it
satisfies condition E in \cite[p.~65]{ref16}.}

\Proof Assume
$\Phi_1,\Phi_2,\ldots\in\prod_{i=1}^{n}\otimes\Hspace_i$,
$\lim_{r,s\to\infty}\lVert\Phi_r-\Phi_s\rVert=0$. For each $\Phi_r$ choose
a $\overline{\Phi}_r^0\in
\mathop{\prod\nolimits'}_{i=1}^{n}\otimes\Hspace_i$ with
$\lVert\Phi_r-\overline{\Phi}_r^0\rVert\leq1/r$; thus
$\lim_{r\to\infty}\lVert\Phi_r-\overline{\Phi}_r^0\rVert=0$. So
$\lim_{r,s\to\infty}\lVert\overline{\Phi}_r^0-
\overline{\Phi}_s^0\rVert=0$, and for
$\Phi\sim(\overline{\Phi}_1^0,\overline{\Phi}_2^0,\ldots)$,
$\lim_{r\to\infty}\lVert\Phi-\overline{\Phi}_r^0\rVert=0$
(by \lemref{2.1.11}), and
thus $\lim_{r\to\infty}\lVert\Phi-\Phi_r\rVert=0$, too.

This discussion can be summed up as follows:

\MainTheorem{I}
{\itshape
The direct product of $\Hspace_1,\ldots,\Hspace_n$,
$\prod_{i=1}^{n}\otimes\Hspace_i$, arises by the topological completion of
the set $\mathop{\prod\nolimits'}_{i=1}^{n}\otimes\Hspace_i$ of all finite
linear aggregates of expressions $\prod_{i=1}^{n}\otimes f_i^0$,
$f_i^0\in\Hspace_i$. It is a set of functionals
$\Phi=\Phi(f_1,\ldots,f_n)$ in the sense of \defref{2.1.1}, thus
\[
 \prod_{i=1}^{n}\otimes\Hspace_i
 \subset\prod_{i=1}^{n}\solidotimes\Hspace_i
\]
and it is a space satisfying conditions A, B, C, E of
\cite[pp.~64--65]{ref16} (cf. our \secref{1.1}, (b)). Further essential
properties are given in Lemmas \hyperref[lem:2.1.1]{2.1.1},
\hyperref[lem:2.1.4]{2.1.4}, \hyperref[lem:2.1.9]{2.1.9},
\hyperref[lem:2.1.10]{2.1.10}, \hyperref[lem:2.1.11]{2.1.11}.}

If $n=1$, $\prod_{i=1}^{n}\otimes\Hspace_i$ coincides not only in our
notation with $\Hspace_1$, but in reality too.

In this case we may write $f_1^0$ for
$\prod_{i=1}^{n}\otimes f_i^0$, or even better $f^0$. The linear aggregates
of $f^0$'s are $f^0$'s again, so
$\mathop{\prod\nolimits'}_{i=1}^{n}\otimes\Hspace_i$ coincides with
$\Hspace_1$. Thus already
$\mathop{\prod\nolimits'}_{i=1}^{n}\otimes\Hspace_i$ is complete, and
$\prod_{i=1}^{n}\otimes\Hspace_i$ contains no further elements.

Note that in this correspondence of the one-variable functionals $\Phi$ to
the $f^0\in\Hspace$, $\Phi(f)$ becomes simply $(f^0,f)$. A well known
theorem of M. Fr\'echet and F. Riesz implies therefore, that in this case the
$\Phi\in\prod_{i=1}^{n}\otimes\Hspace_i$ are characterised by
$\lVert\Phi(f)\rVert\leq C\cdot\lVert f\rVert$ for a suitable constant $C$.
(Cf. for instance \cite[p.~94]{ref16}.)

This latter point is remarkable, because for $n>1$ the condition
\[
 |\Phi(f_1,\ldots,f_n)|\leq C\cdot\prod_{i=1}^{n}\lVert f_i\rVert
\]
is necessary (cf. \lemref{2.1.10}) but not sufficient for the
$\Phi\in\prod_{i=1}^{n}\otimes\Hspace_i$, as we will show after
\lemref{2.2.2}.

\SectionStart{2.2}{Orthogonal expressions}
It is of interest to discuss the relationship between
$\prod_{i=1}^{n}\otimes\Hspace_i$ and a set of complete normalised orthogonal
sets in each $\Hspace_i$.

\Lemma{2.2.1}
{\itshape
Let a complete normalised orthogonal set
$\varphi_{i,1},\varphi_{i,2},\ldots$ be given in each $\Hspace_i$,
$i=1,\ldots,n$. Then the
$\prod_{i=1}^{n}\otimes\varphi_{i,t_i}$,
$t_1,t_2,\ldots,t_n=1,2,\ldots$ form a complete normalised orthogonal set in
$\prod_{i=1}^{n}\otimes\Hspace_i$.}

\Proof By definition
\[
 \left(\prod_{i=1}^{n}\otimes\varphi_{i,t_i},
       \prod_{i=1}^{n}\otimes\varphi_{i,u_i}\right)
 =\prod_{i=1}^{n}(\varphi_{i,t_i},\varphi_{i,u_i})
 =\prod_{i=1}^{n}\delta_{t_i,u_i}
\]
proving the normalised and orthogonal character. The linear aggregates of the
$\varphi_{i,t}$ are dense in $\Hspace_i$, thus those of the
$\prod_{i=1}^{n}\otimes\varphi_{i,t_i}$ are dense in
$\mathop{\prod\nolimits'}_{i=1}^{n}\otimes\Hspace_i$ (by \lemref{2.1.9}),
and therefore in $\prod_{i=1}^{n}\otimes\Hspace_i$. This proves the
completeness too.

\Lemma{2.2.2}
{\itshape
In order that an $n$-variable functional
$\Phi\in\prod_{i=1}^{n}\solidotimes\Hspace_i$ (cf. \defref{2.1.1}) should belong to
$\prod_{i=1}^{n}\otimes\Hspace_i$
\[
 |\Phi(f_1,\ldots,f_n)|\leq C\prod_{i=1}^{n}\lVert f_i\rVert
\]
(for a suitable constant $C$) is necessary. If this is the case, $\Phi$ is
completely characterised by the numerical values
\[
 \Phi(\varphi_{1,t_1},\ldots,\varphi_{n,t_n})
   =a_{t_1,\ldots,t_n}\quad\text{for all }t_1,\ldots,t_n=1,2,\ldots.
\]
The finiteness of
$\sum_{t_1,\ldots,t_n=1}^{\infty}|a_{t_1,\ldots,t_n}|^2$ is then
characteristic for $\Phi\in\prod_{i=1}^{n}\otimes\Hspace_i$. With this
restriction the $a_{t_1,\ldots,t_n}$ can even be prescribed arbitrarily, and
a $\Phi\in\prod_{i=1}^{n}\otimes\Hspace_i$ with these
$a_{t_1,\ldots,t_n}$ will exist.}

\Proof The necessity of the first condition follows from \lemref{2.1.10}. It
implies, owing to the linearity, that $\Phi(f_1,\ldots,f_n)$ is continuous in
the $f_1,\ldots,f_n$. Thus the values
$\Phi(\varphi_{1,t_1},\ldots,\varphi_{n,t_n})$ determine
$\Phi(f_1,\ldots,f_n)$ if every $f_i$ is a limit of linear aggregates of
$\varphi_{i,1},\varphi_{i,2},\ldots$; that is always.

By \lemref{2.1.10},
\[
 \left(\Phi,\prod_{i=1}^{n}\otimes\varphi_{i,t_i}\right)
 =\Phi(\varphi_{1,t_1},\ldots,\varphi_{n,t_n})=a_{t_1,\ldots,t_n},
\]
therefore
$\sum_{t_1,\ldots,t_n=1}^{\infty}|a_{t_1,\ldots,t_n}|^2$ must be finite, by
\lemref{2.2.1} and \cite[p.~66]{ref16}. Conversely if a system
$a_{t_1,\ldots,t_n}$ with a finite
$\sum_{t_1,\ldots,t_n=1}^{\infty}|a_{t_1,\ldots,t_n}|^2$ is given, then
\[
 \Phi=\sum_{t_1,\ldots,t_n=1}^{\infty}a_{t_1,t_2,\ldots,t_n}
       \prod_{i=1}^{n}\otimes\varphi_{i,t_i}
\]
exists and is $\in\prod_{i=1}^{n}\otimes\Hspace_i$ by \lemref{2.2.1} and
\cite[p.~67]{ref16}, and for the same reason
\[
 \Phi(\varphi_{1,t_1},\ldots,\varphi_{n,t_n})
 =\left(\Phi,\prod_{i=1}^{n}\otimes\varphi_{i,t_i}\right)
 =a_{t_1,\ldots,t_n}.
\]

We now see that in the case $n=2$, the preliminary condition on $\Phi$,
\[
 |\Phi(f_1,f_2)|\leq C\lVert f_1\rVert\cdot\lVert f_2\rVert,
\]
means that the matrix
$(a_{t_1,t_2})_{t_1,t_2=1,2,\ldots}$ associated with $\Phi$ is bounded in the
sense of Hilbert's theory; while
$\Phi\in\prod_{i=1}^{n}\otimes\Hspace_i$ means the finiteness of
\[
 \sum_{t_1,t_2=1}^{\infty}|a_{t_1,t_2}|^2,
\]
that is that the matrix belongs to E. Schmidt's class of matrices. Thus while
the preliminary condition is already characteristic for $n=1$, it is not for
$n=2$. If $\Hspace_1=\Hspace_2=\Hspace$ is a Hilbert space, and $I$ a
conjugation (passing to the complex conjugate, cf. \cite[p.~357]{ref15}) in
$\Hspace$, then $\Phi(f_1,f_2)=(If_1,f_2)$ is an example for the opposite;
because with $\varphi_{1,t}=\varphi_{2,t}$ we obtain
$a_{t_1,t_2}=\delta_{t_1,t_2}$.

\Lemma{2.2.3}
{\itshape
If in the notation of \lemref{2.2.2}
$\Phi,\Psi\in\prod_{i=1}^{n}\otimes\Hspace_i$ corresponds to
$a_{t_1,\ldots,t_n}$ resp. $b_{t_1,\ldots,t_n}$, then
\[
 (\Phi,\Psi)=\sum_{t_1,\ldots,t_n=1}^{\infty}
 a_{t_1,\ldots,t_n}\overline{b_{t_1,\ldots,t_n}},
\]
the series being absolutely convergent.}

\Proof The last formula of the proof of \lemref{2.2.2}, together with
\lemref{2.2.1} and \cite[p.~66]{ref16}, give this.

Lemmas \hyperref[lem:2.2.2]{2.2.2} and \hyperref[lem:2.2.3]{2.2.3} show that
the dimension of
$\prod_{i=1}^{n}\otimes\Hspace_i$ is the product of the dimensions of all
$\Hspace_i$, that is: 0 if at least one of them is 0; $\infty$ if none of them
is 0 but at least one $\infty$; the finite product if all are finite. From now
on we assume that all $\Hspace_i$ have dimensions $\ne0$, that is that
$\Hspace_i\ne(0)$.

\SectionStart{2.3}{Effect on rings. Theorem II}
We define now certain operators in $\prod_{i=1}^{n}\otimes\Hspace_i$. We
will denote the ring of all bounded, everywhere defined, linear and closed
operators in $\Hspace_i$ by $\B_i$, $i=1,\ldots,n$ the same in
$\prod_{i=1}^{n}\otimes\Hspace_i$ by $\B\otimes$. (Cf.
\cite[p.~370]{ref17}.)

\Lemma{2.3.1}
{\itshape
If an operator $A_i\in\B_i$ and a
$\Phi\in\prod_{i=1}^{n}\otimes\Hspace_i$ is given, then a unique
$\Phi^*\in\prod_{i=1}^{n}\otimes\Hspace_i$ with
\[
 \Phi^*(f_1,\ldots,f_i,\ldots,f_n)
 =\Phi(f_1,\ldots,A_i^*f_i,\ldots,f_n)
\]
exists. ($A_i^*$ is the adjoint of $A_i$). If
$\lVert A_if_i\rVert\leq D\lVert f_i\rVert$ for all $f_i\in\Hspace_i$, then
$\lVert\Phi^*\rVert\leq D\lVert\Phi\rVert$.}

\Proof We may assume $i=n$. The uniqueness of $\Phi^*$ and its existence in
$\prod_{i=1}^{n}\solidotimes\Hspace_i$ are obvious.

A $C$ with
$|\Phi(f_1,\ldots,f_n)|^2\leq C^2\prod_{i=1}^{n}\lVert f_i\rVert^2$ exists,
therefore there exists for any given $f_1,\ldots,f_{n-1}$ an
$f^0=f^0(f_1,\ldots,f_{n-1})\in\Hspace_n$ with
$\Phi(f_1,\ldots,f_n)=(f^0,f_n)$, and
$\lVert f^0\rVert\leq C\prod_{i=1}^{n-1}\lVert f_i\rVert$ (cf.
\cite[p.~94]{ref16}). Now
\[
 \Phi^*(f_1,\ldots,f_n)
 =\Phi(f_1,\ldots,A_n^*f_n)
 =(f^0,A_n^*f_n)=(A_nf^0,f_n).
\]
Furthermore $A_n$ is bounded, therefore a $D$ with
$\lVert A_nf\rVert\leq D\lVert f\rVert$ exists. So we have:
\begin{align*}
 |\Phi^*(f_1,\ldots,f_n)|
 &=|(A_nf^0,f_n)|
 \leq\lVert A_nf^0\rVert\cdot\lVert f_n\rVert\\
 &\leq D\lVert f^0\rVert\cdot\lVert f_n\rVert
 \leq CD\prod_{i=1}^{n}\lVert f_i\rVert,
\end{align*}
and
\begin{align*}
 &\sum_{t_1,\ldots,t_n=1}^{\infty}
   |\Phi^*(\varphi_{1,t_1},\ldots,\varphi_{n,t_n})|^2\\
 &\quad=\sum_{t_1,\ldots,t_n=1}^{\infty}
   |(A_nf^0(\varphi_{1,t_1},\ldots,\varphi_{n-1,t_{n-1}}),
      \varphi_{n,t_n})|^2\\
 &\quad=\sum_{t_1,\ldots,t_{n-1}=1}^{\infty}
   \lVert A_nf^0(\varphi_{1,t_1},\ldots,
                    \varphi_{n-1,t_{n-1}})\rVert^2\\
 &\quad\leq D^2\sum_{t_1,\ldots,t_{n-1}=1}^{\infty}
   \lVert f^0(\varphi_{1,t_1},\ldots,
               \varphi_{n-1,t_{n-1}})\rVert^2\\
 &\quad\leq D^2\sum_{t_1,\ldots,t_n=1}^{\infty}
   |(f^0(\varphi_{1,t_1},\ldots,\varphi_{n-1,t_{n-1}}),
      \varphi_{n,t_n})|^2\\
 &\quad=D^2\sum_{t_1,\ldots,t_n=1}^{\infty}
   |\Phi(\varphi_{1,t_1},\ldots,\varphi_{n,t_n})|^2.
\end{align*}
Thus $\Phi\in\prod_{i=1}^{n}\otimes\Hspace_i$ implies
$\Phi^*\in\prod_{i=1}^{n}\otimes\Hspace_i$ (by \lemref{2.2.2}), and we
have $\lVert\Phi^*\rVert\leq D\lVert\Phi\rVert$ (by \lemref{2.2.3}).

\Lemma{2.3.2}
{\itshape
If we define an operator $A_i^{(i)}$ by $A_i^{(i)}\Phi=\Phi^*$ in the sense
of \lemref{2.3.1}, then $A_i^{(i)}\in\B\otimes$.}

\Proof $A_i^{(i)}$ is everywhere defined and linear by \lemref{2.3.1}; and
if we choose $D$ with $\lVert A_if\rVert\leq D\lVert f\rVert$, then
$\lVert A_i^{(i)}\Phi\rVert\leq D\lVert\Phi\rVert$ by the same Lemma. Thus
$A_i^{(i)}$ is bounded and therefore it is closed, too.

\Lemma{2.3.3}
{\itshape
\[
 A_i^{(i)}\prod_{j=1}^{n}\otimes f_j
 =\prod_{j=1}^{i-1}\otimes f_j\otimes A_if_i\otimes
  \prod_{j=i+1}^{n}\otimes f_j.
\]}

\Proof Follows immediately from the definition.

\Lemma{2.3.4}
{\itshape
The correspondence $A_i\sim A_i^{(i)}$ is a (one to one) isomorphism for the
operations $\alpha A_i,A_i^*,A_i+B_i$, and $A_iB_i$.}

\Proof This is obvious for $\alpha A_i$ and $A_i+B_i$; for $A_i^*$ and
$A_iB_i$ it follows by applying these operators to the
$\prod_{j=1}^{n}\otimes f_j$ as the limits of linear aggregates of
$\prod_{j=1}^{n}\otimes f_j$ cover all
$\prod_{j=1}^{n}\otimes\Hspace_j$.

That $A_i=B_i$ implies $A_i^{(i)}=B_i^{(i)}$, is clear; the converse is again
seen by considering the $\prod_{j=1}^{n}\otimes f_j$.

\Definition{2.3.1}
Call the set of all $A_i^{(i)}$ if $A_i$ runs over all $\B_i$,
$\B_i^{(i)}$.

\Lemma{2.3.5}
{\itshape
\[
 \B_i^{(i)}=(\B_j^{(j)}; j\ne i)';
\]
that is: $\B_i^{(i)}$ is the set of all operators in $\B\otimes$ which
commute with all elements of all $\B_j^{(j)}$'s $j\ne i$.}

\Proof It is obvious that for $j\ne i$,
$A_i^{(i)}A_j^{(j)}\Phi=A_j^{(j)}A_i^{(i)}\Phi$ if
$\Phi=\prod_{k=1}^{n}\otimes f_k$. Thus this holds for all
$\Phi\in\prod_{k=1}^{n}\otimes\Hspace_k$,
$A_j^{(j)}$ and $A_i^{(i)}$ commute. In other words:
$\B_i^{(i)}\subset(\B_j^{(j)}; j\ne i)'$.

Assume now $A^0\in(\B_j^{(j)}; j\ne i)'$. Then $A^0$ commutes with every
$A_j^{(j)}$, $j\ne i$. Introduce again the complete normalised orthogonal
systems $\varphi_{j,1},\varphi_{j,2},\ldots$ in $\Hspace_j$,
$j=1,\ldots,n$. The projection $P_{[\varphi_{j,t}]}$ (cf.
\cite[p.~74]{ref16}) belongs to $\B_j$, it transforms $\varphi_{j,u}$ into
$\delta_{u,t}\varphi_{j,t}$. Thus $P_{[\varphi_{j,t}]}^{(j)}\in\B_j^{(j)}$
transforms $\prod_{k=1}^{n}\otimes\varphi_{k,t_k}$ into
$\delta_{t_j,t}\prod_{k=1}^{n}\otimes\varphi_{k,t_k}$; and as this is a
complete normalised orthogonal system, this means that
$P_{[\varphi_{j,t}]}^{(j)}$ is the projection of
\[
 \left[\prod_{k=1}^{n}\otimes\varphi_{k,t_k},\ t_j=t\right]
\]
(the smallest linear and closed set containing all
$\prod_{k=1}^{n}\otimes\varphi_{k,t_k}$ with $t_j=t$). Now $A^0$ commutes
with $P_{[\varphi_{j,t}]}^{(j)}$. Therefore
$P_{[\varphi_{j,t}]}^{(j)}\Phi=\Phi$ implies
$P_{[\varphi_{j,t}]}^{(j)}A^0\Phi=A^0\Phi$. The former is the case for
\[
 \Phi=\prod_{k=1}^{i-1}\otimes\varphi_{k,t_k}\otimes f_i\otimes
      \prod_{k=i+1}^{n}\otimes\varphi_{k,t_k}
\]
and $t=t_j$. Thus $A^0\Phi$ when written as a
$\prod_{k=1}^{n}\otimes\varphi_{k,u_k}$ series,
$u_1,\ldots,u_n=1,2,\ldots$ has terms with $u_j=t_j$ only. As this holds for
every $j\ne i$ we have
\[
 A^0\Phi=\prod_{k=1}^{i-1}\otimes\varphi_{k,t_k}\otimes f_i^*\otimes
           \prod_{k=i+1}^{n}\otimes\varphi_{k,t_k}.
\]
Thus $f_i^*$ is determined uniquely by $f_i$ and the $t_k$, $k\ne i$.
Consideration of the operator
$P_{[\varphi_{j,t}+\varphi_{j,s}]}^{(j)}$, $j\ne i$, $t\ne s$, shows however,
that $t_j=t$ and $t_j=s$ give the same. So the value of $t_j$ does not
matter. Thus $f_i^*=A_if_i$ for a fixed operator $A_i$ in $\Hspace_i$. It is
clear that $A_i$ is everywhere defined and linear; and
$\lVert A^0\Phi\rVert\leq D\lVert\Phi\rVert$ implies
$\lVert A_if_i\rVert\leq D\lVert f_i\rVert$ so that $A_i$ is bounded and
therefore closed. So $A_i\in\B_i$. The definition of $A_i$ gives
$A^0\Phi=A_i^{(i)}\Phi$ for all
$\Phi=\prod_{i=1}^{n}\otimes\varphi_{i,t_i}$, thus for all $\Phi$. So we
have $A^0=A_i^{(i)}\in\B_i^{(i)}$.

This proves $\B_i^{(i)}\supset(\B_j^{(j)}; j\ne i)'$, completing the proof
of $\B_i^{(i)}=(\B_j^{(j)}; j\ne i)'$.

\Lemma{2.3.6}
{\itshape
$\B_i^{(i)}$ is a ring and contains the operator 1.}

\Proof This follows from \cite[p.~397]{ref18}, because $\B_i^{(i)}$ has the
form $S'$ (by \lemref{2.3.5}).

\Lemma{2.3.7}
{\itshape
Let $S_i$ be a set $\subset\B_i$, and $S_i^{(i)}$ the set of all
$A_i^{(i)}$, $A_i\in S_i$, $S_i^{(i)}\subset\B\otimes$. Then $S_i$ is a
ring if and only if $S_i^{(i)}$ is one.}

\Proof Define a ring by its characteristics established in
\cite[\S3]{ref22}. The only notions entering there are the operations
$\alpha A,A^*,A+B,AB$, and the strongest topology, as defined in
\cite[\S2]{ref22}. Now all these are invariant under
$A_i\sim A_i^{(i)}$: the algebraic operations by \lemref{2.3.4}, and the
strongest topology by \cite[\S4]{ref22}. Besides closure in $\B_i^{(i)}$ and
in $\B\otimes$ is the same thing, because $\B_i^{(i)}$ is a ring by
\lemref{2.3.6}, thus weakly closed in $\B\otimes$, and a fortiori in the
strongest sense. This completes the proof.

\Lemma{2.3.8}
{\itshape
\[
 R(\B_1^{(1)},\ldots,\B_n^{(n)})=\B\otimes.
\]}

\Proof We could prove this by a direct construction, but the following
indirect proof is quicker: The same considerations as the ones we made in the
proof of \lemref{2.3.5} (but now without the exceptional $i$) show that
\[
 (\B_j^{(j)}; j=1,\ldots,n)'=(\alpha1).
\]
So
\[
 R(\B_1^{(1)},\ldots,\B_n^{(n)})'
 =(\B_1^{(1)},\ldots,\B_n^{(n)})'=(\alpha1),
\]
and by \cite[p.~397]{ref18},
\[
 R(\B_1^{(1)},\ldots,\B_n^{(n)})=(\alpha1)'=\B\otimes.
\]

We will frequently use the above quoted Theorem 8 from \cite[p.~397]{ref18}:
That if $\M$ is a ring which contains 1, then $\M=\M''$. No particular
references will therefore be made to it.

The results of the foregone discussion can be summed up as follows:

\MainTheorem{II}
{\itshape
The sets $\B_i^{(i)}$, $i=1,\ldots,n$, defined in \defref{2.3.1}, are rings
which contain 1; $\B_i^{(i)}$ being isomorphic to $\B_i$. Any two
$\B_i^{(i)}$ commute, and
\[
 R(\B_1^{(1)},\ldots,\B_n^{(n)})=\B\otimes.
\]
Further essential properties are given in Lemmas
\hyperref[lem:2.3.5]{2.3.5}, and \hyperref[lem:2.3.7]{2.3.7}.}

\SectionStart{2.4}{\texorpdfstring{$n\otimes\Hspace=
\Hspace\oplus\Hspace\oplus\cdots\oplus\Hspace$ ($n$-times).
Operator-matrices}{n tensor H equals n copies of H. Operator-matrices}}
The asymmetric aspect of the case $n=2$ is of importance.

Form $\prod_{i=1}^{2}\otimes\Hspace_i=\Hspace_1\otimes\Hspace_2$, and
introduce a (finite or infinite) complete normalised orthogonal system in
$\Hspace_1$: $\varphi_1,\varphi_2,\ldots$. Now define

\Definition{2.4.1}
If $\Phi\in\Hspace_1\otimes\Hspace_2$, form the
$f^0(\varphi_t)\in\Hspace_2$, $t=1,2,\ldots$ as in the proof of
\lemref{2.3.1}. Denote them by $f^{(t)}$, $t=1,2,\ldots$ resp., and write
\[
 \Phi\sim\langle f^{(1)},f^{(2)},\ldots\rangle
       =\langle f^{(t)}\rangle_{t=1,2,\ldots}.
\]

\Lemma{2.4.1}
{\itshape
The $f^{(t)}$, $t=1,2,\ldots$ determine
$\Phi\sim\langle f^{(1)},f^{(2)},\ldots\rangle$ uniquely. If dimension
$\Hspace_1$ is finite, the $f^{(t)}$ can be prescribed arbitrarily; if it is
infinite, then the only restriction is that
$\sum_{t=1}^{\infty}\lVert f^{(t)}\rVert^2$ must be finite.

If $\Phi\sim\langle f^{(1)},f^{(2)},\ldots\rangle$,
$\Psi\sim\langle g^{(1)},g^{(2)},\ldots\rangle$, then
$(\Phi,\Psi)=\sum_t(f^{(t)},g^{(t)})$. The sum is absolutely convergent, if
infinite at all.}

\Proof Let $\psi_1,\psi_2,\ldots$ be a complete normalised orthogonal system
in $\Hspace_2$; form, as in \lemref{2.2.2},
\[
 \Phi(\varphi_{t_1},\psi_{t_2})=(f^{(t_1)},\psi_{t_2})=a_{t_1,t_2}.
\]
$\Phi$ is uniquely determined by the $a_{t_1,t_2}$ and so by the
$f^{(t_1)}$. For given $a_{t_1,t_2}$'s $\Phi$ exists if
$\sum_{t_1,t_2}|a_{t_1,t_2}|^2$ is finite; but if the $a_{t_1,t_2}$ are
defined by the $f^{(t_1)}$: $a_{t_1,t_2}=(f^{(t_1)},\psi_{t_2})$ then
\[
 \sum_{t_2}|a_{t_1,t_2}|^2=\lVert f^{(t_1)}\rVert^2,
\]
therefore the condition is, that
$\sum_{t_1}\lVert f^{(t_1)}\rVert^2$ is finite.

If the dimension of $\Hspace_1$ is finite, the number of the $t_1$'s is, and
the condition is void. If the dimension of $\Hspace_1$ is infinite, we obtain
the condition: $\sum_{t_1=1}^{\infty}\lVert f^{(t_1)}\rVert^2$ is finite.
Finally \lemref{2.2.3} gives:
\[
 (\Phi,\Psi)
 =\sum_{t_1,t_2}(f^{(t_1)},\psi_{t_2})
       \overline{(g^{(t_1)},\psi_{t_2})}
 =\sum_{t_1}(f^{(t_1)},g^{(t_1)}),
\]
all sums being absolutely convergent.

The operators in $\Hspace_1\otimes\Hspace_2$ can be expressed in a simple
normal form, in which only operators from $\Hspace_2$ occur:

\Definition{2.4.2}
If $A^0$ is an operator in $\B\otimes$ (for
$\Hspace_1\otimes\Hspace_2$), then form
\[
 A^0\langle0,\ldots,0,f,0,\ldots\rangle
 =\langle f_1^*,f_2^*,\ldots\rangle
 \quad\text{(the $f$ stands in place no. $t$)}.
\]
Every $f_s^*$ defines an operator $A_{t,s}$ by
$f_s^*=A_{t,s}f_t$. Write
\[
 A^0\sim\langle A_{t,s}\rangle_{t,s=1,2,\ldots}.
\]

\Lemma{2.4.2}
{\itshape
All $A_{t,s}$ belong to $\B_2$ (for $\Hspace_2$), they determine $A^0$
uniquely. If dimension of $\Hspace_1$ is finite, the $A_{t,s}$ can be
prescribed arbitrarily; if it is infinite, they cannot.}

\Proof $A_{t,s}$ is clearly everywhere defined and linear. $A^0$ is bounded,
so a $C$ with $\lVert A^0\Phi\rVert\leq C\lVert\Phi\rVert$ exists. Then
\begin{align*}
 \lVert A_{t,s}f\rVert^2
 &\leq\sum_{s'}\lVert A_{t,s'}f\rVert^2
  =\lVert A^0\langle0,0,\ldots,0,f,0,\ldots\rangle\rVert^2\\
 &\leq C^2\lVert\langle0,\ldots,0,f,0,\ldots\rangle\rVert^2
  =C^2\lVert f\rVert^2,
 \qquad \lVert A_{t,s}f\rVert\leq C\lVert f\rVert.
\end{align*}
Thus $A_{t,s}$ is bounded and closed too. It is clear from
\hyperref[def:2.4.2]{Definition 2.4.3},
that the $A_{t,s}$ determine $A^0$ uniquely for all
$\Phi=\langle0,\ldots,0,f,0,\ldots\rangle$, and therefore for all $\Phi$.

If dimension of $\Hspace_1$ is finite, say $p$, then the formula in
\defref{2.4.2} certainly defines an operator $A^0$, which is everywhere
defined and linear. Each $A_{t,s}$ is bounded, say
$\lVert A_{t,s}f\rVert\leq C_{t,s}\lVert f\rVert$. Then
\[
 A^0\langle f_1,\ldots,f_p\rangle
 =\left\langle\sum_{t=1}^{p}A_{t,1}f_t,\ldots,
                 \sum_{t=1}^{p}A_{t,p}f_t\right\rangle.
\]
\begin{align*}
 \lVert A^0\Phi\rVert^2
 &=\lVert A^0\langle f_1,\ldots,f_p\rangle\rVert^2
  =\sum_{s=1}^{p}\left\lVert\sum_{t=1}^{p}A_{t,s}f_t\right\rVert^2\\
 &\leq\sum_{s=1}^{p}p\sum_{t=1}^{p}\lVert A_{t,s}f_t\rVert^2\\
 &\leq\sum_{s=1}^{p}p\sum_{t=1}^{p}C_{t,s}^2\lVert f_t\rVert^2
  \leq C^2\sum_{t=1}^{p}\lVert f_t\rVert^2=C^2\lVert\Phi\rVert^2,
\end{align*}
where
$C^2=p\max_{t=1,\ldots,p}\sum_{s=1}^{p}C_{t,s}^2$, and so
$\lVert A^0\Phi\rVert\leq C\lVert\Phi\rVert$. Thus $A^0$ is bounded and
closed too.

It would be easy to formulate the restrictions arising if the dimension of
$\Hspace_1$ is infinite, but they are of a very implicit type. In fact, even
if $\Hspace_2$ is one dimensional, they express that the numerical matrix
$\langle a_{t,s}\rangle_{t,s=1,2,\ldots}$ is bounded in the sense of Hilbert's
theory.

The rules of computation for this representation
$A^0\sim\langle A_{t,s}\rangle_{t,s=1,2,\ldots}$ are very similar to those of
common matrix-algebra.

\Lemma{2.4.3}
{\itshape
If $A^0,B^0,C^0\in\B\otimes$, and
$A^0\sim\langle A_{t,s}\rangle_{t,s=1,2,\ldots}$,
$B^0\sim\langle B_{t,s}\rangle_{t,s=1,2,\ldots}$,
$C^0\sim\langle C_{t,s}\rangle_{t,s=1,2,\ldots}$ then the following rules of
computation hold:
\begin{enumerate}[label=(\roman*)]
\item If $C^0=\alpha A^0$, then $C_{t,s}=\alpha A_{t,s}$.
\item If $C^0=(A^0)^*$, then $C_{t,s}=A_{s,t}^*$.
\item If $C^0=A^0+B^0$, then $C_{t,s}=A_{t,s}+B_{t,s}$.
\item If $C^0=A^0B^0$, then
  $C_{t,s}=\sum_{n=1}^{\infty}A_{n,s}B_{t,n}$. The
  $\sum_{n=1}^{\infty}$ converges in the sense of the strong topology of
  $\B_2$, irrespectively of the order in which the $n=1,2,\ldots$ are gone
  through.
\end{enumerate}}

\Proof (i), (iii) are obvious. Denoting
$\langle0,\ldots,0,f,0,\ldots\rangle$ (the $f$ stands in place no. $t$) by
$f^{(t)}$, we see: $(A^0f^{(t)},g^{(s)})=(A_{t,s}f,g)$. Therefore in case
(ii) we have on one hand $(C^0f^{(t)},g^{(s)})=(C_{t,s}f,g)$ and on the
other
\[
 (C^0f^{(t)},g^{(s)})
 =\overline{(A^0g^{(s)},f^{(t)})}
 =\overline{(A_{s,t}g,f)}=(A_{s,t}^*f,g)
\]
thus proving $C_{t,s}=A_{s,t}^*$.

Consider finally case (iv). We have
$C^0f^{(t)}=\langle C_{t,s}f\rangle_{s=1,2,\ldots}$.

On the other hand
\[
 C^0f^{(t)}=A^0B^0f^{(t)}
 =A^0\langle B_{t,n}f\rangle_{n=1,2,\ldots}
 =A^0\sum_{n=1}^{\infty}(B_{t,n}f)^{(n)}
\]
where the $\sum_{n=1}^{\infty}$ is strongly convergent, irrespectively of the
order in which the $n=1,2,\ldots$ are gone through. So we have further (as
$A^0$ is continuous):
\begin{align*}
 C^0f^{(t)}
 &=\sum_{n=1}^{\infty}A^0(B_{t,n}f)^{(n)}
  =\sum_{n=1}^{\infty}
       \langle A_{n,s}B_{t,n}f\rangle_{s=1,2,\ldots}\\
 &=\left\langle\sum_{n=1}^{\infty}A_{n,s}B_{t,n}f
   \right\rangle_{s=1,2,\ldots}
\end{align*}
with the same kind of convergence. This proves
$C_{t,s}=\sum_{n=1}^{\infty}A_{n,s}B_{t,n}$. We now investigate the
correspondences $A_i\sim A_i^{(i)}$.

\Lemma{2.4.4}
{\itshape
If $A\in\B_2$, then $A^{(2)}\in\B\otimes$ is
\[
 A^{(2)}\sim\langle\delta_{t,s}A\rangle_{t,s=1,2,\ldots}.
\]}

\Proof Obviously
$\varphi_t\otimes f\sim\langle0,\ldots,0,f,0,\ldots\rangle$ ($f$ in place
no. $t$), therefore by \lemref{2.3.3}
\[
 A^{(2)}\langle0,\ldots,0,f,0,\ldots\rangle
 =\langle0,\ldots,0,Af,0,\ldots\rangle.
\]
Thus $A_{t,s}=0$ if $s\ne t$ and $=A$ if $s=t$, in other words
$A_{t,s}=\delta_{t,s}A$.

\Lemma{2.4.5}
{\itshape
If $S$ is a set $\subset\B_2$, then the set $S^{(2)}$ of all $A^{(2)}$,
$A^{(2)}\in S$, consists of all
$A^0\sim\langle\delta_{t,s}A\rangle_{t,s=1,2,\ldots}$, $A\in S$ (every
$A\in S$ generates an $A^0$!); and the set $S^{(2)'}$ consists of all
$A^0\sim\langle A_{t,s}\rangle_{t,s=1,2,\ldots}$, $A_{t,s}\in S'$.}

\Proof The first statement follows immediately from \lemref{2.4.4}. As to
the second, observe that
\[
 \langle\delta_{t,s}A\rangle\langle A_{t,s}\rangle
   =\langle AA_{t,s}\rangle;\qquad
 \langle A_{t,s}\rangle\langle\delta_{t,s}A\rangle
   =\langle A_{t,s}A\rangle;
\]
thus $\langle A_{t,s}\rangle$ commutes with
$A^{(2)}\sim\langle\delta_{t,s}A\rangle$ if and only if every $A_{t,s}$
commutes with $A$. Furthermore $A^{(2)*}=A^{*(2)}$ (by \lemref{2.3.4},
$\langle(\delta_{s,t}A)^*\rangle=\langle\delta_{t,s}A^*\rangle$). These
facts together establish the equivalence of
$\langle A_{t,s}\rangle\in S^{(2)'}$ and of all $A_{t,s}\in S'$.

\Lemma{2.4.6}
{\itshape
$\B_2^{(2)}$ is the set of all
$A^0\sim\langle\delta_{t,s}A\rangle_{t,s=1,2,\ldots}$,
$A\in\B_2$ (every $A\in\B_2$ generates an $A^0$!);
$\B_1^{(1)}$ is the set of all
$A^0\sim\langle\alpha_{t,s}1\rangle_{t,s=1,2,\ldots}$.}

\Proof Follows from \lemref{2.4.5} by putting $S=\B_2$, considering that
$\B_2^{(2)'}=\B_1^{(1)}$ (by \lemref{2.3.5}).

In those applications in which the aspect of this \hyperref[sec:2.4]{\S}\
will be used, the
important object will always be $\Hspace_2$, while $\Hspace_1$ will only play
a r\^ole in so far as its dimension may be prescribed. If dimension of
$\Hspace_1$ is $N=1,2,\ldots,\infty$, and if we write $\Hspace$ for
$\Hspace_2$, we will use the abbreviated notation $N\otimes\Hspace$ for
$\Hspace_1\otimes\Hspace_2$.
% --- End of parts/02a-chapters-i-ii.tex ---
\ChapterHeading{3}{III}{\texorpdfstring{Factorisation of $\B$}
{Factorisation of B}}

\SectionStart{3.1}{Factorisations, factors. Problem 1}
We now formulate the two notions which play the central r\^ole in all our
discussions (chiefly the second one). We consider a space $\Hspace$ and the ring
$\B$ of everywhere defined, bounded, linear and closed operators in $\Hspace$.

\Definition{3.1.1}
A system of $n(=1,2,\ldots)$ subrings $\M_1,\ldots,\M_n\neq(0)$ of $\B$ is
called a \emph{factorisation}, if
\begin{enumerate}[label=(\roman*)]
\item $\M_i\subset\M_j'$, that is every element of $\M_i$ commutes with every
  element of $\M_j$, for $i\neq j$,
\item $R(\M_1,\ldots,\M_n)=\B$.
\end{enumerate}

\Definition{3.1.2}
A subring $\M$ of $\B$ is called a \emph{factor}, if
\[
  \M\cdot\M'=(\alpha1),
\]
that is, if the center of $\M$ (the set of all those elements of $\M$ which
commute with every element of $\M$) consists of the numerical multiples of $1$
alone.

The connection between these two notions is given by the following Lemmas:

\Lemma{3.1.1}
\emph{If $\M_1,\ldots,\M_n$ is a factorisation, then every $\M_i$ is a factor.}

\Proof We have
\begin{align*}
 \M_i\cdot\M_i'
 &\subset \left(\prod_{\substack{j=1\\j\neq i}}^n\M_j'\right)\cdot\M_i'
   =\prod_{j=1}^n\M_j'=(\M_1,\ldots,\M_n)'\\
 &=\bigl(R(\M_1,\ldots,\M_n)\bigr)'=\B'=(\alpha1).
\end{align*}
By \cite[p.~393]{ref18}, $\M_i\cdot\M_i'$ contains a projection $E_0$; thus
$E_0=\alpha1$, that is $E_0=0$ or $1$. $E_0=0$ would imply
$\M_i=(0)$ (cf. eod.: if $A\in\M_i$, then $AE_0=A$), thus $E_0=1$.
Therefore $\M_i\cdot\M_i'=(\alpha1)$.

\Lemma{3.1.2}
\emph{$\M$ is a factor if and only if $\M,\M'$ is a factorisation.}

\Proof If $\M,\M'$ is a factorisation then $\M$ is a factor by
\lemref{3.1.1}. If conversely $\M$ is a factor, then
$1\in\M\cdot\M'$, and so $\M,\M'\neq(0)$; further
$1\in\M\subset R(\M,\M')$, and so $\M=\M''$ and
$R(\M,\M')=(R(\M,\M'))''$. Now
\[
 (R(\M,\M'))'=(\M,\M')'=\M'\cdot\M''=\M'\cdot\M=(\alpha1).
\]
$R(\M,\M')=(\alpha1)'=\B$. Finally $\M$ and $\M'$ commute. Thus
$\M,\M'$ is a factorisation.

These Lemmas prove that if $\M$ is a factor, then $\M'$ is one too. If
$\M_1$ is a factor then it is a ring and $1\in\M_1$; thus
$\M_1=\M_1''$; therefore $\M_1'=\M_2$ implies $\M_2'=\M_1$.
Therefore we can define:

\Definition{3.1.3}
If $\M_1,\M_2$ is a factorisation, then $\M_1$ and $\M_2$ are called
\emph{coupled factors} if $\M_1'=\M_2$ or (which is equivalent to it)
$\M_2'=\M_1$. The following problem arises immediately:

\Problem{1}
\emph{Is every factorisation $\M_1,\M_2$ made up of coupled factors?}

We will see that the answer is affirmative in a certain special case
(cf. \lemref{3.2.4}) but negative in general (cf. \secref{13.4}). This
problem is of some interest in the case $n>2$ too, because if
$\M_1,\ldots,\M_n$ is a factorisation, then
\[
 R(\M_1,\ldots,\M_k),\qquad R(\M_{k+1},\ldots,\M_n),
 \qquad 1\leq k\leq n-1,
\]
is one too.

\SectionStart{3.2}{Elementary properties. Problem 2}
A simple example of a factorisation is this one: If
\[
 \Hspace=\prod_{i=1}^n\otimes\Hspace_i
 =\Hspace_1\otimes\cdots\otimes\Hspace_n,
\]
then $\B_1^{(1)},\ldots,\B_n^{(n)}$ form a factorisation by \thmref{II}.
Thus every $\B_i^{(i)}$ is a factor. This suggests the following definitions:

\Definition{3.2.1}
A factorisation $\M_1,\ldots,\M_n$ is a \emph{direct factorisation}, if
$\Hspace$ is isomorphic to a direct product
$\Hspace_1\otimes\cdots\otimes\Hspace_n$, so that
$\M_1,\ldots,\M_n$ becomes $\B_1^{(1)},\ldots,\B_n^{(n)}$ resp.

\Definition{3.2.2}
A factor $\M$ is a \emph{direct factor}, if $\Hspace$ is isomorphic to a
direct product $\Hspace_1\otimes\cdots\otimes\Hspace_n$, so that $\M$
becomes a $\B_i^{(i)}$, $i=1,\ldots,n$.

The following remarks lie near:

\Lemma{3.2.1}
\emph{We can restrict ourselves in \defref{3.2.1} to the case where $n=2$,
$i=1$ (or just as well $i=2$).}

\Proof \hyperref[lem:2.2.2]{Lemmas 2.2.2} and
\hyperref[lem:2.2.3]{2.2.3} show, that
$\prod_{j=1}^n\otimes\Hspace_j$ is isomorphic to
\[
 \Hspace_i\otimes
 \left(\prod_{\substack{j=1\\j\neq i}}^n\otimes\Hspace_j\right)
 \quad\hbox{and to}\quad
 \left(\prod_{\substack{j=1\\j\neq i}}^n\otimes\Hspace_j\right)
 \otimes\Hspace_i.
\]

\Lemma{3.2.2}
\emph{If $\M_1,\ldots,\M_n$ is a direct factorisation, then every $\M_i$ is
a direct factor.}

\Proof Obvious.

\Lemma{3.2.3}
\emph{$\M$ is a direct factor if and only if $\M,\M'$ is a direct
factorisation.}

\Proof If $\M,\M'$ is a direct factorisation, then $\M$ is a direct factor
by \lemref{3.2.2}. If conversely $\M$ is a direct factor, then we may assume
by \lemref{3.2.1} that
$\Hspace=\Hspace_1\otimes\Hspace_2$, $\M=\B_1^{(1)}$. Then by
\lemref{2.3.5} $\M'=(\B_1^{(1)})'=\B_2^{(2)}$, proving the statement.

\hyperref[lem:3.2.2]{Lemmas 3.2.2} and \hyperref[lem:3.2.3]{3.2.3}
are perfect analogs of \hyperref[lem:3.1.1]{Lemmas 3.1.1} and
\hyperref[lem:3.1.2]{3.1.2}. We now solve \probref{1} for direct
factors.

\Lemma{3.2.4}
\emph{If $\M_1,\M_2$ is a factorisation, and either $\M_1$ or $\M_2$ is a
direct factor, then both are it, they are coupled, and $\M_1,\M_2$ is a direct
factorisation.}

\Proof Owing to the symmetry we may assume that $\M_1$ is a direct factor,
and thus that
$\Hspace=\Hspace_1\otimes\Hspace_2$, $\M_1=\B_1^{(1)}$. Now
$\M_2\subset\M_1'=(\B_1^{(1)})'=\B_2^{(2)}$. Besides
\[
 \M_2'\cdot\B_2^{(2)}=\M_2'\cdot\M_1'
 =(\M_1,\M_2)'=(R(\M_1,\M_2))'={\B\otimes}'=(\alpha1).
\]
Now consider the mapping $A_2^{(2)}\leftrightarrow A_2$ of $\B_2^{(2)}$ on
$\B_2$ (belonging to $\Hspace_2$). It carries the ring
$\M_2\subset\B_2^{(2)}$ into a ring $\N\subset\B_2$ by \lemref{2.3.4}
and the ring $\M_2'\cdot\B_2^{(2)}$ into $\N'$ by \lemref{2.3.7}. So
$\N'=(\alpha1)$, and therefore $\N=(\alpha1)'=\B_2$,
$\M_2=\B_2^{(2)}$. Thus $\M_2=\M_1'$, $\M_1=\B_1^{(1)}$,
$\M_2=\B_2^{(2)}$, proving all statements.

We have thus a general picture of the formal properties of factorisations and
factors, general, direct, and coupled. A further point which is of interest in this
connection will be elucidated by \lemref{5.4.1}.

The decisive question is now this:

\Problem{2}
\emph{Is every factorisation direct? Is every factor direct?}

We will see in \secref{8.4} and \secref{13.2}, that the answer to the second
and therefore (by \lemref{3.2.2}) to both questions is no. We will therefore
have to investigate and characterise the non-direct factors.

\ChapterHeading{4}{IV}{Auxiliary Theorems}

\SectionStart{4.1}{\texorpdfstring{Discussion of
$\sum_1^n A_iA_i'=0$. Theorem III}{Discussion of the zero-sum identity.
Theorem III}}
The theorem which follows has some interest of its own. We therefore prove it
in its general form, although we will mostly need its extremely special
Corollary only.

\MainTheorem{III}
\emph{Let $\M$ be a factor and let operators
$A_1,\ldots,A_n\in\M$, $A_1',\ldots,A_n'\in\M'$ be given. We have}
\[
 \sum_{i=1}^n A_iA_i'=0,
\]
\emph{if and only if a matrix $(a_{i,j})_{i,j=1,\ldots,n}$ exists, such that}
\[
 \sum_{i=1}^n a_{i,j}A_i=0,
 \qquad
 \sum_{j=1}^n a_{i,j}A_j'=A_i'.
\]

\Proof The condition is obviously sufficient: If it is valid,
$\sum_{i,j=1}^n a_{i,j}A_iA_j'$ gives, when first summed over $i$, $0$; and
when first summed over $j$, $\sum_{i=1}^n A_iA_i'$. In order to prove its
necessity, assume $\sum_{i=1}^n A_iA_i'=0$.

Form the space $n\otimes\Hspace$ of all systems
$\langle f_1,\ldots,f_n\rangle$ in $\Hspace$, with the inner product
\[
 \bigl(\langle f_1,\ldots,f_n\rangle,
       \langle g_1,\ldots,g_n\rangle\bigr)
 =\sum_{i=1}^n(f_i,g_i)
\]
(cf. \secref{2.4}). Let $\overline{\mathfrak M}$ be the set of all those
$\langle f_1,\ldots,f_n\rangle$, for which
$\sum_{i=1}^n A_iXf_i=0$ for every $X\in\M$.
$\overline{\mathfrak M}$ is obviously linear and closed. Let $\overline E$ be
the projection of $\overline{\mathfrak M}$. $\overline E$ can be written as an
operator matrix $\langle E_{r,s}\rangle_{r,s=1,\ldots,n}$.

Assume now $\langle f_1,\ldots,f_n\rangle\in\overline{\mathfrak M}$. If
$X_0\in\M$, then for every $X\in\M$
\[
 \sum_{i=1}^n A_iXX_0f_i=0,
\]
thus $\langle X_0f_1,\ldots,X_0f_n\rangle\in\overline{\mathfrak M}$. If
$X_0'\in\M'$, then for every $X\in\M$,
\[
 \sum_{i=1}^n A_iXX_0'f_i
 =X_0'\sum_{i=1}^n A_iXf_i=0,
\]
thus $\langle X_0'f_1,\ldots,X_0'f_n\rangle\in\overline{\mathfrak M}$.
Thus $\Phi\in\overline{\mathfrak M}$ and $X_0\in\M$ or $\M'$ imply for
$X_0^{(2)}=\langle\delta_{r,s}X_0\rangle_{r,s=1,\ldots,n}$,
$X_0^{(2)}\Phi\in\overline{\mathfrak M}$. As every
$\overline E\Phi\in\overline{\mathfrak M}$, we have
$X_0^{(2)}\overline E\Phi\in\overline{\mathfrak M}$,
\[
 \overline E X_0^{(2)}\overline E\Phi=X_0^{(2)}\overline E\Phi,
\]
that is $\overline E X_0^{(2)}\overline E=X_0^{(2)}\overline E$.
Replacing $X_0$ by $X_0^*$, and thus $X_0^{(2)}$ by
$X_0^{*(2)}=X_0^{(2)*}$, and applying $*$ gives
$\overline E X_0^{(2)}\overline E=\overline E X_0^{(2)}$. Thus
$X_0^{(2)}\overline E=\overline E X_0^{(2)}$, or in other words
$X_0E_{r,s}=E_{r,s}X_0$. Thus $E_{r,s}$ commutes with every $X_0\in\M$ or
$\M'$, $E_{r,s}\in\M'\cdot\M''=\M'\cdot\M=(\alpha1)$, that is
$E_{r,s}=a_{r,s}1$.

Consider $\langle A_1'f,\ldots,A_n'f\rangle$ for an arbitrary
$f\in\Hspace$. If $X\in\M$, then
\[
 \sum_{i=1}^n A_iXA_i'f
 =\sum_{i=1}^n A_iA_i'Xf=0.
\]
Thus $\langle A_1'f,\ldots,A_n'f\rangle\in\overline{\mathfrak M}$,
\[
 \overline E\langle A_1'f,\ldots,A_n'f\rangle
 =\langle A_1'f,\ldots,A_n'f\rangle.
\]
That is:
\[
 A_i'f=\sum_{j=1}^n E_{i,j}A_j'f
 =\sum_{j=1}^n a_{i,j}A_j'f,
 \qquad
 \sum_{j=1}^n a_{i,j}A_j'=A_i'.
\]

Consider next $\langle A_1^*f,\ldots,A_n^*f\rangle$ for an arbitrary
$f\in\Hspace$. Further choose an
$\langle f_1,\ldots,f_n\rangle\in\overline{\mathfrak M}$. Then
\begin{align*}
 \bigl(\langle A_1^*f,\ldots,A_n^*f\rangle,
       \langle f_1,\ldots,f_n\rangle\bigr)
 &=\sum_{i=1}^n(A_i^*f,f_i)
   =\sum_{i=1}^n(f,A_if_i)\\
 &=\left(f,\sum_{i=1}^n A_if_i\right)=0.
\end{align*}
Thus $\langle A_1^*f,\ldots,A_n^*f\rangle$ is orthogonal to
$\overline{\mathfrak M}$, and therefore
\[
 \overline E\langle A_1^*f,\ldots,A_n^*f\rangle=0.
\]
That is:
\[
 0=\sum_{j=1}^nE_{i,j}A_j^*f
  =\sum_{j=1}^na_{i,j}A_j^*f,
 \qquad
 \sum_{j=1}^na_{i,j}A_j^*=0.
\]
Applying $*$ gives $\sum_{j=1}^n\overline{a}_{i,j}A_j=0$.
$\overline E$ is Hermitian, this implies obviously
$E_{r,s}^*=E_{s,r}$, $\overline a_{r,s}=a_{s,r}$. Thus the last equation
becomes $\sum_{j=1}^n a_{j,i}A_j=0$. So the proof is completed.

\par\smallskip\noindent\textbf{Corollary:}
\emph{For $A\in\M$, $A'\in\M'$ we have $AA'=0$ if and only if $A=0$ or
$A'=0$.}

\Proof Put $n=1$, $a_{1,1}=\alpha$: $AA'=0$ means
$\alpha A=0$, $(1-\alpha)A'=0$. This means for
$\alpha=0$, $\alpha=1$, $\alpha\neq0,1$, that $A'=0$ resp. $A=0$ resp.
both.

\SectionStart{4.2}{\texorpdfstring{The notion $\eta\,\M$}
{The notion eta M}}
We define

\Definition{4.2.1}
Let $\M$ be a ring which contains $1$. A subset $\mathfrak S$ of $\Hspace$ or
an arbitrary operator $Z$ in $\Hspace$ (not necessarily $\in\B$!) are said to
\emph{belong} to the ring $\M$ in symbols $\mathfrak S\mathrel\eta\M$ resp.
$Z\mathrel\eta\M$ (we use the $\eta$ in order to distinguish this from the
element-relation $\in$), if they are invariant under every unitary
$U\in\M'$: that is, $U\mathfrak S=\mathfrak S$ resp.
$UZU^{-1}=Z$ for all $U\in\M'$.

The relationship between $\eta$ and $\in$ is easily established.

\Lemma{4.2.1}
\emph{If $\mathfrak M$ is a linear closed set, then
$\mathfrak M\mathrel\eta\M$ is equivalent to $P_{\mathfrak M}\in\M$
(cf. \cite[p.~74]{ref16}). If $Z\in\B$, then $Z\mathrel\eta\M$ is equivalent
to $Z\in\M$.}

\Proof $\mathfrak M\mathrel\eta\M$ means
$UP_{\mathfrak M}U^{-1}=P_{U\mathfrak M}=P_{\mathfrak M}$,
$UP_{\mathfrak M}=P_{\mathfrak M}U$ if $U$ is unitary and $\in\M'$, thus
$P_{\mathfrak M}\mathrel\eta\M$. Thus it suffices to prove the second
statement.

If $Z\in\B$ then $Z\mathrel\eta\M$ means $UZ=ZU$ for all unitary
$U\in\M'$ that is for $U\in\M'^u$ (cf. \cite[p.~392]{ref18}). That is:
\[
 Z\in(\M'^u)'=(R(\M'^u))'=\M''=\M.
\]

\Lemma{4.2.2}
\emph{The conditions of \defref{4.2.1} can be replaced by
$U\mathfrak S\subset\mathfrak S$ and $UZ\subset ZU$. The latter means that
$Z$ commutes with $U$ in the sense of \cite[p.~404]{ref18}.}

\Proof Replacing $U$ by $U^{-1}$ (which is unitary and $\in\M'$ too,
$U^{-1}=U^*$) and multiplying left resp. on both sides with $U$ gives
$\mathfrak S\subset U\mathfrak S$ and $ZU\subset UZ$. Thus
$\mathfrak S=U\mathfrak S$ and $UZU^{-1}=Z$, $UZ=ZU$.

\SectionStart{4.3}{Partially isometric operators}
We will now define a class of operators which is a generalisation of the
unitary ones.

\Definition{4.3.1}
An operator $U\in\B$ is \emph{partially isometric} if it maps a linear closed
set $\mathfrak E$ isometrically on another linear closed set $\mathfrak F$,
while it maps $\Hspace-\mathfrak E$ on $0$. That is: $f\in\mathfrak E$
implies $Uf\in\mathfrak F$ and $\lVert Uf\rVert=\lVert f\rVert$, the range of
$U$ is all $\mathfrak F$, $f\in\Hspace-\mathfrak E$ implies $Uf=0$.
$\mathfrak E$ and $\mathfrak F$ are the \emph{initial} resp. \emph{final} sets
of $U$.

A unitary operator $U$ is obviously such a partially isometric one, with
$\Hspace$ as initial and final set.

\Lemma{4.3.1}
\emph{$P_{\mathfrak E}=U^*U$, $P_{\mathfrak F}=UU^*$, $U^*$ is partially
isometric too, with interchanged initial and final sets: $\mathfrak F$,
$\mathfrak E$; as a mapping of $\mathfrak F$ on $\mathfrak E$, $U^*$ is the
inverse of the mapping of $\mathfrak E$ on $\mathfrak F$, $U$.}

\Proof If $f,g\in\mathfrak E$, then substitution of
$h=(f\pm g)/2,(f\pm ig)/2$ into $\lVert Uh\rVert^2=\lVert h\rVert^2$
gives $(Uf,Ug)=(f,g)$, that is
$(Uf,Ug)=(P_{\mathfrak E}f,P_{\mathfrak E}g)$
(cf. \cite[p.~71]{ref16}). If $f$ or $g\in\Hspace-\mathfrak E$, the latter
equation remains valid, as both sides vanish. Owing to the linearity, it holds
therefore for all $f,g$. We can write it as
$(U^*Uf,g)=(P_{\mathfrak E}f,g)$. So $U^*U=P_{\mathfrak E}$. If the
character of $U^*$ were already established, substitution of $U^*$ would give
$UU^*=P_{\mathfrak F}$. So we need only to discuss $U^*$.

For $f\in\mathfrak E$, $U^*Uf=P_{\mathfrak E}f=f$, so $U^*$ maps
$\mathfrak F$ to $\mathfrak E$ inversely to $U$. Thus this mapping is
isometric in $\mathfrak F$. Assume now $f\in\Hspace-\mathfrak F$. Then
$(f,Ug)=0$ whenever $Ug\in\mathfrak F$, but this is the case for
$g\in\mathfrak E$ and for $g\in\Hspace-\mathfrak E$ (then $Ug=0$), thus for
all $g$. So $(U^*f,g)=(f,Ug)=0$, $U^*f=0$. This completes the proof of
$U^*$'s being partially isometric with the initial and final sets
$\mathfrak F$, $\mathfrak E$ resp., and its being inverse to $U$ there.

\Lemma{4.3.2}
\emph{Any of the four following conditions is characteristic for partially
isometric operators. $UU^*U=U$, $U^*UU^*=U^*$, $U^*U$ is a projection,
$UU^*$ is a projection.}

\Proof Owing to the symmetry between $U$ and $U^*$ it suffices to consider
the first and the third conditions. If $UU^*U=U$ then
$U^*UU^*U=U^*U$, $(U^*U)^2=U^*U$, thus $U^*U$ is a projection; so we need
only to prove that the first condition is necessary, and that the third one is
sufficient.

If $U$ is partially isometric, we have always $Uf\in\mathfrak F$, so
$P_{\mathfrak F}Uf=Uf$, $P_{\mathfrak F}U=U$; but
$UU^*=P_{\mathfrak F}$, so $UU^*U=U$.

If $U^*U$ is a projection, put $U^*U=P_{\mathfrak E}$. Thus for
$f\in\mathfrak E$,
$\lVert Uf\rVert^2=(U^*Uf,f)=(P_{\mathfrak E}f,f)=\lVert f\rVert^2$,
$Uf$ is isometric in $\mathfrak E$. The image $\mathfrak F$ of
$\mathfrak E$ is therefore isomorphic to $\mathfrak E$, thus a linear closed
set (that is: a linear and topologically complete space) too. For
$f\in\Hspace-\mathfrak E$,
$\lVert Uf\rVert^2=(U^*Uf,f)=(P_{\mathfrak E}f,f)=0$. Thus $U$ is partially
isometric.

\Lemma{4.3.3}
\emph{Let $U_1,U_2,\ldots$ be a (finite or infinite) sequence of partially
isometric operators, with the resp. initial sets
$\mathfrak E_1,\mathfrak E_2,\ldots$ and the resp. final sets
$\mathfrak F_1,\mathfrak F_2,\ldots$; let the $\mathfrak E_i$ be mutually
orthogonal and let the $\mathfrak F_i$ be mutually orthogonal. Then the sum
$\sum_iU_i$ is strongly convergent (if infinite at all), it is partially
isometric, and its initial and final sets are
$[\mathfrak E_i;i=1,2,\ldots]$ and $[\mathfrak F_i;i=1,2,\ldots]$.}

\Proof The strong convergence of $\sum_iU_i$ means the same thing as that of
each $\sum_iU_if$. As $U_if\in\mathfrak F_i$, the addends are mutually
orthogonal, thus the above convergence is equivalent to the finiteness of
$\sum_i\lVert U_if\rVert^2$ (cf. \cite[p.~67]{ref16}). Now
\[
 \sum_i\lVert U_if\rVert^2
 =\sum_i(U_i^*U_if,f)
 =\sum_i(P_{\mathfrak E_i}f,f)
 =\bigl(P_{[\mathfrak E_i;i=1,2,\ldots]}f,f\bigr),
\]
and thus finite. Replacing $U_i$ by $U_i^*$, we see that $\sum_iU_i^*$ is
strongly convergent too.

So
\[
 \left(\sum_iU_i\right)^*\left(\sum_iU_i\right)
 =\sum_iU_i^*\sum_jU_j=\sum_{i,j}U_i^*U_j.
\]
If $i\neq j$, then $U_jf\in\mathfrak F_j\subset\Hspace-\mathfrak F_i$,
$U_i^*U_jf=0$, $U_i^*U_j=0$. So our above expression is equal to
\[
 \sum_iU_i^*U_i=\sum_iP_{\mathfrak E_i}
 =P_{[\mathfrak E_i;i=1,2,\ldots]}.
\]
Thus $\sum_iU_i$ is a partially isometric operator by \lemref{4.3.2} and its
initial set is $[\mathfrak E_i;i=1,2,\ldots]$, by \lemref{4.3.1}. Replacing
$U_i$ by $U_i^*$, we see that the final set is
\[
 [\mathfrak F_i;i=1,2,\ldots].
\]

\SectionStart{4.4}{\texorpdfstring{The canonical decomposition $A=WB$}
{The canonical decomposition A=WB}}
A construction which is described in \cite[p.~307, Theorem~7]{ref21}, will be
of importance in all our discussions. This Theorem applies to all linear,
closed operators $A$ with an everywhere dense domain, and leads to some
operators of which we will consider $B$ and $W$. $B$ is self adjoint, and
definite, and $W$ is obviously partially isometric, its initial and final sets
being
\[
 [\operatorname{Range}B]
 =\Hspace-(f;Bf=0)
 =\Hspace-(f;Af=0)
 =[\operatorname{Range}A^*]
\]
and
\[
 \Hspace-(f;A^*f=0)=[\operatorname{Range}A]
\]
resp. We have
\[
 A=WB,\qquad A^*=BW^*,\qquad B^2=A^*A.
\]
Cf. loc. cit. \cite{ref21}.

We define:

\Definition{4.4.1}
If $A$ is linear, closed, and has an everywhere dense domain, then the above
described $W,B$ are its \emph{canonical decomposition}.

We are interested in the relation between rings and the canonical
decomposition.

\Lemma{4.4.1}
\emph{If $\M$ is a ring containing $1$, $A\mathrel\eta\M$, and $W,B$ the
canonical decomposition of $A$, then $W\in\M$, $B\mathrel\eta\M$.}

\Proof $W,B$ are uniquely determined by their properties which follow:
\begin{gather*}
 B\text{ is self adjoint and definite};\qquad B^2=A^*A.\\
 A=WB;\qquad Bf=0\text{ implies }Wf=0.
\end{gather*}
In fact: The two first properties determine $B$ (cf. \cite[p.~303]{ref21});
and the two last ones determine $W$ on Range $B$ and on $(f;Bf=0)$,
therefore on
\[
 [[\operatorname{Range}B],(f;Bf=0)]
 =[\Hspace-(f;Af=0),(f;Af=0)]=\Hspace,
\]
that is they determine $W$ completely.

Now if $U$ is unitary and $W,B$ is the canonical decomposition of $A$, then
$UWU^{-1},UBU^{-1}$ is the one of $UAU^{-1}$. Thus $UAU^{-1}=A$ implies
$UWU^{-1}=W$, $UBU^{-1}=B$. If $A\mathrel\eta\M$ then let $U$ run over all
$\M'$, this gives $W,B\mathrel\eta\M$. As $W\in\B$, we even have
$W\in\M$.

\ChapterHeading{5}{V}{Direct Factors. Ring Isomorphisms}

\SectionStart{5.1}{\texorpdfstring{$\mathfrak M_f^{\M'}$,
$\mathfrak M_f^{\M}$. The notion of minimality. The minimality criterion}
{Two cyclic subspaces. The notion of minimality. The minimality criterion}}
We define

\Definition{5.1.1}
If $\M$ is a ring containing $1$, denote $[Af;A\in\M']$ by
$\mathfrak M_f^{\M'}$. Write $E_f^{\M'}$ for
$P_{\mathfrak M_f^{\M'}}$.

\Lemma{5.1.1}
\emph{$f\in\mathfrak M_f^{\M'}\mathrel\eta\M$; whenever
$f\in\mathfrak M\mathrel\eta\M$ then
$\mathfrak M_f^{\M'}\subset\mathfrak M$.}

\Proof $A=1$ gives $f\in\mathfrak M_f^{\M'}$. If $U\in\M'$, then
\[
 U(Af;A\in\M')=(UAf;A\in\M')\subset(Bf;B\in\M');
\]
that is $U\mathfrak M_f^{\M'}\subset\mathfrak M_f^{\M'}$. Thus
$\mathfrak M_f^{\M'}\mathrel\eta\M$ (by \lemref{4.2.2}). Conversely: Assume
$f\in\mathfrak M\mathrel\eta\M$, and consider those $A$ for which $Af$ and
$A^*f$ both $\in\mathfrak M$ for all $f\in\mathfrak M$. They clearly form a
ring (use the strong topology; for instance by \cite[\S3]{ref22}), and contain
all unitary $U\in\M'$, that is all $(\M')^u$. Thus they contain all
$R((\M')^u)=\M'$ (cf. \cite[p.~392]{ref18}), and so
\[
 (Af;A\in\M')\subset\mathfrak M,\qquad
 [Af;A\in\M']\subset\mathfrak M,\qquad
 \mathfrak M_f^{\M'}\subset\mathfrak M.
\]

\Definition{5.1.2}
$\mathfrak M\mathrel\eta\M$ is \emph{minimal} (with respect to $\M$), if
$\mathfrak M\neq(0)$, and $\mathfrak N\mathrel\eta\M$,
$\mathfrak N\subset\mathfrak M$ implies $\mathfrak N=(0)$ or
$\mathfrak M$. Replacing $\mathfrak M,\mathfrak N$ by
$E=P_{\mathfrak M}$, $F=P_{\mathfrak N}$ we can say (cf. \lemref{4.2.1} and
\cite[p.~76]{ref16}): projection $E\in\M$ is minimal (with respect to $\M$),
if $E\neq0$, and if for every projection $F\in\M$, $F\leq E$ implies $F=0$
or $F=E$.

\Lemma{5.1.2}
\emph{$\mathfrak M\mathrel\eta\M$ is minimal if and only if
$\mathfrak M\neq(0)$, and for every $f\in\mathfrak M$, $f\neq0$,
$\mathfrak M_f^{\M'}=\mathfrak M$.}

\Proof Assume first that $\mathfrak M\mathrel\eta\M$ is minimal. Then
$\mathfrak M\neq(0)$. If $f\in\mathfrak M$, $f\neq0$, then
$\mathfrak M_f^{\M'}\mathrel\eta\M$; as
$f\in\mathfrak M_f^{\M'}$, $\mathfrak M_f^{\M'}\neq0$ and as
$f\in\mathfrak M$, $\mathfrak M_f^{\M'}\subset\mathfrak M$. Thus
$\mathfrak M_f^{\M'}=\mathfrak M$ by definition.

Assume now conversely that our condition is fulfilled. If
$\mathfrak N\subset\mathfrak M$, $\mathfrak N\mathrel\eta\M$, then either
$\mathfrak N=(0)$, or an $f\in\mathfrak N$, $f\neq0$ exists. Then
$\mathfrak M_f^{\M'}\subset\mathfrak N$; but as $f\in\mathfrak M$,
$f\neq0$, $\mathfrak M_f^{\M'}=\mathfrak M$; so
$\mathfrak M\subset\mathfrak N\subset\mathfrak M$,
$\mathfrak N=\mathfrak M$.

\Lemma{5.1.3}
\emph{If $\M$ is a factor, then $\mathfrak M_f^{\M'}$ is minimal if and only
if $f\neq0$ and
$\mathfrak M_f^{\M'}\cdot\mathfrak M_f^{\M}=(\alpha f)$.}

\Proof Assume first that $f\neq0$ and
$\mathfrak M_f^{\M'}\cdot\mathfrak M_f^{\M}=(\alpha f)$. Assume that
$\mathfrak M_f^{\M'}$ is not minimal. As $f\neq0$,
$\mathfrak M_f^{\M'}\neq(0)$, there must exist an
$\mathfrak N\mathrel\eta\M$ with
$\mathfrak N\subset\mathfrak M_f^{\M'}$,
$\mathfrak N\neq0,\mathfrak M_f^{\M'}$. Put
$\mathfrak P=\mathfrak M_f^{\M'}-\mathfrak N$; then $\mathfrak N,\mathfrak P$
are orthogonal, both $\subset\mathfrak M_f^{\M'}$, both $\neq0$.

Thus for $F=P_{\mathfrak N}$, $G=P_{\mathfrak P}$ we have:
$F,G\in\M$, $FG=0$, $F,G\neq0$. As
$E_f^{\M}=P_{\mathfrak M_f^{\M}}\in\M'$ and $\neq0$ too, the Corollary to
\thmref{III} gives
\[
 P_{\mathfrak N}P_{\mathfrak M_f^{\M}}
 =F\cdot P_{\mathfrak M_f^{\M}}\neq0,
 \qquad
 P_{\mathfrak P}P_{\mathfrak M_f^{\M}}
 =G\cdot P_{\mathfrak M_f^{\M}}\neq0.
\]
So $\mathfrak N\cdot\mathfrak M_f^{\M}$ and
$\mathfrak P\cdot\mathfrak M_f^{\M}$ are both $\neq(0)$. But they are
$\subset\mathfrak M_f^{\M'}\cdot\mathfrak M_f^{\M}=(\alpha f)$, so they
are both $=(\alpha f)$. Thus
$f\in\mathfrak N\cdot\mathfrak M_f^{\M}$ and
$\mathfrak P\cdot\mathfrak M_f^{\M}$ and a fortiori
$f\in\mathfrak N$ and $\mathfrak P$. As $\mathfrak N,\mathfrak P$ are
orthogonal, this would imply $f=0$, contradicting our original assumption.

Assume now conversely that $\mathfrak M_f^{\M'}$ is minimal. As
$\mathfrak M_f^{\M'}\neq(0)$, we have at any rate $f\neq0$.

Put $E=P_{\mathfrak M_f^{\M'}}$, $P_{\mathfrak M_f^{\M}}=E'$. As
$E\in\M$, $E'\in\M'$, therefore
$E,E'$ commute; therefore the $E$-image of $\mathfrak M_f^{\M}$ is
$\mathfrak M_f^{\M'}\cdot\mathfrak M_f^{\M}$. Thus (remember $Ef=f$!)
\begin{align*}
 \mathfrak M_f^{\M'}\cdot\mathfrak M_f^{\M}
 &=[EAf,A\in\M]=[EAEf,A\in\M]\\
 &=[Bf,B\in\overline{\M}],
\end{align*}
where $\overline{\M}$ is the set of all $B\in\M$ with $EB=BE=B$
($\overline{\M}$ is obviously a ring, thus
$\overline{\M}=R(\overline{\M}^P)$ (cf. \cite[p.~392]{ref18})).
$F'\in\overline{\M}^P$ means: $F'$ is a projection, and
$\in\overline{\M}$, and so
$EF'=F'E=F'$, $F'\in\M$. So $F'\in\M$, $F'\leq E$, and as $E$ is minimal,
$F'=0$ or $E$. So $\overline{\M}=R(0,E)=(\alpha E)$. Therefore
\[
 \mathfrak M_f^{\M'}\cdot\mathfrak M_f^{\M}
 =[Bf,B\in(\alpha E)]=[\alpha Ef]=[\alpha f]=(\alpha f).
\]

\Lemma{5.1.4}
\emph{If $\M$ is a factor, then $\mathfrak M_f^{\M'}$ is minimal (with
respect to $\M$) if and only if $\mathfrak M_f^{\M}$ is minimal (with
respect to $\M'$).}

\Proof The criterium of \lemref{5.1.3} is obviously symmetric in $\M$ and
$\M'$, as $\M=\M''$.

\Definition{5.1.3}
An $f\neq0$ is \emph{minimal} (with respect to $\M,\M'$) if its
$\mathfrak M_f^{\M'}$, or, which is equivalent, its
$\mathfrak M_f^{\M}$ is minimal.

\Lemma{5.1.5}
\emph{If $\M$ is a factor, and if $f$ is minimal, then every $g\neq0$,
$g\in\mathfrak M_f^{\M'}$ or $\mathfrak M_f^{\M}$ is minimal.}

\Proof $\mathfrak M_f^{\M'}$ is minimal, and since $g\neq0$,
$g\in\mathfrak M_f^{\M'}$ implies by \lemref{5.1.2},
$\mathfrak M_g^{\M'}=\mathfrak M_f^{\M'}$, $g$ is too. $g\neq0$,
$g\in\mathfrak M_f^{\M}$ is taken care of by symmetry.

\SectionStart{5.2}{Ring Isomorphisms}
Various sorts of ring isomorphisms are possible. They are defined as follows:

\Definition{5.2.1}
Let $\Hspace_{\mathrm I}$ and $\Hspace_{\mathrm{II}}$ be two spaces,
$\B_{\mathrm I}$ and $\B_{\mathrm{II}}$ their operator rings,
$\M_{\mathrm I}\subset\B_{\mathrm I}$ and
$\M_{\mathrm{II}}\subset\B_{\mathrm{II}}$ rings in $\Hspace_{\mathrm I}$
resp. $\Hspace_{\mathrm{II}}$ which contain $1$. If $V$ is any set of
operations and properties defined in both $\B_{\mathrm I}$ and
$\B_{\mathrm{II}}$, then $\M_{\mathrm I}$ and $\M_{\mathrm{II}}$ are called
$V$-\emph{ring-isomorphic}, if a one-to-one mapping of $\M_{\mathrm I}$ on
$\M_{\mathrm{II}}$ exists which leaves the operations and properties of $V$
invariant.

If $V=(\alpha A,A^*,A+B,AB,\text{strongest closure})$, we speak of
\emph{full ring-isomorphism}; if $V=(\alpha A,A^*,A+B,AB)$, of
\emph{algebraic ring-isomorphism}.

It is clear that every isomorphism of the spaces $\Hspace_{\mathrm I}$,
$\Hspace_{\mathrm{II}}$ generates full ring isomorphisms (in particular one
between $\B_{\mathrm I}$ and $\B_{\mathrm{II}}$). But the really interesting
cases are such ring isomorphisms between an $\M_{\mathrm I}$ and an
$\M_{\mathrm{II}}$, which do not originate from spatial isomorphisms.

As our methods are invariant with respect to the spaces $\Hspace$, all notions
we considered are invariant under spatial isomorphisms. But it is not so with
ring isomorphisms. For instance the operation $\M'$ is not invariant even under
full ring isomorphisms. (Let $\Hspace_{\mathrm I}$,
$\Hspace_{\mathrm{II}}$ be two finite dimensional Euclidean spaces of different
dimensions. Put $\M_{\mathrm I}=\M_{\mathrm{II}}=(\alpha1)$, then
$\M_{\mathrm I}'=\B_{\mathrm I}$,
$\M_{\mathrm{II}}'=\B_{\mathrm{II}}$; thus $\M_{\mathrm I}$,
$\M_{\mathrm{II}}$ are fully ring isomorphic, while $\M_{\mathrm I}'$,
$\M_{\mathrm{II}}'$ are not.) It is therefore of importance to show that some
fundamental notions are invariant under certain types of ring isomorphism.

\Lemma{5.2.1}
\emph{Every $(A^*,AB)$-ring-isomorphism leaves the following notions
invariant: To be $0$; to be $1$; to be a projection; to be partially isometric;
for two operators: to commute; for two projections $E,F$: $F\leq E$; for a
projection: to be minimal.}

\Proof $A=0$ means $AX=A$ for all $X\in\M_{\mathrm I}$; $A=1$ means
$AX=X$ for all $X\in\M_{\mathrm I}$; $A$ is a projection if $A^*=A$,
$AA=A$; $A$ is partially isometric if $AA^*A=A$; commutativity means
$AB=BA$; $F\leq E$ means $EF=F$; minimality can be expressed as a
combination of the foregoing statements.

\Lemma{5.2.2}
\emph{Every $(\alpha A,A^*,AB)$-ring-isomorphism leaves the notion of a
factor (as applied to the rings $\M_{\mathrm I}$, $\M_{\mathrm{II}}$
themselves) invariant.}

\Proof As the rings $\M$ under consideration contain $1$, being a factor means
$\M\cdot\M'=(\alpha1)$, that is: If $A\in\M$ is such that for every
$B\in\M$, $AB=BA$, then $A=\alpha1$. The first half is invariant under
$(A^*,AB)$-ring-isomorphisms, and so is the $1$ (cf. \lemref{5.2.1}), while
the equation $A=\alpha1$ is then invariant under
$(\alpha A)$-ring-isomorphisms.

\SectionStart{5.3}{Characterisation of direct factors by the existence of
minimal elements. Theorems IV and V}
We are now able to give a simple criterium for direct factors.

\Lemma{5.3.1}
\emph{The three following conditions are equivalent for a factor $\M$:}
\begin{enumerate}[label=(\roman*)]
\item \emph{$\M$ contains a minimal projection $E$ (with respect to $\M$).}
\item \emph{$\M'$ contains a minimal projection $E'$ (with respect to
  $\M'$).}
\item \emph{There exists a minimal $f\in\Hspace$ (with respect to
  $\M,\M'$).}
\end{enumerate}

\Proof (iii) implies (i). (i) implies (iii) by \lemref{5.1.2}. Similarly (ii)
and (iii) are equivalent.

\Lemma{5.3.2}
\emph{The conditions of \lemref{5.3.1} are fulfilled for every direct factor
$\M$.}

\Proof If $\M$ is a direct factor in the $\B$ of $\Hspace$ then $\Hspace$ is
isomorphic to $\Hspace_1\otimes\Hspace_2$, so that $\M$ becomes
$\B_1^{(1)}$. $\B_1^{(1)}$ is algebraically (even fully) ring-isomorphic to
$\B_1$ of $\Hspace_1$ by \lemref{2.3.4} (resp. \cite[\S4]{ref22}); and the
condition of \lemref{5.3.1}, in its form (i), invariant under algebraic
ring-isomorphisms by \lemref{5.2.1}. Thus we need only to consider $\B_1$ in
$\Hspace_1$. Now if $\varphi$ is any element $\neq0$ of $\Hspace_1$, then
$P_{[\varphi]}\in\B_1$ and obviously minimal.

We begin now to discuss the sufficiency of these conditions.

\Lemma{5.3.3}
\emph{If a factor $\M$ fulfills the conditions of \lemref{5.3.1}, then every
$\mathfrak M\mathrel\eta\M$, $\mathfrak M\neq0$, contains a minimal $f$ with
$\lVert f\rVert=1$.}

\Proof A minimal $f^0$ exists, of course $f^0\neq0$. Thus
$E_{f^0}^{\M}\neq0$; besides $E_{f^0}^{\M}\in\M'$. Now
$P_{\mathfrak M}\neq0$, $P_{\mathfrak M}\in\M$, so by the Corollary of
\thmref{III},
\[
 P_{\mathfrak M_{f^0}^{\M}\cdot\mathfrak M}
 =P_{\mathfrak M}\cdot E_{f^0}^{\M}\neq0,
 \qquad
 \mathfrak M_{f^0}^{\M}\cdot\mathfrak M\neq(0).
\]
Thus an $f\in\mathfrak M_{f^0}^{\M}\cdot\mathfrak M$ with $f\neq0$ exists,
multiplying it with $1/\lVert f\rVert$ makes $\lVert f\rVert=1$.
As $f\in\mathfrak M_{f^0}^{\M}$, $f\neq0$ this $f$ is minimal by
\lemref{5.1.5}; and besides $f\in\mathfrak M$.

\Lemma{5.3.4}
\emph{If a factor $\M$ fulfills the conditions of \lemref{5.3.1}, then a
finite or infinite normalised orthogonal system $\varphi_1,\varphi_2,\ldots$
exists, such that}
\begin{enumerate}[label=(\roman*)]
\item \emph{every $\varphi_n$ is minimal;}
\item \emph{$\mathfrak M_{\varphi_1}^{\M'},
  \mathfrak M_{\varphi_2}^{\M'},\ldots$ are mutually orthogonal;}
\item \emph{$[\mathfrak M_{\varphi_1}^{\M'},
  \mathfrak M_{\varphi_2}^{\M'},\ldots]=\Hspace$.}
\end{enumerate}

\Proof Let $\Omega$ be the first uncountable Cantor ordinal number. Define for
all $\alpha<\Omega$, a $\varphi_\alpha$ as follows: If $\alpha<\Omega$, and
all $\varphi_\beta$, $\beta<\alpha$ are already defined then form
\[
 \Hspace-[\mathfrak M_{\varphi_\beta}^{\M'};\ \beta<\alpha].
\]
As this is obviously $\mathrel\eta\M$, it will contain a minimal $f$ with
$\lVert f\rVert=1$ if it is $\neq(0)$. Let then $\varphi_\alpha$ be some such
$f$; otherwise let $\varphi_\alpha$ (and all $\varphi_\gamma$,
$\gamma\geq\alpha$) be undefined.

Let the first $\alpha$ for which $\varphi_\alpha$ is undefined, be
$\bar\alpha\leq\Omega$. If $\bar\alpha<\Omega$, we must have
\[
 \Hspace-[\mathfrak M_{\varphi_\beta}^{\M'};\ \beta<\bar\alpha]=0,
 \qquad
 [\mathfrak M_{\varphi_\beta}^{\M'};\ \beta<\bar\alpha]=\Hspace.
\]
Consider now all $\varphi_\alpha$, $\alpha<\bar\alpha$. Always
$\lVert\varphi_\alpha\rVert=1$. If $\alpha,\beta<\bar\alpha$, assume
$\alpha>\beta$. Then $A',B'\in\M'$ imply $A'^*B'\in\M'$,
$A'^*B'\varphi_\beta\in\mathfrak M_{\varphi_\beta}^{\M'}$ orthogonal to
$\varphi_\alpha$, thus $B'\varphi_\beta$ orthogonal to
$A'\varphi_\alpha$ (because
$(A'^*B'\varphi_\beta,\varphi_\alpha)
=(B'\varphi_\beta,A'\varphi_\alpha)$). So
$(A'\varphi_\alpha;A'\in\M')$ is orthogonal to
$(B'\varphi_\beta;B'\in\M')$,
$[A'\varphi_\alpha;A'\in\M']$ to
$[B'\varphi_\beta;B'\in\M']$ and so
$\mathfrak M_{\varphi_\alpha}^{\M'}$ to
$\mathfrak M_{\varphi_\beta}^{\M'}$. This is symmetric in $\alpha,\beta$ so
it holds for $\alpha<\beta$ too, that is whenever $\alpha\neq\beta$.

Thus $\varphi_\alpha,\varphi_\beta$ are orthogonal if
$\alpha\neq\beta$, that is the $\varphi_\alpha$, $\alpha<\bar\alpha$, form a
normalised orthogonal set. It must therefore be finite or countably infinite
(cf. \cite[p.~66]{ref16}), thus excluding $\bar\alpha=\Omega$. Thus
$\bar\alpha<\Omega$, and so
$[\mathfrak M_{\varphi_\alpha}^{\M'};\alpha<\bar\alpha]=\Hspace$.

If $\bar\alpha$ is finite, $\Hspace\neq(0)$ excludes $\bar\alpha=1$, so
$\bar\alpha=n+1$, $n=1,2,\ldots$, and $\alpha$ runs over $1,\ldots,n$.
If $\bar\alpha$ is infinite then $\alpha<\bar\alpha$ has the same aleph as
$\alpha<\omega$, and may therefore (by a re-indexing of the $\varphi_\alpha$)
be replaced by it; so we may assume that $\alpha$ runs over $1,2,\ldots$.

This completes the proof.

\Lemma{5.3.5}
\emph{If a factor $\M$ fulfills the conditions of \lemref{5.3.1} then a
complete normalised orthogonal set $f_{m,n}$, $m=1,2,\ldots$;
$n=1,2,\ldots$ exists (both ranges for $m$ and for $n$ may be finite or
infinite), such that:}
\begin{enumerate}[label=(\roman*)]
\item \emph{every $f_{m,n}$ is minimal,}
\item \emph{$\mathfrak M_{f_{m,n}}^{\M'}=[f_{m,1},f_{m,2},\ldots]$,
  $\mathfrak M_{f_{m,n}}^{\M}=[f_{1,n},f_{2,n},\ldots]$.}
\end{enumerate}

\Proof If we replace $\M$ by $\M'$, condition (i) of \lemref{5.3.1} becomes
(ii), so $\M'$ fulfills it too. Apply therefore \lemref{5.3.4} to $\M$ and
to $\M'$, obtaining the two normalised orthogonal systems
$\varphi_1,\varphi_2,\ldots$ and $\psi_1,\psi_2,\ldots$. As
$E_{\varphi_m}^{\M'}\in\M$ and $\neq0$,
$E_{\psi_n}^{\M}\in\M'$ and $\neq0$, we have by the Corollary to
\thmref{III}
\[
 P_{\mathfrak M_{\varphi_m}^{\M'}\cdot
    \mathfrak M_{\psi_n}^{\M}}
 =E_{\varphi_m}^{\M'}\cdot E_{\psi_n}^{\M}\neq0,
 \qquad
 \mathfrak M_{\varphi_m}^{\M'}\cdot
 \mathfrak M_{\psi_n}^{\M}\neq(0).
\]
Choose an
$f_{m,n}\in\mathfrak M_{\varphi_m}^{\M'}\cdot
\mathfrak M_{\psi_n}^{\M}$ which is $\neq0$; multiplying it with
$1/\lVert f_{m,n}\rVert$ makes $\lVert f_{m,n}\rVert=1$.

If $m\neq p$, $f_{m,n}\in\mathfrak M_{\varphi_m}^{\M'}$,
$f_{p,q}\in\mathfrak M_{\varphi_p}^{\M'}$; if $n\neq q$,
$f_{m,n}\in\mathfrak M_{\psi_n}^{\M}$,
$f_{p,q}\in\mathfrak M_{\psi_q}^{\M}$ both imply
$(f_{m,n},f_{p,q})=0$. Thus the $f_{m,n}$ form a normalised orthogonal system.

$\mathfrak M_{\varphi_m}^{\M'}$ is minimal and $f_{m,n}$ belongs to it and is
$\neq0$, therefore \lemref{5.1.5} applies: $f_{m,n}$ is minimal too, and
$\mathfrak M_{\varphi_m}^{\M'}=\mathfrak M_{f_{m,n}}^{\M'}$. Similarly
$\mathfrak M_{\psi_n}^{\M}=\mathfrak M_{f_{m,n}}^{\M}$. As $f_{m,n}$ is
minimal, \lemref{5.1.3} gives
$\mathfrak M_{f_{m,n}}^{\M'}\cdot
\mathfrak M_{f_{m,n}}^{\M}=(\alpha f_{m,n})$, that is
\[
 \mathfrak M_{\varphi_m}^{\M'}\cdot
 \mathfrak M_{\psi_n}^{\M}=(\alpha f_{m,n}).
\]
Thus
\[
 \sum_{m,n}P_{[f_{m,n}]}
 =\sum_{m,n}E_{\varphi_m}^{\M'}\cdot E_{\psi_n}^{\M}
 =\sum_mE_{\varphi_m}^{\M'}\sum_nE_{\psi_n}^{\M}
 =1\cdot1=1;
\]
that is: the normalised orthogonal system of all $f_{m,n}$ is complete.

$f_{m,n}\in\mathfrak M_{f_{m,n}}^{\M'}
=\mathfrak M_{\varphi_m}^{\M'}$, and if $p\neq m$, then
$f_{p,n}\in\mathfrak M_{\varphi_p}^{\M'}$ is orthogonal to
$\mathfrak M_{\varphi_m}^{\M'}$. This proves
$\mathfrak M_{\varphi_m}^{\M'}=[f_{m,1},f_{m,2},\ldots]$. Similarly
$\mathfrak M_{\psi_n}^{\M}=[f_{1,n},f_{2,n},\ldots]$.

Thus all parts of our Lemma are proved.

\Lemma{5.3.6}
\emph{If a factor $\M$ fulfills the conditions of \lemref{5.3.1} we can
choose the $f_{m,n}$ of \lemref{5.3.5} so that we have:}

\emph{There is a unique partially isometric operator $U\in\M$ which has the
initial and final sets $[f_{m,1},f_{m,2},\ldots]$ resp.
$[f_{p,1},f_{p,2},\ldots]$ and for which $Uf_{m,n}=f_{p,n}$,
$Uf_{r,n}=0$, if $r\neq m$. Denote it by $U_{m,p}$.}

\emph{If an $A\in\M$ has the properties}
\[
 P_{[f_{p,1},f_{p,2},\ldots]}A
 =AP_{[f_{m,1},f_{m,2},\ldots]}=A
\]
\emph{then $A=\alpha U_{m,p}$.}

\Proof Assume first $A\in\M$ and
\[
 P_{[f_{p,1},f_{p,2},\ldots]}A
 =AP_{[f_{m,1},f_{m,2},\ldots]}=A
\]
and the same for $B$. Then applying a $*$ gives:
\[
 P_{[f_{m,1},f_{m,2},\ldots]}B^*
 =B^*P_{[f_{p,1},f_{p,2},\ldots]}=B^*
\]
and so
\[
 P_{[f_{m,1},f_{m,2},\ldots]}B^*A
 =B^*AP_{[f_{m,1},f_{m,2},\ldots]}=B^*A.
\]
Now $B^*A\in\M$ and
$P_{[f_{m,1},f_{m,2},\ldots]}=E_{f_{m,n}}^{\M'}$ is minimal with respect to
$\M$, so the argument made at the end of the proof of \lemref{5.1.3} applies
\[
 B^*A=\beta P_{[f_{m,1},f_{m,2},\ldots]}.
\]

Put first $A=B$, then
$B^*B=\beta_1P_{[f_{m,1},f_{m,2},\ldots]}$. As the left side is Hermitian
and definite, $\beta_1\geq0$. If $B\neq0$, then even $\beta_1\neq0$ and we
have
\[
 (\beta_1^{-1/2}B)^*(\beta_1^{-1/2}B)
 =P_{[f_{m,1},f_{m,2},\ldots]}.
\]
Thus $\beta_1^{-1/2}B$ is a partially isometric operator with the initial set
$[f_{m,1},f_{m,2},\ldots]$ (cf. \hyperref[lem:4.3.1]{Lemmas 4.3.1} and
\hyperref[lem:4.3.2]{4.3.2}). Replacing $B$ by $B^*$ we interchange $m$ and $p$, so
$\beta_2^{-1/2}B^*$ is partially isometric with the initial set
$[f_{p,1},f_{p,2},\ldots]$ and $\beta_2^{-1/2}B$ with the final set
$[f_{p,1},f_{p,2},\ldots]$. As $\beta_1^{-1/2}B$,
$\beta_2^{-1/2}B$ are both partially isometric, we have
$|\beta_1^{-1/2}|=|\beta_2^{-1/2}|$, that is $\beta_1=\beta_2$. So
$\beta_1^{-1/2}B$ has the initial and final sets
$[f_{m,1},f_{m,2},\ldots]$ resp. $[f_{p,1},f_{p,2},\ldots]$.

Return now to the pair $A,B$. If $B\neq0$, then we have
\[
 A=P_{[f_{p,1},f_{p,2},\ldots]}A
 =\frac1{\beta_1}BB^*A
 =\frac\beta{\beta_1}BP_{[f_{m,1},f_{m,2},\ldots]}
 =\frac\beta{\beta_1}B.
\]
So we see: either $B=0$ or $A=\alpha B$.

Consider now $f_{m,n}$ and $f_{p,n}$. As
$\mathfrak M_{f_{m,n}}^{\M}=[f_{1,n},f_{2,n},\ldots]$ we have
$f_{p,n}\in\mathfrak M_{f_{m,n}}^{\M}$. So an $A\in\M$ with
$\lVert f_{p,n}-Af_{m,n}\rVert<1$ exists. A fortiori
\[
 \left\lVert
 P_{[f_{p,1},f_{p,2},\ldots]}f_{p,n}
 -P_{[f_{p,1},f_{p,2},\ldots]}Af_{m,n}
 \right\rVert<1
\]
or as
\[
 P_{[f_{p,1},f_{p,2},\ldots]}f_{p,n}=f_{p,n},
 \qquad
 P_{[f_{m,1},f_{m,2},\ldots]}f_{m,n}=f_{m,n}
\]
we have $\lVert f_{p,n}-A_0f_{m,n}\rVert<1$, where
\[
 A_0=P_{[f_{p,1},f_{p,2},\ldots]}A
 P_{[f_{m,1},f_{m,2},\ldots]}.
\]
By what we prove above $A_0=0$ or $A_0=\beta_0U$, where $\beta_0\neq0$ and
$U$ is partially isometric with the initial resp. final set
$[f_{m,1},f_{m,2},\ldots]$ resp. $[f_{p,1},f_{p,2},\ldots]$. Now
$\lVert f_{p,n}\rVert=1$ implies $A_0f_{m,n}\neq0$ and so the latter
alternative holds.

As $A_0\in\M$,
$A_0f_{m,n}\in\mathfrak M_{f_{m,n}}^{\M}
=\mathfrak M_{f_{p,n}}^{\M}$, as
$A_0=P_{[f_{p,1},f_{p,2},\ldots]}A_0$ therefore
\[
 A_0f_{m,n}\in[f_{p,1},f_{p,2},\ldots]
 =\mathfrak M_{f_{p,n}}^{\M'}.
\]
Thus
$A_0f_{m,n}\in\mathfrak M_{f_{p,n}}^{\M'}\cdot
\mathfrak M_{f_{p,n}}^{\M}=(\alpha f_{p,n})$,
$A_0f_{m,n}=\alpha_0f_{p,n}$,
$Uf_{m,n}=(\alpha_0/\beta_0)f_{p,n}$. As
$\lVert f_{m,n}\rVert=\lVert f_{p,n}\rVert=1$ and $f_{m,n}$ is in the
initial set of $U$, we have $|\alpha_0/\beta_0|=1$. Replacing $U$ by
$(\beta_0/\alpha_0)U$ does not affect its other properties, but makes
$Uf_{m,n}=f_{p,n}$. So we have:
\[
 P_{[f_{p,1},f_{p,2},\ldots]}U
 =UP_{[f_{m,1},f_{m,2},\ldots]}=U,
 \qquad Uf_{m,n}=f_{p,n}.
\]

As $U$ depends on $m,n,p$ write $U=U_{m,p}^{(n)}$. The first equations
determine it uniquely up to a constant factor, the last one makes this factor
equal to $1$. Comparing $U_{m,p}^{(n)}$, $U_{m,p}^{(q)}$ we see that the first
equations still make them agree up to a constant factor:
\[
 U_{m,p}^{(n)}=\theta_{m,p}^{(n,q)}U_{m,p}^{(q)}
\]
and $\theta_{m,p}^{(n,q)}$ is obviously of absolute value $1$. On the other
hand the uniqueness property gives
\[
 U_{m,p}^{(n)}\cdot U_{p,q}^{(n)}=U_{m,r}^{(n)}.
\]

Replace every $f_{p,q}$, $p\neq1$, by
$\theta_{1,p}^{(1,q)}f_{p,q}$. This makes every
$\theta_{1,p}^{(1,q)}=1$. Now obviously
\[
 \theta_{1,p}^{(1,n)}\cdot\theta_{1,p}^{(n,q)}
 =\theta_{1,p}^{(1,q)},
\]
therefore $\theta_{1,p}^{(n,q)}=1$. Furthermore it is easy to see, that
\[
 \theta_{1,m}^{(n,q)}\cdot\theta_{m,p}^{(n,q)}
 =\theta_{1,p}^{(n,q)},
\]
therefore $\theta_{m,p}^{(n,q)}=1$. Thus $U_{m,p}^{(n)}$ does not depend on
$n$, $U_{m,p}^{(n)}=U_{m,p}$.

$U_{m,p}$ is partially isometric, its initial and final sets are
$[f_{m,1},f_{m,2},\ldots]$ resp. $[f_{p,1},f_{p,2},\ldots]$, and
$U_{m,p}f_{m,n}=f_{p,n}$. If $r\neq m$, then $f_{r,n}$ is orthogonal to
$[f_{m,1},f_{m,2},\ldots]$, and so $U_{m,p}f_{r,n}=0$. $U_{m,p}$ is uniquely
determined by these properties, because already the $U_{m,p}^{(n)}$ were.

Thus we have proved the first statement of our Lemma. The second results by
replacing in our result on $A,B$ at the beginning of this proof $A,B$ by
$A,U_{m,p}$.

\Lemma{5.3.7}
\emph{If a factor $\M$ fulfills the conditions of \lemref{5.3.1}, then we
have with the $U_{m,p}$ of \lemref{5.3.6}}
\[
 \M=R(U_{m,p};m,p=1,2,\ldots).
\]

\Proof As all $U_{m,p}\in\M$, we have
$R(U_{m,p};m,p=1,2,\ldots)\subset\M$. Assume now conversely $A\in\M$.
Put
\[
 A_{m,p}=P_{[f_{p,1},f_{p,2},\ldots]}
 A P_{[f_{m,1},f_{m,2},\ldots]}.
\]
Then $A_{m,p}\in\M$,
\[
 P_{[f_{p,1},f_{p,2},\ldots]}A_{m,p}
 =A_{m,p}P_{[f_{m,1},f_{m,2},\ldots]}=A_{m,p},
\]
so by \lemref{5.3.6} $A_{m,p}=a_{m,p}U_{m,p}$. Thus
$A_{m,p}\in R(U_{m',p'};m',p'=1,2,\ldots)$. But
\begin{align*}
 \sum_{m,p}A_{m,p}
 &=\sum_{m,p}P_{[f_{p,1},f_{p,2},\ldots]}
 A P_{[f_{m,1},f_{m,2},\ldots]}\\
 &=\left(\sum_pP_{[f_{p,1},f_{p,2},\ldots]}\right)
 A\left(\sum_mP_{[f_{m,1},f_{m,2},\ldots]}\right)\\
 &=1\cdot A\cdot1=A.
\end{align*}
If these sums are infinite at all, they are strongly convergent,
cf. \cite[pp.~77--78]{ref16}, so
$A\in R(U_{m',p'};m',p'=1,2,\ldots)$. This holds for any $A\in\M$, so
\[
 \M\subset R(U_{m,p};m,p=1,2,\ldots).
\]
So we have proved
\[
 \M=R(U_{m,p};m,p=1,2,\ldots).
\]

\Lemma{5.3.8}
\emph{If a factor $\M$ fulfills the conditions of \lemref{5.3.1}, then it is
a direct factor.}

\Proof Consider the $f_{m,n}$ and the $U_{m,p}$ of \lemref{5.3.7}; $m,p$
have the same finite or infinite range, $n$ has another finite of infinite
range. Take two spaces $\Hspace_1$ and $\Hspace_2$ with the same dimensions as
there are elements in the resp. ranges of $m$ and $n$; and let
$\varphi_m^0$, $m=1,2,\ldots$, and $\psi_n^0$, $n=1,2,\ldots$, be complete,
normalised orthogonal sets in $\Hspace_1$ resp. $\Hspace_2$. Then
$\varphi_m^0\otimes\psi_n^0$, $m=1,2,\ldots$, $n=1,2,\ldots$, is one in
$\Hspace_1\otimes\Hspace_2$. If we let correspond
$\varphi_m^0\otimes\psi_n^0$ to $f_{m,n}$, we obtain an isomorphism of
$\Hspace_1\otimes\Hspace_2$ and $\Hspace$.

Define in $\Hspace_1$ operators $V_{m,p}$ by
\[
 V_{m,p}\varphi_m^0=\varphi_p^0,
 \qquad
 V_{m,p}\varphi_r^0=0\quad\text{if }r\neq m.
\]
Then (cf. \lemref{2.3.3})
\[
 V_{m,p}^{(1)}\varphi_m^0\otimes\varphi_r^0
 =\varphi_p^0\otimes\varphi_r^0,
\]
\[
 V_{m,p}^{(1)}\varphi_r^0\otimes\varphi_r^0=0
 \quad\text{if }r\neq m.
\]
Thus the above spatial isomorphism, of $\Hspace_1\otimes\Hspace_2$ and
$\Hspace$ transforms $V_{m,p}^{(1)}$ into $U_{m,p}$, and thus
$R(V_{m,p}^{(1)};m,p=1,2,\ldots)$ into
$R(U_{m,p};m,p=1,2,\ldots)=\M$.

Now consider any $A_1\in\B_1$. Then
\[
 P_{[\varphi_p^0]}A_1P_{[\varphi_m^0]}
 =(A_1\varphi_m^0,\varphi_p^0)V_{m,p}
 =\alpha_{m,p}V_{m,p}.
\]
So
$P_{[\varphi_p^0]}A_1P_{[\varphi_m^0]}
\in R(V_{m',p'};m',p'=1,2,\ldots)$. But
\[
 \sum_{m,p}P_{[\varphi_p^0]}A_1P_{[\varphi_m^0]}
 =\left(\sum_pP_{[\varphi_p^0]}\right)
 A_1\left(\sum_mP_{[\varphi_m^0]}\right)
 =1\cdot A_1\cdot1=A_1,
\]
(if the sums are infinite at all, they are strongly convergent), so
\[
 A_1\in R(V_{m,p};m,p=1,2,\ldots).
\]
Thus $R(V_{m,p};m,p=1,2,\ldots)=\B_1$. Now
\hyperref[lem:2.3.6]{Lemmas 2.3.6} and \hyperref[lem:2.3.7]{2.3.7} give
$R(V_{m,p}^{(1)};m,p=1,2,\ldots)=\B_1^{(1)}$. Thus our special isomorphism
transforms $\B_1^{(1)}$ into $\M$, proving that $\M$ is a direct factor.

The results of \hyperref[lem:5.3.1]{Lemmas 5.3.1},
\hyperref[lem:5.3.2]{5.3.2} and \hyperref[lem:5.3.8]{5.3.8}
give together:

\MainTheorem{IV}
\emph{A factor $\M$ is direct if and only if it fulfills the (equivalent)
conditions of \lemref{5.3.1}.}

The condition (i) in \lemref{5.3.1} gives, together with Lemmas
\hyperref[lem:5.2.1]{5.2.1} and \hyperref[lem:5.2.2]{5.2.2}, the

\MainTheorem{V}
\emph{Every algebraic ring-isomorphism leaves invariant the notions of a
factor and of a direct factor (as applied to the rings $\M_{\mathrm I}$,
$\M_{\mathrm{II}}$ themselves).}

\SectionStart{5.4}{Direct factors and direct factorisations}
\thmref{V} makes it possible to prove the following Lemma, which belongs
naturally to \secref{3.2}.

\Lemma{5.4.1}
\emph{If $\M_1,\ldots,\M_n$ is a factorisation, and at least $n-1$ of its
factors $\M_i$ are direct, then the factorisation is a direct one too.}

\Proof For $n=1$, $R(\M_1)=\B$ means $\M_1=\B_1$, thus $\M_1$ is a direct
factor: Put $\Hspace_1$ a one dimensional space, $\Hspace_2=\Hspace$; then
the results of \secref{2.4} show that
\[
 \Hspace_1\otimes\Hspace_2=\Hspace_2=\Hspace,
 \qquad \B_2^{(2)}=\B_2=\B=\M_1.
\]
So the case $n=1$ is settled, and we must only prove the Lemma for
$n=2,3,\ldots$, if it is already established for $n-1$.

There is an $i$ such that all $\M_j$, $j\neq i$, are direct. As a permutation
of all $\M_j$'s does not matter, we may assume that this $i$ is $\neq1$, that
is, that $\M_1$ is direct.

Then $\Hspace$ is isomorphic to $\Hspace_1\otimes\Hspace_2$, so that $\M_1$
becomes $\B_1^{(1)}$. Let $\M_2,\ldots,\M_n$ become
$\N_2,\ldots,\N_n$. If $j\neq1$, then $\M_j\subset\M_1'$ so
$\N_j\subset(\B_1^{(1)})'=\B_2^{(2)}$. The correspondence
$A^{(2)}\longrightarrow A$ maps therefore the $\N_j$'s on rings
$\N_j^0\subset\B_2$ (cf. \lemref{2.3.7}). The $\N_j$'s are algebraically
(even fully) ring-isomorphic to the resp. $\N_j^0$ (cf. \lemref{2.3.4}, resp.
\cite[\S4]{ref22}). Now all $\M_2,\ldots,\M_n$ are factors, and at least
$n-2$ of them direct, so the same holds for $\N_2,\ldots,\N_n$ and, by
\thmref{V}, for $\N_2^0,\ldots,\N_n^0$.

As the $\N_j,\N_k$ commute, the $\N_2^0,\ldots,\N_n^0$ do too; as
$R(\N_2,\ldots,\N_n)=\B_2^{(2)}$,
$R(\N_2^0,\ldots,\N_n^0)=\B_2$ by \lemref{2.3.7}. Thus
$\N_2^0,\ldots,\N_n^0$ is a factorisation in $\B_2$ of $\Hspace_2$.

So our Lemma for $n-1$ applies: $\N_2^0,\ldots,\N_n^0$ is a direct
factorisation in $\B_2$ of $\Hspace_2$. So $\Hspace_2$ is isomorphic to
$\Hspace_2^*\otimes\cdots\otimes\Hspace_n^*$, and
$\N_2^0,\ldots,\N_n^0$ become
$\B_2^{*(2)},\ldots,\B_n^{*(n)}$ resp.

Now it is obvious, that $\Hspace$ is isomorphic to
$\Hspace_1\otimes\Hspace_2^*\otimes\cdots\otimes\Hspace_n^*$ and that
$\M_1,\M_2,\ldots,\M_n$ become
$\B_1^{(1)},\B_2^{*(2)},\ldots,\B_n^{*(n)}$ (the $n-1$ last ones have to be
taken now in the product
$\Hspace_1\otimes\Hspace_2^*\otimes\cdots\otimes\Hspace_n^*$!) resp. Thus
$\M_1,\ldots,\M_n$ is a direct factorisation.

Note that now \lemref{3.2.2} shows, that all $\M_i$'s are direct. Note
furthermore, that our Lemma would even then not have been obvious, if all $n$
factors $\M_i$ had been known to be direct. $n-1$ can not be replaced in our
Lemma by $n-2$: Because this would give for every factor $\M$, that the
factorisation $\M,\M'$ ($n=2$) is direct, and thus $\M$ is direct; and then by
our Lemma every factorisation $\M_1,\ldots,\M_n$ would be direct. Thus the
answer to both parts of \probref{2} at the end of \secref{3.2} would be
affirmative, which is not the case (cf. there and \secref{8.4} and
\secref{13.2}).
% --- End of parts/02b-chapters-iii-v.tex ---
\PartHeading{II}{The General Problem}

\ChapterHeading{6}{VI}{Relative Dimensionality}

\SectionStart{6.1}{Equality of relative dimensions. Additivity. Equivalence theorem}
In what follows \(\M\) will always denote a factor in \(\B\) of the space
\(\Hspace\). Thus \(\M'\) is a factor too.

\Definition{6.1.1} We write \(\mathfrak M\sim\mathfrak N\ (\cdots\M)\),
and for \(E=P_{\mathfrak M}\), \(F=P_{\mathfrak N}\),
\(E\sim F\ (\cdots\M)\), if a partially isometric \(U\in\M\) exists, the
initial and final sets of which are \(\mathfrak M\) resp. \(\mathfrak N\).
(If no misunderstanding is possible we will omit the \((\cdots\M)\).) We say
that \(\mathfrak M,\mathfrak N\) have \emph{the same relative dimension}
(with respect to \(\M\)).

In discussing this notion \(\sim\), we have always the choice to speak about
the linear closed sets \(\mathfrak M\), or their projections
\(E=P_{\mathfrak M}\). We will mostly use the first mentioned terminology,
remembering however the equivalence of the two.

\Lemma{6.1.1} {\itshape
\(\mathfrak M\sim\mathfrak N\) implies
\(\mathfrak M,\mathfrak N\mathbin{\eta}\M\). If
\(\mathfrak M,\mathfrak N,\mathfrak P\mathbin{\eta}\M\) then we have:
\begin{gather*}
  \mathfrak M\sim\mathfrak M,\\
  \mathfrak M\sim\mathfrak N\ \text{implies}\
    \mathfrak N\sim\mathfrak M,\\
  \mathfrak M\sim\mathfrak N,\ \mathfrak N\sim\mathfrak P\
    \text{ imply }\mathfrak M\sim\mathfrak P.
\end{gather*}
}

\Proof If \(U\) is the partially isometric operator in question, we have
\(P_{\mathfrak M}=U^*U\in\M\), \(P_{\mathfrak N}=UU^*\in\M\), and so
\(\mathfrak M,\mathfrak N\mathbin{\eta}\M\).

\(\mathfrak M\sim\mathfrak M\) results by putting
\(U=P_{\mathfrak M}\); \(\mathfrak M\sim\mathfrak N\) gives
\(\mathfrak N\sim\mathfrak M\) by replacing \(U\) by \(U^*\);
\(\mathfrak M\sim\mathfrak N\) with \(U\) and
\(\mathfrak N\sim\mathfrak P\) with \(V\) gives
\(\mathfrak M\sim\mathfrak P\) by considering \(VU\).

Thus \(\sim\) is naturally applied to linear closed sets
\(\mathfrak M\mathbin{\eta}\M\), and has there the properties of an
equivalence. We may say: Two linear, closed sets \(\mathfrak M,\mathfrak N\)
which are invariant under all unitary elements of \(\M'\), are of the same
dimension in the usual sense of the word, if an isometric mapping \(U\) of
\(\mathfrak M\) on \(\mathfrak N\) exists, that is a partially isometric
operator \(U\) with the initial and final sets \(\mathfrak M\) resp.
\(\mathfrak N\). But they have the same relative dimension, if even the
mapping \(U\) can be chosen so as to be invariant under all unitary elements
of \(\M'\). (Cf. \defref{4.2.1} and \lemref{4.2.1}.)

We now prove two Lemmas, the first of which expresses an additivity property
of our \(\sim\), while the second one is the analogon of the Cantor--Bernstein
``equivalence theorem'' of general set-theory. (In fact, it is a special case
of the Kuratowski--Tarski generalisation of that theorem.)

\Lemma{6.1.2} {\itshape
Let \(\mathfrak M_1,\mathfrak M_2,\ldots\) and
\(\mathfrak N_1,\mathfrak N_2,\ldots\) be two (finite or infinite) sequences
of the same length; let the \(\mathfrak M_i\) be mutually orthogonal and let
the \(\mathfrak N_i\) be mutually orthogonal. If we have
\(\mathfrak M_i\sim\mathfrak N_i\) for all \(i\), then
\[
  \clos{\mathfrak M_1,\mathfrak M_2,\ldots}
  \sim
  \clos{\mathfrak N_1,\mathfrak N_2,\ldots}.
\]
}

\Proof Let \(U_i\in\M\) be the partially isometric mapping of
\(\mathfrak M_i\) on \(\mathfrak N_i\), then \(U=\sum_iU_i\) does the
desired thing for \(\clos{\mathfrak M_1,\mathfrak M_2,\ldots}\) and
\(\clos{\mathfrak N_1,\mathfrak N_2,\ldots}\) by \lemref{4.3.3}.

\Lemma{6.1.3} {\itshape
If \(\mathfrak M\sim\mathfrak N'\subset\mathfrak N\) and
\(\mathfrak N\sim\mathfrak M'\subset\mathfrak M\), then
\(\mathfrak M\sim\mathfrak N\).
}

\Proof Our relations imply
\(\mathfrak M,\mathfrak N,\mathfrak M',\mathfrak N'\mathbin{\eta}\M\).
Now let \(U\in\M\) be the partially isomorphic mapping of \(\mathfrak N\)
on \(\mathfrak M'\); it maps \(\mathfrak N'\subset\mathfrak N\) on some
\(\mathfrak M''\subset\mathfrak M'\). Then consideration of
\(UP_{\mathfrak N'}\) proves that \(\mathfrak N'\sim\mathfrak M''\), and
so \(\mathfrak M\sim\mathfrak M''\).

Let \(V\in\M\) be the partially isomorphic mapping of \(\mathfrak M\) on
\(\mathfrak M''\). Form for every \(n=0,1,2,\ldots\) the \(V^n\)-images
of \(\mathfrak M\) and of \(\mathfrak M'\), and call them
\(\mathfrak M^{(2n)}\) resp. \(\mathfrak M^{(2n+1)}\)
(\(\mathfrak M,\mathfrak M',\mathfrak M''\) coincide with
\(\mathfrak M^{(0)},\mathfrak M^{(1)},\mathfrak M^{(2)}\) resp.). We have
\(\mathfrak M^{(0)}\supset\mathfrak M^{(1)}\supset\mathfrak M^{(2)}\),
thus (apply \(V^n\)!)
\[
 \mathfrak M^{(2n)}\supset\mathfrak M^{(2n+1)}
 \supset\mathfrak M^{(2n+2)},
\]
therefore
\[
 \mathfrak M^{(0)}\supset\mathfrak M^{(1)}\supset
 \mathfrak M^{(2)}\supset\cdots.
\]

\(V^n\) is partially isometric with the initial and final sets
\(\mathfrak M^{(0)}\) resp. \(\mathfrak M^{(2n)}\),
\(V^nP_{\mathfrak M^{(1)}}\) similarly with \(\mathfrak M^{(1)}\) and
\(\mathfrak M^{(2n+1)}\). Thus all
\(\mathfrak M^{(p)}\mathbin{\eta}\M\), \(n=0,1,2,\ldots\), and
\[
 \mathfrak M^{(0)}\sim\mathfrak M^{(2n)},\qquad
 \mathfrak M^{(1)}\sim\mathfrak M^{(2n+1)},
\]
that is: \(\mathfrak M^{(p)}\sim\mathfrak M\) or \(\mathfrak M'\) if
\(p\) is even resp. odd.

\(V\) maps \(\mathfrak M^{(p)}\) on \(\mathfrak M^{(p+2)}\),
\(\mathfrak M^{(p+1)}\) on \(\mathfrak M^{(p+3)}\), therefore
\(\mathfrak M^{(p)}-\mathfrak M^{(p+1)}\) on
\(\mathfrak M^{(p+2)}-\mathfrak M^{(p+3)}\). Using
\(V(P_{\mathfrak M^{(p)}}-P_{\mathfrak M^{(p+1)}})\) shows therefore,
that
\[
 \mathfrak M^{(p)}-\mathfrak M^{(p+1)}
 \sim
 \mathfrak M^{(p+2)}-\mathfrak M^{(p+3)}.
\]
Now apply \lemref{6.1.2} to the sequences
\begin{multline*}
 \mathfrak M^{(0)}\cdot\mathfrak M^{(1)}\cdot\mathfrak M^{(2)}\cdots,
 \ \mathfrak M^{(0)}-\mathfrak M^{(1)},
 \ \mathfrak M^{(1)}-\mathfrak M^{(2)},
 \ \mathfrak M^{(2)}-\mathfrak M^{(3)},\\
 \mathfrak M^{(3)}-\mathfrak M^{(4)},\ldots,
\end{multline*}
and
\begin{multline*}
 \mathfrak M^{(0)}\cdot\mathfrak M^{(1)}\cdot\mathfrak M^{(2)}\cdots,
 \ \mathfrak M^{(2)}-\mathfrak M^{(3)},
 \ \mathfrak M^{(1)}-\mathfrak M^{(2)},
 \ \mathfrak M^{(4)}-\mathfrak M^{(5)},\\
 \mathfrak M^{(3)}-\mathfrak M^{(4)},\ldots.
\end{multline*}
It gives \(\mathfrak M^{(0)}\sim\mathfrak M^{(1)}\), that is
\(\mathfrak M\sim\mathfrak M'\), and thus
\(\mathfrak M\sim\mathfrak N\).

\SectionStartTOC{6.2}
 {}
 {\texorpdfstring{$(\operatorname{Range}X)\sim(\operatorname{Range}X^*)$}{(Range X) similar to (Range X*)}. Comparability theorem}
The following Lemma is an important means of establishing equality of
relative dimension.

\Lemma{6.2.1} {\itshape
If \(X\) is a linear, closed operator with an everywhere dense domain, and
\(X\mathbin{\eta}\M\), then
\[
 \clos{\operatorname{Range}X}\sim\clos{\operatorname{Range}X^*}.
\]
}

\Proof Form the canonical decomposition of \(X\) in the sense of
\defref{4.4.1}. \(X\mathbin{\eta}\M\) implies by \lemref{4.4.1}
\(W\in\M\). \(W\) is partially isometric its initial and final sets being
\(\clos{\operatorname{Range}X^*}\) resp.
\(\clos{\operatorname{Range}X}\). So
\[
 \clos{\operatorname{Range}X^*}\sim\clos{\operatorname{Range}X},
 \qquad
 \clos{\operatorname{Range}X}\sim\clos{\operatorname{Range}X^*}.
\]

We now prove in two steps the analogue of the ``comparability theorem'' of
general set-theory.

\Lemma{6.2.2} {\itshape
If \(\mathfrak M,\mathfrak N\mathbin{\eta}\M\) and \(\ne0\), then two
\(\mathfrak P,\mathfrak Q\ne0\),
\(\mathfrak P\subset\mathfrak M\), \(\mathfrak Q\subset\mathfrak N\),
with \(\mathfrak P\sim\mathfrak Q\) exist.
}

\Proof Choose an \(f\in\mathfrak M\), \(f\ne0\). Then
\(E_f^{\M}\in\M'\) and \(\ne0\), \(P_{\mathfrak N}\in\M\) and
\(\ne0\), so by the Corollary to \thmref{III}.
\[
 P_{\mathfrak M_f^{\M}\cdot\mathfrak N}
   =P_{\mathfrak N}\cdot E_f^{\M}\ne0,
 \qquad
 \mathfrak M_f^{\M}\cdot\mathfrak N\ne(0).
\]
Choose \(g\in\mathfrak M_f^{\M}\cdot\mathfrak N\), \(g\ne0\),
multiplying with \(1/\lVert g\rVert\) makes \(\lVert g\rVert=1\). As
\(g\in\mathfrak M_f^{\M}\), therefore an \(A\in\M\) with
\(\lVert g-Af\rVert<1\) exists. So
\(\lVert P_{\mathfrak N}g-P_{\mathfrak N}Af\rVert<1\). Now
\(P_{\mathfrak N}g=g\), \(P_{\mathfrak M}f=f\) have the effect that
\[
 \lVert g-P_{\mathfrak N}AP_{\mathfrak M}f\rVert<1.
\]
As \(\lVert g\rVert=1\) this implies
\(P_{\mathfrak N}AP_{\mathfrak M}f\ne0\),
\(P_{\mathfrak N}AP_{\mathfrak M}\ne0\).

Now apply \lemref{6.2.1}:
\[
 \clos{\operatorname{Range}P_{\mathfrak N}AP_{\mathfrak M}}
 \sim
 \clos{\operatorname{Range}(P_{\mathfrak N}AP_{\mathfrak M})^*}.
\]
As \(P_{\mathfrak N}AP_{\mathfrak M}\ne0\) and so
\((P_{\mathfrak N}AP_{\mathfrak M})^*\ne0\), both above sets are
\(\ne(0)\). Furthermore
\[
 \clos{\operatorname{Range}P_{\mathfrak N}AP_{\mathfrak M}}
 \subset
 \clos{\operatorname{Range}P_{\mathfrak N}}=\mathfrak N,
\]
\[
 \clos{\operatorname{Range}(P_{\mathfrak N}AP_{\mathfrak M})^*}
 =\clos{\operatorname{Range}P_{\mathfrak M}A^*P_{\mathfrak N}}
 \subset
 \clos{\operatorname{Range}P_{\mathfrak M}}=\mathfrak M.
\]
So
\[
 \mathfrak P=
 \clos{\operatorname{Range}(P_{\mathfrak N}AP_{\mathfrak M})^*},
 \qquad
 \mathfrak Q=
 \clos{\operatorname{Range}P_{\mathfrak N}AP_{\mathfrak M}}
\]
meet our requirements.

\Lemma{6.2.3} {\itshape
If \(\mathfrak M,\mathfrak N\mathbin{\eta}\M\), then either
\(\mathfrak M\sim\mathfrak M'\subset\mathfrak N\) or
\(\mathfrak N\sim\mathfrak N'\subset\mathfrak M\).
}

\Proof Let \(\Omega\) be the first uncountable Cantor ordinal number. Define
for all \(\alpha<\Omega\) a pair \(\mathfrak M_\alpha,\mathfrak N_\alpha\)
as follows: All \(\mathfrak M_\alpha,\mathfrak N_\alpha\) will be linear,
closed sets, \(\mathbin{\eta}\M\), \(\ne(0)\), and
\[
 \mathfrak M_\alpha\subset\mathfrak M,\qquad
 \mathfrak N_\alpha\subset\mathfrak N,\qquad
 \mathfrak M_\alpha\sim\mathfrak N_\alpha.
\]
If \(\alpha<\Omega\) and all \(\mathfrak M_\beta,\mathfrak N_\beta\),
\(\beta<\alpha\) are already defined, then form
\[
 \mathfrak M-\clos{\mathfrak M_\beta,\ \beta<\alpha}
 \quad\hbox{and}\quad
 \mathfrak N-\clos{\mathfrak N_\beta,\ \beta<\alpha}.
\]
As \(\mathfrak M,\mathfrak N,\mathfrak M_\beta,\mathfrak N_\beta\) all
\(\mathbin{\eta}\M\), these two sets are \(\mathbin{\eta}\M\) too. If
both are \(\ne0\), then there exist two \(\mathfrak P,\mathfrak Q\ne(0)\)
with
\[
 \mathfrak P\subset
   \mathfrak M-\clos{\mathfrak M_\beta,\ \beta<\alpha},
 \qquad
 \mathfrak Q\subset
   \mathfrak N-\clos{\mathfrak N_\beta,\ \beta<\alpha},
 \qquad
 \mathfrak P\sim\mathfrak Q.
\]
Let them be \(\mathfrak M_\alpha=\mathfrak P\),
\(\mathfrak N_\alpha=\mathfrak Q\) for some such pair
\(\mathfrak P,\mathfrak Q\). As we see
\(\mathfrak M_\alpha,\mathfrak N_\alpha\ne(0)\),
\(\mathfrak M_\alpha\subset\mathfrak M\),
\(\mathfrak N_\alpha\subset\mathfrak N\),
\(\mathfrak M_\alpha\sim\mathfrak N_\alpha\) are maintained. If the above
condition is not fulfilled, let \(\mathfrak M_\alpha,\mathfrak N_\alpha\)
(and all \(\mathfrak M_\gamma,\mathfrak N_\gamma\),
\(\gamma\geq\alpha\)) be undefined.

Let the first \(\alpha\) for which \(\mathfrak M_\alpha,\mathfrak N_\alpha\)
are undefined, be \(\bar\alpha\leq\Omega\). If
\(\bar\alpha<\Omega\), we must have
\[
 \mathfrak M-\clos{\mathfrak M_\alpha,\ \alpha<\bar\alpha}=(0)
 \quad\hbox{or}\quad
 \mathfrak N-\clos{\mathfrak N_\alpha,\ \alpha<\bar\alpha}=(0).
\]
Consider now all \(\mathfrak M_\alpha,\mathfrak N_\alpha\),
\(\alpha<\bar\alpha\). Always \(\mathfrak M_\alpha\ne(0)\). If
\(\alpha,\beta<\bar\alpha\) assume \(\alpha>\beta\). Then
\[
 \mathfrak M_\alpha\subset
 \mathfrak M-\clos{\mathfrak M_{\beta'},\ \beta'<\alpha},
\]
so \(\mathfrak M_\alpha\) is orthogonal to \(\mathfrak M_\beta\). This is
symmetric in \(\alpha,\beta\), so it holds for \(\alpha<\beta\) too, that is
whenever \(\alpha\ne\beta\). Thus all \(\mathfrak M_\alpha\),
\(\alpha<\bar\alpha\), are \(\ne(0)\) and mutually orthogonal. The same is
of course true for the \(\mathfrak N_\alpha\),
\(\alpha<\bar\alpha\).

Choose an \(f_\alpha\in\mathfrak M_\alpha\), \(f_\alpha\ne0\), multiplying
with \(1/\lVert f_\alpha\rVert\) makes \(\lVert f_\alpha\rVert=1\). So the
\(f_\alpha\), \(\alpha<\bar\alpha\) form a normalised orthogonal set. It
must therefore be finite or countably infinite (cf. \cite[p.~66]{ref16}), thus
excluding \(\bar\alpha=\Omega\).

If \(\bar\alpha\) is finite, we have \(\bar\alpha=n+1\),
\(n=0,1,2,\ldots\), and \(\alpha\) runs over \(1,\ldots,n\). If
\(\bar\alpha\) is infinite, then \(\alpha<\bar\alpha\) has the same aleph as
\(\alpha<\omega\) and may therefore (by a re-indexing of the
\(\mathfrak M_\alpha\)) be replaced by it; so we may assume that \(\alpha\)
runs over \(1,2,\ldots\).

Writing \(i\) for \(\alpha\), we see that we have a finite or infinite
sequence of pairs \(\mathfrak M_i,\mathfrak N_i\), \(i=1,2,\ldots\); where
the \(\mathfrak M_i\) are mutually orthogonal, the \(\mathfrak N_i\) are
mutually orthogonal; \(\mathfrak M_i\subset\mathfrak M\),
\(\mathfrak N_i\subset\mathfrak N\); \(\mathfrak M_i\sim\mathfrak N_i\)
and either
\[
 \clos{\mathfrak M_i,\ i=1,2,\ldots}=\mathfrak M
 \quad\hbox{or}\quad
 \clos{\mathfrak N_i,\ i=1,2,\ldots}=\mathfrak N.
\]
Now \lemref{6.1.2} applies:
\[
 \clos{\mathfrak M_i,\ i=1,2,\ldots}
 \sim
 \clos{\mathfrak N_i,\ i=1,2,\ldots}.
\]
Thus if we define
\(\mathfrak N'=\clos{\mathfrak N_i,\ i=1,2,\ldots}\) resp.
\(\mathfrak M'=\clos{\mathfrak M_i,\ i=1,2,\ldots}\) we will have
\(\mathfrak M\sim\mathfrak N'\subset\mathfrak N\) resp.
\(\mathfrak N\sim\mathfrak M'\subset\mathfrak M\). So the proof is
completed.

\SectionStart{6.3}{The ordering of relative dimensions. Theorem VI}
Now we are in the position to define:

\Definition{6.3.1} We write \(\mathfrak M\lesssim\mathfrak N\) or
\(\mathfrak N\gtrsim\mathfrak M\), if
\(\mathfrak M\sim\mathfrak N'\subset\mathfrak N\). We write
\(\mathfrak M<\mathfrak N\) or \(\mathfrak N>\mathfrak M\), if
\(\mathfrak M\lesssim\mathfrak N\) but not
\(\mathfrak M\sim\mathfrak N\). (As to replacing
\(\mathfrak M,\mathfrak N\) by \(E=P_{\mathfrak M}\),
\(F=P_{\mathfrak N}\), and the explanatory symbol \((\cdots\M)\), cf.
\defref{6.1.1}.)

And we can prove:

\Lemma{6.3.1} {\itshape
If \(\mathfrak M,\mathfrak N,\mathfrak P,\mathfrak Q\mathbin{\eta}\M\),
then the following statements hold:
\begin{enumerate}[label=(\roman*)]
 \item \(\mathfrak M\lesssim\mathfrak N\) means
   \(\mathfrak M<\mathfrak N\) or \(\mathfrak M\sim\mathfrak N\).
 \item If \(\mathfrak M\sim\mathfrak P\),
   \(\mathfrak N\sim\mathfrak Q\), then
   \(\mathfrak M<\mathfrak N\) is equivalent to
   \(\mathfrak P<\mathfrak Q\).
 \item One and only one of the three relations
   \(\mathfrak M\mathrel{\substack{>\\[-0.7ex]\sim\\[-0.7ex]<}}\mathfrak N\)
   holds.
 \item \(\mathfrak M<\mathfrak N\), \(\mathfrak N<\mathfrak P\) imply
   \(\mathfrak M<\mathfrak P\).
\end{enumerate}
}

\Proof We prove the statements in a somewhat changed order.

Ad (i): Obvious.

Ad (iv): \(\mathfrak M\lesssim\mathfrak N\),
\(\mathfrak N\lesssim\mathfrak P\) imply
\(\mathfrak M\sim\mathfrak N'\subset\mathfrak N\),
\(\mathfrak N\sim\mathfrak P'\subset\mathfrak P\). If \(U\in\M\) is
the isometric mapping of \(\mathfrak N\) on \(\mathfrak P'\), and
\(\mathfrak P''\) the \(U\)-image of \(\mathfrak N'\), then we have
\[
 \mathfrak M\sim\mathfrak N'\sim\mathfrak P''
 \subset\mathfrak P'\subset\mathfrak P,
\]
so \(\mathfrak M\lesssim\mathfrak P\).

Now if we had \(\mathfrak M\sim\mathfrak P\), then we could write
\(\mathfrak P\sim\mathfrak M\subset\mathfrak M\), that is
\(\mathfrak P\lesssim\mathfrak M\); together with
\(\mathfrak M\lesssim\mathfrak N\) this gives
\(\mathfrak P\lesssim\mathfrak N\). \(\mathfrak N\lesssim\mathfrak P\)
and \(\mathfrak P\lesssim\mathfrak N\) gives by \lemref{6.1.3}
\(\mathfrak N\sim\mathfrak P\), contradicting
\(\mathfrak N<\mathfrak P\). Thus we must have
\(\mathfrak M<\mathfrak P\).

Ad (ii): Assume \(\mathfrak M<\mathfrak N\). As
\(\mathfrak M\sim\mathfrak P\), \(\mathfrak N\sim\mathfrak Q\), we can
write \(\mathfrak P\lesssim\mathfrak N\),
\(\mathfrak M\lesssim\mathfrak N\),
\(\mathfrak N\lesssim\mathfrak Q\), implying
\(\mathfrak P\lesssim\mathfrak Q\) (cf. the beginning of the proof of
(iv)). \(\mathfrak P\sim\mathfrak Q\) would imply
\(\mathfrak M\sim\mathfrak P\), \(\mathfrak P\sim\mathfrak Q\),
\(\mathfrak Q\sim\mathfrak N\) so \(\mathfrak M\sim\mathfrak N\),
contradicting \(\mathfrak M<\mathfrak N\). Thus we must have
\(\mathfrak P<\mathfrak Q\). In the same way
\(\mathfrak P<\mathfrak Q\) implies \(\mathfrak M<\mathfrak N\).

Ad (iii): By \lemref{6.2.3} we have
\(\mathfrak M\lesssim\mathfrak N\) or
\(\mathfrak N\lesssim\mathfrak M\), by (i) this means
\(\mathfrak M\mathrel{\substack{>\\[-0.7ex]\sim\\[-0.7ex]<}}\mathfrak N\).
\(\sim\) and \(>\) at once imply \(\mathfrak M<\mathfrak M\) by (ii);
\(<\) and \(\sim\) similarly; \(<\) and \(>\) imply again
\(\mathfrak M<\mathfrak M\) by (iv). Now
\(\mathfrak M<\mathfrak M\) is impossible as
\(\mathfrak M\sim\mathfrak M\). So precisely one of the three relations
holds.

Summing up Lemmas \hyperref[lem:6.1.1]{6.1.1} and
\hyperref[lem:6.3.1]{6.3.1} we have:

\MainTheorem{VI} {\itshape
The notions \(\mathfrak M\sim\mathfrak N\) and
\(\mathfrak M<\mathfrak N\) (that is \(\mathfrak N>\mathfrak M\)) for
\(\mathfrak M,\mathfrak N\mathbin{\eta}\M\) (cf. \defref{6.1.1} and
\hyperref[def:6.3.1]{6.3.1}) have the properties of an equivalence and of a complete
ordering.

Further essential properties are given in Lemmas
\hyperref[lem:6.1.2]{6.1.2} and \hyperref[lem:6.3.1]{6.3.1}.
}

Some immediate consequences:

\Lemma{6.3.2} {\itshape
Assume \(\mathfrak M,\mathfrak N\mathbin{\eta}\M\). Then we have:
\begin{enumerate}[label=(\roman*)]
 \item \(\mathfrak M\subset\mathfrak N\) implies
   \(\mathfrak M\lesssim\mathfrak N\).
 \item Always \((0)\lesssim\mathfrak M\lesssim\Hspace\).
 \item \(\mathfrak M\sim(0)\) implies \(\mathfrak M=(0)\).
\end{enumerate}
}

\Proof Ad (i): Write
\(\mathfrak M\sim\mathfrak M\subset\mathfrak N\).

Ad (ii): By \((0)\subset\mathfrak M\subset\Hspace\) and (i).

Ad (iii): \(\mathfrak M\) is the image of \((0)\) by some linear \(U\) so
\(\mathfrak M=(0)\).

\ChapterHeading{7}{VII}{Finite and Infinite Sets}

\SectionStart{7.1}{Division. Finiteness and infiniteness}
We continue along the lines which are familiar from set-theory:

\Definition{7.1.1} An \(\mathfrak M\mathbin{\eta}\M\) is \emph{finite} if
\(\mathfrak M\sim\mathfrak N\subset\mathfrak M\) implies
\(\mathfrak N=\mathfrak M\); it is \emph{infinite}, if this is not the case.
(As to replacing \(\mathfrak M\) by \(E=P_{\mathfrak M}\), and the
explanatory symbol \((\cdots\M)\), cf. \defref{6.1.1}.)

Some immediate consequences:

\Lemma{7.1.1} {\itshape
Assume \(\mathfrak M,\mathfrak N\mathbin{\eta}\M\). Then we have:
\begin{enumerate}[label=(\roman*)]
 \item If \(\mathfrak M\lesssim\mathfrak N\) and \(\mathfrak M\) is
   infinite, then \(\mathfrak N\) is too; if \(\mathfrak N\) is finite,
   \(\mathfrak M\) is too.
 \item \((0)\) is finite.
 \item If infinite \(\mathfrak M\)'s exist at all, then \(\Hspace\) is
   infinite.
\end{enumerate}
}

\Proof Ad (i): The second statement follows from the first one, so we
consider the first one only. As \(\mathfrak M\) is infinite, a
\(\mathfrak P\) with \(\mathfrak M\sim\mathfrak P\subsetneqq\mathfrak M\)
exists. Then
\[
 \clos{\mathfrak M,\mathfrak N-\mathfrak M}
 \sim
 \clos{\mathfrak P,\mathfrak N-\mathfrak M}
 \subsetneqq
 \clos{\mathfrak M,\mathfrak N-\mathfrak M},
\]
\[
 \mathfrak N\sim\clos{\mathfrak P,\mathfrak N-\mathfrak M}
 \subsetneqq\mathfrak N.
\]
So \(\mathfrak N\) is infinite too.

Ad (ii): \(\mathfrak N\subset(0)\) alone implies
\(\mathfrak N=(0)\).

Ad (iii): If \(\mathfrak M\) is infinite,
\(\mathfrak M\lesssim\Hspace\) implies by (i), that \(\Hspace\) is infinite
too.

We now begin the preparations for a quantitative comparison of sets
\(\mathfrak M,\mathfrak N\mathbin{\eta}\M\).

\Lemma{7.1.2} {\itshape
Assume \(\mathfrak M,\mathfrak N\mathbin{\eta}\M\),
\(\mathfrak M\ne(0)\). Then there exists a finite or infinite sequence
\(\mathfrak P_1,\mathfrak P_2,\ldots\) (its length can be 0 too) and a
\(\mathfrak Q\), such that we have:

\(\mathfrak P_1,\mathfrak P_2,\ldots,\mathfrak Q\) are mutually orthogonal
linear, closed sets, all \(\mathbin{\eta}\M\),
\[
 \clos{\mathfrak P_1,\mathfrak P_2,\ldots,\mathfrak Q}=\mathfrak N,
 \qquad
 \mathfrak M\sim\mathfrak P_1\sim\mathfrak P_2\sim\cdots,
 \qquad
 \mathfrak Q<\mathfrak M,
\]
and if the sequence is infinite, we may even assume \(\mathfrak Q=(0)\).
}

\Proof Let \(\Omega\) be the first uncountable Cantor ordinal number. Define
for all \(\alpha<\Omega\) a \(\mathfrak P_\alpha\) as follows: All
\(\mathfrak P_\alpha\) will be linear, closed sets,
\[
 \mathfrak P_\alpha\mathbin{\eta}\M,
 \qquad \mathfrak P_\alpha\subset\mathfrak N,
 \qquad \mathfrak M\sim\mathfrak P_\alpha.
\]
If \(\alpha<\Omega\) and all \(\mathfrak P_\beta\), \(\beta<\alpha\) are
already defined, then form
\(\mathfrak N-\clos{\mathfrak P_\beta,\ \beta<\alpha}\). As
\(\mathfrak N,\mathfrak P_\beta\) all \(\mathbin{\eta}\M\), this set is
\(\mathbin{\eta}\M\) too. If it is \(\gtrsim\mathfrak M\), then there
exists a \(\mathfrak P\mathbin{\eta}\M\) with
\[
 \mathfrak P\subset
 \mathfrak N-\clos{\mathfrak P_\beta,\ \beta<\alpha},
 \qquad \mathfrak M\sim\mathfrak P.
\]
Let then \(\mathfrak P_\alpha\) be some such \(\mathfrak P\). If the above
condition is not fulfilled, let \(\mathfrak P_\alpha\) (and all
\(\mathfrak P_\gamma\), \(\gamma\geq\alpha\)) be undefined.

Let the first \(\alpha\) for which \(\mathfrak P_\alpha\) is undefined, be
\(\bar\alpha\leq\Omega\). If \(\bar\alpha<\Omega\), we must have
\[
 \mathfrak N-\clos{\mathfrak P_\alpha,\ \alpha<\bar\alpha}<\mathfrak M.
\]
We see in the same way as in the proof of \lemref{6.2.3}, that the
\(\mathfrak P_\alpha\), \(\alpha<\bar\alpha\), are mutually orthogonal; and
as \(\mathfrak P_\alpha\sim\mathfrak M\ne(0)\), therefore
\(\mathfrak P_\alpha\ne(0)\). This implies, as above, that their set is finite
or countably infinite, thus excluding \(\bar\alpha=\Omega\).

If \(\bar\alpha\) is finite, we have \(\bar\alpha=n+1\),
\(n=0,1,2,\ldots\), and \(\alpha\) runs over \(1,\ldots,n\). If
\(\bar\alpha\) is infinite, then \(\alpha<\bar\alpha\) has the same aleph as
\(\alpha<\omega\), and may therefore (by a re-indexing of the
\(\mathfrak P_\alpha\)) be replaced by it; so we may assume that \(\alpha\)
runs over \(1,2,\ldots\).

Writing \(i\) for \(\alpha\) we have a finite or infinite sequence
\(\mathfrak P_i\), \(i=1,2,\ldots\); where the \(\mathfrak P_i\) are mutually
orthogonal; \(\mathfrak M\sim\mathfrak P_1\sim\mathfrak P_2\sim\cdots\);
and
\[
 \mathfrak Q=\mathfrak N-\clos{\mathfrak P_1,\mathfrak P_2,\ldots}
 <\mathfrak M.
\]
Thus \(\mathfrak P_1,\mathfrak P_2,\ldots,\mathfrak Q\) are mutually
orthogonal and
\[
 \mathfrak N=\clos{\mathfrak P_1,\mathfrak P_2,\ldots,\mathfrak Q}.
\]

So the case of a finite sequence is completely settled. In the case of an
infinite sequence we must still get rid of \(\mathfrak Q\). Assume therefore
that the \(\mathfrak P_1,\mathfrak P_2,\ldots\) form an infinite sequence.

We have \(\mathfrak Q<\mathfrak M\sim\mathfrak P_1\), so
\(\mathfrak Q\sim\mathfrak Q'\subset\mathfrak P_1\). Apply
\lemref{6.1.2} to the two sequences
\(\mathfrak Q,\mathfrak P_1,\mathfrak P_2,\ldots\) and
\(\mathfrak Q',\mathfrak P_2,\mathfrak P_3,\ldots\). It gives
\begin{align*}
 \mathfrak N
 &=\clos{\mathfrak Q,\mathfrak P_1,\mathfrak P_2,\ldots}
 \sim\clos{\mathfrak Q',\mathfrak P_2,\mathfrak P_3,\ldots}\\
 &\subset\clos{\mathfrak P_1,\mathfrak P_2,\ldots}
 \subset\clos{\mathfrak Q,\mathfrak P_1,\mathfrak P_2,\ldots}
 =\mathfrak N.
\end{align*}
So
\[
 \mathfrak N\lesssim\clos{\mathfrak P_1,\mathfrak P_2,\ldots}
 \lesssim\mathfrak N,
 \qquad
 \clos{\mathfrak P_1,\mathfrak P_2,\ldots}\sim\mathfrak N.
\]
Let \(U\in\M\) be an isometric mapping of
\(\clos{\mathfrak P_1,\mathfrak P_2,\ldots}\) on \(\mathfrak N\),
\(\mathfrak P_i'\) the \(U\)-image of \(\mathfrak P_i\). Then the
\(\mathfrak P_i'\) are mutually orthogonal and
\(\clos{\mathfrak P_1',\mathfrak P_2',\ldots}=\mathfrak N\). Using
\(UP_{\mathfrak P_i}\) shows \(\mathfrak P_i\sim\mathfrak P_i'\); so
\(\mathfrak M\sim\mathfrak P_1'\sim\mathfrak P_2'\sim\cdots\). Thus we
can replace
\[
 \mathfrak P_1,\mathfrak P_2,\ldots,\mathfrak Q
 \quad\hbox{by}\quad
 \mathfrak P_1',\mathfrak P_2',\ldots,(0).
\]
This completes the proof.

\Lemma{7.1.3} {\itshape
Assume \(\mathfrak M,\mathfrak N\mathbin{\eta}\M\),
\(\mathfrak M\ne(0)\), \(\mathfrak N\) finite. Then the sequence
\(\mathfrak P_1,\mathfrak P_2,\ldots\) from \lemref{7.1.2} must be finite,
and its length is uniquely determined by \(\mathfrak M,\mathfrak N\). Call
this number
\[
 \left[\frac{\mathfrak N}{\mathfrak M}\right]=0,1,2,\ldots.
\]
}

\Proof If the sequence \(\mathfrak P_1,\mathfrak P_2,\ldots\) was infinite,
we could apply \lemref{6.1.2} to
\(\mathfrak P_1,\mathfrak P_2,\ldots\) and
\(\mathfrak P_2,\mathfrak P_3,\ldots\), giving
\[
 \mathfrak N=\clos{\mathfrak P_1,\mathfrak P_2,\ldots}
 \sim\clos{\mathfrak P_2,\mathfrak P_3,\ldots}
 \subsetneqq\clos{\mathfrak P_1,\mathfrak P_2,\ldots}=\mathfrak N
\]
(\(\subsetneqq\) because \(\mathfrak P_1\ne(0)\), as
\(\mathfrak P_1\sim\mathfrak M\ne(0)\)), thus \(\mathfrak N\) infinite,
contradicting the assumption.

Assume now we had two representations
\[
 \mathfrak N
 =\clos{\mathfrak P_1,\ldots,\mathfrak P_m,\mathfrak Q}
 =\clos{\mathfrak P_1',\ldots,\mathfrak P_n',\mathfrak Q'}
\]
in the sense of \lemref{7.1.2}, with \(m\ne n\). Owing to the symmetry we
may assume \(m<n\), \(m+1\leq n\).

Now \(\mathfrak Q<\mathfrak M\sim\mathfrak P_{m+1}'\), so
\(\mathfrak Q\sim\mathfrak Q^*\subset\mathfrak P_{m+1}'\), and
\(\mathfrak Q^*\ne\mathfrak P_{m+1}'\). So by \lemref{6.1.2},
\begin{align*}
 \mathfrak N
 &=\clos{\mathfrak P_1,\ldots,\mathfrak P_m,\mathfrak Q}
 \sim\clos{\mathfrak P_1',\ldots,\mathfrak P_m',\mathfrak Q^*}\\
 &\subsetneqq
 \clos{\mathfrak P_1',\ldots,\mathfrak P_m',\mathfrak P_{m+1}'}
 \subset
 \clos{\mathfrak P_1',\ldots,\mathfrak P_n',\mathfrak Q'}
 =\mathfrak N,
\end{align*}
and again \(\mathfrak N\) would be infinite, contradicting the assumption.

\SectionStart{7.2}{Properties of infinite}
We are now in the position to determine the chief characteristics of infinity.

\Lemma{7.2.1} {\itshape
Assume \(\mathfrak M,\mathfrak N\mathbin{\eta}\M\), \(\mathfrak N\)
infinite. Then \(\mathfrak M\lesssim\mathfrak N\).
}

\Proof We have \(\mathfrak N\sim\mathfrak N'\subsetneqq\mathfrak N\).
Let \(U\in\M\) be the partially isometric mapping of \(\mathfrak N\) on
\(\mathfrak N'\). Form for every \(n=0,1,2,\ldots\) the \(U^n\)-image of
\(\mathfrak N\), and call them \(\mathfrak N^{(n)}\)
(\(\mathfrak N,\mathfrak N'\) coincide with
\(\mathfrak N^{(0)},\mathfrak N^{(1)}\) resp.). We have
\(\mathfrak N^{(0)}\supset\mathfrak N^{(1)}\), thus (apply \(U^n\)!)
\(\mathfrak N^{(n)}\supset\mathfrak N^{(n+1)}\); therefore
\[
 \mathfrak N^{(0)}\supset\mathfrak N^{(1)}\supset
 \mathfrak N^{(2)}\supset\cdots.
\]
\(U^n\) is partially isometric with the final set \(\mathfrak N^{(n)}\)
therefore \(\mathfrak N^{(n)}\mathbin{\eta}\M\). \(U\) maps
\(\mathfrak N^{(n)}\) on \(\mathfrak N^{(n+1)}\),
\(\mathfrak N^{(n+1)}\) on \(\mathfrak N^{(n+2)}\), therefore
\(\mathfrak N^{(n)}-\mathfrak N^{(n+1)}\) on
\(\mathfrak N^{(n+1)}-\mathfrak N^{(n+2)}\). Using
\[
 U(P_{\mathfrak N^{(n)}}-P_{\mathfrak N^{(n+1)}})
\]
shows therefore, that
\[
 \mathfrak N^{(n)}-\mathfrak N^{(n+1)}
 \sim
 \mathfrak N^{(n+1)}-\mathfrak N^{(n+2)}.
\]
Put
\(\mathfrak P_n^*=\mathfrak N^{(n)}-\mathfrak N^{(n+1)}\), then we see:
The \(\mathfrak P_1^*,\mathfrak P_2^*,\ldots\) form an infinite sequence,
are mutually orthogonal, all \(\subset\mathfrak N^{(0)}=\mathfrak N\) and
\(\mathfrak P_1^*\sim\mathfrak P_2^*\sim\cdots\). As
\[
 \mathfrak P_1^*=\mathfrak N^{(0)}-\mathfrak N^{(1)}
 =\mathfrak N-\mathfrak N'\ne(0),
\]
all \(\mathfrak P_n^*\ne0\).

Apply now the construction of \lemref{7.1.2} to
\(\mathfrak P_1^*,\mathfrak N\) (for \(\mathfrak M,\mathfrak N\)) and
choose \(\mathfrak P_\alpha\) for \(\alpha=1,2,\ldots\) equal to
\(\mathfrak P_1^*,\mathfrak P_2^*,\ldots\) resp. Then an infinite
\(\bar\alpha\) will result, which will be replaced, as described in the proof
of \lemref{7.1.2}, by \(\bar\alpha=\omega\). And then the \(\mathfrak Q\)
will be eliminated. So
\[
 \mathfrak N=\clos{\mathfrak P_1,\mathfrak P_2,\ldots},
\]
where the sequence is infinite, and
\(\mathfrak P_1^*\sim\mathfrak P_1\sim\mathfrak P_2\sim\cdots\).

Apply next the result of \lemref{7.1.2} to
\(\mathfrak P_1^*,\mathfrak M\) (for \(\mathfrak M,\mathfrak N\)), then
\[
 \mathfrak M=\clos{\mathfrak P_1',\mathfrak P_2',\ldots,\mathfrak Q'},
\]
where the sequence is finite or infinite, and
\[
 \mathfrak P_1^*\sim\mathfrak P_1'\sim\mathfrak P_2'\sim\cdots,
 \qquad \mathfrak Q'<\mathfrak P_1^*.
\]

If it is infinite, we may put \(\mathfrak Q'=(0)\), and so, using
\lemref{6.1.2},
\[
 \mathfrak M=\clos{\mathfrak P_1',\mathfrak P_2',\ldots}
 \sim\clos{\mathfrak P_1,\mathfrak P_2,\ldots}=\mathfrak N.
\]

If it is finite, we have
\[
 \mathfrak M=\clos{\mathfrak P_1',\ldots,\mathfrak P_n',\mathfrak Q'}
 \quad\hbox{and}\quad
 \mathfrak Q'<\mathfrak P_1^*\sim\mathfrak P_{n+1},
 \quad
 \mathfrak Q'\sim\mathfrak Q''\subset\mathfrak P_{n+1}.
\]
Thus \lemref{6.1.2} gives
\begin{align*}
 \mathfrak M
 &=\clos{\mathfrak P_1',\ldots,\mathfrak P_n',\mathfrak Q'}
 \sim\clos{\mathfrak P_1,\ldots,\mathfrak P_n,\mathfrak Q''}\\
 &\subset\clos{\mathfrak P_1,\ldots,\mathfrak P_n,\mathfrak P_{n+1}}
 \subset\clos{\mathfrak P_1,\mathfrak P_2,\ldots}=\mathfrak N.
\end{align*}
So we have \(\mathfrak M\lesssim\mathfrak N\) in any event.

\Lemma{7.2.2} {\itshape
There are two possibilities:
\begin{enumerate}[label=(\alph*)]
 \item An \(\mathfrak M\mathbin{\eta}\M\) is infinite if and only if
   \(\mathfrak M\sim\Hspace\) and finite if and only if
   \(\mathfrak M<\Hspace\).
 \item All \(\mathfrak M\mathbin{\eta}\M\) are finite. In this case
   \(\mathfrak M\sim\Hspace\) implies \(\mathfrak M=\Hspace\).
\end{enumerate}
}

\Proof If \(\Hspace\) is finite, every \(\mathfrak M\) is so, by
\lemref{7.1.1}, (iii). If \(\mathfrak M\ne\Hspace\), then
\(\mathfrak M\subsetneqq\Hspace\), and \(\Hspace\sim\mathfrak M\) would
imply that \(\Hspace\) is infinite. Thus if \(\Hspace\) is finite, case (b)
holds.

If \(\Hspace\) is infinite, then \(\mathfrak M\sim\Hspace\) implies that
\(\mathfrak M\) too is infinite by \lemref{7.1.1}, (i). If conversely
\(\mathfrak M\) is infinite, \lemref{7.2.1} gives if applied to
\(\Hspace,\mathfrak M\) (for \(\mathfrak M,\mathfrak N\))
\(\Hspace\lesssim\mathfrak M\). By \lemref{6.3.2}, (i),
\(\mathfrak M\lesssim\Hspace\), so \(\mathfrak M\sim\Hspace\). And as we
have always \(\mathfrak M\lesssim\Hspace\) (cf. as above), therefore
\(\mathfrak M<\Hspace\) is characteristic for finite \(\mathfrak M\)'s.
Thus case (a) holds.

A convenient characterisation of infinity is this:

\Lemma{7.2.3} {\itshape
An \(\mathfrak M\mathbin{\eta}\M\) is infinite if and only if
\(\mathfrak M\ne(0)\) and an \(\mathfrak N\subset\mathfrak M\) with
\(\mathfrak M\sim\mathfrak N\sim\mathfrak M-\mathfrak N\) exists.
}

\Proof The condition is sufficient: If \(\mathfrak N\) is such, then
\[
 \mathfrak M-\mathfrak N\sim\mathfrak M\ne(0),\qquad
 \mathfrak M-\mathfrak N\ne(0),\qquad
 \mathfrak N\ne\mathfrak M,
\]
so \(\mathfrak M\sim\mathfrak N\subsetneqq\mathfrak M\).

It is necessary: If \(\mathfrak M\) is infinite, then the proof of
\lemref{7.2.1} shows that an infinite sequence
\(\mathfrak P_1,\mathfrak P_2,\ldots\) exists such that all
\(\mathfrak P_i\) are mutually orthogonal, \(\ne(0)\),
\[
 \mathfrak M=\clos{\mathfrak P_1,\mathfrak P_2,\ldots},
 \qquad
 \mathfrak P_1\sim\mathfrak P_2\sim\mathfrak P_3\sim\cdots.
\]
Now \lemref{6.1.2} gives
\[
 \clos{\mathfrak P_1,\mathfrak P_2,\ldots}
 \sim\clos{\mathfrak P_1,\mathfrak P_3,\ldots}
 \sim\clos{\mathfrak P_2,\mathfrak P_4,\ldots}
\]
so for \(\mathfrak N=\clos{\mathfrak P_1,\mathfrak P_3,\ldots}\),
\(\mathfrak N\subset\mathfrak M\),
\(\mathfrak M\sim\mathfrak N\sim\mathfrak M-\mathfrak N\). As
\(\mathfrak P_1\ne(0)\), necessarily \(\mathfrak M\ne(0)\).

\SectionStart{7.3}{Additivity of finiteness}
Continuing our discussion along the lines which are customary in general
set-theory, it is natural to establish the additivity of the notion of
finiteness. We will prove it in five successive steps, the first Lemma having
a certain interest of its own.

\Lemma{7.3.1} {\itshape Assume
\(\mathfrak M,\mathfrak N,\mathfrak P\mathbin{\eta}\M\);
\(\mathfrak M,\mathfrak N\) orthogonal;
\(\mathfrak P\subset[\mathfrak M,\mathfrak N]\). Then
\(\mathfrak P\) can be represented in the following form (all sets which
follow being linear and closed):

There exist two orthogonal sets
\(\mathfrak M',\mathfrak M''\mathbin{\eta}\M\) and
\(\subset\mathfrak M\); two orthogonal sets
\(\mathfrak N',\mathfrak N''\mathbin{\eta}\M\) and
\(\subset\mathfrak N\); and a linear, closed operator
\(A\mathbin{\eta}\M\), with the following properties; Domain \(A\) is a
part of and dense in \(\mathfrak M''\); Range \(A\) is a part of and dense
in \(\mathfrak N''\); \(Af=0\) implies \(f=0\). The set
\(\mathfrak P^*\) of all \(f+Af\) (\(f\in\operatorname{Domain}A\)) is
linear, closed, and \(\mathbin{\eta}\M\).
\[
 \mathfrak P^*\sim\mathfrak M''\sim\mathfrak N'',
\]
\(\mathfrak M',\mathfrak N',\mathfrak P^*\) are mutually orthogonal and
\[
 \mathfrak P=[\mathfrak M',\mathfrak N',\mathfrak P^*].
\]}

\Proof Put \(\mathfrak M'=\mathfrak M\mathbin{\cdot}\mathfrak P\),
\(\mathfrak N'=\mathfrak N\mathbin{\cdot}\mathfrak P\), then
\(\mathfrak M',\mathfrak N'\mathbin{\eta}\M\),
\(\mathfrak M'\subset\mathfrak M\),
\(\mathfrak N'\subset\mathfrak N\), thus
\(\mathfrak M',\mathfrak N'\) are orthogonal; and
\(\mathfrak M',\mathfrak N'\subset\mathfrak P\). Thus we can form
\(\mathfrak P^*=\mathfrak P-[\mathfrak M',\mathfrak N']\), then
\(\mathfrak P^*\mathbin{\eta}\M\); \(\mathfrak P^*\) is orthogonal to
both \(\mathfrak M',\mathfrak N'\) and
\(\mathfrak P=[\mathfrak M',\mathfrak N',\mathfrak P^*]\). Thus we must
only give now the characterisation of \(\mathfrak P^*\) given in our Lemma
in terms of \(\mathfrak M''\), \(\mathfrak N''\), and \(A\).

An \(f\in\mathfrak P^*\) is
\(f\in[\mathfrak M,\mathfrak N]\), thus \(f=g+h\),
\(g\in\mathfrak M\), \(h\in\mathfrak N\) (cf. \cite[p.~76]{ref16}), \(g,h\)
being clearly uniquely determined. If in such a decomposition \(g=0\),
then \(f=h\in\mathfrak N\), \(f\in\mathfrak P^*\subset\mathfrak P\), so
\(f\in\mathfrak N\mathbin{\cdot}\mathfrak P
=\mathfrak N'\subset\mathfrak P-\mathfrak P^*\). As
\(f\in\mathfrak P^*\) this implies \(f=0\), \(h=0\). Similarly \(h=0\)
implies \(g=0\). Owing to the linearity this means, that \(g\) determines
\(h\) uniquely and conversely (as far as they belong to any
\(f\in\mathfrak P^*\) at all). Thus we can define a one-valued operator
\(A\) with a one-valued inverse \(A^{-1}\) by \(Ag=h\), whenever \(g,h\)
belong to an \(f\in\mathfrak P^*\) in the above way. Thus
\(\mathfrak P^*\) is the set of all \(f+Af\). As it is
\(\mathbin{\eta}\M\), linear, and closed, the same is true for \(A\).

Put now \(\mathfrak M''=[\operatorname{Domain}A]\),
\(\mathfrak N''=[\operatorname{Range}A]\). All we must prove is, that
\(\mathfrak M''\) is \(\subset\mathfrak M\) and orthogonal to
\(\mathfrak M'\); that \(\mathfrak N''\subset\mathfrak N\) and orthogonal
to \(\mathfrak N'\); and that
\(\mathfrak P^*\sim\mathfrak N''\sim\mathfrak M''\).

If \(g\in\operatorname{Domain}A\), then \(f=g+h\),
\(f\in\mathfrak P^*\), \(g\in\mathfrak M\), \(h\in\mathfrak N\). So
\(g\in\mathfrak M\). \(f\) is orthogonal to \(\mathfrak M'\), and
\(h\in\mathfrak N\) is orthogonal to \(\mathfrak M\), so to
\(\mathfrak M'\) too. Thus \(g=f-h\) is orthogonal to \(\mathfrak M'\),
\(g\in\mathfrak M-\mathfrak M'\). Thus
\(\operatorname{Domain}A\subset\mathfrak M-\mathfrak M'\),
\(\mathfrak M''=[\operatorname{Domain}A]
\subset\mathfrak M-\mathfrak M'\). Similarly
\(\mathfrak N''\subset\mathfrak N-\mathfrak N'\).

Consider now the operator
\(A_1=(1+A)P_{\mathfrak M''}\) (the linear operator, which agrees with
\(1+A\) in \(\mathfrak M''\) and with 0 in
\(\Hspace-\mathfrak M''\)). \(A_1\) is clearly
\(\mathbin{\eta}\M\), linear, and closed, and its domain is everywhere
dense. Now
\([\operatorname{Range}A_1]
=[\operatorname{Range}(1+A)]=\mathfrak P^*\). On the other hand
\begin{align*}
 [\operatorname{Range}A_1^*]
 &=\Hspace-\setof{f}{A_1f=0}\\
 &=\Hspace-\setof{f}{(1+A)P_{\mathfrak M''}f=0}\\
 &=\Hspace-\setof{f}{P_{\mathfrak M''}f=0}\\
 &=\Hspace-(\Hspace-\mathfrak M'')=\mathfrak M''.
\end{align*}
Now \lemref{6.2.1} gives, applied to \(A_1\),
\(\mathfrak P^*\sim\mathfrak M''\). Similarly
\(\mathfrak P^*\sim\mathfrak N''\).

Thus all parts of our Lemma are proved.

\Lemma{7.3.2} \emph{Assume
\(\mathfrak M,\mathfrak N,\mathfrak P\mathbin{\eta}\M\);
\(\mathfrak M,\mathfrak N\) orthogonal;
\(\mathfrak P\subset[\mathfrak M,\mathfrak N]\). Then either
\(\mathfrak P\lesssim\mathfrak M\), or
\([\mathfrak M,\mathfrak N]-\mathfrak P\lesssim\mathfrak N\).}

\Proof Apply \lemref{7.3.1} to \(\mathfrak M,\mathfrak N,\mathfrak P\)
obtaining the sets
\(\mathfrak M',\mathfrak M'',\mathfrak N',\mathfrak N'',\mathfrak P^*\)
and the operator \(A\); and to
\(\mathfrak M,\mathfrak N,[\mathfrak M,\mathfrak N]-\mathfrak P\)
obtaining the sets
\(\overline{\mathfrak M}',\overline{\mathfrak M}'',
\overline{\mathfrak N}',\overline{\mathfrak N}'',
\overline{\mathfrak P}^*\) and the operator \(B\). As
\(\mathfrak P,[\mathfrak M,\mathfrak N]-\mathfrak P\) are orthogonal,
\(\mathfrak M'=\mathfrak M\mathbin{\cdot}\mathfrak P\),
\(\overline{\mathfrak M}'
=\mathfrak M\mathbin{\cdot}([\mathfrak M,\mathfrak N]-\mathfrak P)\)
are too; similarly \(\mathfrak N',\overline{\mathfrak N}'\). Form
\[
 \mathfrak M^0=\mathfrak M-[\mathfrak M',\overline{\mathfrak M}'],
 \qquad
 \mathfrak N^0=\mathfrak N-[\mathfrak N',\overline{\mathfrak N}'].
\]
If \(f\in\operatorname{Domain}A\) then
\(f+Af\in\mathfrak P^*\subset\mathfrak P\) is orthogonal to
\(\overline{\mathfrak M}'\subset[\mathfrak M,\mathfrak N]-\mathfrak P\).
\(Af\in\mathfrak N\) is orthogonal too to
\(\overline{\mathfrak M}'\subset\mathfrak M\). Thus \(f\) is
orthogonal to \(\overline{\mathfrak M}'\). Thus all
\(\mathfrak M''=[\operatorname{Domain}A]\) is orthogonal to
\(\overline{\mathfrak M}'\). Besides we know that it is orthogonal to
\(\mathfrak M'\), and \(\subset\mathfrak M\); so
\(\mathfrak M''\subset\mathfrak M-[\mathfrak M',
\overline{\mathfrak M}']=\mathfrak M^0\). Similarly
\(\overline{\mathfrak M}''\subset\mathfrak M^0\). Similarly
\(\mathfrak N'',\overline{\mathfrak N}''\subset\mathfrak N^0\).

Now consider \(\mathfrak N',\overline{\mathfrak M}'\). We have
\(\mathfrak N'\gtreqless\overline{\mathfrak M}'\) (by \thmref{VI}),
that is either
\(\mathfrak N'\lesssim\overline{\mathfrak M}'\),
\(\mathfrak N'\sim\mathfrak M^*\subset\overline{\mathfrak M}'\), or
\(\mathfrak N'\gtrsim\overline{\mathfrak M}'\),
\(\overline{\mathfrak M}'\sim\mathfrak M^*\subset\mathfrak N'\). In the
first case (by \lemref{6.1.2})
\[
 \mathfrak P=[\mathfrak M',\mathfrak N',\mathfrak P^*]
 \sim[\mathfrak M',\mathfrak M^*,\mathfrak M'']
 \subset[\mathfrak M',\overline{\mathfrak M}',\mathfrak M^0]
 =\mathfrak M;
\]
in the second case (as above)
\[
 [\mathfrak M,\mathfrak N]-\mathfrak P
 =[\overline{\mathfrak M}',\overline{\mathfrak N}',
   \overline{\mathfrak P}^*]
 \sim[\mathfrak M^*,\overline{\mathfrak N}',
   \overline{\mathfrak N}'']
 \subset[\mathfrak N',\overline{\mathfrak N}',\mathfrak N^0]
 =\mathfrak N.
\]
Thus \(\mathfrak P\lesssim\mathfrak M\) resp.
\([\mathfrak M,\mathfrak N]-\mathfrak P\lesssim\mathfrak N\).

\Lemma{7.3.3} \emph{Assume
\(\mathfrak M,\mathfrak N\mathbin{\eta}\M\);
\(\mathfrak M,\mathfrak N\) orthogonal. If \(\mathfrak M,\mathfrak N\)
are finite, then \([\mathfrak M,\mathfrak N]\) is finite too.}

\Proof Assume that \([\mathfrak M,\mathfrak N]\) is infinite. Then there
exists by \lemref{7.2.3} a \(\mathfrak P\subset[\mathfrak M,\mathfrak N]\)
such that
\([\mathfrak M,\mathfrak N]\sim\mathfrak P
\sim[\mathfrak M,\mathfrak N]-\mathfrak P\). By \lemref{7.3.2} we have
\(\mathfrak P\lesssim\mathfrak M\) or
\([\mathfrak M,\mathfrak N]-\mathfrak P\lesssim\mathfrak N\), so
\(\mathfrak M\) or \(\mathfrak N\gtrsim[\mathfrak M,\mathfrak N]\).
By \lemref{7.1.1}, (i), then \(\mathfrak M\) or \(\mathfrak N\) must be
infinite, contradicting our assumption.

\Lemma{7.3.4} \emph{If
\(\mathfrak M,\mathfrak N\mathbin{\eta}\M\), then
\([\mathfrak M,\mathfrak N]-\mathfrak N\lesssim\mathfrak M\).}

\Proof \([\mathfrak M,\mathfrak N]-\mathfrak N\) is obviously the set of
all \(P_{\Hspace-\mathfrak N}f\) where \(f\) runs over
\([\mathfrak M,\mathfrak N]\). So it is equal to
\begin{align*}
 &[P_{\Hspace-\mathfrak N}f; f\in[\mathfrak M,\mathfrak N]]\\
 &=[P_{\Hspace-\mathfrak N}(g+h); g\in\mathfrak M, h\in\mathfrak N]\\
 &=[P_{\Hspace-\mathfrak N}g; g\in\mathfrak M, h\in\mathfrak N]\\
 &=[P_{\Hspace-\mathfrak N}g; g\in\mathfrak M]\\
 &=[P_{\Hspace-\mathfrak N}P_{\mathfrak M}g'; g'\in\Hspace]
 =[\operatorname{Range}P_{\Hspace-\mathfrak N}P_{\mathfrak M}].
\end{align*}
At the same time
\[
 [\operatorname{Range}(P_{\Hspace-\mathfrak N}P_{\mathfrak M})^*]
 =[\operatorname{Range}P_{\mathfrak M}P_{\Hspace-\mathfrak N}]
 \subset[\operatorname{Range}P_{\mathfrak M}]=\mathfrak M.
\]
Now applying \lemref{6.2.1} to
\(P_{\Hspace-\mathfrak N}P_{\mathfrak M}\) gives:
\[
 [\mathfrak M,\mathfrak N]-\mathfrak N
 =[\operatorname{Range}P_{\Hspace-\mathfrak N}P_{\mathfrak M}]
 \sim[\operatorname{Range}(P_{\Hspace-\mathfrak N}P_{\mathfrak M})^*]
 \subset\mathfrak M
\]
and therefore
\([\mathfrak M,\mathfrak N]-\mathfrak N\lesssim\mathfrak M\).

\Lemma{7.3.5} \emph{If
\(\mathfrak M_1,\ldots,\mathfrak M_n\mathbin{\eta}\M\), and all
\(\mathfrak M_i\) are finite, then
\([\mathfrak M_1,\ldots,\mathfrak M_n]\) is finite too.}

\Proof The statement is obvious for \(n=1\). Consider now \(n=2\). By
\lemref{7.3.4} and \lemref{7.1.1}, (i),
\([\mathfrak M_1,\mathfrak M_2]-\mathfrak M_2\) is finite. As
\([\mathfrak M_1,\mathfrak M_2]-\mathfrak M_2\) and
\(\mathfrak M_2\) are orthogonal, \lemref{7.3.3} applies:
\([[\mathfrak M_1,\mathfrak M_2]-\mathfrak M_2,\mathfrak M_2]\), that
is \([\mathfrak M_1,\mathfrak M_2]\), is finite. Thus the statement holds
for \(n=2\), too.

Assume now \(n=3,4,\ldots\), and that the statement is already established
for \(n-1\). Then \([\mathfrak M_1,\ldots,\mathfrak M_{n-1}]\) is finite,
and as our statement holds for \(n=2\),
\([[\mathfrak M_1,\ldots,\mathfrak M_{n-1}],\mathfrak M_n]\), that is
\([\mathfrak M_1,\ldots,\mathfrak M_n]\), is finite too. This completes
the proof.
% --- End of parts/03a-chapters-vi-vii.tex ---
\ChapterHeading{8}{VIII}{Numerical Dimensionality}

\SectionStart{8.1}{Fundamental sequences}
We will use the information obtained in the foregoing about the relative
dimensionality to define a quantitative measurement for it. We will do this by
refining the number $\reldiv{\mathfrak N}{\mathfrak M}$ in \lemref{7.1.3},
which gave an integral and therefore only approximate measure of the ratio of
the sets.

Our task is identical with that one which has to be met if one desires to pass
from the purely geometrical calculus with intervals to the analytical one with
numerical lengths: the problem which was solved by Euclid's famous algorithm.
Our \lemref{7.1.3} is indeed the first step in this algorithm.

It is, however, technically preferable to use a somewhat different procedure.
We will have to distinguish for a short time between the two possibilities that
there is or is not a minimal $\mathfrak M\mathbin\eta\M$. The first case is not
really interesting, because if a minimal $\mathfrak M\mathbin\eta\M$ exists, we
possess already the complete characterization of $\M$ (as a direct factor) by
\thmref{IV}. But as it makes no difference for the possibility of a quantitative
measurement of the relative dimensionality, we will preserve the completeness
of our discussion by including this case.

\Definition{8.1.1}
Assume $\mathfrak M,\mathfrak N\mathbin\eta\M$, finite, and $\ne(0)$. If we
have
\[
 \mathfrak N=[\mathfrak P_1,\ldots,\mathfrak P_n,\mathfrak D],
\]
where $n=0,1,2,\ldots$ (finite!); $\mathfrak P_1,\ldots,\mathfrak P_n,
\mathfrak D$ are mutually orthogonal; $\mathfrak P_i\sim\mathfrak M$; then we
use the abbreviated notation
\[
 \mathfrak N\simeq n\odot\mathfrak M\mathbin{\#}\mathfrak D.
\]
If $\mathfrak D=(0)$, we omit it.

\Lemma{8.1.1}
{\itshape
If $\mathfrak M\mathbin\eta\M$, finite, $\ne(0)$, and not minimal, then an
$\mathfrak N\mathbin\eta\M$, finite, $\ne(0)$, with
\[
 \reldiv{\mathfrak M}{\mathfrak N}\geq2
\]
exists.}

\Proof As $\mathfrak M$ is not minimal, an $\mathfrak N\mathbin\eta\M$,
$\mathfrak N\subset\mathfrak M$, $\mathfrak N\ne(0),\mathfrak M$ exists. So
$\mathfrak N,\mathfrak M-\mathfrak N\ne(0)$. We have
$\mathfrak N\gtreqless\mathfrak M-\mathfrak N$, that is
$\mathfrak N\succsim\mathfrak M-\mathfrak N$ or
$\mathfrak N\precsim\mathfrak M-\mathfrak N$. In the second case replace
$\mathfrak N$ by $\mathfrak M-\mathfrak N$, so we have at any rate
$\mathfrak N\precsim\mathfrak M-\mathfrak N$.

$\mathfrak N\subset\mathfrak M$ implies the finiteness of $\mathfrak N$. As
$\mathfrak N\sim\mathfrak N'\subset\mathfrak M-\mathfrak N$, and
$\mathfrak N,\mathfrak N'$ are orthogonal and $\subset\mathfrak M$, we have
with $\mathfrak P_1=\mathfrak N$, $\mathfrak P_2=\mathfrak N'$,
$\mathfrak D=\mathfrak M-[\mathfrak N,\mathfrak N']$, obviously
\[
 \mathfrak M\simeq2\odot\mathfrak N\mathbin{\#}\mathfrak D.
\]
Apply now \lemref{7.1.3} to $\mathfrak N,\mathfrak D$;
\[
 \mathfrak D\simeq p\odot\mathfrak N\mathbin{\#}\mathfrak D',
 \qquad p=0,1,2,\ldots,\quad \mathfrak D'\prec\mathfrak N.
\]
Then
\[
 \mathfrak M\simeq(2+p)\odot\mathfrak N\mathbin{\#}\mathfrak D',
 \qquad \mathfrak D'\prec\mathfrak N,
\]
and so $\reldiv{\mathfrak M}{\mathfrak N}=2+p\geq2$.

\Lemma{8.1.2}
{\itshape
Every minimal $\mathfrak M\mathbin\eta\M$ is finite and $\ne(0)$.}

\Proof $\mathfrak M\ne(0)$ is part of the definition. $\mathfrak M$ infinite
would imply $\mathfrak M\sim\mathfrak N\subsetneqq\mathfrak M$,
so $\mathfrak N\mathbin\eta\M$, and as $\mathfrak M$ is minimal,
$\mathfrak N=(0)$. Then $\mathfrak M\sim(0)$, $\mathfrak M=(0)$ which is
contradictory.

\Definition{8.1.2}
Certain (finite or infinite) sequences
$\mathfrak S=(\overline{\mathfrak M}_1,\overline{\mathfrak M}_2,\ldots)$ of
elements which are $\mathbin\eta\M$, finite, $\ne(0)$, will be called
\emph{fundamental sequences}. They are these:
\begin{enumerate}[leftmargin=2.4em]
\item[$(\alpha)$] Any minimal $\overline{\mathfrak M}_1\mathbin\eta\M$ is a fundamental
  sequence by itself (length one!).
\item[$(\beta)$] Any infinite sequence
  $\overline{\mathfrak M}_1,\overline{\mathfrak M}_2,\ldots$ of
  $\mathbin\eta\M$, finite, $\ne(0)$ elements is a fundamental sequence if
  \[
   \reldiv{\overline{\mathfrak M}_i}{\overline{\mathfrak M}_{i+1}}\geq2
   \qquad\text{for }i=1,2,\ldots.
  \]
\end{enumerate}

\Lemma{8.1.3}
{\itshape
If sets $\mathfrak M\mathbin\eta\M$ which are finite and $\ne(0)$ do exist at
all, then there exists at least one fundamental sequence.}

\Proof If a minimal $\mathfrak M\mathbin\eta\M$ exists, put
$\overline{\mathfrak M}_1=\mathfrak M$. By \lemref{8.1.2} this meets the
requirements of case $(\alpha)$. If no minimal $\mathfrak M\mathbin\eta\M$ exists,
choose a finite $\mathfrak M\mathbin\eta\M$, $\mathfrak M\ne(0)$, and put
$\overline{\mathfrak M}_1=\mathfrak M$; and obtain
$\overline{\mathfrak M}_{i+1}$ from $\overline{\mathfrak M}_i$ by means of
\lemref{8.1.1}. This meets the requirements of case $(\beta)$.

We now prove two important evaluations.

\Lemma{8.1.4}
{\itshape
Assume $\mathfrak M,\mathfrak N,\mathfrak P\mathbin\eta\M$ finite, and
$\ne(0)$. Then
\[
 \reldiv{\mathfrak N}{\mathfrak M}
 \reldiv{\mathfrak P}{\mathfrak N}
 \leq \reldiv{\mathfrak P}{\mathfrak M}
 <\left(\reldiv{\mathfrak N}{\mathfrak M}+1\right)
  \left(\reldiv{\mathfrak P}{\mathfrak N}+1\right)
\]
and, if $\mathfrak N,\mathfrak P$ are orthogonal,
\[
 \reldiv{\mathfrak N}{\mathfrak M}
 +\reldiv{\mathfrak P}{\mathfrak M}
 \leq\reldiv{[\mathfrak N,\mathfrak P]}{\mathfrak M}
 \leq\reldiv{\mathfrak N}{\mathfrak M}
 +\reldiv{\mathfrak P}{\mathfrak M}+1.
\]}

\Proof First put
\[
 \begin{aligned}
 \reldiv{\mathfrak N}{\mathfrak M}&=m,
 \quad \mathfrak N\simeq m\odot\mathfrak M\mathbin{\#}\mathfrak M',
 \quad \mathfrak M'\prec\mathfrak M,\\
 \text{and}\quad \reldiv{\mathfrak P}{\mathfrak N}&=n,
 \quad \mathfrak P\simeq n\odot\mathfrak N\mathbin{\#}\mathfrak N',
 \quad \mathfrak N'\prec\mathfrak N.
 \end{aligned}
\]
This clearly implies
$\mathfrak P\simeq mn\odot\mathfrak M\mathbin{\#}\mathfrak M''$, from which
$\reldiv{\mathfrak P}{\mathfrak M}\geq mn$ follows as at the end of the proof
of \lemref{8.1.1}.

If we had however
$\reldiv{\mathfrak P}{\mathfrak M}\geq(m+1)(n+1)$ that would mean
\[
 \mathfrak P\simeq
 \reldiv{\mathfrak P}{\mathfrak M}\odot\mathfrak M
 \mathbin{\#}\mathfrak M'''
 \simeq(m+1)(n+1)\odot\mathfrak M\mathbin{\#}\mathfrak M^{\mathrm{IV}},
\]
where
\[
 \mathfrak M^{\mathrm{IV}}\simeq
 \left(\reldiv{\mathfrak P}{\mathfrak M}-(m+1)(n+1)\right)
 \odot\mathfrak M\mathbin{\#}\mathfrak M'''.
\]
Or $(n+1)\odot\mathfrak N^\circ\simeq\mathfrak P^\circ\subset\mathfrak P$
where $\mathfrak N^\circ\simeq(m+1)\odot\mathfrak M$. Now clearly
\[
 \mathfrak N\simeq m\odot\mathfrak M\mathbin{\#}\mathfrak M'
 \precsim\mathfrak N^\circ,
\]
so we may as well write
$(n+1)\odot\mathfrak N\simeq\mathfrak P^{\circ\circ}
\subset\mathfrak P^\circ\subset\mathfrak P$, that is
\[
 \mathfrak P\simeq(n+1)\odot\mathfrak N\mathbin{\#}\mathfrak N''.
\]
This gives, as above,
$\reldiv{\mathfrak P}{\mathfrak N}\geq n+1$, contradicting
$\reldiv{\mathfrak P}{\mathfrak N}=n$. So we must have
\[
 mn\leq\reldiv{\mathfrak P}{\mathfrak M}<(m+1)(n+1),
\]
which is the first set of inequalities.

Next assume that $\mathfrak N,\mathfrak P$ are orthogonal, and put
\[
 \reldiv{\mathfrak N}{\mathfrak M}=m,
 \quad \mathfrak N\simeq m\odot\mathfrak M\mathbin{\#}\mathfrak M',
 \quad \mathfrak M'\prec\mathfrak M;
\]
and
\[
 \reldiv{\mathfrak P}{\mathfrak M}=p,
 \quad \mathfrak P\simeq p\odot\mathfrak M\mathbin{\#}\mathfrak M^{\mathrm V},
 \quad \mathfrak M^{\mathrm V}\prec\mathfrak M.
\]
Then
$[\mathfrak N,\mathfrak P]\simeq(m+p)\odot\mathfrak M
\mathbin{\#}[\mathfrak M',\mathfrak M^{\mathrm V}]$ and so, as above,
\[
 \reldiv{[\mathfrak N,\mathfrak P]}{\mathfrak M}\geq m+p.
\]
If we had, however,
$\reldiv{[\mathfrak N,\mathfrak P]}{\mathfrak M}\geq m+p+2$, that would mean
\[
 [\mathfrak N,\mathfrak P]\simeq
 \reldiv{[\mathfrak N,\mathfrak P]}{\mathfrak M}\odot\mathfrak M
 \mathbin{\#}\mathfrak M^{\mathrm{VI}}
 \simeq(m+p+2)\odot\mathfrak M\mathbin{\#}\mathfrak M^{\mathrm{VII}},
\]
where
\[
 \mathfrak M^{\mathrm{VII}}\simeq
 \left(\reldiv{[\mathfrak N,\mathfrak P]}{\mathfrak M}-(m+p+2)\right)
 \odot\mathfrak M\mathbin{\#}\mathfrak M^{\mathrm{VI}}.
\]
Thus
$[\mathfrak N,\mathfrak P]=[\overline{\mathfrak N},
\overline{\mathfrak P},\mathfrak M^{\mathrm{VII}}]$ where
$\overline{\mathfrak N},\overline{\mathfrak P},\mathfrak M^{\mathrm{VII}}$
are mutually orthogonal, and
$\overline{\mathfrak N}\simeq(m+1)\odot\mathfrak M$,
$\overline{\mathfrak P}\simeq(n+1)\odot\mathfrak M$. This implies
$\mathfrak N\sim\overline{\mathfrak N}^{\circ}
\subsetneqq\overline{\mathfrak N}$,
$\mathfrak P\sim\overline{\mathfrak P}^{\circ}
\subsetneqq\overline{\mathfrak P}$, and so
\[
 [\mathfrak N,\mathfrak P]
 \sim[\overline{\mathfrak N}^{\circ},\overline{\mathfrak P}^{\circ}]
 \subsetneqq[\overline{\mathfrak N},\overline{\mathfrak P}]
 \subset[\overline{\mathfrak N},\overline{\mathfrak P},
 \mathfrak M^{\mathrm{VII}}]=[\mathfrak N,\mathfrak P],
\]
contradicting the finiteness of $[\mathfrak N,\mathfrak P]$, which follows from
the finiteness of $\mathfrak N$ and of $\mathfrak P$ (by \lemref{7.3.3} or
\hyperref[lem:7.3.5]{7.3.5}). So
$\reldiv{[\mathfrak N,\mathfrak P]}{\mathfrak M}<m+p+2$ and therefore
\[
 m+p\leq\reldiv{[\mathfrak N,\mathfrak P]}{\mathfrak M}\leq m+p+1,
\]
which is the second set of inequalities.

These evaluations put us in position to establish the quantitative ratios of
relative dimensionalities as follows.

\Lemma{8.1.5}
{\itshape
If $\mathfrak S=(\overline{\mathfrak M}_1,\overline{\mathfrak M}_2,\ldots)$
is a fundamental sequence, and $\mathfrak M,\mathfrak N\mathbin\eta\M$ are
finite and $\ne(0)$, then
\begin{equation*}
 \phantomsection\label{eq:8.1.1}
 \lim_i
 \frac{\reldiv{\mathfrak N}{\overline{\mathfrak M}_i}}
      {\reldiv{\mathfrak M}{\overline{\mathfrak M}_i}}
 =\left(\frac{\mathfrak N}{\mathfrak M}\right)_{\mathfrak S}
\end{equation*}
exists. (In case (I) we mean by $\lim$ the value at $i=1$, in case (II) we mean
by it $\lim_{i\to\infty}$.) It is a finite and positive real number.}

\Proof In case (I) we must only prove that
$\reldiv{\mathfrak N}{\overline{\mathfrak M}_1}$,
$\reldiv{\mathfrak M}{\overline{\mathfrak M}_1}$ are both $\ne0$.
$\reldiv{\mathfrak N}{\overline{\mathfrak M}_1}=0$ would imply
$\mathfrak N\sim\mathfrak D\prec\overline{\mathfrak M}_1$ so
$\mathfrak N\sim\mathfrak N'\subset\overline{\mathfrak M}_1$ excluding
$\mathfrak N'=\overline{\mathfrak M}_1$. As
$\mathfrak N'\mathbin\eta\M$ and $\overline{\mathfrak M}_1$ minimal, this
necessitates $\mathfrak N'=(0)$, $\mathfrak N\sim\mathfrak N'=(0)$,
$\mathfrak N=(0)$, thus contradicting $\mathfrak N\ne(0)$. Therefore
$\reldiv{\mathfrak N}{\overline{\mathfrak M}_1}\ne0$, similarly
$\reldiv{\mathfrak M}{\overline{\mathfrak M}_1}\ne0$.

Consider now case (II). We have by \lemref{8.1.4}
\[
 \reldiv{\mathfrak N}{\overline{\mathfrak M}_{i+1}}
 \geq\reldiv{\mathfrak N}{\overline{\mathfrak M}_i}
       \reldiv{\overline{\mathfrak M}_i}{\overline{\mathfrak M}_{i+1}}
 \geq2\reldiv{\mathfrak N}{\overline{\mathfrak M}_i},
\]
so either always $\reldiv{\mathfrak N}{\overline{\mathfrak M}_i}=0$, or
$\lim_{i\to\infty}\reldiv{\mathfrak N}{\overline{\mathfrak M}_i}=\infty$.
The former would mean $\mathfrak N\prec\overline{\mathfrak M}_i$ for all
$i=1,2,\ldots$. Now \lemref{8.1.4} gives
\[
 \reldiv{\overline{\mathfrak M}_i}{\overline{\mathfrak M}_{i+j}}
 \geq
 \reldiv{\overline{\mathfrak M}_i}{\overline{\mathfrak M}_{i+1}}
 \cdots
 \reldiv{\overline{\mathfrak M}_{i+j-1}}{\overline{\mathfrak M}_{i+j}}
 \geq2^j,
\]
and so in particular
$\reldiv{\overline{\mathfrak M}_1}{\overline{\mathfrak M}_i}\geq2^{i-1}$
(write $1,i-1$ for $i,j$). So
$\overline{\mathfrak M}_1\simeq2^{i-1}\odot\overline{\mathfrak M}_i
\mathbin{\#}\mathfrak M_i'$, and \emph{a fortiori}
$\overline{\mathfrak M}_1\simeq2^{i-1}\odot\mathfrak N
\mathbin{\#}\mathfrak N_i'$. This would imply, as we frequently concluded in
the proofs of Lemmas \hyperref[lem:8.1.1]{8.1.1} and
\hyperref[lem:8.1.4]{8.1.4},
$\reldiv{\overline{\mathfrak M}_1}{\mathfrak N}\geq2^{i-1}$ for all
$i=1,2,\ldots$, which is impossible. So we have proved
\[
 \lim_{i\to\infty}\reldiv{\mathfrak N}{\overline{\mathfrak M}_i}=\infty;
\]
similarly
$\lim_{i\to\infty}\reldiv{\mathfrak M}{\overline{\mathfrak M}_i}=\infty$.

Another application of \lemref{8.1.4} gives
\[
 \reldiv{\mathfrak N}{\overline{\mathfrak M}_{i+j}}
 \leq
 \left(\reldiv{\mathfrak N}{\overline{\mathfrak M}_i}+1\right)
 \left(\reldiv{\overline{\mathfrak M}_i}{\overline{\mathfrak M}_{i+j}}+1\right),
\]
\[
 \reldiv{\mathfrak M}{\overline{\mathfrak M}_{i+j}}
 \geq
 \reldiv{\mathfrak M}{\overline{\mathfrak M}_i}
 \reldiv{\overline{\mathfrak M}_i}{\overline{\mathfrak M}_{i+j}},
\]
so
\[
 \frac{\reldiv{\mathfrak N}{\overline{\mathfrak M}_{i+j}}}
      {\reldiv{\mathfrak M}{\overline{\mathfrak M}_{i+j}}}
 \leq
 \frac{\reldiv{\mathfrak N}{\overline{\mathfrak M}_i}+1}
      {\reldiv{\mathfrak M}{\overline{\mathfrak M}_i}}
 \frac{\reldiv{\overline{\mathfrak M}_i}{\overline{\mathfrak M}_{i+j}}+1}
      {\reldiv{\overline{\mathfrak M}_i}{\overline{\mathfrak M}_{i+j}}}
 \leq
 \frac{\reldiv{\mathfrak N}{\overline{\mathfrak M}_i}+1}
      {\reldiv{\mathfrak M}{\overline{\mathfrak M}_i}}
 \frac{2^j+1}{2^j}.
\]
Now $j\to\infty$ gives
\[
 \limsup_{j\to\infty}
 \frac{\reldiv{\mathfrak N}{\overline{\mathfrak M}_{i+j}}}
      {\reldiv{\mathfrak M}{\overline{\mathfrak M}_{i+j}}}
 \leq
 \frac{\reldiv{\mathfrak N}{\overline{\mathfrak M}_i}+1}
      {\reldiv{\mathfrak M}{\overline{\mathfrak M}_i}},
\]
that is
\[
 \limsup_{i\to\infty}
 \frac{\reldiv{\mathfrak N}{\overline{\mathfrak M}_i}}
      {\reldiv{\mathfrak M}{\overline{\mathfrak M}_i}}
 \leq
 \frac{\reldiv{\mathfrak N}{\overline{\mathfrak M}_i}+1}
      {\reldiv{\mathfrak M}{\overline{\mathfrak M}_i}}.
\]
For a sufficiently large $i$,
$\reldiv{\mathfrak M}{\overline{\mathfrak M}_i}\ne0$, so the expression on
the right side is finite, and thus the $\limsup_{i\to\infty}$ on the left side is
finite too. Now $i\to\infty$ gives, considering
$\lim_{i\to\infty}\reldiv{\mathfrak M}{\overline{\mathfrak M}_i}=\infty$,
\[
 \limsup_{i\to\infty}
 \frac{\reldiv{\mathfrak N}{\overline{\mathfrak M}_i}}
      {\reldiv{\mathfrak M}{\overline{\mathfrak M}_i}}
 \leq
 \liminf_{i\to\infty}
 \frac{\reldiv{\mathfrak N}{\overline{\mathfrak M}_i}+1}
      {\reldiv{\mathfrak M}{\overline{\mathfrak M}_i}}
 =
 \liminf_{i\to\infty}
 \frac{\reldiv{\mathfrak N}{\overline{\mathfrak M}_i}}
      {\reldiv{\mathfrak M}{\overline{\mathfrak M}_i}}.
\]
Thus \hyperref[eq:8.1.1]{\(\displaystyle
\lim_{i\to\infty}
\frac{\reldiv{\mathfrak N}{\overline{\mathfrak M}_i}}
     {\reldiv{\mathfrak M}{\overline{\mathfrak M}_i}}\)}
exists and is finite. It is non negative by
its nature, and interchanging $\mathfrak M,\mathfrak N$ proves that it cannot
be $0$. This completes the proof.

\Lemma{8.1.6}
{\itshape
If $\mathfrak S=(\overline{\mathfrak M}_1,\overline{\mathfrak M}_2,\ldots)$
is a fundamental sequence, and $\mathfrak M,\mathfrak N,\mathfrak P
\mathbin\eta\M$ are finite and $\ne(0)$, then the following statements hold:
\begin{enumerate}[label=(\roman*)]
\item If $\mathfrak N\sim\mathfrak P$, then
 $\left(\frac{\mathfrak N}{\mathfrak M}\right)_{\mathfrak S}
 =\left(\frac{\mathfrak P}{\mathfrak M}\right)_{\mathfrak S}$.
\item If $\mathfrak M\sim\mathfrak N$, then
 $\left(\frac{\mathfrak P}{\mathfrak M}\right)_{\mathfrak S}
 =\left(\frac{\mathfrak P}{\mathfrak N}\right)_{\mathfrak S}$.
\item
 $\left(\frac{\mathfrak M}{\mathfrak M}\right)_{\mathfrak S}=1$;
 $\left(\frac{\mathfrak N}{\mathfrak M}\right)_{\mathfrak S}
 =\dfrac{1}{\left(\frac{\mathfrak M}{\mathfrak N}\right)_{\mathfrak S}}$;
 $\left(\frac{\mathfrak P}{\mathfrak M}\right)_{\mathfrak S}
 =\left(\frac{\mathfrak N}{\mathfrak M}\right)_{\mathfrak S}
  \cdot\left(\frac{\mathfrak P}{\mathfrak N}\right)_{\mathfrak S}$.
\item If $\mathfrak N,\mathfrak P$ are orthogonal, then
\[
 \left(\frac{[\mathfrak N,\mathfrak P]}{\mathfrak M}\right)_{\mathfrak S}
 =\left(\frac{\mathfrak N}{\mathfrak M}\right)_{\mathfrak S}
 +\left(\frac{\mathfrak P}{\mathfrak M}\right)_{\mathfrak S}.
\]
\end{enumerate}}

\Proof We prove these statements in a somewhat changed order:

Ad (iii): Obvious from the definition.

Ad (i): Clearly
$\reldiv{\mathfrak N}{\overline{\mathfrak M}_i}
=\reldiv{\mathfrak P}{\overline{\mathfrak M}_i}$; this implies
$\left(\frac{\mathfrak N}{\mathfrak M}\right)_{\mathfrak S}
=\left(\frac{\mathfrak P}{\mathfrak M}\right)_{\mathfrak S}$.

Ad (ii): Follows from (i), (iii).

Ad (iv): Consider first case $(\alpha)$. Then
\[
 \mathfrak N\simeq
 \reldiv{\mathfrak N}{\overline{\mathfrak M}_1}
 \odot\overline{\mathfrak M}_1\mathbin{\#}\mathfrak N',
 \qquad \mathfrak N'\prec\overline{\mathfrak M}_1.
\]
So $\mathfrak N'\sim\mathfrak N''\subset\overline{\mathfrak M}_1$, excluding
$\mathfrak N''=\overline{\mathfrak M}_1$. As
$\mathfrak N''\mathbin\eta\M$ and $\overline{\mathfrak M}_1$ minimal, this
necessitates $\mathfrak N''=(0)$, $\mathfrak N'\sim\mathfrak N''=(0)$,
$\mathfrak N'=(0)$. So
\[
 \mathfrak N\simeq
 \reldiv{\mathfrak N}{\overline{\mathfrak M}_1}
 \odot\overline{\mathfrak M}_1,
\]
similarly
$\mathfrak P\simeq\reldiv{\mathfrak P}{\overline{\mathfrak M}_1}
\odot\overline{\mathfrak M}_1$, and therefore
\[
 [\mathfrak N,\mathfrak P]\simeq
 \left(
   \reldiv{\mathfrak N}{\overline{\mathfrak M}_1}
  +\reldiv{\mathfrak P}{\overline{\mathfrak M}_1}
 \right)\odot\overline{\mathfrak M}_1.
\]
Thus
\[
 \reldiv{[\mathfrak N,\mathfrak P]}{\overline{\mathfrak M}_1}
 =\reldiv{\mathfrak N}{\overline{\mathfrak M}_1}
 +\reldiv{\mathfrak P}{\overline{\mathfrak M}_1}.
\]
This means
\[
 \frac{\reldiv{[\mathfrak N,\mathfrak P]}{\overline{\mathfrak M}_1}}
      {\reldiv{\mathfrak M}{\overline{\mathfrak M}_1}}
 =
 \frac{\reldiv{\mathfrak N}{\overline{\mathfrak M}_1}}
      {\reldiv{\mathfrak M}{\overline{\mathfrak M}_1}}
 +
 \frac{\reldiv{\mathfrak P}{\overline{\mathfrak M}_1}}
      {\reldiv{\mathfrak M}{\overline{\mathfrak M}_1}},
\]
and under the conditions of case $(\alpha)$,
\[
 \left(\frac{[\mathfrak N,\mathfrak P]}{\mathfrak M}\right)_{\mathfrak S}
 =\left(\frac{\mathfrak N}{\mathfrak M}\right)_{\mathfrak S}
 +\left(\frac{\mathfrak P}{\mathfrak M}\right)_{\mathfrak S}.
\]

Consider next case $(\beta)$. Then \lemref{8.1.4} gives
\[
 \frac{
   \reldiv{\mathfrak N}{\overline{\mathfrak M}_i}
  +\reldiv{\mathfrak P}{\overline{\mathfrak M}_i}}
  {\reldiv{\mathfrak M}{\overline{\mathfrak M}_i}}
 \leq
 \frac{\reldiv{[\mathfrak N,\mathfrak P]}{\overline{\mathfrak M}_i}}
      {\reldiv{\mathfrak M}{\overline{\mathfrak M}_i}}
 \leq
 \frac{
   \reldiv{\mathfrak N}{\overline{\mathfrak M}_i}
  +\reldiv{\mathfrak P}{\overline{\mathfrak M}_i}+1}
  {\reldiv{\mathfrak M}{\overline{\mathfrak M}_i}}.
\]
Now if $i\to\infty$ this becomes, considering
$\lim_{i\to\infty}\reldiv{\mathfrak M}{\overline{\mathfrak M}_i}=\infty$
(cf. the proof of \lemref{8.1.5}),
\[
 \left(\frac{\mathfrak N}{\mathfrak M}\right)_{\mathfrak S}
 +\left(\frac{\mathfrak P}{\mathfrak M}\right)_{\mathfrak S}
 \leq
 \left(\frac{[\mathfrak N,\mathfrak P]}{\mathfrak M}\right)_{\mathfrak S}
 \leq
 \left(\frac{\mathfrak N}{\mathfrak M}\right)_{\mathfrak S}
 +\left(\frac{\mathfrak P}{\mathfrak M}\right)_{\mathfrak S},
\]
that is
\[
 \left(\frac{[\mathfrak N,\mathfrak P]}{\mathfrak M}\right)_{\mathfrak S}
 =\left(\frac{\mathfrak N}{\mathfrak M}\right)_{\mathfrak S}
 +\left(\frac{\mathfrak P}{\mathfrak M}\right)_{\mathfrak S}.
\]

\SectionStart{8.2}{\texorpdfstring{$D(\mathfrak M)$. Uniqueness of
$D(\mathfrak M)$ and
$\left(\frac{\mathfrak N}{\mathfrak M}\right)_{\mathfrak S}$}
{D(M). Uniqueness of D(M) and (N/M) sub S}}
We introduce the numerical relative dimensions by an implicit definition.

\Definition{8.2.1}
A real-valued function $D(\mathfrak M)$, which is defined for all linear, closed
sets $\mathfrak M\mathbin\eta\M$, is a \emph{relative dimension function}, if
it has the following properties:
\begin{enumerate}[label=(\roman*)]
\item $D(\mathfrak M)$ is $0$ if $\mathfrak M=(0)$; it is finite and positive if
  $\mathfrak M$ is finite and $\ne(0)$; it is $\infty$ if $\mathfrak M$ is
  infinite.
\item If $\mathfrak M\sim\mathfrak N$, then
  $D(\mathfrak M)=D(\mathfrak N)$.
\item If $\mathfrak M,\mathfrak N$ are orthogonal, then
  $D([\mathfrak M,\mathfrak N])=D(\mathfrak M)+D(\mathfrak N)$.
\end{enumerate}
(As to replacing $\mathfrak M$ by $E=P_{\mathfrak M}$, and the use of the
explanatory symbol $(\cdots\M)$, cf. \defref{6.1.1}. If we want to make the
relationship between $\M$ and the function $D(\mathfrak M)$ explicit, we will
write $D_{\M}(\mathfrak M)$.)

The results of \secref{8.1} make it possible to deal exhaustively with the
existence question.

\Lemma{8.2.1}
{\itshape
There always exists at least one relative dimension function $D(\mathfrak M)$.}

\Proof If no finite $\mathfrak M\mathbin\eta\M$, $\mathfrak M\ne(0)$,
exists, put
\[
 D(\mathfrak M)=
 \begin{cases}
  0&\text{for }\mathfrak M=(0),\\
  \infty&\text{otherwise}.
 \end{cases}
\]
This meets all requirements. If finite $\mathfrak M\mathbin\eta\M$,
$\mathfrak M\ne(0)$ do exist, choose one, say $\mathfrak M_0$. Furthermore
choose a fundamental sequence $\mathfrak S$ by \lemref{8.1.3}. Now define
\[
 D(\mathfrak M)=
 \begin{cases}
  0,&\mathfrak M=(0),\\[2pt]
  \left(\dfrac{\mathfrak M}{\mathfrak M_0}\right)_{\mathfrak S},
     &\mathfrak M\text{ finite and }\ne(0),\\[6pt]
  \infty,&\mathfrak M\text{ infinite}.
 \end{cases}
\]
Remember that $[\mathfrak M,\mathfrak N]$ is infinite if and only if
$\mathfrak M$ or $\mathfrak N$ is infinite, by Lemmas
\hyperref[lem:7.1.1]{7.1.1}, (i), and \hyperref[lem:7.3.3]{7.3.3} or
\hyperref[lem:7.3.5]{7.3.5}. Now (i) is obviously fulfilled, and (ii),
(iii) follow from \lemref{8.1.6}, (i), (iv), resp.

\Lemma{8.2.2}
{\itshape
If $D(\mathfrak M)$ is a relative dimension function and $\mathfrak S$ a
fundamental sequence, then we have, whenever
$\mathfrak M,\mathfrak N\mathbin\eta\M$, finite and $\ne(0)$ the relation
\begin{equation*}
 \phantomsection\label{eq:8.2.1}
 \frac{D(\mathfrak N)}{D(\mathfrak M)}
 =\left(\frac{\mathfrak N}{\mathfrak M}\right)_{\mathfrak S}.
\end{equation*}}

\Proof Distinguish the two cases $(\alpha)$ and $(\beta)$ of \defref{8.1.2}.

In the case $(\alpha)$ we have (cf. the proof of \lemref{8.1.6}, (iv)),
\[
 \mathfrak N\simeq\reldiv{\mathfrak N}{\overline{\mathfrak M}_1}
 \odot\overline{\mathfrak M}_1,
 \qquad
 \mathfrak M\simeq\reldiv{\mathfrak M}{\overline{\mathfrak M}_1}
 \odot\overline{\mathfrak M}_1,
\]
therefore by \defref{8.2.1}
\[
 D(\mathfrak N)=\reldiv{\mathfrak N}{\overline{\mathfrak M}_1}
 D(\overline{\mathfrak M}_1),
 \qquad
 D(\mathfrak M)=\reldiv{\mathfrak M}{\overline{\mathfrak M}_1}
 D(\overline{\mathfrak M}_1),
\]
and $D(\overline{\mathfrak M}_1)$ finite and positive, so
\[
 \frac{D(\mathfrak N)}{D(\mathfrak M)}
 =\frac{\reldiv{\mathfrak N}{\overline{\mathfrak M}_1}}
       {\reldiv{\mathfrak M}{\overline{\mathfrak M}_1}}
 =\left(\frac{\mathfrak N}{\mathfrak M}\right)_{\mathfrak S}.
\]

In the case $(\beta)$ we have
\[
 \mathfrak N\simeq\reldiv{\mathfrak N}{\overline{\mathfrak M}_i}
 \odot\overline{\mathfrak M}_i\mathbin{\#}\mathfrak M_i',
 \quad \mathfrak M_i'\prec\overline{\mathfrak M}_i;
\]
\[
 \mathfrak M\simeq\reldiv{\mathfrak M}{\overline{\mathfrak M}_i}
 \odot\overline{\mathfrak M}_i\mathbin{\#}\mathfrak M_i'',
 \quad \mathfrak M_i''\prec\overline{\mathfrak M}_i.
\]
So $\mathfrak M_i'\sim\mathfrak M_i'''\subset\overline{\mathfrak M}_i$,
implying (using \defref{8.2.1} here and below as we did above)
\[
 D(\overline{\mathfrak M}_i)
 =D(\mathfrak M_i''')+D(\overline{\mathfrak M}_i-\mathfrak M_i''')
 \geq D(\mathfrak M_i''')=D(\mathfrak M_i').
\]
Therefore
\[
 D(\mathfrak N)
 =\reldiv{\mathfrak N}{\overline{\mathfrak M}_i}
   D(\overline{\mathfrak M}_i)+D(\mathfrak M_i')
 \leq\left(\reldiv{\mathfrak N}{\overline{\mathfrak M}_i}+1\right)
   D(\overline{\mathfrak M}_i),
\]
\[
 D(\mathfrak M)
 =\reldiv{\mathfrak M}{\overline{\mathfrak M}_i}
   D(\overline{\mathfrak M}_i)+D(\mathfrak M_i'')
 \geq\reldiv{\mathfrak M}{\overline{\mathfrak M}_i}
   D(\overline{\mathfrak M}_i),
\]
and so
\[
 \frac{D(\mathfrak N)}{D(\mathfrak M)}
 \leq
 \frac{\reldiv{\mathfrak N}{\overline{\mathfrak M}_i}+1}
      {\reldiv{\mathfrak M}{\overline{\mathfrak M}_i}}.
\]
Now $i\to\infty$ gives, remembering
$\lim_{i\to\infty}\reldiv{\mathfrak M}{\overline{\mathfrak M}_i}=\infty$
from the proof of \lemref{8.1.5},
\[
 \frac{D(\mathfrak N)}{D(\mathfrak M)}
 \leq\lim_{i\to\infty}
 \frac{\reldiv{\mathfrak N}{\overline{\mathfrak M}_i}}
      {\reldiv{\mathfrak M}{\overline{\mathfrak M}_i}}
 =\left(\frac{\mathfrak N}{\mathfrak M}\right)_{\mathfrak S}.
\]
Interchanging $\mathfrak M,\mathfrak N$ gives, considering \lemref{8.1.6},
(iii), that $\geq$ holds too; so we have equality again.

\Lemma{8.2.3}
{\itshape
All relative dimension functions $D(\mathfrak M)$ are obtained by taking one
of them, say $D_0(\mathfrak M)$, and forming all $aD_0(\mathfrak M)$, for all
finite and positive constants $a$.}

\Proof It is clear, that every $aD_0(\mathfrak M)$ is a relative dimension
function, along with $D_0(\mathfrak M)$.

Consider now two relative dimension functions $D_0(\mathfrak M)$ and
$D(\mathfrak M)$. If no finite $\mathfrak M\mathbin\eta\M$,
$\mathfrak M\ne(0)$ exist, we have
$D_0(\mathfrak M)=D(\mathfrak M)=0$ resp. $\infty$ for $\mathfrak M=(0)$
resp. $\ne(0)$, so we have $D(\mathfrak M)=aD_0(\mathfrak M)$ with $a=1$.
Let us therefore assume, that finite $\mathfrak M\mathbin\eta\M$,
$\mathfrak M\ne(0)$ do exist, and choose a fundamental sequence $\mathfrak S$
by \lemref{8.1.3}. Then \lemref{8.2.2} gives for
$\mathfrak M,\mathfrak N\mathbin\eta\M$, finite and $\ne0$,
\[
 \frac{D(\mathfrak N)}{D(\mathfrak M)}
 =\frac{D_0(\mathfrak N)}{D_0(\mathfrak M)}
 =\left(\frac{\mathfrak N}{\mathfrak M}\right)_{\mathfrak S},
\]
that is
\[
 \frac{D(\mathfrak N)}{D_0(\mathfrak N)}
 =\frac{D(\mathfrak M)}{D_0(\mathfrak M)}.
\]
Thus $D(\mathfrak M)/D_0(\mathfrak M)$ is, for these $\mathfrak M$,
independent of $\mathfrak M$; therefore it is a finite, positive constant $a$.
So we have $D(\mathfrak M)=aD_0(\mathfrak M)$ for every
$\mathfrak M\mathbin\eta\M$ which is finite and $\ne(0)$. But this holds too,
if $\mathfrak M$ is $=(0)$ or infinite, as then both sides become $=0$ resp.
$=\infty$. So the desired relation holds without exceptions. This completes
the proof.

We may remark, on the other hand, that \lemref{8.2.2} shows too, that
$\left(\frac{\mathfrak M}{\mathfrak N}\right)_{\mathfrak S}$ is independent
of the choice of the fundamental sequence $\mathfrak S$. It is nevertheless
dependent on its existence, whereas $D(\mathfrak M)$ is not.

\SectionStart{8.3}{\texorpdfstring{Formulation of the essential characteristics
of $D(\mathfrak M)$. Total additivity and continuity of $D(\mathfrak M)$.
Theorem VII}{Formulation of the essential characteristics of D(M). Total
additivity and continuity of D(M). Theorem VII}}
The connection between the ordering and the relative dimension functions is
given by this Lemma.

\Lemma{8.3.1}
{\itshape
If $D(\mathfrak M)$ is a relative dimension function and
$\mathfrak M,\mathfrak N\mathbin\eta\M$, then the relations
$\mathfrak M\sim\mathfrak N$, $\mathfrak M\prec\mathfrak N$,
$\mathfrak M\succ\mathfrak N$ are equivalent to
$D(\mathfrak M)=D(\mathfrak N)$, $D(\mathfrak M)<D(\mathfrak N)$,
 $D(\mathfrak M)>D(\mathfrak N)$, resp.}

\Proof $\mathfrak M\sim\mathfrak N$ implies
$D(\mathfrak M)=D(\mathfrak N)$ by \defref{8.2.1}, (ii). Consider now
$\mathfrak M\prec\mathfrak N$. Then $\mathfrak M$ is finite, because
$\mathfrak M$ infinite would imply $\mathfrak M\succsim\mathfrak N$; so
$D(\mathfrak M)$ is finite too. Now
$\mathfrak M\sim\mathfrak N'\subset\mathfrak N$, and
$\mathfrak N'=\mathfrak N$ is excluded. So
$\mathfrak N-\mathfrak N'\ne(0)$,
$D(\mathfrak N-\mathfrak N')>0$. (Use both times \defref{8.2.1}, (i).) Now
\defref{8.2.1}, (ii) and (iii) give
\[
 D(\mathfrak N)=D(\mathfrak N')+D(\mathfrak N-\mathfrak N')
 >D(\mathfrak N')=D(\mathfrak M),
\]
$D(\mathfrak M)<D(\mathfrak N)$. Interchanging $\mathfrak M,\mathfrak N$
shows that $\mathfrak M\succ\mathfrak N$ implies
$D(\mathfrak M)>D(\mathfrak N)$. As we have complete disjunctions on both
sides, the converse is true too.

Summing up Lemmas \hyperref[lem:8.2.1]{8.2.1},
\hyperref[lem:8.2.3]{8.2.3}, and \hyperref[lem:8.3.1]{8.3.1}, we have:

\MainTheorem{VII}
{\itshape
A relative dimension function $D(\mathfrak M)$ as characterized by
\defref{8.2.1}, does always exist, and it is unique, except for an arbitrary
constant (real, finite, and positive) factor $a$. A particular choice of
$D(\mathfrak M)$, that is of $a$, is a normalization of the relative dimension.

If $\mathfrak M,\mathfrak N\mathbin\eta\M$ then the relations
$\mathfrak M\sim\mathfrak N$, $\mathfrak M\prec\mathfrak N$,
$\mathfrak M\succ\mathfrak N$ are equivalent to
$D(\mathfrak M)=D(\mathfrak N)$, $D(\mathfrak M)<D(\mathfrak N)$,
 $D(\mathfrak M)>D(\mathfrak N)$, resp.}

We will use $D(\mathfrak M)$ to obtain a detailed characterization of $\M$
itself. But first we establish some continuity properties of $D(\mathfrak M)$.

\Lemma{8.3.2}
{\itshape
Let $\mathfrak M_1,\mathfrak M_2,\ldots$ be a (finite or infinite) sequence,
all $\mathfrak M_i\mathbin\eta\M$ and mutually orthogonal. Then
\begin{equation*}
 \phantomsection\label{eq:8.3.1}
 D([\mathfrak M_1,\mathfrak M_2,\ldots])
 =\sum_iD(\mathfrak M_i).
\end{equation*}}

\Proof Those $\mathfrak M_i$ which are $=(0)$ have no effect on either side
of this equation, and we can omit them. So we may assume that all
$\mathfrak M_i\ne(0)$, $D(\mathfrak M_i)>0$.

Assume first that the sequence is finite, say
$\mathfrak M_1,\ldots,\mathfrak M_n$. For $n=0,1$ the statement is obvious,
for $n=2$ it coincides with \defref{8.2.1}, (iii). So we need only to consider
$n=3,4,\ldots$, assuming that the statement is already proved for $n-1$. Now
as
\[
 [\mathfrak M_1,\ldots,\mathfrak M_n]
 =[[\mathfrak M_1,\ldots,\mathfrak M_{n-1}],\mathfrak M_n]
\]
and as the statement holds for $n-1$ and for two addends, it holds for $n$ too.

Assume next that the sequence is infinite. Then we have
\[
 [\mathfrak M_1,\mathfrak M_2,\ldots]
 =[\mathfrak M_1,\ldots,\mathfrak M_n,
   [\mathfrak M_{n+1},\mathfrak M_{n+2},\ldots]],
\]
so
\[
 D([\mathfrak M_1,\mathfrak M_2,\ldots])
 =\sum_{i=1}^nD(\mathfrak M_i)
 +D([\mathfrak M_{n+1},\mathfrak M_{n+2},\ldots])
 \geq\sum_{i=1}^nD(\mathfrak M_i).
\]
Thus
$\sum_{i=1}^nD(\mathfrak M_i)\leq D([\mathfrak M_1,\mathfrak M_2,\ldots])$
for every $n=1,2,\ldots$, and therefore
\[
 \sum_{i=1}^\infty D(\mathfrak M_i)
 \leq D([\mathfrak M_1,\mathfrak M_2,\ldots]).
\]

Assume now, that $<$ does hold. Then $\sum_{i=1}^\infty D(\mathfrak M_i)$ is
finite. Therefore $\lim_{i\to\infty}D(\mathfrak M_i)=0$. So there exists for
every $\epsilon>0$ a finite $\mathfrak N\mathbin\eta\M$ (in fact an
$\mathfrak M_i$) with $0<D(\mathfrak N)<\epsilon$. Choose such an
$\mathfrak N$ for
\[
 \epsilon=D([\mathfrak M_1,\mathfrak M_2,\ldots])
 -\sum_{i=1}^\infty D(\mathfrak M_i)>0.
\]
Consider the expression
\[
 D([\mathfrak M_1,\mathfrak M_2,\ldots])
 -\sum_{i=1}^\infty D(\mathfrak M_i).
\]
It is equal to
\[
 \{D(\mathfrak M_1)+D([\mathfrak M_2,\mathfrak M_3,\ldots])\}
 -\sum_{i=1}^\infty D(\mathfrak M_i)
 =D([\mathfrak M_2,\mathfrak M_3,\ldots])
 -\sum_{i=2}^\infty D(\mathfrak M_i).
\]
Thus it is unchanged if we replace
$\mathfrak M_1,\mathfrak M_2,\ldots$ by
$\mathfrak M_2,\mathfrak M_3,\ldots$. Applying this $i_0$ times, we see, that
it is unchanged even if we pass to
$\mathfrak M_{i_0+1},\mathfrak M_{i_0+2},\ldots$. Now
$\sum_{i=1}^\infty D(\mathfrak M_{i_0+i})=
\sum_{i_0+1}^\infty D(\mathfrak M_i)$ is arbitrarily small if $i_0$ is
sufficiently great. So we can make it $\leq D(\mathfrak N)$. Now write again
$\mathfrak M_1,\mathfrak M_2,\ldots$ for
$\mathfrak M_{i_0+1},\mathfrak M_{i_0+2},\ldots$. Then we see:
$\sum_{i=1}^\infty D(\mathfrak M_i)\leq D(\mathfrak N)$ and still
\[
 D(\mathfrak N)<D([\mathfrak M_1,\mathfrak M_2,\ldots])
 -\sum_{i=1}^\infty D(\mathfrak M_i),
\]
and \emph{a fortiori}
$D(\mathfrak N)<D([\mathfrak M_1,\mathfrak M_2,\ldots])$.

We will now construct a sequence $\mathfrak N_1,\mathfrak N_2,\ldots$ of
mutually orthogonal sets $\mathfrak N_i\mathbin\eta\M$,
$\mathfrak N_i\subset\mathfrak N$ with $\mathfrak M_i\sim\mathfrak N_i$. We
proceed by induction. Assume that for an $i=1,2,\ldots$ all
$\mathfrak N_1,\ldots,\mathfrak N_{i-1}$ have already been constructed
fulfilling these conditions, then we construct $\mathfrak N_i$ as follows.
$D(\mathfrak N_j)=D(\mathfrak M_j)$, $j=1,\ldots,i-1$, is finite, so
\[
\begin{aligned}
 D(\mathfrak N-[\mathfrak N_1,\ldots,\mathfrak N_{i-1}])
 &=D(\mathfrak N)-\sum_{j=1}^{i-1}D(\mathfrak N_j)\\
 &\geq\sum_{j=1}^\infty D(\mathfrak M_j)
       -\sum_{j=1}^{i-1}D(\mathfrak N_j)\\
 &=\sum_{j=i}^\infty D(\mathfrak M_j)
 \geq D(\mathfrak M_i).
\end{aligned}
\]
Thus
$\mathfrak M_i\precsim
\mathfrak N-[\mathfrak N_1,\ldots,\mathfrak N_{i-1}]$, and thus
$\mathfrak M_i\sim\mathfrak N_i\subset
\mathfrak N-[\mathfrak N_1,\ldots,\mathfrak N_{i-1}]$. This definition of
$\mathfrak N_i$ meets all requirements.

Now $\mathfrak M_i\sim\mathfrak N_i$, all $\mathfrak M_i$ are mutually
orthogonal, and all $\mathfrak N_i$ are too, so \lemref{6.1.2} applies:
\[
 [\mathfrak M_1,\mathfrak M_2,\ldots]
 \sim[\mathfrak N_1,\mathfrak N_2,\ldots]\subset\mathfrak N,
 \qquad
 [\mathfrak M_1,\mathfrak M_2,\ldots]\precsim\mathfrak N,
\]
$D([\mathfrak M_1,\mathfrak M_2,\ldots])\leq D(\mathfrak N)$. But on the
other hand we had
$D(\mathfrak N)<D([\mathfrak M_1,\mathfrak M_2,\ldots])$ which is a
contradiction. So our original assumption concerning the $<$ must have been
wrong, and there must be for the original sequence
\[
 \sum_{i=1}^\infty D(\mathfrak M_i)
 =D([\mathfrak M_1,\mathfrak M_2,\ldots]).
\]
This completes the proof.

\Lemma{8.3.3}
{\itshape
Let $\mathfrak M_1,\mathfrak M_2,\ldots$ be an infinite sequence, all
$\mathfrak M_i\mathbin\eta\M$ and
$\mathfrak M_1\subset\mathfrak M_2\subset\cdots$. Then
\[
 \lim_{i\to\infty}D(\mathfrak M_i)
 =D([\mathfrak M_1,\mathfrak M_2,\ldots]).
\]}

\Proof Put $\mathfrak N_1=\mathfrak M_1$,
$\mathfrak N_i=\mathfrak M_i-\mathfrak M_{i-1}$ for $i=2,3,\ldots$. Then all
$\mathfrak N_i\mathbin\eta\M$ and they are mutually orthogonal. So
\lemref{8.3.2} applies to them and gives
\[
 \sum_{i=1}^\infty D(\mathfrak N_i)
 =D([\mathfrak N_1,\mathfrak N_2,\ldots]).
\]
Now $[\mathfrak N_1,\ldots,\mathfrak N_i]=\mathfrak M_i$ and so
$[\mathfrak N_1,\mathfrak N_2,\ldots]=
[\mathfrak M_1,\mathfrak M_2,\ldots]$ and furthermore
\[
 \sum_{i=1}^\infty D(\mathfrak N_i)
 =\lim_{n\to\infty}\sum_{i=1}^nD(\mathfrak N_i)
 =\lim_{n\to\infty}D([\mathfrak N_1,\ldots,\mathfrak N_n])
 =\lim_{n\to\infty}D(\mathfrak M_n).
\]
Therefore we have
$\lim_{i\to\infty}D(\mathfrak M_i)=
D([\mathfrak M_1,\mathfrak M_2,\ldots])$.

\Lemma{8.3.4}
{\itshape
Let $\mathfrak M_1,\mathfrak M_2,\ldots$ be an infinite sequence, all
$\mathfrak M_i\mathbin\eta\M$ and
$\mathfrak M_1\supset\mathfrak M_2\supset\cdots$. If not all $\mathfrak M_i$
are infinite, then
\[
 \lim_{i\to\infty}D(\mathfrak M_i)
 =D(\mathfrak M_1\cdot\mathfrak M_2\cdot\mathfrak M_3\cdots).
\]}

\Proof Assume that $\mathfrak M_{i_0}$ is finite. Put
$\mathfrak M_i'=\mathfrak M_{i_0}-\mathfrak M_{i_0+i}$. Then all
$\mathfrak M_i'\mathbin\eta\M$ and
$\mathfrak M_1'\subset\mathfrak M_2'\subset\cdots$, so \lemref{8.3.3} applies
to them and gives
\[
 \lim_{i\to\infty}D(\mathfrak M_i')
 =D([\mathfrak M_1',\mathfrak M_2',\ldots]).
\]
But
\[
 D(\mathfrak M_i')=D(\mathfrak M_{i_0})-D(\mathfrak M_{i_0+i})
\]
and
\begin{align*}
 [\mathfrak M_1',\mathfrak M_2',\ldots]
 &=[\mathfrak M_{i_0}-\mathfrak M_{i_0+1},
     \mathfrak M_{i_0}-\mathfrak M_{i_0+2},\ldots]\\
 &=\mathfrak M_{i_0}-(\mathfrak M_{i_0+1}\cdot
     \mathfrak M_{i_0+2}\cdots)\\
 &=\mathfrak M_{i_0}-(\mathfrak M_1\cdot\mathfrak M_2\cdots).
\end{align*}
Therefore we have
\[
 D(\mathfrak M_{i_0})-\lim_{i\to\infty}D(\mathfrak M_{i_0+i})
 =D(\mathfrak M_{i_0})-D(\mathfrak M_1\cdot\mathfrak M_2\cdots),
\]
so that
\[
 \lim_{i\to\infty}D(\mathfrak M_{i_0+i})
 =D(\mathfrak M_1\cdot\mathfrak M_2\cdots),
\]
or, which is the same thing,
$\lim_{i\to\infty}D(\mathfrak M_i)=
D(\mathfrak M_1\cdot\mathfrak M_2\cdots)$.

The following criterion is occasionally useful, because it contains a sufficient
condition for relative dimension functions which makes no explicit use of the
notions of finiteness and infiniteness for sets $\mathfrak M\mathbin\eta\M$.

\Lemma{8.3.5}
{\itshape
If a real-valued function $D'(\mathfrak M)$ which is defined for all linear,
closed sets $\mathfrak M\mathbin\eta\M$ assumes only values
$\geq0$, $\leq\infty$ and if some of its values are actually $\ne0,\infty$ then
the conditions (ii), (iii) in \defref{8.2.1} are sufficient in order that
$D'(\mathfrak M)$ be a relative dimension function.}

\Proof We must derive (i) in \defref{8.2.1} from these conditions. Choose an
$\overline{\mathfrak M}\mathbin\eta\M$ with
$D'(\overline{\mathfrak M})>0$, $<\infty$.

As $(0)$ is orthogonal to $\overline{\mathfrak M}$ and
$[\overline{\mathfrak M},(0)]=\overline{\mathfrak M}$; (iii) gives
$D'(\overline{\mathfrak M})+D'(0)=D'(\overline{\mathfrak M})$, as
$D'(\overline{\mathfrak M})$ is finite, this means $D'(0)=0$.

If any $\mathfrak M$ is infinite, then $\Hspace$ is it, too, and
$\mathfrak M\sim\Hspace$. Then by (ii) we must only prove
$D'(\Hspace)=\infty$. By \lemref{7.2.3} applied to $\Hspace$ we have
$2D'(\Hspace)=D'(\Hspace)$, $D'(\Hspace)=0$ or $\infty$. But
\[
 D'(\Hspace)=D'(\overline{\mathfrak M})
 +D'(\Hspace-\overline{\mathfrak M})
 \geq D'(\overline{\mathfrak M})>0,
\]
so $D'(\Hspace)=\infty$. This disposes of all infinite $\mathfrak M$'s.

If $\mathfrak M$ is finite and $\ne(0)$ then apply \lemref{7.1.3} to
$\mathfrak M,\overline{\mathfrak M}$ and to
$\overline{\mathfrak M},\mathfrak M$ (for its $\mathfrak M,\mathfrak N$).
Then (ii), (iii) give
\[
 D'(\mathfrak M)\leq
 \left(\reldiv{\mathfrak M}{\overline{\mathfrak M}}+1\right)
 D'(\overline{\mathfrak M}).
\]
Thus $D'(\overline{\mathfrak M})<\infty$ necessitates
$D'(\mathfrak M)<\infty$. Next
\[
 D'(\overline{\mathfrak M})\leq
 \left(\reldiv{\overline{\mathfrak M}}{\mathfrak M}+1\right)
 D'(\mathfrak M).
\]
Thus $D'(\overline{\mathfrak M})>0$ necessitates $D'(\mathfrak M)>0$. So
$0<D'(\mathfrak M)<\infty$, completing the proof.

\SectionStart{8.4}{\texorpdfstring{Range of $D(\mathfrak M)$. The cases
(I)--(III). Theorem VIII. Problems 3--7}{Range of D(M). The cases (I)--(III).
Theorem VIII. Problems 3--7}}
We now undertake to find invariant characteristics of $\M$ in terms of its
relative dimension functions $D(\mathfrak M)$.

\Lemma{8.4.1}
{\itshape
The range $\Delta$ of $D(\mathfrak M)$ (that is the set of all values of a given
$D(\mathfrak M)$ for all $\mathfrak M\mathbin\eta\M$) has these properties:
\begin{enumerate}[label=(\roman*)]
\item The elements of $\Delta$ are real numbers, $\geq0$, $\leq\infty$.
\item $0\in\Delta$, $\Delta$ contains a greatest element $\overline\alpha>0$.
\item $\alpha,\beta\in\Delta$ and $\alpha>\beta$ imply
  $\alpha-\beta\in\Delta$.
\item $\alpha_1,\alpha_2,\ldots\in\Delta$ and
  $\sum_{i=1}^{\infty}\alpha_i\leq\overline\alpha$ imply
  $\sum_{i=1}^{\infty}\alpha_i\in\Delta$.
\end{enumerate}}

\Proof Ad (i): Obvious.

Ad (ii): $D((0))=0$; $D(\Hspace)=\overline\alpha$,
$\Hspace\ne(0)$ therefore $D(\Hspace)>0$.

Ad (iii): $\beta$ must be finite. Choose
$\mathfrak M,\mathfrak N\mathbin\eta\M$ with
$D(\mathfrak M)=\alpha$, $D(\mathfrak N)=\beta$. Then
$\mathfrak N\precsim\mathfrak M$,
$\mathfrak N\sim\mathfrak N'\subset\mathfrak M$,
$D(\mathfrak N')=D(\mathfrak N)=\beta$ finite and so
$D(\mathfrak M-\mathfrak N')=D(\mathfrak M)-D(\mathfrak N')
=\alpha-\beta$.

Ad (iv): If any $\alpha_j=\infty$ then
$\sum_{i=1}^{\infty}\alpha_i=\infty$, so
$\sum_{i=1}^{\infty}\alpha_i=\alpha_j\in\Delta$, so we may assume that all
$\alpha_i$ are finite. Define a sequence $\mathfrak M_1,\mathfrak M_2,\ldots$
of mutually orthogonal sets $\mathfrak M_i\mathbin\eta\M$,
$D(\mathfrak M_i)=\alpha_i$ by induction. Assume that for an
$i=1,2,\ldots$, $\mathfrak M_1,\ldots,\mathfrak M_{i-1}$ have already been
constructed fulfilling these conditions, then construct $\mathfrak M_i$ as
follows.
\begin{align*}
 D(\Hspace-[\mathfrak M_1,\ldots,\mathfrak M_{i-1}])
 &=D(\Hspace)-\sum_{j=1}^{i-1}D(\mathfrak M_j)\\
 &=\overline\alpha-\sum_{j=1}^{i-1}\alpha_j
 \geq\sum_{j=1}^{\infty}\alpha_j-\sum_{j=1}^{i-1}\alpha_j
 =\sum_{j=i}^{\infty}\alpha_j\geq\alpha_i.
\end{align*}
Choose an $\mathfrak M_i'\mathbin\eta\M$ with
$D(\mathfrak M_i')=\alpha_i$; then
$\mathfrak M_i'\precsim\Hspace-[\mathfrak M_1,\ldots,\mathfrak M_{i-1}]$,
$\mathfrak M_i'\sim\mathfrak M_i\subset
\Hspace-[\mathfrak M_1,\ldots,\mathfrak M_{i-1}]$. This definition of
$\mathfrak M_i$ meets all requirements. Now by \lemref{8.3.2}
\[
 D([\mathfrak M_1,\mathfrak M_2,\ldots])
 =\sum_{i=1}^{\infty}D(\mathfrak M_i)
 =\sum_{i=1}^{\infty}\alpha_i.
\]

\Lemma{8.4.2}
{\itshape
The only sets $\Delta$ which fulfill the requirements (i)--(iv) of
\lemref{8.4.1} are the following ones ($\bar\epsilon,\overline\alpha$ are finite):
\begin{description}[style=multiline,leftmargin=4.4em]
\item[$(\mathrm I_n)$] $\bar\epsilon>0$, $n=1,2,\ldots$;
  $\Delta$ consists of all $i\bar\epsilon$, $i=0,1,\ldots,n$.
\item[$(\mathrm I_\infty)$] $\bar\epsilon>0$;
  $\Delta$ consists of all $i\bar\epsilon$, $i=0,1,\ldots,\infty$.
\item[$(\mathrm{II}_1)$] $\overline\alpha>0$;
  $\Delta$ consists of all $\alpha$, $0\leq\alpha\leq\overline\alpha$.
\item[$(\mathrm{II}_\infty)$] $\Delta$ consists of all $\alpha$,
  $0\leq\alpha\leq\infty$.
\item[$(\mathrm{III}_\infty)$] $\Delta$ consists of $0,\infty$.
\end{description}}

\Proof $\Delta$ contains $0$ by (ii) and some element $>0$ by (i), (ii). If it
contains no element $>0$, $<\infty$ then it consists of $0,\infty$ which is
precisely case $(\mathrm{III}_\infty)$. Assume now that it does contain
elements $>0$, $<\infty$.

Assume first, that this set contains a smallest element $\bar\epsilon$. If
$\alpha\in\Delta$ and finite, then there exists an $i=0,1,2,\ldots$ with
$i\bar\epsilon\leq\alpha<(i+1)\bar\epsilon$. If
$i\bar\epsilon\ne\alpha$ then by (iii) all
$\alpha,\alpha-\bar\epsilon,\ldots,\alpha-i\bar\epsilon$ belong to $\Delta$.
But $0<\alpha-i\bar\epsilon<\bar\epsilon$, which is impossible. So
$\alpha=i\bar\epsilon$. In particular
$\overline\alpha=n\bar\epsilon$, $n=0,1,2,\ldots$ if it is finite, and as
$\overline\alpha>0$, $n\ne0$. If $\overline\alpha$ is infinite we may write
$\overline\alpha=\infty\bar\epsilon$, so at any rate
$\overline\alpha=n\bar\epsilon$, $n=1,2,\ldots,\infty$.

For all other $\alpha\in\Delta$, $\alpha<\overline\alpha$ so $\alpha$ finite,
$\alpha=i\bar\epsilon$ and as $\alpha<\overline\alpha$, $i<n$,
$i=0,1,2,\ldots,n-1$. Conversely, if $i=0,1,\ldots,n-1$, then
$i\bar\epsilon<n\bar\epsilon=\overline\alpha$ and (iv) with
$\alpha_1=\cdots=\alpha_i=\bar\epsilon$,
$\alpha_{i+1}=\alpha_{i+2}=\cdots=0$ gives $i\bar\epsilon\in\Delta$. Thus
$\Delta$ consists of all $i\bar\epsilon$, $i=0,1,\ldots,n$, which is case
$(\mathrm I_n)$ if $n=1,2,\ldots$, and case $(\mathrm I_\infty)$ if
$n=\infty$.

Assume next that there is no smallest element $>0$, $<\infty$ in $\Delta$.
Let $\epsilon'$ be the g.r.l.b. of the elements $\epsilon>0$; if we had
$\epsilon'>0$ then there would exist an $\alpha\in\Delta$, $\alpha>0$ with
$\alpha<2\epsilon'$ and as $\alpha=\epsilon'$ would imply that $\alpha$ is the
smallest element $>0$, $<\infty$ in $\Delta$, therefore $\alpha>\epsilon'$.
Now there exists a $\beta\in\Delta$, $\beta>0$ with $\beta<\alpha$ and
$\beta>\epsilon'$. So $\alpha,\beta\in\Delta$, $\alpha>\beta$ and by (iii)
$\alpha-\beta\in\Delta$, but $\alpha-\beta>0$,
$\alpha-\beta<2\epsilon'-\epsilon'=\epsilon'$ which is impossible. Thus
$\epsilon'=0$ and so for every $\epsilon>0$ an $\alpha\in\Delta$,
$0<\alpha<\epsilon$ with $0<\alpha<\epsilon$ exists.

Now assume $0\leq\beta_1<\beta_2\leq\overline\alpha$. Choose an
$\alpha\in\Delta$ with $0<\alpha<\beta_2-\beta_1$. Then an
$i=0,1,2,\ldots$ with $i\alpha\leq\beta_1<(i+1)\alpha$ exists, so
\[
 (i+1)\alpha=i\alpha+\alpha<\beta_1+\beta_2-\beta_1=\beta_2.
\]
As $(i+1)\alpha<\overline\alpha$, therefore (iv) with
$\alpha_1=\cdots=\alpha_{i+1}=\alpha$,
$\alpha_{i+2}=\alpha_{i+3}=\cdots=0$ gives $(i+1)\alpha\in\Delta$. So we see:
There exists a $\beta\in\Delta$ with $\beta_1<\beta<\beta_2$.

Now consider any $\alpha$ with $0<\alpha<\overline\alpha$ and form a sequence
$\beta_1,\beta_2,\ldots$ with
$0<\beta_1<\beta_2<\cdots<\alpha$;
$\lim_{i\to\infty}\beta_i=\alpha$. For each $i$ choose a
$\beta_i'\in\Delta$ with $\beta_i<\beta_i'<\beta_{i+1}$. So
$\beta_i'>\beta_{i-1}'$, $\beta_i'-\beta_{i-1}'\in\Delta$, $i=2,3,\ldots$.
Put $\alpha_1=\beta_1'$, $\alpha_i=\beta_i'-\beta_{i-1}'$ for
$i=2,3,\ldots$, and apply (iv):
$\sum_{i=1}^\infty\alpha_i=\lim_{i\to\infty}\beta_i'=\alpha
<\overline\alpha$ and so $\alpha\in\Delta$.

Thus $\Delta$ consists of all $\alpha$ with
$0\leq\alpha\leq\overline\alpha$ which is case $(\mathrm{II}_1)$ if
$\overline\alpha$ is finite, and case $(\mathrm{II}_\infty)$ if
$\overline\alpha=\infty$. These considerations exhaust all possibilities.

We now apply Lemmas \hyperref[lem:8.4.1]{8.4.1} and
\hyperref[lem:8.4.2]{8.4.2} simultaneously, remembering
that we can normalise $D(\mathfrak M)$ so as to make $\bar\epsilon=1$ (cases
$(\mathrm I_n)$ and $(\mathrm I_\infty)$) resp. $\overline\alpha=1$ (case
$(\mathrm{II}_1)$). Then we have:

\MainTheorem{VIII}
{\itshape
The range of $D(\mathfrak M)$ (that is the set of all values of a given
$D(\mathfrak M)$ for all $\mathfrak M\mathbin\eta\M$) is one of the following
sets:
\begin{description}[style=multiline,leftmargin=4.7em]
\item[$(\mathrm I_n)$] $n=1,2,\ldots$. The set $0,1,\ldots,n$.
\item[$(\mathrm I_\infty)$] The set $0,1,\ldots,\infty$.
\item[$(\mathrm{II}_1)$] The set of all $\alpha$, $0\leq\alpha\leq1$.
\item[$(\mathrm{II}_\infty)$] The set of all $\alpha$,
  $0\leq\alpha\leq\infty$.
\item[$(\mathrm{III}_\infty)$] The set $0,\infty$.
\end{description}}

In the cases $(\mathrm I_n)$, $(\mathrm I_\infty)$, $(\mathrm{II}_1)$ we have
included a normalisation of $D(\mathfrak M)$, the \emph{standard
normalisation}. Case $(\mathrm{II}_\infty)$ is not normalised; in case
$(\mathrm{III}_\infty)$ no normalisation is needed, because $0,\infty$ are
unaltered by multiplication with a finite positive number.

We call the cases $(\mathrm I_n)$, $(\mathrm{II}_1)$ the \emph{finite cases},
the cases $(\mathrm I_\infty)$, $(\mathrm{II}_\infty)$,
$(\mathrm{III}_\infty)$, the \emph{infinite cases}. On the other hand we call
the cases (I) (that is $(\mathrm I_n)$, $(\mathrm I_\infty)$) the
\emph{discrete cases}, the cases (II), (that is $(\mathrm{II}_1)$,
$(\mathrm{II}_\infty)$) the \emph{continuous cases}, and the case (III), (that
is $(\mathrm{III}_\infty)$) the \emph{purely infinite case}.

The following problems arise immediately:

\Problem{3}\emph{Which ones of the classes
$(\mathrm I_n)$--$(\mathrm{III}_\infty)$ do really exist?}

\Problem{4}\emph{Which combinations of these classes do occur for coupled
factors $\M,\M'$? For general factorisations $\M_1,\ldots,\M_n$?}

\Problem{5}\emph{To which classes do the direct factors $\M$ belong?}

\Problem{6}\emph{Are all factors $\M$ belonging to the same class
ring-isomorphic, or do there exist further invariant characteristics?}

\Problem{7}\emph{Are all factorisations $\M_1,\ldots,\M_n$ in which each
$\M_i$ belongs to a given class ring-isomorphic or even spatially isomorphic,
or do there exist further invariant characteristics?}

We will be able to contribute considerable material to clarify these questions,
but the only one to which we can give a complete answer is \probref{5}.

\SectionStart{8.5}{Invariance under ring-isomorphisms. Theorem IX}
We establish the invariance of the subjects of the last
\hyperref[sec:8.3]{\S}\hyperref[sec:8.4]{\S}\ under
ring-isomorphisms.

\Lemma{8.5.1}
{\itshape
Replace the $\mathfrak M,\mathfrak N\mathbin\eta\M$ by their
$E=P_{\mathfrak M}$, $F=P_{\mathfrak N}$ which are $\in\M$
(Cf. Definitions \hyperref[def:6.1.1]{6.1.1},
\hyperref[def:7.1.1]{7.1.1}, and \hyperref[def:8.2.1]{8.2.1}). The notions
$E\sim F$ $(\cdots\M)$, finiteness and infiniteness of an $E$ with respect to
$\M$, $D_{\M}(E)$ (apart from its normalisation) are invariant under algebraic
ring-isomorphisms.}

\Proof $E\sim F$ means that a $U\in\M$ with
$UU^*U=U$, $U^*U=E$, $UU^*=F$ exists. $E$ is infinite, if an $F\in\M$ with
$F^2=F^*=F$, $EF=F$, $F\ne E$, $F\sim E$ exists. $D(E)$ is characterised in
\defref{8.2.1} by the sole use of the notions $0$, finiteness, infiniteness,
$E\sim F$ and $E+F$ with $EF=0$. Thus all these constructions are invariant
under algebraic ring-isomorphisms.

Now we can state, combining \lemref{8.5.1} and \thmref{VIII}:

\MainTheorem{IX}
{\itshape
The cases $(\mathrm I_n)$--$(\mathrm{III}_\infty)$ of \thmref{VIII} are
invariant under algebraic ring-isomorphisms. So is $D_{\M}(E)$ (we write
$E=P_{\mathfrak M}\in\M$ for $\mathfrak M\mathbin\eta\M$), apart from its
normalisation, and even the standard normalisation (in the cases
$(\mathrm I_n)$--$(\mathrm{II}_1)$).}

\SectionStart{8.6}{The discrete cases}
We now locate the direct factors in this classification.

\Lemma{8.6.1}
{\itshape
$\M$ is a direct factor if and only if it belongs to the discrete cases:
$(\mathrm I_n)$ or $(\mathrm I_\infty)$. If we then use the spatial isomorphism
of $\Hspace$ with $\Hspace_1\otimes\cdots\otimes\Hspace_m$ where $\M$ becomes
$\B_i^{(i)}$ then the corresponding $\Hspace_i$ will be an $n$-dimensional
Euclidean space resp. a Hilbert space. If $E_i=P_{\mathfrak M_i}$ is a
projection in $\Hspace_i$ then $D_{\M}(E_i^{(i)})$ is, in the standard
normalisation, simply the common notion of dimension of $\mathfrak M_i$.}

\Proof In the cases (I) an $\mathfrak M\mathbin\eta\M$ with
$D(\mathfrak M)=1$ exists. If $\mathfrak N\mathbin\eta\M$,
$\mathfrak N\subset\mathfrak M$ then $D(\mathfrak N)\leq D(\mathfrak M)$,
$D(\mathfrak N)=0,1$. The former implies $\mathfrak N=(0)$, the latter
$\mathfrak N\sim\mathfrak M$ and as $\mathfrak M$ is finite, it excludes
$\mathfrak N\subsetneqq\mathfrak M$, so it implies
$\mathfrak N=\mathfrak M$. Obviously $\mathfrak M\ne(0)$. Thus
$\mathfrak M$ is minimal, and therefore $\M$ is a direct factor
(cf. \thmref{IV}). So our condition is sufficient.

All other statements assume that $\M$ is a direct factor, and thus
algebraically (even fully) ring-isomorphic to $\B_i$. Then we may replace (by
\thmref{IX}) $\Hspace,\M$, by $\Hspace_i,\B_i$; that is, we may assume
$\M=\B$.

Every one-dimensional $\mathfrak M=[\varphi]$, $\varphi\ne0$ (now all
$\mathfrak M\mathbin\eta\M$) is $\ne(0)$ and finite
($\mathfrak N\subsetneqq\mathfrak M$ implies $\mathfrak N=(0)$ excluding
$\mathfrak N\sim\mathfrak M$) so
$D_{\B}([\varphi])>0$, $<\infty$. Besides
$[\varphi]\sim[\psi]$ so
$D_{\B}([\varphi])=D_{\B}([\psi])$; $D_{\B}([\varphi])$ is a constant.
Normalise $D_{\B}(\mathfrak M)$ so as to have $D_{\B}([\varphi])=1$. Then
\lemref{8.3.2} gives $D_{\B}(\mathfrak M)=$ common notion of dimension of
$\mathfrak M$.

Thus if $\Hspace$ is an $n$-dimensional Euclidean space, this
$D_{\B}(\mathfrak M)$ has the range $0,1,\ldots,n$; and if it is a Hilbert
space, it has the range $0,1,\ldots,\infty$. This proves that $\M$ is in the
cases $(\mathrm I_n)$ resp. $(\mathrm I_\infty)$, and that this normalisation of
$D_{\B}(\mathfrak M)$ is the standard one. This completes the proof.

\lemref{8.6.1} makes it possible to solve Problems
\hyperref[prob:3]{3}--\hyperref[prob:7]{7} in the
discrete cases (that is for direct factors). The details are these:

The discrete cases $(\mathrm I_n)$, for every $n=1,2,\ldots$ and
$(\mathrm I_\infty)$ do occur; and for factorisations
$\M_1,\ldots,\M_m$ we can prescribe any combination of the discrete cases
$(\mathrm I_n)$ and $(\mathrm I_\infty)$ for the $\M_i$'s. We need only form
$\Hspace=\Hspace_1\otimes\cdots\otimes\Hspace_m$ where the $\Hspace_i$ are
the corresponding $n$-dimensional Euclidean or Hilbert spaces, and put
$\M_i=\B_i^{(i)}$. Furthermore if $m-1$ of the factors $\M_i$ are in discrete
cases, then all are (\lemref{5.4.1}); and so in particular two factors
$\M_1,\M_2$ (and \emph{a fortiori} two coupled factors) are either both in
discrete cases, or none of them is. The direct factors are identical with those
in discrete cases. If in a factorisation $\M_1,\ldots,\M_m$ all $\M_i$ are in
discrete cases, then they are all direct factors, so $\M_1,\ldots,\M_m$ is a
direct factorisation (\lemref{5.4.1}). Thus a spatial isomorphism carries
$\Hspace$ into $\Hspace_1\otimes\cdots\otimes\Hspace_m$ and each $\M_i$ into
$\B_i^{(i)}$; therefore $\Hspace_i$ is determined by the class of $\M_i$
(cf. above). Thus if two factors $\M,\N$ have the same discrete class, they are
fully ring-isomorphic (\lemref{2.3.4} and \cite[\S4]{ref22}, compare the
factorisations $\M,\M'$ and $\N,\N'$); and if in two factorisations
$\M_1,\ldots,\M_m$ and $\N_1,\ldots,\N_m$ the corresponding $\M_i,\N_i$
have the same discrete classes for $i=1,\ldots,m$, then we have even spatial
isomorphism.
% --- End of parts/03b-chapter-viii.tex ---
\PartHeadingTOC{III}{Pairs of factors}{Pairs of Factors}

\ChapterHeadingTOC{9}{IX}{Qualitative Comparison of
$\mathfrak M_f^{\M'}$ and $\mathfrak M_f^{\M}$}
{\texorpdfstring{Qualitative comparison of $\mathfrak M_f^{\M'}$ and
 $\mathfrak M_f^{\M}$}{Qualitative comparison of two cyclic subspaces}}

\SectionStart{9.1}{\texorpdfstring{Definite series
$\sum_{i=0}^{\infty}A_i$. The operator $B^{1/2}$}{Definite series. The
operator B to the one-half}}
We need first a discussion of infinite sums of definite operators. The
considerations which follow are based on a construction of K. Friedrichs,
\cite[pp.~472, 476]{ref7}. The situation which we will investigate is
described as follows:

\Definition{9.1.1}
Let $A_1,A_2,\ldots$ be a sequence of everywhere defined, bounded,
Hermitian, (semi-) definite operators. In other words; For each
$i=1,2,\ldots$ an $\alpha_i$ exists, so that we have identically
\[
 0\leq(A_if,f)\leq\alpha_i\lVert f\rVert^2
\]
for all $f\in\Hspace$ (Cf. \hyperlink{cite.ref16}{(16)} pp.~73--74).

Put $A_0=1$. Define $\mathfrak A$ as the set of $f\in\Hspace$ for which
$\sum_{i=0}^{\infty}(A_if,f)$ is finite. For any two $f,g\in\Hspace$ for
which $\sum_{i=0}^{\infty}(A_if,g)$ is absolutely convergent, call this
$Q(f,g)$.

\Lemma{9.1.1}
{\itshape
$\mathfrak A$ is a linear set. If $f,g\in\mathfrak A$ then $Q(f,g)$ is
defined. With $\alpha f,f+g$ as linear operations and $Q(f,g)$ as inner
product $\mathfrak A$ fulfills the conditions A, B, E of
\hyperlink{cite.ref16}{(16)} pp.~64--66. (We may however have
$\mathfrak A=(0)$.)}

\Proof As $(A_i\alpha f,\alpha f)=|\alpha|^2(A_if,f)$ therefore
$f\in\mathfrak A$ implies $\alpha f\in\mathfrak A$. As
\begin{align*}
 (A_i(f+g),f+g)
 &\leq(A_i(f+g),f+g)+(A_i(f-g),f-g)\\
 &=2(A_if,f)+2(A_ig,g)
\end{align*}
therefore $f,g\in\mathfrak A$ imply $f+g\in\mathfrak A$. Thus
$\mathfrak A$ is a linear set. We have
\[
 |(A_if,g)|\leq\tfrac12(A_if,f)+\tfrac12(A_ig,g),
\]
(this is proved literally as Schwarz's inequality cf. \cite[p.~64]{ref16}),
thus if $f,g\in\mathfrak A$ then $\sum_{i=0}^{\infty}(A_if,g)$ is absolutely
convergent, and so $Q(f,g)$ is defined.

We now verify A, B, E, loc. cit. All parts of A, B are obvious, except B, 4.
This holds, as $f\ne0$ implies
\[
 Q(f,f)=\sum_{i=0}^{\infty}(A_if,f)\geq(A_0f,f)=\lVert f\rVert^2>0.
\]
Consider now E. We must prove: If $f_1,f_2,\ldots\in\mathfrak A$ and
\[
 \lim_{m,n\to\infty}Q(f_m-f_n,f_m-f_n)=0
\]
then there exists an $f\in\mathfrak A$ with
$\lim_{n\to\infty}Q(f_n-f,f_n-f)=0$.

As
$Q(f_m-f_n,f_m-f_n)\geq\lVert f_m-f_n\rVert^2$ we have
$\lim_{m,n\to\infty}\lVert f_m-f_n\rVert=0$. So an $f\in\Hspace$ with
$\lim_{n\to\infty}\lVert f_n-f\rVert=0$ exists. As
\[
 \lim_{m,n\to\infty}Q(f_m-f_n,f_m-f_n)=0
\]
there exists for every $\epsilon>0$ an $n_0=n_0(\epsilon)$ so that
$m,n\geq n_0$ imply
\[
 Q(f_m-f_n,f_m-f_n)\leq\epsilon
\]
and a fortiori
\[
 \sum_{i=0}^{j}(A_i(f_m-f_n),f_m-f_n)\leq\epsilon.
\]
Let now $n\to\infty$; as all $A_i$ are continuous, this gives
\[
 \sum_{i=0}^{j}(A_i(f_m-f),f_m-f)\leq\epsilon.
\]
This holds for every $j=0,1,2,\ldots$, therefore
\[
 \sum_{i=0}^{\infty}(A_i(f_m-f),f_m-f)\leq\epsilon.
\]
This proves $f_m-f\in\mathfrak A$,
$f=f_m-(f_m-f)\in\mathfrak A$, and then
$Q(f_m-f,f_m-f)\leq\epsilon$. So we have $f\in\mathfrak A$ and
$\lim_{m\to\infty}Q(f_m-f,f_m-f)=0$. This completes the proof.

\Lemma{9.1.2}
{\itshape
For every $f\in\Hspace$ there exists one and only one $f^*\in\mathfrak A$ so
that $(f,g)=Q(f^*,g)$ for every $g\in\mathfrak A$. We have
$\lVert f^*\rVert\leq\lVert f\rVert$.}

\Proof Consider $L(g)=(f,g)$ for $g\in\mathfrak A$ as a functional in
$\mathfrak A$. We have obviously 1) $L(\alpha g)=\overline\alpha L(g)$,
2) $L(g_1+g_2)=L(g_1)+L(g_2)$. Furthermore we have
\[
 |L(g)|=|(f,g)|\leq\lVert f\rVert\cdot\lVert g\rVert
 \leq\lVert f\rVert\sqrt{Q(g,g)},
\]
that is 3) $|L(g)|\leq a_0\sqrt{Q(g,g)}$, where
$a_0=\lVert f\rVert$. Under these conditions the existence of an
$f^*\in\mathfrak A$ is certain, for which $L(g)=Q(f^*,g)$ for all
$g\in\mathfrak A$ and this $f^*$ is unique. (Cf. \cite[p.~11]{ref11}, and
\cite[p.~34]{ref13a}.) In other words: $(f,g)=Q(f^*,g)$, $g=f^*$ gives in
particular:
\[
 Q(f^*,f^*)=(f,f^*)\leq\lVert f\rVert\cdot\lVert f^*\rVert;
\]
but as $Q(f^*,f^*)\geq\lVert f^*\rVert^2$ this implies
$\lVert f^*\rVert\leq\lVert f\rVert$.

\Definition{9.1.2}
Define an operator $B$ by $Bf=f^*$, in the sense of \lemref{9.1.2}.

\Lemma{9.1.3}
{\itshape
$B$ is everywhere defined, bounded, Hermitian, (semi-) definite (in
$\Hspace$); and $\operatorname{Range}B\subset\mathfrak A$.}

\Proof $B$ is obviously everywhere defined.
$(Bf,f)=(f^*,f)=Q(f^*,f^*)$ and
\[
 0\leq Q(f^*,f^*)\leq\lVert f^*\rVert\cdot\lVert f\rVert
 \leq\lVert f\rVert^2,
\]
thus $B$ is bounded, Hermitian, and (semi-) definite. As
$Bf=f^*\in\mathfrak A$ we have $\operatorname{Range}B\subset\mathfrak A$.

\Lemma{9.1.4}
{\itshape
Form $B^{1/2}$ in the sense of F. Riesz (cf. in the proof). $B^{1/2}$ is
everywhere defined, bounded, Hermitean, (semi-) definite.
$\operatorname{Range}B^{1/2}=\mathfrak A$. For
$f\in\Hspace-[\mathfrak A]$, $B^{1/2}f=0$, while for
$f,g\in[\mathfrak A]$, $(f,g)=Q(B^{1/2}f,B^{1/2}g)$.}

\Proof As to the definition and character of $B^{1/2}$ cf. for instance
\cite[p.~113]{ref16}.

$B^{1/2}f=0$ implies $Bf=B^{1/2}(B^{1/2}f)=0$; $Bf=0$ implies
$\lVert B^{1/2}f\rVert^2=(B^{1/2}f,B^{1/2}f)=(Bf,f)=0$,
$B^{1/2}f=0$; $Bf=0$ means that the $f^*$ of \lemref{9.1.2} is 0, that is
$(f,g)=0$ for $g\in\mathfrak A$. This means $f\in\Hspace-[\mathfrak A]$.
So we have:
\[
 (f;B^{1/2}f=0)=(f;Bf=0)=\Hspace-[\mathfrak A].
\]
\[
 [\operatorname{Range}B^{1/2}]=[\operatorname{Range}B]=[\mathfrak A].
\]

Consider an $f\in[\mathfrak A]$. Then $f$ is a condensation point of
$\operatorname{Range}B^{1/2}$ so a sequence $f_1,f_2,\ldots$ with
$\lim_{n\to\infty}\lVert B^{1/2}f_n-f\rVert=0$ exists. Now all
$Bf_n\in\mathfrak A$ and
\begin{align*}
 Q(Bf_m-Bf_n,Bf_m-Bf_n)
 &=(f_m-f_n,Bf_m-Bf_n)\\
 &=(B^{1/2}f_m-B^{1/2}f_n,
    B^{1/2}f_m-B^{1/2}f_n)\\
 &=\lVert B^{1/2}f_m-B^{1/2}f_n\rVert^2,
\end{align*}
so $\lim_{m,n\to\infty}Q(Bf_m-Bf_n,Bf_m-Bf_n)=0$. But
$\mathfrak A$ fulfills the completeness postulate E (for $Q(\ldots,\ldots)$
cf. \lemref{9.1.1}), so an $f^0\in\mathfrak A$ with
\[
 \lim_{n\to\infty}Q(Bf_n-f^0,Bf_n-f^0)=0
\]
exists. Now $Q(g,g)\geq\lVert g\rVert^2$ so
$\lim_{n\to\infty}\lVert Bf_n-f^0\rVert=0$. On the other hand the
continuity of $B^{1/2}$ implies
\[
 \lim_{n\to\infty}\lVert Bf_n-B^{1/2}f\rVert
 =\lim_{n\to\infty}\lVert B^{1/2}(B^{1/2}f_n-f)\rVert=0.
\]
So $f^0=B^{1/2}f$, proving $B^{1/2}f\in\mathfrak A$. Besides we found
\[
 \lim_{n\to\infty}Q(Bf_n-B^{1/2}f,Bf_n-B^{1/2}f)=0.
\]

Consider two $f,g\in[\mathfrak A]$. Then
$B^{1/2}f,B^{1/2}g\in\mathfrak A$ and forming the above sequence
$f_1,f_2,\ldots$ for $f$ and its analogue $g_1,g_2,\ldots$ for $g$, the
continuity of the inner product of $\mathfrak A$ in the metric of
$\mathfrak A$ (all $Q(\ldots,\ldots)$ cf. \cite[p.~65]{ref16}) necessitates
\begin{align*}
 Q(B^{1/2}f,B^{1/2}g)
 &=\lim_{n\to\infty}Q(Bf_n,Bg_n)
  =\lim_{n\to\infty}(f_n,Bg_n)\\
 &=\lim_{n\to\infty}(B^{1/2}f_n,B^{1/2}g_n)=(f,g).
\end{align*}

Thus $B^{1/2}$ maps $[\mathfrak A]$ on some subset $\mathfrak A'$ of
$\mathfrak A$ and this mapping is isomorphic, if we use the inner product
$(f,g)$ in $[\mathfrak A]$ but $Q(f,g)$ in $\mathfrak A'$.

Therefore it is one-to-one. Now $[\mathfrak A]$ is a closed subset in the
topologically complete set $\Hspace$ (topology of the metric
$\lVert f-g\rVert=(f-g,f-g)^{1/2}$), therefore $\mathfrak A'$ is complete
too, and so it must be a closed subset of $\mathfrak A$ (topology of the
metric $[Q(f-g,f-g)]^{1/2}$). In this sense $\mathfrak A'$ is a closed linear
set in $\mathfrak A$. Consider therefore $\mathfrak A-\mathfrak A'$ (for the
inner product $Q(f,g)$). $f\in\mathfrak A-\mathfrak A'$ means
$f\in\mathfrak A$ and $Q(f,g)=0$ for every $g\in\mathfrak A'$ that is
$Q(f,B^{1/2}g')=0$ for every $g'\in[\mathfrak A]$. For every $g''$,
$g'=B^{1/2}g''\in[\mathfrak A]$ so we have a fortiori $Q(f,Bg'')=0$,
$(f,g'')=0$ for every $g''$, and so $f=0$. This proves
$\mathfrak A-\mathfrak A'=(0)$. Now every $f\in\mathfrak A$ is the sum of
an element of $\mathfrak A'$ and one of $\mathfrak A-\mathfrak A'$ (this is
the fundamental theorem on projections, which applies to spaces which
fulfill, like $\mathfrak A$, only A., B., E., by \cite[p.~14]{ref11}, or
\cite[p.~34]{ref13a}); as $\mathfrak A-\mathfrak A'=(0)$ this means
$f\in\mathfrak A'$. Thus we have proved $\mathfrak A'=\mathfrak A$. We see
that $B^{1/2}$ maps $[\mathfrak A]$ on $\mathfrak A$. In
$\Hspace-[\mathfrak A]$ (we use again the inner product $(f,g)$ of
$\Hspace$) however it is $=0$. Therefore
$\operatorname{Range}B^{1/2}=\mathfrak A$.

Thus we have proved all statements of our Lemma.

Another essential property of $B^{1/2}$ is this:

\Lemma{9.1.5}
{\itshape
Let $\M$ be a ring which contains 1. If the above $A_1,A_2,\ldots$ are
elements of $\M$ then $B^{1/2}\in\M$.}

\Proof $A_n\in\M$, $A_n\mathbin\eta\M$ so $A_n$ is invariant under every
unitary $U'\in\M'$. Therefore $\mathfrak A$ and $Q(f,g)$ are invariant under
these and with them $B$, which was uniquely characterised with their help.
So $B\mathbin\eta\M$ and by \lemref{4.2.1} $B\in\M$. This implies
$B^{1/2}\in\M$ for instance by \cite[p.~213]{ref19}.

\SectionStart{9.2}{\texorpdfstring{Application to
$\mathfrak M_{f_0}^{\M'}$}{Application to a cyclic subspace}}
Consider now the sets $\mathfrak M_{f_0}^{\M'}$ of \defref{5.1.1}.

\Lemma{9.2.1}
{\itshape
Let $\M$ be a ring which contains 1. If
$g_0\in\mathfrak M_{f_0}^{\M'}$ then two linear, closed operators
$X',Y'\mathbin\eta\M'$ with everywhere dense domains can be found such that
$g_0=X'Y'f_0$. (Of course $f_0\in\operatorname{Domain}Y'$,
$Y'f_0\in\operatorname{Domain}X'$. $X'$ is even bounded.)}

\Proof $g_0$ is a condensation point of the set of all $A'f_0$,
$A'\in\M'$. So we can find for every $n=1,2,\ldots$ an $A_n'\in\M'$ with
$\lVert A_n'f_0-g_0\rVert\leq 1/2^n$. Thus
\[
 \lVert A_{n+1}'f_0-A_n'f_0\rVert
 \leq 1/2^{n+1}+1/2^n=3/2^{n+1},
\]
and so
\[
\begin{aligned}
 &\bigl(2^n(A_{n+1}'-A_n')^*(A_{n+1}'-A_n')f_0,f_0\bigr)\\
 &\qquad=2^n\bigl((A_{n+1}'-A_n')f_0,
                  (A_{n+1}'-A_n')f_0\bigr)\\
 &\qquad=2^n\lVert(A_{n+1}'-A_n')f_0\rVert^2
 \leq 2^n(3/2^{n+1})^2=9/2^{n+2}.
\end{aligned}
\]
We now perform the construction of \defref{9.1.1} with
$2^n(A_{n+1}'-A_n')^*(A_{n+1}'-A_n')$ in place of $A_n$,
$n=1,2,\ldots$. We thus obtain $\mathfrak A$ and $Q(f,g)$, and the
$B,B^{1/2}$ of \defref{9.1.2} and \lemref{9.1.4}, which we will denote by
$B',B'^{1/2}$. \lemref{9.1.5} gives $B'^{1/2}\in\M'$. Note that our
above evaluation secures $f_0\in\mathfrak A$.

$B'^{1/2}$ is a one-to-one mapping of $[\mathfrak A]$ on $\mathfrak A$
(even isomorphic!, cf. \lemref{9.1.4}). Denote its inverse (in
$\mathfrak A$ only!) by $Y_0'$. As $B'^{1/2}\in\M'$ we have
$\mathfrak A=\operatorname{Range}B'^{1/2}\mathbin\eta\M$ and so
$Y_0'\mathbin\eta\M'$; by its very nature $Y_0'$ is linear and closed. But
$\operatorname{Domain}Y_0'=\mathfrak A$ is only dense in $[\mathfrak A]$.
If we put $Y'=Y_0'P_{[\mathfrak A]}$, then $Y'$ is clearly
$\mathbin\eta\M'$, linear, closed, and has an everywhere dense domain.

Consider an $f\in\Hspace$. Then
$B'^{1/2}f=B'^{1/2}P_{[\mathfrak A]}f\in\mathfrak A$ and
$Q(B'^{1/2}f,B'^{1/2}f)=\lVert P_{[\mathfrak A]}f\rVert^2\leq\lVert f\rVert^2$.
But
\begin{align*}
 Q(B'^{1/2}f,B'^{1/2}f)
 &=\bigl(B'^{1/2}f,B'^{1/2}f\bigr)\\
 &\quad+\sum_{n=1}^{\infty}
   \bigl(2^n(A_{n+1}'-A_n')^*(A_{n+1}'-A_n')B'^{1/2}f,
         B'^{1/2}f\bigr)\\
 &=\lVert B'^{1/2}f\rVert^2
   +\sum_{n=1}^{\infty}2^n
     \lVert(A_{n+1}'-A_n')B'^{1/2}f\rVert^2\\
 &\geq 2^n\lVert(A_{n+1}'-A_n')B'^{1/2}f\rVert^2
\end{align*}
for every $n=1,2,\ldots$; so
\[
 \lVert(A_{n+1}'-A_n')B'^{1/2}f\rVert
 \leq\frac1{\sqrt{2^n}}\lVert f\rVert.
\]
Therefore
\begin{align*}
 \lVert(A_{n+p}'-A_n')B'^{1/2}f\rVert
 &\leq\left(\frac1{\sqrt{2^n}}+\cdots+
              \frac1{\sqrt{2^{n+p-1}}}\right)\lVert f\rVert\\
 &\leq\frac1{\sqrt{2^n}}
       \left(\frac1{1-\frac1{\sqrt2}}\right)\lVert f\rVert
  =\frac{2+\sqrt2}{\sqrt{2^n}}\lVert f\rVert,
\end{align*}
or, which is the same thing,
\[
 \lVert(A_m'-A_n')B'^{1/2}f\rVert
 \leq\frac{2+\sqrt2}{\sqrt{2^{\min(m,n)}}}\lVert f\rVert.
\]
\begingroup\rightskip=0pt plus .8em
Thus $\lim_{m,n\to\infty}\lVert(A_m'-A_n')B'^{1/2}f\rVert=0$ and
therefore a unique $f^*$ with
$\lim_{n\to\infty}\lVert A_n'B'^{1/2}f-f^*\rVert=0$ exists. We define an
operator $X'$ by $X'f=f^*$; then $X'$ is clearly the uniform limit of the
sequence $A_1'B'^{1/2},A_2'B'^{1/2},\ldots$ (cf.
\cite[p.~384]{ref18}), thus bounded and $\in\M'$ too.\par\endgroup

Now $f_0\in\mathfrak A$, therefore $Y'f_0=Y_0'f_0\in[\mathfrak A]$ and
$B'^{1/2}Y'f_0=f_0$; $g_0$ is the limit of the sequence
$A_n'f_0=A_n'B'^{1/2}Y'f_0$, so $g_0=X'Y'f_0$. This completes the proof.

Note that we cannot make much use of the operator-product $X'Y'$ itself.
It is not closed in general, and it is uncertain whether it possesses a
single valued closure (cf. \cite[p.~300]{ref21}). We will see in
\hyperref[sec:16.3]{\S16.3}--\hyperref[sec:16.4]{16.4} how the notion of finiteness helps to bridge this gap; for
the moment this is unimportant, because we are able to deal with the
$X',Y'$ separately.

\SectionStart{9.3}{\texorpdfstring{Equivalence of
$\mathfrak M_{f_0}^{\M'}\gtreqless\mathfrak M_{g_0}^{\M'}\ (\cdots\M)$
with $\mathfrak M_{f_0}^{\M}\gtreqless\mathfrak M_{g_0}^{\M}\ (\cdots\M')$,
resp.}{Equivalence of the two comparisons of cyclic subspaces, resp.}}
We will compare the $\mathfrak M_f^{\M'}$ and the $\mathfrak M_f^{\M}$.
From now on we assume that $\M$ is a factor.

\Lemma{9.3.1}
{\itshape
If $g_0=X'f_0$ where $X'\mathbin\eta\M'$ is a linear, closed operator with
an everywhere dense domain, then
\[
 \mathfrak M_{g_0}^{\M}\precsim\mathfrak M_{f_0}^{\M}
 \quad(\cdots\M').
\]}

\Proof $X'\mathbin\eta\M'$ means that $X'$ commutes with all unitary
$U\in\M''=\M$, that is with all elements of $\M^{(U)}$ (cf.
\cite[p.~392]{ref18}); as it is linear and closed however, it will even
commute with those of $\mathrm R(\M^{(U)})=\M$ (cf.
\cite[pp.~392, 405]{ref18}). So we have
\begin{align*}
 \mathfrak M_{g_0}^{\M}
 &=[Ag_0;A\in\M]=[AX'f_0;A\in\M]\\
 &=[X'Af_0;A\in\M]=[X'f;f\in(Af_0,A\in\M)]\\
 &\subset[X'f;f\in\operatorname{Domain}X'\cdot
              \mathfrak M_{f_0}^{\M}].
\end{align*}
Put
$[\operatorname{Domain}X'\cdot\mathfrak M_{f_0}^{\M}]=\mathfrak N'$;
clearly $\mathfrak N'\subset\mathfrak M_{f_0}^{\M}$ and
$X'\mathbin\eta\M'$, $\mathfrak M_{f_0}^{\M}\mathbin\eta\M'$ imply
$\mathfrak N'\mathbin\eta\M'$. It is clear that
$X'\cdot P_{\mathfrak N'}$ is a linear, closed operator with an everywhere
dense domain, and $\mathbin\eta\M$. Our above relation proves
\[
 \mathfrak M_{g_0}^{\M}\subset
 [\operatorname{Range}(X'P_{\mathfrak N'})].
\]
On the other hand
\begin{align*}
 [\operatorname{Range}(X'P_{\mathfrak N'})^*]
 &=\Hspace-(f;X'P_{\mathfrak N'}f=0)\\
 &\subset\Hspace-(f;P_{\mathfrak N'}f=0)\\
 &=\Hspace-(\Hspace-\mathfrak N')=\mathfrak N'
 \subset\mathfrak M_{f_0}^{\M}.
\end{align*}
Using \lemref{6.2.1} we obtain
\[
 \mathfrak M_{g_0}^{\M}\subset
 [\operatorname{Range}(X'P_{\mathfrak N'})]
 \sim[\operatorname{Range}(X'P_{\mathfrak N'})^*]
 \subset\mathfrak M_{f_0}^{\M},
\]
(the $\sim$ is $(\cdots\M')$), and so
$\mathfrak M_{g_0}^{\M}\precsim\mathfrak M_{f_0}^{\M}
(\cdots\M')$.

\Lemma{9.3.2}
{\itshape
$\mathfrak M_{g_0}^{\M'}\precsim\mathfrak M_{f_0}^{\M'}
(\cdots\M)$ implies
$\mathfrak M_{g_0}^{\M}\precsim\mathfrak M_{f_0}^{\M}
(\M')$.}

\Proof $\mathfrak M_{g_0}^{\M'}\precsim\mathfrak M_{f_0}^{\M'}
(\cdots\M)$ means
$\mathfrak M_{g_0}^{\M'}\sim\mathfrak N\ (\cdots\M)$,
$\mathfrak N\subset\mathfrak M_{f_0}^{\M'}$. So a partially isometric
$U\in\M$ maps $\mathfrak M_{g_0}^{\M'}$ on $\mathfrak N$ while $U^*$ maps
$\mathfrak N$ on $\mathfrak M_{g_0}^{\M'}$.

Now $Ug_0\in\mathfrak N\subset\mathfrak M_{f_0}^{\M'}$, so by
\lemref{9.2.1} $Ug_0=X'Y'f_0$, $X',Y'$ linear, closed, and with everywhere
dense domains. By \lemref{9.3.1}
\[
 \mathfrak M_{f_0}^{\M}\succsim
 \mathfrak M_{Y'f_0}^{\M}\succsim
 \mathfrak M_{X'Y'f_0}^{\M}=\mathfrak M_{Ug_0}^{\M}
 \quad(\cdots\M').
\]
But $U^*Ug_0=g_0$ and so
\[
 \mathfrak M_{g_0}^{\M}=[Ag_0;A\in\M]
 =[AU^*Ug_0;A\in\M]
 \subset[BUg_0;B\in\M]=\mathfrak M_{Ug_0}^{\M}.
\]
So
$\mathfrak M_{g_0}^{\M}\subset\mathfrak M_{Ug_0}^{\M}\precsim
\mathfrak M_{f_0}^{\M}\ (\cdots\M')$, and hence
$\mathfrak M_{g_0}^{\M}\precsim\mathfrak M_{f_0}^{\M}
(\cdots\M')$.

\Lemma{9.3.3}
{\itshape
$\mathfrak M_{f_0}^{\M'}\gtreqless\mathfrak M_{g_0}^{\M'}
(\cdots\M)$ is equivalent to
$\mathfrak M_{f_0}^{\M}\gtreqless\mathfrak M_{g_0}^{\M}
 (\cdots\M')$, resp.}

\Proof
$\mathfrak M_{f_0}^{\M'}\succsim\mathfrak M_{g_0}^{\M'}
(\cdots\M)$ implies
$\mathfrak M_{f_0}^{\M}\succsim\mathfrak M_{g_0}^{\M}$ by
\lemref{9.3.2}. The converse follows by replacing $\M$ by $\M'$ (and so
$\M'$ by $\M''=\M$); thus
$\mathfrak M_{f_0}^{\M'}\succsim\mathfrak M_{g_0}^{\M'}
(\cdots\M)$ and
$\mathfrak M_{f_0}^{\M}\succsim\mathfrak M_{g_0}^{\M}
(\cdots\M')$ are equivalent. The same follows for $\precsim$ by
interchanging $f_0$ and $g_0$.

The same follows for $\sim$ because it means that $\precsim$ and
$\succsim$ both hold; it follows for $\prec$ because this means that
$\precsim$ holds, but not $\sim$; finally it follows for $\succ$ by
interchanging again $f_0$ and $g_0$. Thus all statements of our Lemma are
proved.

\ChapterHeading{10}{X}{\texorpdfstring{Ratio of
$D_{\M}(\mathfrak M_f^{\M'})$ and $D_{\M'}(\mathfrak M_f^{\M})$}
{Ratio of the two relative dimensions}}

\SectionStart{10.1}{\texorpdfstring{The range of
$D_{\M}(\mathfrak M_f^{\M'})$ and $D_{\M'}(\mathfrak M_f^{\M})$}
{The range of the two relative dimensions}}
The last \S\S\ give the basis for a quantitative comparison of
$\mathfrak M_f^{\M'}$ with $\mathfrak M_f^{\M}$. $\M,\M'$ will again be
factors; $D_{\M}(\mathfrak M)$, $D_{\M'}(\mathfrak M')$ arbitrarily
 normalised relative dimension functions of $\M,\M'$ resp.;
 $\Delta,\Delta'$ their resp. ranges. We define:

\Definition{10.1.1}
Denote the range of $D_{\M}(\mathfrak M_f^{\M'})$ (for all
$f\in\Hspace$) by $\Delta_0$ and the range of
$D_{\M'}(\mathfrak M_f^{\M})$ (for all $f\in\Hspace$) by $\Delta_0'$.

\Lemma{10.1.1}
{\itshape
If we let $D_{\M'}(\mathfrak M_f^{\M})$ correspond to
$D_{\M}(\mathfrak M_f^{\M'})$ a one-to-one mapping of $\Delta_0$ on
$\Delta_0'$ results. Describe this mapping by
$\alpha'=\phi(\alpha)$, $\alpha\in\Delta_0$, $\alpha'\in\Delta_0'$.
$\phi(\alpha)$ is a monotone increasing function of $\alpha$.}

\Proof \lemref{9.3.3} can be formulated like this:
\[
 D_{\M}(\mathfrak M_f^{\M'})\gtreqless
 D_{\M}(\mathfrak M_g^{\M'})
\quad\hbox{is equivalent to}\quad
 D_{\M'}(\mathfrak M_f^{\M})\gtreqless
 D_{\M'}(\mathfrak M_g^{\M}),
\]
 resp. This implies all parts of our Lemma.

We will now determine $\Delta_0,\Delta_0',\phi(\alpha)$ and the set of all
pairs $\mathfrak M_f^{\M'},\mathfrak M_f^{\M}$.

\Lemma{10.1.2}
{\itshape
If $A\in\M$, then
$[Ag;g\in\mathfrak M_f^{\M'}]=\mathfrak M_{Af}^{\M'}$.}

\Proof We have
\begin{align*}
 [Ag;g\in\mathfrak M_f^{\M'}]
 &=[Ag;g\in[B'f;B'\in\M']]\\
 &=[Ag;g\in(B'f;B'\in\M')]\\
 &=[AB'f;B'\in\M']=[B'Af;B'\in\M']
 =\mathfrak M_{Af}^{\M'}.
\end{align*}

\Lemma{10.1.3}
{\itshape
The necessary and sufficient conditions for the existence of an
$f\in\Hspace$ with
$\mathfrak M_f^{\M'}=\mathfrak M$, $\mathfrak M_f^{\M}=\mathfrak N$ are
these: $\mathfrak M\mathbin\eta\M$, $\mathfrak N\mathbin\eta\M'$,
$D_{\M}(\mathfrak M)\in\Delta_0$,
$D_{\M'}(\mathfrak N)=\phi(D_{\M}(\mathfrak M))$.}

\Proof The necessity being obvious, let us consider the sufficiency. Our
assumptions imply the existence of a $g$ with
$D_{\M}(\mathfrak M_g^{\M'})=D_{\M}(\mathfrak M)$,
$D_{\M'}(\mathfrak M_g^{\M})=D_{\M'}(\mathfrak N)$. So
$\mathfrak M_g^{\M'}\sim\mathfrak M\ (\cdots\M)$,
$\mathfrak M_g^{\M}\sim\mathfrak N\ (\cdots\M')$. Let $U\in\M$ be
partially isometric with the initial and final sets
$\mathfrak M_g^{\M'},\mathfrak M$ and $U'\in\M'$ similarly with
$\mathfrak M_g^{\M},\mathfrak N$.

Consider now $\mathfrak M_{U'g}^{\M'}$. It is equal to
$[A'U'g;A'\in\M']$, thus certainly a subset of
$[B'g;B'\in\M']=\mathfrak M_g^{\M'}$. But on the other hand for every
$B'\in\M'$ we have for $A'=B'U'^*\in\M'$,
$A'U'g=B'U'^*U'g=B'P_{\mathfrak M_g^{\M}}g=B'g$. So our set is equal to
$\mathfrak M_g^{\M'}$. Now \lemref{10.1.2} gives
\[
 \mathfrak M_{UU'g}^{\M'}=[Uh;h\in\mathfrak M_{U'g}^{\M'}]
 =[Uh;h\in\mathfrak M_g^{\M'}]=\mathfrak M.
\]
Similarly $\mathfrak M_{U'Ug}^{\M}=\mathfrak N$. Thus
$f=UU'g=U'Ug$ meets all our requirements.

\Lemma{10.1.4}
{\itshape
$\alpha_1,\alpha_2,\ldots\in\Delta_0$,
$\sum_{i=1}^{\infty}\alpha_i\in\Delta$,
$\sum_{i=1}^{\infty}\phi(\alpha_i)\in\Delta'$ imply
\[
 \sum_{i=1}^{\infty}\alpha_i\in\Delta_0,
 \qquad
 \sum_{i=1}^{\infty}\phi(\alpha_i)
 =\phi\left(\sum_{i=1}^{\infty}\alpha_i\right).
\]}

\Proof We wish to construct a sequence of mutually orthogonal
\[
 \mathfrak M_1,\mathfrak M_2,\ldots\mathbin\eta\M
\]
so that $D_{\M}(\mathfrak M_i)=\alpha_i$. We proceed by induction. Assume
that $\mathfrak M_1,\mathfrak M_2,\ldots,\mathfrak M_{i-1}$ have already
been constructed, fulfilling these requirements as well as
\[
 D_{\M}\bigl(\Hspace-[\mathfrak M_1,\ldots,\mathfrak M_{i-1}]\bigr)
 \geq\sum_{j=i}^{\infty}\alpha_j,
\]
where $i=1,2,\ldots$. (For $i=1$, this is the case, as
$D_{\M}(\Hspace)\geq\sum_{j=1}^{\infty}\alpha_j$, because
$\sum_{j=1}^{\infty}\alpha_j\in\Delta$.) We then construct
$\mathfrak M_i$ as follows: If $\alpha_i=\infty$ then
$\sum_{j=i}^{\infty}\alpha_j=\infty$ and
$\Hspace-[\mathfrak M_1,\ldots,\mathfrak M_{i-1}]$ is infinite. Apply
\lemref{7.2.3} to it, and put $\mathfrak M_i=\mathfrak N$ for the resulting
$\mathfrak N$ which obviously meets all requirements. If $\alpha_i$ is
finite, choose an $\mathfrak M_i'\mathbin\eta\M$ with
$D_{\M}(\mathfrak M_i')=\alpha_i$. Then
\[
 \mathfrak M_i'\precsim
 \Hspace-[\mathfrak M_1,\ldots,\mathfrak M_{i-1}],
 \qquad
 \mathfrak M_i'\sim\mathfrak M_i\subset
 \Hspace-[\mathfrak M_1,\ldots,\mathfrak M_{i-1}],
\]
and use this $\mathfrak M_i$. Then
\begin{align*}
 D_{\M}(\mathfrak M_i)&=\alpha_i,\\
 D_{\M}\bigl(\Hspace-[\mathfrak M_1,\ldots,\mathfrak M_i]\bigr)
 &=D_{\M}\bigl(\Hspace-[\mathfrak M_1,\ldots,\mathfrak M_{i-1}]\bigr)
   -D_{\M}(\mathfrak M_i)\\
 &\geq\sum_{j=i}^{\infty}\alpha_j-\alpha_i
   =\sum_{j=i+1}^{\infty}\alpha_j,
\end{align*}
and again all conditions are fulfilled.

Choose similarly a sequence
$\mathfrak N_1,\mathfrak N_2,\ldots\mathbin\eta\M'$ so that
$D_{\M'}(\mathfrak N_i)=\phi(\alpha_i)$. (Now
$\sum_{i=1}^{\infty}\phi(\alpha_i)\in\Delta'$ must be used.) By
\lemref{10.1.3} an $f_i\in\Hspace$ with
\[
 \mathfrak M_{f_i}^{\M'}=\mathfrak M_i,
 \qquad \mathfrak M_{f_i}^{\M}=\mathfrak N_i
\]
exists. As we can replace $f_i$ by any $\alpha_i'f_i$, we may assume
$\lVert f_i\rVert\leq1/i$.

As we see, the $\mathfrak M_{f_i}^{\M'}$, $i=1,2,\ldots$, are mutually
orthogonal and so are the $\mathfrak M_{f_i}^{\M}$,
$i=1,2,\ldots$. So the $f_i$ are mutually orthogonal, and as
$\sum_{i=1}^{\infty}\lVert f_i\rVert^2\leq
\sum_{i=1}^{\infty}1/i^2$ is finite, we can form
$f=\sum_{i=1}^{\infty}f_i$ (cf. \cite[p.~67]{ref16}).

As $f_i\in\mathfrak M_{f_i}^{\M'}$, therefore
$f\in[\mathfrak M_{f_i}^{\M'};i=1,2,\ldots]$ and considering the mutual
orthogonality of the $\mathfrak M_{f_i}^{\M'}$,
$i=1,2,\ldots$, we have $f_i=E_{f_i}^{\M'}f$. Similarly
$f\in[\mathfrak M_{f_i}^{\M};i=1,2,\ldots]$ and
$f_i=E_{f_i}^{\M}f$.

Now
\begin{align*}
 \mathfrak M_f^{\M'}
 &=[A'f;A'\in\M']
 =\left[\sum_{i=1}^{\infty}A'f_i;A'\in\M'\right]\\
 &\subset[[A'f_i;A'\in\M'];i=1,2,\ldots]
 =[\mathfrak M_{f_i}^{\M'};i=1,2,\ldots]
\end{align*}
and
\[
 \mathfrak M_f^{\M'}=[A'f;A'\in\M']
 \supset[B'E_{f_i}^{\M'}f;B'\in\M']
 =[B'f_i;B'\in\M']=\mathfrak M_{f_i}^{\M'}.
\]
Thus
\[
 \mathfrak M_f^{\M'}=[\mathfrak M_{f_i}^{\M'};i=1,2,\ldots],
\]
similarly
$\mathfrak M_f^{\M}=[\mathfrak M_{f_i}^{\M};i=1,2,\ldots]$. So we have
\begin{align*}
 D_{\M}(\mathfrak M_f^{\M'})
 &=\sum_{i=1}^{\infty}D_{\M}(\mathfrak M_{f_i}^{\M'})
 =\sum_{i=1}^{\infty}\alpha_i,\\
 D_{\M'}(\mathfrak M_f^{\M})
 &=\sum_{i=1}^{\infty}D_{\M'}(\mathfrak M_{f_i}^{\M})
 =\sum_{i=1}^{\infty}\phi(\alpha_i).
\end{align*}
Therefore
$\sum_{i=1}^{\infty}\alpha_i\in\Delta_0$ and
$\sum_{i=1}^{\infty}\phi(\alpha_i)
=\phi(\sum_{i=1}^{\infty}\alpha_i)$.

\Lemma{10.1.5}
{\itshape
$0\in\Delta_0$; an $\alpha\in\Delta_0$ with $\alpha>0$ exists; if
$\alpha\in\Delta$, $\alpha\leq\beta\in\Delta_0$ then
$\alpha\in\Delta_0$.}

\Proof For $f=0$,
$D_{\M}(\mathfrak M_f^{\M'})=D_{\M}((0))=0$; for $f\ne0$,
$\mathfrak M_f^{\M'}\ne(0)$,
$D_{\M}(\mathfrak M_f^{\M'})>0$.
$\alpha\leq\beta\in\Delta_0$, $\alpha\in\Delta$ implies that an $f$ with
$D_{\M}(\mathfrak M_f^{\M'})=\beta$ and an
$\mathfrak N\mathbin\eta\M$ with $D_{\M}(\mathfrak N)=\alpha$ exist. So
$\mathfrak N\precsim\mathfrak M_f^{\M'}$,
$\mathfrak N\sim\mathfrak N'\subset\mathfrak M_f^{\M'}$, and by
\lemref{10.1.2}
\[
 \mathfrak M_{P_{\mathfrak N'}f}^{\M'}
 =[P_{\mathfrak N'}g;g\in\mathfrak M_f^{\M'}]
 =\mathfrak N'\sim\mathfrak N,
\]
so
$D_{\M}(\mathfrak M_{P_{\mathfrak N'}f}^{\M'})
=D_{\M}(\mathfrak N)=\alpha$.

\SectionStartTOC{10.2}{}
{Proof of the ratio formula and of the range alternatives. Theorem X}%
With the help of Lemmas \hyperref[lem:10.1.4]{10.1.4},
\hyperref[lem:10.1.5]{10.1.5} the discussion
of $\Delta_0,\Delta_0'$ and $\phi(\alpha)$ can now be completed. We will use
\thmref{VIII}, and discuss the three cases (I), (II), (III) for $\M$
separately.

Assume first case (I) for $\M$. Then we have case (I) for $\M'$ too, by
Lemmas \hyperref[lem:3.2.3]{3.2.3},
\hyperref[lem:3.2.4]{3.2.4}, and \hyperref[lem:8.6.1]{8.6.1}. Use
$D_{\M}(\mathfrak M)$, $D_{\M'}(\mathfrak M)$ in their standard
normalisations. By \lemref{10.1.5} $0\in\Delta_0$; and an
$\alpha\in\Delta_0$ with $\alpha>0$, so $\alpha\geq1$ exists, which
implies, as $1\in\Delta$, that $1\in\Delta_0$. Similarly
$0,1\in\Delta_0'$. As $\phi(\alpha)$ is monotone increasing (by
\lemref{10.1.1}) and as $0$ is the smallest element of both $\Delta_0$ and
$\Delta_0'$, so $\phi(0)=0$, and as $1$ is the smallest element $\ne0$ of
both $\Delta_0$ and $\Delta_0'$, so $\phi(1)=1$.

Denote the case of $\M$ more precisely by $\mathrm I_{m'}$, where
$m'=1,2,\ldots$ or $\infty$; similarly $\M'$'s by $\mathrm I_{n'}$,
where $n'=1,2,\ldots$ or $\infty$. If
$p=0,1,2,\ldots,\min(m',n')$, then apply \lemref{10.1.4} with
$\alpha_1=\cdots=\alpha_p=1$,
$\alpha_{p+1}=\alpha_{p+2}=\cdots=0$ if $p$ is finite, and
$\alpha_1=\alpha_2=\cdots=1$ if $p$ is infinite. Then
$p\in\Delta_0$, $\phi(p)=p$ obtains.

If $\Delta_0$ had further elements, we would have
\[
 q\in\Delta_0,
 \qquad q\ne0,1,2,\ldots,\min(m',n').
\]
But $q\in\Delta$, $q=0,1,2,\ldots,m'$, therefore
$\min(m',n')<q\leq m'$. Thus
$\phi(\min(m',n'))<\phi(q)$ (by \lemref{10.1.1}), but
$\phi(\min(m',n'))=\min(m',n')$ and
$\phi(q)\in\Delta_0'\subset\Delta'$, so
$\min(m',n')<\phi(q)\leq n'$. So $\min(m',n')<m'$ and $<n'$, which is
impossible. Thus $\Delta_0$ is precisely the set
$0,1,2,\ldots,\min(m',n')$ and in it $\phi(\alpha)=\alpha$; therefore
$\Delta_0'$ is the same set. We see that we have either
$\Delta_0=\Delta$, or $\Delta_0'=\Delta'$.

Assume next case (III) for $\M$. $\Delta$ has precisely two elements:
$0,\infty$, so that $\Delta_0=\Delta$ by \lemref{10.1.5}. So
$\Delta_0=(0,\infty)$,
$\Delta_0'=(\phi(0),\phi(\infty))=(0,\phi(\infty))$.
\lemref{10.1.5} applied to $\M'$ excludes the existence of an
$\alpha'\in\Delta'$ with $0<\alpha'<\phi(\infty)$. This excludes case
(II) for $\M'$. Case (I) for $\M'$ is ruled out as above (interchange
$\M,\M'$), so we have case (III) for $\M'$ too. Now we have for the same
reasons as in $\M$, $\Delta_0'=\Delta'=(0,\infty)$,
$\phi(\infty)=\infty$. So we have again $\phi(\alpha)=\alpha$, but now
$\Delta_0=\Delta$ and $\Delta_0'=\Delta'$.

Assume finally case (II) for $\M$. The above arguments exclude cases (I)
and (III) for $\M'$ (interchange $\M$, $\M'$,!), so we have case (II) for
$\M'$ too. We have four possibilities: $\M$ in a case
$(\mathrm{II}_{m'})$ and $\M'$ in a case $(\mathrm{II}_{n'})$ with
$m',n'=1,\infty$.

Consider a pair $\alpha\in\Delta_0$, $\alpha'\in\Delta_0'$ with
$\alpha'=\phi(\alpha)$ and a $p=1,2,\ldots$. Then
$(\alpha/p)\in\Delta_0$, $(\alpha'/p)\in\Delta_0'$. Assume
$(\alpha'/p)>\phi(\alpha/p)$. Obviously
$p(\alpha/p)=\alpha\in\Delta_0$. Further
$p\phi(\alpha/p)<\alpha'\in\Delta_0'\subset\Delta'$; as we have case (II)
for $\M'$, this implies $p\phi(\alpha/p)\in\Delta'$. Thus
\lemref{10.1.4} applies for
$\alpha_1=\cdots=\alpha_p=(\alpha/p)$,
$\alpha_{p+1}=\alpha_{p+2}=\cdots=0$, giving
$p\phi(\alpha/p)=\phi(\alpha)=\alpha'$, $(\alpha'/p)=\phi(\alpha/p)$.
This contradicts our assumption $(\alpha'/p)>\phi(\alpha/p)$.
Interchanging of $\M,\M'$ excludes
$(\alpha'/p)<\phi(\alpha/p)$ too. So we have proved
$\alpha'/p=\phi(\alpha/p)$. In other words:
\[
 \phi(\alpha/p)=(1/p)\phi(\alpha).
\]

Take now a $q=0,1,\ldots,p$. We have
$(q/p)\alpha\leq\alpha\in\Delta_0\subset\Delta$, so (owing to case (II)
for $\M$ and for $\M'$) $(q/p)\alpha\in\Delta$,
$q\phi(\alpha/p)\in\Delta'$. Therefore \lemref{10.1.4} applies for
$\alpha_1=\cdots=\alpha_q=\alpha/p$,
$\alpha_{q+1}=\alpha_{q+2}=\cdots=0$, giving
$(q/p)\alpha\in\Delta_0$ and
$\phi((q/p)\alpha)=q\phi(\alpha/p)=(q/p)\phi(\alpha)$. In other words: If
$0\leq\beta\leq\alpha$ then
$\phi(\beta)=(\beta/\alpha)\phi(\alpha)$ if $\beta/\alpha$ is rational.

By \lemref{10.1.5} $0\leq\beta\leq\alpha$ implies
$\beta\in\Delta_0$ so $\phi(\beta)$ is defined in all this interval. As it
is monotone increasing (\lemref{10.1.1}), we have
$\phi(\beta)=(\beta/\alpha)\phi(\alpha)$ for all these $\beta$, that is
$\phi(\alpha)/\alpha=\phi(\beta)/\beta$ (assuming
$\alpha,\beta\ne0$).

Assume $\alpha,\beta\in\Delta_0$ and $\ne0$. If $\alpha\geq\beta$ then
$0\leq\beta\leq\alpha$ and we know that
$\phi(\alpha)/\alpha=\phi(\beta)/\beta$. By symmetry this equation holds
for $\alpha\leq\beta$ too, that is, it holds always. So
$\phi(\alpha)/\alpha$ is constant, say $C$. $C$ is finite and positive by
its nature. We have $\phi(\alpha)=C\alpha$ if
$\alpha\in\Delta_0$, $\alpha\ne0$ and of course for $\alpha=0$ too.

Let the l.u.b. of $\Delta_0$ be $m''$, and that of $\Delta_0'$, $n''$. We
have $0<m''\leq m'$, $0<n''\leq n'$, and of course $n''=Cm''$.
$m''$ belongs to $\Delta_0$: Choose $\alpha_1,\alpha_2,\ldots$ with
$\alpha_i\geq0$, $\sum_1^i\alpha_j<m''$,
$\sum_{1}^{\infty}\alpha_j=m''$ (for instance
$\alpha_i=m''/2^i$ if $m''$ is finite, and $\alpha_i=1$ if $m''$ is
infinite), and apply \lemref{10.1.4}. Similarly $n''$ belongs to
$\Delta_0'$.

Assume now $m''<m'$, $n''<n'$. Then $m'',n''$ are both finite, and a
$\vartheta>1$, $<2$ with $\vartheta m''<m'$,
$\vartheta n''<n'$ exists. Then $(\vartheta/2)m''<m'$ and so
$(\vartheta/2)m''\in\Delta_0$ while $2(\vartheta/2)m''=\vartheta m''\leq m'$ is in
$\Delta$, $2\phi((\vartheta/2)m'')=2C(\vartheta/2)m''
=\vartheta n''\leq n'$ is in $\Delta'$. So \lemref{10.1.4} applies for
$\alpha_1=\alpha_2=(\vartheta/2)m''$,
$\alpha_3=\alpha_4=\cdots=0$ giving $\vartheta m''\in\Delta_0$. This is
impossible, as $\vartheta m''>m''$.

So we have $m''=m'$ and $n''\leq n'$ or $m''\leq m'$ and $n''=n'$, and
besides $n''=Cm''$. Let us now consider the four possible combinations of
$m',n'=1,\infty$.

If $m'=n'=1$ then $C\geq1$ necessitates $n''=n'=1$,
$m''=(1/C)n'=1/C\leq1$ so $\Delta_0'=\Delta'$, while $C\leq1$ implies
$m''=m'=1$, $n''=Cm''=C\leq1$ so $\Delta_0=\Delta$. As both
$\Delta,\Delta'$ (that is $D_{\M}(\mathfrak M)$,
$D_{\M'}(\mathfrak M)$) are normalised, $C$ has an invariant meaning.
If $m'=1$, $n'=\infty$ then $\Delta'$ is not normalised. So we can make
$C=1$. Now obviously $m''=m'=1$, $n''=m''=1<\infty$. So
$\Delta_0=\Delta$. If $m'=\infty$, $n'=1$ then $\Delta$ is not
normalised. So we can make $C=1$ again. Now obviously
$n''=n'=1$, $m''=n''=1<\infty$. So $\Delta_0'=\Delta'$. If finally
$m'=n'=\infty$ then both $\Delta$ and $\Delta'$ are not normalised. So we
can make $C=1$ again, and then $m''=n''=\infty$ and thus
$\Delta_0=\Delta$, $\Delta_0'=\Delta'$.

All these results together with those of Lemmas
\hyperref[lem:10.1.1]{10.1.1} and \hyperref[lem:10.1.3]{10.1.3} give the

\MainTheorem{X}
{\itshape
In a factorisation $\M,\M'$ the factors $\M,\M'$ must both belong to class
(I), or both to class (II), or both to class (III).

The function $\phi(\alpha)$ defined in \lemref{10.1.1} always has the form
$\phi(\alpha)=C\alpha$ where $C$ is a finite, positive constant. In case
(I), using the standard normalisations of $\Delta,\Delta'$ necessarily
$C=1$; in case (II), if both $\M,\M'$ are in case $(\mathrm{II}_1)$ and
using the standard normalisations of $\Delta,\Delta'$ the value of $C$ is
uniquely defined; in case (II), if at least one of $\M,\M'$ is in case
$(\mathrm{II}_\infty)$ and if we normalise the corresponding $\Delta$ or
$\Delta'$ suitably, then we can make $C=1$; in case (III) the value of $C$
is arbitrary (as $\alpha=0,\infty$) so we can take $C=1$.

Always either $\Delta_0=\Delta$ or $\Delta_0'=\Delta'$; details are given
in the discussion above.

An $f$ with $\mathfrak M_f^{\M'}=\mathfrak M$,
$\mathfrak M_f^{\M}=\mathfrak N$ exists if and only if
$\mathfrak M\mathbin\eta\M$, $\mathfrak N\mathbin\eta\M'$ and if
$D_{\M}(\mathfrak M)\in\Delta_0$,
$D_{\M'}(\mathfrak N)\in\Delta_0'$,
$D_{\M'}(\mathfrak N)=CD_{\M}(\mathfrak M)$. (Owing to the last condition
each one of the two preceding conditions implies the other.)}

This Theorem brings us nearer to the answers of \hyperref[prob:4]{Problems 4} and \hyperref[prob:7]{7}: It
excludes certain combinations of cases for $\M,\M'$ and it shows that
$\M,\M'$ possess a further invariant, if $\M,\M'$ are both of class
$(\mathrm{II}_1)$, the number $C$. This latter statement is of course only
then of value, if such $\M,\M'$ both of class $(\mathrm{II}_1)$, and with
various values of $C$, do really exist. We will see that this is the case,
and besides get further information on all these points in
\hyperref[sec:13.2]{\S\S13.2} and \hyperref[sec:13.3]{13.3}.

\ChapterHeading{11}{XI}{Composition and Decomposition of Factors}

\SectionStart{11.1}{Normalcy and composition. Problems 8--9}
A natural question to ask is this:

\Problem{8}
{\itshape
Under what conditions can two given $\M,\N$ belong both to one
factorisation $\M_1,\ldots,\M_n$?}

The following Lemma gives a necessary and sufficient condition but thereby
raises a new problem.

\Lemma{11.1.1}
{\itshape
$\M,\N$ belong both to one (suitably chosen) factorisation if and only if
(i) $\M,\N$ are factors, (ii) $\M,\N$ commute (that is: $\M\subset\N'$),
(iii) $\mathrm R(\M,\N)$ is a factor.}

\Proof If $\M,\N$ belong to a factorisation
$\M_1,\ldots,\M_n$, then $\M=\M_i$, $\N=\M_j$, $i\ne j$. Thus (i) is
fulfilled by \lemref{3.1.1}; (ii) by definition; and (iii) by
\lemref{3.1.1} again, as we may make $i=1$, $j=2$ by a permutation of
$\M_1,\ldots,\M_n$ and then remember (cf. the remark at the end of
\secref{3.1}) that $\mathrm R(\M_1,\M_2)$,
$\mathrm R(\M_3,\ldots,\M_n)$ too is a factorisation if $n>3$ (for $n=2$,
$\mathrm R(\M_1,\M_2)=\B$ is trivially a factor). Thus our condition is
necessary.

Assume now that (i)--(iii) are fulfilled. $(\mathrm R(\M,\N))'=\M'\N'$
commutes with $\M$ and $\N$; $\M,\N$ commute with each other; and
\[
 \mathrm R(\M,\N,(\mathrm R(\M,\N))')
 =\mathrm R(\mathrm R(\M,\N),(\mathrm R(\M,\N))')=\B
\]
as $\mathrm R(\M,\N)$ is a factor. So
$\M,\N,(\mathrm R(\M,\N))'$ is a factorisation and it contains $\M,\N$.
Thus our condition is sufficient.

The new problem therefore is:

\Problem{9}
{\itshape
If $\M,\N$ are two factors which commute, under what conditions will then
$\mathrm R(\M,\N)$ be a factor? How does the class of $\mathrm R(\M,\N)$
(in the sense of \thmref{VIII}) connect with the classes of $\M$ and $\N$?}

We will only succeed in giving a partial answer. This answer could be
obtained by using somewhat less formal machinery than we are going to use.
But it seems reasonable to make use of it, because it makes the algebraic
side of these questions clearer.

We define:

\Definition{11.1.1}
Let $\M,\P$ be rings which contain 1. If $\M\subset\P$ we define
$\M_{\P}'=\P\cdot\M'$ ($\M_{\P}'$ too is a ring which contains 1 and
$\subset\P$).

This notion shares many formal properties with $\M'$ with which it
coincides for $\P=\B$. As the decisive property of $\M'$ is $\M''=\M$
(cf. \cite[p.~397]{ref18}), it is reasonable to define:

\Definition{11.1.2}
$\P$ is normal, if $\M\subset\P$ implies $\M_{\P\P}''=\M$.

\Lemma{11.1.2}
{\itshape
If $\P$ is normal, then it is a factor; and every factorisation
$\Q,\P'$ is a coupled one.}

\Proof Put $\M=(\alpha1)$. Then $\M\subset\P$ so
$\M_{\P\P}''=\M$; now
$(\alpha1)_{\P}'=\P(\alpha1)'=\P\B=\P$,
$\P_{\P}'=\P\cdot\P'$. So we have
$\P\cdot\P'=(\alpha1)$. So $\P$ is a factor. That $\Q,\P'$ is a
factorisation, means $\Q\subset\P''=\P$ and
$\mathrm R(\Q,\P')=\B$, that is
$(\mathrm R(\Q,\P'))'=\B'$,
$\Q_{\P}'=\Q'\P=(\mathrm R(\Q,\P'))'=\B'=(\alpha1)$. So
$\Q=\Q_{\P\P}''=(\alpha1)_{\P}'=\P$.

Our notion of normalcy could easily be expressed in terms of the
conditions which characterise the ``B-lattices'' of G. Birkhoff (cf.
\cite[p.~445]{ref1}) but we will not enter on this aspect of the subject at
present.

The connection with \probref{9} is established by the following Lemma:

\Lemma{11.1.3}
{\itshape
The factors $\M,\N$ are commutative and $\mathrm R(\M,\N)$ is a factor too
(that is: the conditions of \lemref{11.1.1} are fulfilled), if a $\P$ with
$\M\subset\P$, $\N\subset\P'$ exists, so that $\P,\P'$ are both normal.

This is in particular the case, if $\M,\N$ commute, and besides
$\M,\M'$ or $\N,\N'$ are both normal.}

\Proof $\M\subset\P\subset\N'$ shows that $\M,\N$ commute. We have
\[
 (\mathrm R(\M,\M_{\P}'))_{\P}'
 =\P\cdot(\mathrm R(\M,\M_{\P}'))'
 =\P\M'(\M_{\P}')'
 =\M'\M_{\P\P}''
\]
and as $\M\subset\P$, $\P$ normal, this is
$\M'\M=(\alpha1)$. Therefore, considering
$\mathrm R(\M,\M_{\P}')\subset\P$, $\P$ normal, we have
\[
 \mathrm R(\M,\M_{\P}')
 =(\mathrm R(\M,\M_{\P}'))_{\P\P}''
 =(\alpha1)_{\P}'=\P.
\]
Similarly $\mathrm R(\N,\N_{\P'}')=\P'$. Therefore
\begin{align*}
 \mathrm R(\mathrm R(\M,\N),\mathrm R(\M_{\P}',\N_{\P'}'))
 &=\mathrm R(\M,\N,\M_{\P}',\N_{\P'}')\\
 &=\mathrm R(\mathrm R(\M,\M_{\P}'),
               \mathrm R(\N,\N_{\P'}'))\\
 &=\mathrm R(\P,\P')=\B. \qquad(\P\text{ is a factor!})
\end{align*}
Now $\M$ commutes with $\M_{\P}'$ because
$\M_{\P}'\subset\M'$, and with $\N_{\P'}'$ because
$\M\subset\P$, $\N_{\P'}'\subset\P'$. So $\M$ commutes with
$\mathrm R(\M_{\P}',\N_{\P'}')$. Similarly $\N$ commutes with
$\mathrm R(\M_{\P}',\N_{\P'}')$. Therefore
$\mathrm R(\M_{\P}',\N_{\P'}')$ commutes with $\mathrm R(\M,\N)$.
Thus $\mathrm R(\M,\N)$, $\mathrm R(\M_{\P}',\N_{\P'}')$ is a
factorisation, and so (by \lemref{3.1.1}) $\mathrm R(\M,\N)$ is a factor.

 The second half of our Lemma obtains by putting $\P=\M$ resp.
$\P=\N'$.

\SectionStart{11.2}{Properties of normalcy. Problem 10}
We derive some properties of normalcy.

\Lemma{11.2.1}
{\itshape
The operation $\M_{\P}'$ (for $\M\subset\P$) and the normalcy of $\P$ are
invariants under $(A^*,AB)$-ring-isomorphisms of $\P$.}

\Proof The first statement is obvious, the second follows immediately
from the first one.

\Lemma{11.2.2}
{\itshape
If $\P$ is a factor of class (I), then $\P,\P'$ are both normal.}

\Proof As $\P'$ is of class (I) too (by Lemmas
\hyperref[lem:3.2.4]{3.2.4} and \hyperref[lem:8.6.1]{8.6.1}, or by
\thmref{X}), it suffices to consider $\P$ alone.
Now \lemref{8.6.1} and \thmref{II} imply that $\P$ is algebraically (even
fully) ring-isomorphic to the $\B$ of some space $\Hspace$. Therefore
\lemref{11.2.1} permits us to assume $\P=\B$.

But then $\M_{\P}'$ becomes $\M_{\B}'=\M'$ and as $\M''=\M$ (cf.
\cite[p.~397]{ref18}) therefore $\P=\B$ is normal.

Note that Lemmas \hyperref[lem:11.1.2]{11.1.2} and
\hyperref[lem:11.2.2]{11.2.2} give together the
statement of \lemref{3.2.4} again: Every direct factorisation $\M,\N$ is
 coupled (let $\M,\N$ correspond to $\Q,\P'$ resp.).

As we will find non-coupled factorisations of class (II) (Cf.
\secref{13.4}), \lemref{11.1.2} implies that non-normal factors of class
(II) do exist. We did not succeed however in showing that no factor of
class (II) is normal. Nevertheless the only non-normal $\P$'s we know are
of class (II).

Thus the following question is only partially answered:

\Problem{10}
{\itshape Which factors are normal?}

The second half of \probref{9} remains to be discussed: Determining the
class of $\mathrm R(\M,\N)$ in terms of the classes of $\M$ and $\N$. If
$\M$ or $\N$ belongs to case (I) (by symmetry we may assume that $\N$
does), then the results of the next two \S\S\ permit us to solve this
problem. In all other cases our information is incomplete.

\SectionStart{11.3}{\texorpdfstring{Decomposition by an
$\mathfrak M\mathbin\eta\M$}{Decomposition by an M eta bold M}}
We will now consider an operation which is in a certain sense inverse to
the ``composition'' $\mathrm R(\M,\N)$. It is based on the following Lemma:

\Lemma{11.3.1}
{\itshape
Let $\M$ be a ring which contains 1, $E$ a projection $\in\M$. If
$A'\in\M'$ then form $EA'=A'E=A_E'$. Then these $A_E'$ are identical with
those (bounded) operators $\bar A$ for which
$E\bar A=\bar AE=\bar A$ and which commute with all $EBE$, $B\in\M$.}

\Proof The necessity is clear: If $A'\in\M'$ then $A'$ commutes with
$E\in\M$ so we can define
\[
 EA'=A'E=A_E';\qquad
 EA_E'=EEA'=EA'=A_E',\qquad
 A_E'E=A'EE=A'E=A_E',
\]
and if $B\in\M$ then
\[
 EBEA_E'=EBA_E'=EBEA'=A'EBE=A_E'BE=A_E'EBE.
\]

Let us now consider the sufficiency. Assume therefore
$E\bar A=\bar AE=\bar A$ and that $\bar A$ commutes with all $EBE$,
$B\in\M$.

As $\bar A$ is bounded, there exists an $\alpha>0$ so that
$\alpha^2\cdot1-\bar A^*\bar A$ is (semi-) definite, and then there exists
a unique (semi-) definite $\bar C$ with
$\bar C^2=\alpha^2\cdot1-\bar A^*\bar A$. (That is:
$\bar C=(\alpha^2\cdot1-\bar A^*\bar A)^{1/2}$, cf.
\cite[p.~303]{ref21}.) The $EBE$, $B\in\M$ commute with $\bar A$,
therefore with $\bar A^*,\bar A^*\bar A$ and $\bar C$ too. In particular
$E=EEE$ does so.

Assume now $B_1,\ldots,B_i\in\M$, $f_1,\ldots,f_i$ arbitrary. We then
form:
\begin{align*}
 &\left\lVert\sum_{j=1}^{i}B_j\bar C E f_j\right\rVert^2
 +\left\lVert\sum_{j=1}^{i}B_j\bar A E f_j\right\rVert^2\\
 &=\left(\sum_{j=1}^{i}B_j\bar C E f_j,
          \sum_{k=1}^{i}B_k\bar C E f_k\right)
  +\left(\sum_{j=1}^{i}B_j\bar A E f_j,
          \sum_{k=1}^{i}B_k\bar A E f_k\right)\\
 &=\left(\sum_{j,k=1}^{i}
    \{E\bar C B_k^*B_j\bar C E+E\bar A^*B_k^*B_j\bar A E\}f_j,f_k\right)\\
 &=\left(\sum_{j,k=1}^{i}
    \{\bar CEB_k^*B_jE\bar C+\bar A^*EB_k^*B_jE\bar A\}f_j,f_k\right)\\
 &=\left(\sum_{j,k=1}^{i}
    \{EB_k^*B_jE(\bar C^2+\bar A^*\bar A)\}f_j,f_k\right)\\
 &=\left(\sum_{j,k=1}^{i}\{EB_k^*B_jE(\alpha^21)\}f_j,f_k\right)\\
 &=\alpha^2\left(\sum_{j=1}^{i}B_jEf_j,
                  \sum_{k=1}^{i}B_kEf_k\right)
 =\alpha^2\left\lVert\sum_{j=1}^{i}B_jEf_j\right\rVert^2,
\end{align*}
and therefore
\[
 \left\lVert\sum_{j=1}^{i}B_j\bar A E f_j\right\rVert
 \leq\alpha\left\lVert\sum_{j=1}^{i}B_jE f_j\right\rVert.
\]
Thus $\sum_{j=1}^{i}B_jEf_j=0$ implies
$\sum_{j=1}^{i}B_j\bar A Ef_j=0$ and so, owing to the linearity,
\[
 \overline{\overline A}\left(\sum_{j=1}^{i}B_jEf_j\right)
 =\sum_{j=1}^{i}B_j\bar A Ef_j
\]
($i=0,1,2,\ldots$; $B_j\in\M$, $f_j$ arbitrary for $j=1,\ldots,i$)
defines a one-valued linear operator $\overline{\overline A}$. We can now
write $\lVert\overline{\overline A}f\rVert\leq\alpha\lVert f\rVert$ so
that $\overline{\overline A}$ is continuous. Therefore we can form the
closure of $\overline{\overline A}$, $\widetilde{\bar A}$. This will be
one-valued, again fulfilling
$\lVert\widetilde{\bar A}f\rVert\leq\alpha\lVert f\rVert$ and
 $\operatorname{Domain}\widetilde{\bar A}'
=[\operatorname{Domain}\overline{\overline A}]$. (All this results by
direct application of the criterium of \cite[p.~300]{ref21}.)

If $B\in\M$ then
\begin{align*}
 \overline{\overline A}\left(B\sum_{j=1}^{i}B_jEf_j\right)
 &=\overline{\overline A}\left(\sum_{j=1}^{i}BB_jEf_j\right)
 =\sum_{j=1}^{i}(BB_j)\bar A Ef_j\\
 &=B\left(\sum_{j=1}^{i}B_j\bar A Ef_j\right)
 =B\left(\overline{\overline A}
           \left(\sum_{j=1}^{i}B_jEf_j\right)\right),
\end{align*}
so $\overline{\overline A}\mathbin\eta\M'$ (cf. \lemref{4.2.2}).
Therefore $\widetilde{\bar A}\mathbin\eta\M'$ and
$\mathfrak M_0=\operatorname{Domain}\widetilde{\bar A}
\mathbin\eta\M'$. Thus $F_0=P_{\mathfrak M_0}\in\M'$. Note that
$\operatorname{Domain}\widetilde{\bar A}\supset
\operatorname{Domain}\overline{\overline A}\supset\mathfrak M$, if
$E=P_{\mathfrak M}$ (put $i=1$, $B_1=1$), so $F_0\geq E$.

Our definition of $F_0$ shows, that $\widetilde{\bar A}F_0$ is everywhere
defined. It is bounded too, and as
$\widetilde{\bar A}\mathbin\eta\M'$, $F_0\in\M'$, therefore
$\widetilde{\bar A}F_0\mathbin\eta\M'$. These facts imply together
$\widetilde{\bar A}F_0\in\M'$. Put
$A'=\widetilde{\bar A}F_0$.

Now $\widetilde{\bar A}Ef=\bar AEf$ (put $i=1$, $B_1=1$ in the
definition), and so $\widetilde{\bar A}E=\bar AE=\bar A$. Therefore
$A'E=\widetilde{\bar A}F_0E=\widetilde{\bar A}E=\bar A$. In other words:
$EA'=A'E=A_E'=\bar A$, completing the proof.

We can now carry out the constructions which are necessary for our
``decomposition.''

\Definition{11.3.1}
Let $\M$ be a ring which contains 1 and $\mathfrak M$ a closed linear set
$\ne(0)$. Consider those operators $A\in\M$ which are reduced by
$\mathfrak M$ and form their parts in $\mathfrak M$, $A_{(\mathfrak M)}$
(cf. \cite[p.~78]{ref16}). These are bounded operators in $\mathfrak M$.
Denote the set of all these $A_{(\mathfrak M)}$ ($A\in\M$ and reduced by
$\mathfrak M$) by $\M_{(\mathfrak M)}$.

\Lemma{11.3.2}
{\itshape
Let $\M$ be a ring which contains 1 and $\mathfrak M$ a closed linear set
$\ne(0)$, $\mathfrak M\mathbin\eta\M$, $E=P_{\mathfrak M}$. Then
$(\M_{(\mathfrak M)})'=\M'_{(\mathfrak M)}$. (These are rings of
operators in the space $\mathfrak M$.)}

\Proof An operator $A_0$ in $\mathfrak M$ is at the same time an operator
in $\Hspace$ but its domain is $\subset\mathfrak M$. If $A_0$ is bounded,
then $A_0E$ is everywhere defined in $\Hspace$ and so (being bounded)
$\in\B$. As its range is $\subset\mathfrak M$,
$E(A_0E)=A_0E$. Obviously $(A_0E)E=A_0E$ so $A_0E$ commutes with
$E=P_{\mathfrak M}$ and is reduced by $\mathfrak M$. Its parts in
 $\mathfrak M$, $\Hspace-\mathfrak M$ are $A_0,0$ resp.

Thus $A_0,B_0$ (in $\mathfrak M$) commute if and only if $A_0E,B_0E$ (in
$\Hspace$) commute. $A_0\in(\M_{(\mathfrak M)})'$ means therefore, that
$A_0E$ commutes with every $B_0E$, $B_0\in\M_{(\mathfrak M)}$.

The statement on $B_0$ means $B_0=B_{(\mathfrak M)}$, $B\in\M$ but then
\[
 EBE=EB=BE=B_0E.
\]
(Note, that $B_0E$ differs in general from $EB_0=B_0$, their domains being
 $\Hspace$ resp. $\mathfrak M$.) So $A_0E$ must commute with every
$EBE$ where $B\in\M$, $B$ commutes with $E$. Replacing $B$ by $EBE$
(which commutes with $E$ and leaves $EBE$ unaltered) shows that any
$B\in\M$ will do. Now \lemref{11.3.1} applies to $A_0E$:
$A_0E=EA'=A'E$ for some $A'\in\M'$. This means,
$A_0=(A_0E)_{(\mathfrak M)}=A'_{(\mathfrak M)}$,
$A_0\in\M'_{(\mathfrak M)}$.

Conversely: If $A_0\in\M'_{(\mathfrak M)}$,
$B_0\in\M_{(\mathfrak M)}$ then
$A_0=A'_{(\mathfrak M)}$, $B_0=B_{(\mathfrak M)}$ where
$A'\in\M'$, $B\in\M$ so $A',B$ commute, and $A_0,B_0$ too. So
$A_0\in\M'_{(\mathfrak M)}$ implies
$A_0\in(\M_{(\mathfrak M)})'$. So we have completely proved
$(\M_{(\mathfrak M)})'=\M'_{(\mathfrak M)}$.

\Lemma{11.3.3}
{\itshape
Let $\M,\mathfrak M,E$ be as above. Consider the correspondence
$X\longleftrightarrow X_{(\mathfrak M)}$. This is a one-to-one mapping and
even algebraic ring-isomorphism of the following rings on each other:

(i) At any rate of the ring of all $A\in\M$ with $EA=AE=A$ on all
$\M_{(\mathfrak M)}$.

(ii) If $\M$ is a factor, of all $\M'$ on all
$\M'_{(\mathfrak M)}$.}

\Proof The invariance of $\alpha A,A^*,A+B,AB$ is obvious, only the one to
one character must therefore be proved.

Ad (i): If $A_0\in\M_{(\mathfrak M)}$, form the $A_0E$ from the proof of
\lemref{11.3.2}. If $A_0=A_{(\mathfrak M)}$, $A\in\M$, then as we saw,
$A_0E=EAE$, so $A_0E\in\M$ too; and besides
$A_0=(A_0E)_{(\mathfrak M)}$, $E(A_0E)=(A_0E)E=A_0E$. Furthermore $A_0$
and $A_0E$ determine each other uniquely. This completes the proof.

Ad (ii): Every $A'\in\M'$ commutes with $E=P_{\mathfrak M}\in\M$ so it
is reduced by $\mathfrak M$ and $A'_{(\mathfrak M)}$ can be formed.
$A'_{(\mathfrak M)}=B'_{(\mathfrak M)}$ means that $A'f=B'f$ for all
$f\in\mathfrak M$, that is for all $f=Eg$ so $A'E=B'E$. By the corollary
to \thmref{III} this implies $A'=B'$ as $E\ne0$. This completes the proof.

\Lemma{11.3.4}
{\itshape
Let $\M,\mathfrak M,E$ be as above. If $\M$ is a factor then
$\M_{(\mathfrak M)}$ is one too.}

\Proof Apply \lemref{11.3.2}:
\[
 \M_{(\mathfrak M)}\cdot(\M_{(\mathfrak M)})'
 =\M_{(\mathfrak M)}\cdot\M'_{(\mathfrak M)}
 =(\M\cdot\M')_{(\mathfrak M)}
 =(\alpha1)_{(\mathfrak M)}=(\alpha1_{(\mathfrak M)})
\]
and $1_{(\mathfrak M)}$ is the unity operator in $\mathfrak M$.

\Lemma{11.3.5}
{\itshape
Let $\M$ be a factor, $\mathfrak M,E$ as above. Then we have:

(i) If $F$ runs over all projections $F\in\M$, $F\leq E$ (cf.
\cite[p.~76]{ref16}) then $F_{(\mathfrak M)}$ runs over all projections
$\in\M_{(\mathfrak M)}$.
$F_{(\mathfrak M)}\sim G_{(\mathfrak M)}
(\cdots\M_{(\mathfrak M)})$, ($F,G\in\M$, $\leq E$), is equivalent to
$F\sim G\ (\cdots\M)$.

(ii) If $F'$ runs over all projections $F'\in\M'$ then
$F'_{(\mathfrak M)}$ runs over all projections
$\in\M'_{(\mathfrak M)}$.
$F'_{(\mathfrak M)}\sim G'_{(\mathfrak M)}
(\cdots\M'_{(\mathfrak M)})$, ($F',G'\in\M'$), is equivalent to
$F'\sim G'\ (\cdots\M')$.}

\Proof Ad (i): The first part is a restatement of \lemref{11.3.3}, (i) as
$EF=FE=F$ means $F\leq E$ (cf. \cite[p.~76]{ref16}). As to the second
part, note that $F\sim G\ (\cdots\M)$ means $F=U^*U$, $G=UU^*$ for some
$U\in\M$ (cf. \lemref{4.3.1} and \defref{6.1.1}) and
$F_{(\mathfrak M)}\sim G_{(\mathfrak M)}
(\cdots\M_{(\mathfrak M)})$ means similarly
\[
 F_{(\mathfrak M)}=(U_{(\mathfrak M)})^*U_{(\mathfrak M)},
 \qquad
 G_{(\mathfrak M)}=U_{(\mathfrak M)}(U_{(\mathfrak M)})^*
\]
(cf. as above) for some $U\in\M$, $EU=UE=U$ (cf.
\lemref{11.3.3}, (i)). The latter equation means again $F=U^*U$,
$G=UU^*$ (by \lemref{11.3.3}, (i), owing to the algebraic
ring-isomorphism), so that all we need to prove is that the $U$ of the
first case fulfill automatically $EU=UE=U$. As the initial and final sets
of that $U$ have the projections $F,G\leq E$ they are both
$\subset\mathfrak M$ and this implies $EU=UE=U$.

Ad (ii): This follows immediately from \lemref{8.5.1} as $\M'$ and
$\M'_{(\mathfrak M)}$ are algebraically ring-isomorphic.

\Lemma{11.3.6}
{\itshape
Let $\M$ be a factor, $\mathfrak M,E$ as above. Let
 $D_{\M}(F)$, $D_{\M'}(F')$ ($F\in\M$ resp. $F'\in\M'$) be relative
 dimension functions in $\M$ resp. $\M'$. Then
\[
 D_{\M}^{(\mathfrak M)}(F_{(\mathfrak M)})=D_{\M}(F)
 \quad\hbox{and}\quad
 D_{\M'}^{(\mathfrak M)}(F'_{(\mathfrak M)})=D_{\M'}(F')
\]
 are relative dimension functions in $\M_{(\mathfrak M)}$ resp.
$\M'_{(\mathfrak M)}$.}

\Proof Owing to \defref{9.1.1} this follows immediately from
\lemref{11.3.5}.

\Lemma{11.3.7}
{\itshape
Let $\M$ be a factor, $\mathfrak M,E$ as above. Let
\[
 D_{\M}(F),\quad D_{\M'}(F'),\quad
 D_{\M}^{(\mathfrak M)}(F),\quad D_{\M'}^{(\mathfrak M)}(F')
\]
be the relative dimension functions of \lemref{11.3.6} and
$\Delta,\Delta',\Delta_{(\mathfrak M)},\Delta'_{(\mathfrak M)}$ their
 resp. ranges. Then $\Delta_{(\mathfrak M)}$ is the set of all
$\alpha\in\Delta$, $\alpha\leq D_{\M}(E)$ while
$\Delta'_{(\mathfrak M)}=\Delta'$.}

\Proof $\Delta_{(\mathfrak M)}$ is by \lemref{11.3.5}, (i) and
\lemref{11.3.6}, the set of all $D_{\M}(F)$, $F\leq E$, that is the set of
all $D_{\M}(\mathfrak N)$, $\mathfrak N\subset\mathfrak M$. Replacing
these $\mathfrak N$ by all
$\mathfrak N\sim\mathfrak N'\subset\mathfrak M$, that is by all
$\mathfrak N\precsim\mathfrak M$, obviously does not affect the set of
their $D_{\M}(\mathfrak N)$. In other words: we have all
$D_{\M}(\mathfrak N)$ with
$D_{\M}(\mathfrak N)\leq D_{\M}(\mathfrak M)$, that is all
$\alpha\in\Delta$ with
$\alpha\leq D_{\M}(\mathfrak M)=D_{\M}(E)$.

$\Delta'_{(\mathfrak M)}=\Delta'$ follows immediately from
\lemref{11.3.5}, (ii), and \lemref{11.3.6}.

\SectionStart{11.4}{\texorpdfstring{{\protect\small Decomposition by two
$\mathfrak M\mathbin\eta\M$, $\mathfrak M'\mathbin\eta\M'$.
Determination of their classes. The invariant $C$}}{Decomposition by two
subspaces. Determination of their classes. The invariant C}}
The treatment of \secref{11.3} was asymmetric, insofar as the
$\mathfrak M$ occurring in it was assumed to be $\mathbin\eta\M$. We will
now obtain symmetric results by considering simultaneously an
$\mathfrak M\mathbin\eta\M$ and an
$\mathfrak M'\mathbin\eta\M'$ and applying the preceding Lemmas twice.

\Lemma{11.4.1}
{\itshape
Let $\M$ be a factor; $\mathfrak M,\mathfrak M'$ closed linear sets both
$\ne(0)$; $\mathfrak M\mathbin\eta\M$,
$\mathfrak M'\mathbin\eta\M'$, $E=P_{\mathfrak M}$,
$E'=P_{\mathfrak M'}$. Perform the operations of \defref{11.3.1} for the
closed linear set $\mathfrak M\cdot\mathfrak M'$. The correspondence
$X\longleftrightarrow X_{(\mathfrak M\cdot\mathfrak M')}$ is then a
one-to-one mapping and even an algebraic ring-isomorphism of the following
rings on each other:

(i) Of the ring of all $A\in\M$ with $EA=AE=A$ on all
$\M_{(\mathfrak M\cdot\mathfrak M')}$. 

(ii) Of the ring of all $A'\in\M'$ with $E'A'=A'E'=A'$ on all
$\M'_{(\mathfrak M\cdot\mathfrak M')}$. 

Besides we have

(iii) $(\M_{(\mathfrak M\cdot\mathfrak M')})'
=\M'_{(\mathfrak M\cdot\mathfrak M')}$,
$\M_{(\mathfrak M\cdot\mathfrak M')}$ being a factor.

(These are rings of operators in the space
$\mathfrak M\cdot\mathfrak M'\ne(0)$.)}

\Proof As $E\in\M$, $E'\in\M'$ therefore $E,E'$ commute, and so
$EE'=P_{\mathfrak M\cdot\mathfrak M'}$. The corollary to \thmref{III}
gives, owing to $E,E'\ne0$ that $EE'\ne0$,
$\mathfrak M\cdot\mathfrak M'\ne(0)$. As $E'\in\M'$ therefore
$E'_{(\mathfrak M)}\in\M'_{(\mathfrak M)}$ but
$E'_{(\mathfrak M)}=EE'=P_{\mathfrak M\cdot\mathfrak M'}$, so
$\mathfrak M\cdot\mathfrak M'\mathbin\eta\M_{(\mathfrak M)}$.

Now (i) results by applying first \lemref{11.3.3} (i) to
$\M,\M',\mathfrak M$ and then \lemref{11.3.3} (ii) to
$\M'_{(\mathfrak M)},\M_{(\mathfrak M)},
\mathfrak M\cdot\mathfrak M'$. (ii) follows from (i) by interchanging
$\M,\mathfrak M$ with $\M',\mathfrak M'$. (iii) follows by applying
Lemmas \hyperref[lem:11.3.2]{11.3.2} and
\hyperref[lem:11.3.4]{11.3.4} first to
$\M,\M',\mathfrak M$ and then to
$\M'_{(\mathfrak M)},\M_{(\mathfrak M)},
\mathfrak M\cdot\mathfrak M'$.

\Lemma{11.4.2}
{\itshape
Let $\M,\mathfrak M,\mathfrak M',E,E'$ be as above. Then we have:

(i) If $F$ runs over all projections $F\in\M$, $F\leq E$ then
$F_{(\mathfrak M\cdot\mathfrak M')}$ runs over all projections
$\in\M_{(\mathfrak M\cdot\mathfrak M')}$. If $D_{\M}(F)$ is a relative
dimension function in $\M$ then
\[
 D_{\M}^{(\mathfrak M\cdot\mathfrak M')}
 (F_{(\mathfrak M\cdot\mathfrak M')})=D_{\M}(F)
\]
is one in $\M_{(\mathfrak M\cdot\mathfrak M')}$. Denote their ranges by
 $\Delta$ resp. $\Delta_{(\mathfrak M\cdot\mathfrak M')}$, then
$\Delta_{(\mathfrak M\cdot\mathfrak M')}$ is the set of all
$\alpha\in\Delta$, $\alpha\leq D_{\M}(E)$.

(ii) If $F'$ runs over all projections $F'\in\M'$, $F'\leq E'$, then
$F'_{(\mathfrak M\cdot\mathfrak M')}$ runs over all projections
$\in\M'_{(\mathfrak M\cdot\mathfrak M')}$. If $D_{\M'}(F')$ is a
relative dimension function in $\M'$, then
\[
 D_{\M'}^{(\mathfrak M\cdot\mathfrak M')}
 (F'_{(\mathfrak M\cdot\mathfrak M')})=D_{\M'}(F')
\]
is one in $\M'_{(\mathfrak M\cdot\mathfrak M')}$. Denote their ranges by
 $\Delta'$ resp. $\Delta'_{(\mathfrak M\cdot\mathfrak M')}$, then
$\Delta'_{(\mathfrak M\cdot\mathfrak M')}$ is the set of all
$\alpha'\in\Delta'$, $\alpha'\leq D_{\M'}(E')$.}

\Proof (i) results by applying Lemmas \hyperref[lem:11.3.5]{11.3.5} (i)
 resp. (ii), \hyperref[lem:11.3.6]{11.3.6} and
\hyperref[lem:11.3.7]{11.3.7} first to
$\M,\M',\mathfrak M$ and then to
$\M'_{(\mathfrak M)},\M_{(\mathfrak M)},
\mathfrak M\cdot\mathfrak M'$. (ii) follows from (i) by interchanging
$\M,\mathfrak M$ with $\M',\mathfrak M'$.

We see that the coupled factorisation $\M,\M'$ in $\Hspace$ generates a
coupled factorisation
$\M_{(\mathfrak M\cdot\mathfrak M')},
\M'_{(\mathfrak M\cdot\mathfrak M')}$ in
$\mathfrak M\cdot\mathfrak M'\ne(0)$. \lemref{11.4.2} gives us the
means to determine the classes of the latter factors in terms of those of
the former ones.

\Lemma{11.4.3}
{\itshape
Let $\M,\M',\mathfrak M,\mathfrak M',E,E'$ be as above. Then
$\M_{(\mathfrak M\cdot\mathfrak M')}$,
$\M'_{(\mathfrak M\cdot\mathfrak M')}$ belong to the same ones of the
classes (I), (II), (III), as $\M,\M'$ (cf. \thmref{X}). In particular:

(i) If $\M,\M'$ are in class (I), and if $D_{\M}(F)$,
$D_{\M'}(F')$ are given in their standard normalisations, that the same is
true for
$D_{\M}^{(\mathfrak M\cdot\mathfrak M')}
(F_{(\mathfrak M\cdot\mathfrak M')})$,
$D_{\M'}^{(\mathfrak M\cdot\mathfrak M')}
(F'_{(\mathfrak M\cdot\mathfrak M')})$. We have class
$(\mathrm I_n)$ or $(\mathrm I_\infty)$ for
 $\M_{(\mathfrak M\cdot\mathfrak M')}$ (resp.
$\M'_{(\mathfrak M\cdot\mathfrak M')}$) corresponding to whether
 $\mathfrak M$ (resp. $\mathfrak M'$) is finite-$n$-dimensional or
infinite-dimensional.

(ii) If $\M,\M'$ are in class (II), then
 $\M_{(\mathfrak M\cdot\mathfrak M')}$ (resp.
$\M'_{(\mathfrak M\cdot\mathfrak M')}$) are in class
$(\mathrm{II}_1)$ or $(\mathrm{II}_\infty)$ corresponding to whether
 $\mathfrak M$ (resp. $\mathfrak M'$) is finite or infinite. So if
$\mathfrak M,\mathfrak M'$ are both finite, the $C$ of \thmref{X}
(referred to the standard normalisations) has an invariant meaning. If we
then vary $\mathfrak M,\mathfrak M'$ then $C$ will be proportional to
$D_{\M}(\mathfrak M)/D_{\M'}(\mathfrak M')$.

(iii) If $\M,\M'$ are in class (III), no further comment is necessary.}

\Proof Let first $\M,\M'$ be in class (I). Let $D_{\M}(F)$,
$D_{\M'}(F')$ be in standard normalisation. Put
$D_{\M}(E)=m$, $D_{\M'}(E')=n$ then $m,n>0$, that is
$m,n=1,2,\ldots,\infty$. So the ranges of
$D_{\M}^{(\mathfrak M\cdot\mathfrak M')}
(F_{(\mathfrak M\cdot\mathfrak M')})$ and
$D_{\M'}^{(\mathfrak M\cdot\mathfrak M')}
(F'_{(\mathfrak M\cdot\mathfrak M')})$ are
 $0,1,\ldots,m$ resp. $0,1,\ldots,n$. Thus
$\M_{(\mathfrak M\cdot\mathfrak M')}$,
$\M'_{(\mathfrak M\cdot\mathfrak M')}$ belong to class (I), and (i) is
true.

Let second, $\M,\M'$ be in class (II). Choose $D_{\M}(F)$ and then
$D_{\M'}(F')$ in any normalisations. Put
$D_{\M}(E)=\alpha_0$, $D_{\M'}(E')=\alpha_0'$, then
$\alpha_0,\alpha_0'>0$ and the ranges of
$D_{\M}^{(\mathfrak M\cdot\mathfrak M')}
(F_{(\mathfrak M\cdot\mathfrak M')})$ and
$D_{\M'}^{(\mathfrak M\cdot\mathfrak M')}
(F'_{(\mathfrak M\cdot\mathfrak M')})$ consist of all $\alpha$ with
 $0\leq\alpha\leq\alpha_0$ resp. of all $\alpha'$ with
$0\leq\alpha'\leq\alpha_0'$. Thus
$\M_{(\mathfrak M\cdot\mathfrak M')}$,
$\M'_{(\mathfrak M\cdot\mathfrak M')}$ belong to class (II), and all
parts of (ii), except those referring to $C$ are proved.

Concerning $C$ (cf. \thmref{X}) observe this:
$\mathfrak M_f^{\M_{(\mathfrak M\cdot\mathfrak M')}'}$ for
$f\in\mathfrak M\cdot\mathfrak M'$ is
\[
 [A_{(\mathfrak M\cdot\mathfrak M')}f;A\in\M],
\]
that is $[AEE'f;A\in\M, A\text{ commutes with }E]$. Then we may write as
well $EAEE'f$ and replacing $A$ by $EAE$ (this is $\in\M$, commutes with
$E$, and leaves $EAEE'f$ unaltered), we can admit all $A\in\M$. Now
$EAEE'f=EAf$ as $f\in\mathfrak M\cdot\mathfrak M'$ so we have
\[
 [EAf;A\in\M]=[Eg;g\in\mathfrak M_f^{\M}]
 =[EE_f^{\M}h;h\in\Hspace].
\]
But $E\in\M$, $E_f^{\M}\in\M'$ commute, so
$E\cdot E_f^{\M}=P_{\mathfrak M\cdot\mathfrak M_f^{\M}}$; therefore
\[
 \mathfrak M_f^{\M_{(\mathfrak M\cdot\mathfrak M')}'}
 =\mathfrak M\cdot\mathfrak M_f^{\M},
 \qquad
 E_f^{\M_{(\mathfrak M\cdot\mathfrak M')}'}
 =(E_f^{\M})_{(\mathfrak M\cdot\mathfrak M')}.
\]
Besides $f\in\mathfrak M'\mathbin\eta\M'$ implies
$\mathfrak M_f^{\M}\subset\mathfrak M'$,
$\mathfrak M\cdot\mathfrak M_f^{\M}\subset
\mathfrak M\cdot\mathfrak M'$; therefore we may even write
\[
 E_f^{\M_{(\mathfrak M\cdot\mathfrak M')}'}
 =(E_f^{\M})_{(\mathfrak M\cdot\mathfrak M')}.
\]
Similarly
\[
 E_f^{(\M'_{(\mathfrak M\cdot\mathfrak M')})'}
 =(E_f^{\M'})_{(\mathfrak M\cdot\mathfrak M')}.
\]

This proves that
$D_{\M}^{(\mathfrak M\cdot\mathfrak M')}
(F_{(\mathfrak M\cdot\mathfrak M')})$,
$D_{\M'}^{(\mathfrak M\cdot\mathfrak M')}
(F'_{(\mathfrak M\cdot\mathfrak M')})$ have the same $C$ as
$D_{\M}(F)$, $D_{\M'}(F')$. If $\mathfrak M,\mathfrak M'$ are finite, we
must pass to the standard normalisations, that is divide by
\begin{align*}
 D_{\M}^{(\mathfrak M\cdot\mathfrak M')}
 (1_{(\mathfrak M\cdot\mathfrak M')})
 &=D_{\M}^{(\mathfrak M\cdot\mathfrak M')}
 (E_{(\mathfrak M\cdot\mathfrak M')})
  =D_{\M}(E)=D_{\M}(\mathfrak M),\quad\text{resp.}\\
 D_{\M'}^{(\mathfrak M\cdot\mathfrak M')}
 (1_{(\mathfrak M\cdot\mathfrak M')})
 &=D_{\M'}^{(\mathfrak M\cdot\mathfrak M')}
 (E'_{(\mathfrak M\cdot\mathfrak M')})
 =D_{\M'}(E')=D_{\M'}(\mathfrak M').
\end{align*}
This multiplies $C$ by
$D_{\M}(\mathfrak M)/D_{\M'}(\mathfrak M')$, proving the rest of (ii).

Let finally $\M,\M'$ be in class (III). Then $D_{\M}(F)$,
$D_{\M'}(F')$ have both the range $0,\infty$. The ranges of
$D_{\M}^{(\mathfrak M\cdot\mathfrak M')}
(F_{(\mathfrak M\cdot\mathfrak M')})$,
$D_{\M'}^{(\mathfrak M\cdot\mathfrak M')}
(F'_{(\mathfrak M\cdot\mathfrak M')})$ must therefore be $0$ alone or
$0,\infty$. The former is impossible (because
$\mathfrak M\cdot\mathfrak M'\ne(0)$), so we have $0,\infty$ again, that
is, $\M_{(\mathfrak M\cdot\mathfrak M')}$,
$\M'_{(\mathfrak M\cdot\mathfrak M')}$ are in class (III). (iii) is
void.

\lemref{11.4.3} show that if case (II) occurs at all, then combinations
$\M,\M'$ of two cases $(\mathrm{II}_1)$ with an arbitrarily prescribed
$C$ (cf. \thmref{X}) exist. In fact: Choose $\M$ of class (II), then
$\M'$ is of class (II) too. Choose $D_{\M}(\mathfrak M)$,
$D_{\M'}(\mathfrak M')$ in such a normalisation that
$D_{\M}(\Hspace)\geq1$, $D_{\M'}(\Hspace)\geq1$. Then
$\mathfrak M\mathbin\eta\M$, $\mathfrak M'\mathbin\eta\M'$ can be chosen
with arbitrarily prescribed
$D_{\M}(\mathfrak M)=\alpha_0$,
$D_{\M'}(\mathfrak M')=\alpha_0'$ if only
$0\leq\alpha_0,\alpha_0'\leq1$. Thus $\alpha_0/\alpha_0'$ can be prescribed
arbitrarily, and with it $C$. Thus the invariant $C$ exists effectively,
as discussed at the end of \secref{10.2}, if case (II) exists at all. That
this is the case will be seen in \hyperref[chap:13]{Chapter XIII}.

\SectionStart{11.5}{\texorpdfstring{Determination of the resulting class for
composition $\mathrm R(\M,\N)$, $\N$ of class I}{Determination of the resulting
class for composition R(M,N), N of class I}}
We now return to the question formulated at the end of \secref{11.2}. We
will determine the class of $\mathrm R(\M,\N)$ if $\M,\N$ are two factors
which commute, $\N$ being of class (I).

\Lemma{11.5.1}
{\itshape
Let $\M,\N$ be two factors which commute, form the factor
$\mathrm R(\M,\N)$. Let $\varphi$ be minimal with respect to $\N$ (cf.
\defref{5.1.3}). Then
$\mathfrak M_{\varphi}^{\N'}\mathbin\eta\N$ and so
$\mathbin\eta\mathrm R(\M,\N)$, and
\[
 X\longleftrightarrow X_{(\mathfrak M_{\varphi}^{\N'})}
\]
is an algebraic ring-isomorphism of $\M$ and
$(\mathrm R(\M,\N))_{(\mathfrak M_{\varphi}^{\N'})}$.}

\Proof $\mathfrak M_{\varphi}^{\N'}\mathbin\eta\N$ is obvious; as
$\N\subset\mathrm R(\M,\N)$,
$\N'\supset(\mathrm R(\M,\N))'$ it implies
$\mathfrak M_{\varphi}^{\N'}\mathbin\eta\mathrm R(\M,\N)$. So
$E_{\varphi}^{\N'}\in\N$ and $\mathrm R(\M,\N)$; thus it commutes with
every $X\in\M$.

$X\longleftrightarrow X_{(\mathfrak M_{\varphi}^{\N'})}$ is an algebraic
ring-isomorphism if it is a one to one mapping; besides it is a one to one
mapping of all $Y\in\mathrm R(\M,\N)$ with
$E_{\varphi}^{\N'}Y=YE_{\varphi}^{\N'}=Y$ on
$(\mathrm R(\M,\N))_{(\mathfrak M_{\varphi}^{\N'})}$ (by
\lemref{11.3.3}, (i)). So we must only prove this: The correspondence
\[
 X_{(\mathfrak M_{\varphi}^{\N'})}
 =Y_{(\mathfrak M_{\varphi}^{\N'})}
\]
generates a one to one mapping of all $X\in\M$ on all
$Y\in\mathrm R(\M,\N)$ with
$E_{\varphi}^{\N'}Y=YE_{\varphi}^{\N'}=Y$. If $X\in\M$ is given, then
$Y=E_{\varphi}^{\N'}X$ has the above properties, and clearly
$Y_{(\mathfrak M_{\varphi}^{\N'})}
=X_{(\mathfrak M_{\varphi}^{\N'})}$. Besides if $X$ is given, our
correspondence determines $Y$ uniquely, because such a $Y$ is uniquely
determined by its $Y_{(\mathfrak M_{\varphi}^{\N'})}$, just because the
correspondence
$Y\longleftrightarrow Y_{(\mathfrak M_{\varphi}^{\N'})}$,
$Y\in\mathrm R(\M,\N)$,
$E_{\varphi}^{\N'}Y=YE_{\varphi}^{\N'}=Y$ is one-to-one (by
\lemref{11.3.3}, (i)). So
$X_{(\mathfrak M_{\varphi}^{\N'})}
=Y_{(\mathfrak M_{\varphi}^{\N'})}$ is a one-to-one mapping of all
$X\in\M$ on a part of the $Y\in\mathrm R(\M,\N)$,
$E_{\varphi}^{\N'}Y=YE_{\varphi}^{\N'}=Y$. All we must prove now is that
all these $Y$ correspond to some $X$, that is, that they all have the form
$Y=E_{\varphi}^{\N'}X$, $X\in\M$.

We may as well prove: If $Y\in\mathrm R(\M,\N)$ then
$E_{\varphi}^{\N'}YE_{\varphi}^{\N'}=E_{\varphi}^{\N'}X$ for some
$X\in\M$. Let $\P$ be the set of those $Y_0$ for which this holds for all
$Y=AY_0B$, $A,B\in\N$. That is: For which these
$E_{\varphi}^{\N'}AY_0BE_{\varphi}^{\N'}
\in\M_{(\mathfrak M_{\varphi}^{\N'})}$, cf. \lemref{11.3.3} (i).
(Observe $E_{\varphi}^{\N'}\in\N\subset\M'$.) Clearly $\P$ is weakly
closed (along with $\M_{(\mathfrak M_{\varphi}^{\N'})}$) and
$Y_1,Y_2\in\P$ imply $\alpha Y_1,Y_1^*,Y_1+Y_2\in\P$ and
$Y_1Y_2\in\P$, the last one because
\begin{align*}
 &\sum_{p=1}^{\infty}E_{\varphi}^{\N'}AY_1U_{1,p}E_{\varphi}^{\N'}
 \cdot E_{\varphi}^{\N'}U_{p,1}Y_2BE_{\varphi}^{\N'}\\
 &\quad=\sum_{p=1}^{\infty}E_{\varphi}^{\N'}AY_1U_{1,p}
 P_{[f_{1,1},f_{1,2},\ldots]}U_{p,1}Y_2BE_{\varphi}^{\N'}\\
 &\quad=\sum_{p=1}^{\infty}E_{\varphi}^{\N'}AY_1
 P_{[f_{p,1},f_{p,2},\ldots]}Y_2BE_{\varphi}^{\N'}
 =E_{\varphi}^{\N'}AY_1Y_2BE_{\varphi}^{\N'}.
\end{align*}
(Choose the $U_{m,p},f_{m,p}$ by \lemref{5.3.6} for $\N$ with
$f_{1,1}=\varphi$.) So $\P$ is a ring.

If $Y_0=XZ$, $X\in\M$, $Z\in\N$ then for $A,B\in\N$
\[
 E_{\varphi}^{\N'}AXZBE_{\varphi}^{\N'}
 =E_{\varphi}^{\N'}AZBE_{\varphi}^{\N'}\cdot X.
\]
As $E_{\varphi}^{\N'}$ is minimal with respect to $\N$, the last part of
the proof of \lemref{5.1.3} applies:
\[
 E_{\varphi}^{\N'}AZBE_{\varphi}^{\N'}=\alpha E_{\varphi}^{\N'}
\]
(replace there $\M,E=E_f^{\M'},B$ by
$\N,E_{\varphi}^{\N'},E_{\varphi}^{\N'}AZBE_{\varphi}^{\N'}$). Thus the
above expression is
$=\alpha E_{\varphi}^{\N'}X\in
\M_{(\mathfrak M_{\varphi}^{\N'})}$. So $Y_0=XZ\in\P$. This implies
$\P\supset\M$ and $\N$, and so $\P\supset\mathrm R(\M,\N)$. Application
to $Y_0\in\mathrm R(\M,\N)$, $A=B=1$ completes the proof.

\Lemma{11.5.2}
{\itshape
Let $\M,\N$ be two factors which commute, $\N$ being of class (I). Then
$\mathrm R(\M,\N)$ belongs to the same class (I), (II), (III) as $\M$.
In particular:

(i) If $\M$ is in a class $(\mathrm I_m)$,
$m=1,2,\ldots,\infty$, and $\N$ is in a class $(\mathrm I_n)$,
$n=1,2,\ldots,\infty$ then $\mathrm R(\M,\N)$ is in class
$(\mathrm I_{mn})$.

(ii) If $\M$ is in class $(\mathrm{II}_m)$, $m=1,\infty$ and $\N$ is in
class $(\mathrm I_n)$, $n=1,2,\ldots,\infty$ then
$\mathrm R(\M,\N)$ is in class $(\mathrm{II}_p)$, where $p=1$ if $m,n$
are both finite, and $p=\infty$ if $m$ or $n$ is infinite.

(iii) If $\M$ is in class (III), no further comment is necessary.}

\Proof As $\N$ is of class (I), minimal $\varphi$'s (with respect to
$\N$) exist by \thmref{IV}. Choose one, and put
$\mathfrak M=\mathfrak M_{\varphi}^{\N'}$, then $\M$ is algebraically
ring-isomorphic to
$(\mathrm R(\M,\N))_{(\mathfrak M)}$ by \lemref{11.5.1}. Thus it has the
same class by \thmref{IX}. Now \lemref{11.4.3} shows (replacing there
$\M,\mathfrak M$ by $\mathrm R(\M,\N),\mathfrak M$), that $\M$ and
$\mathrm R(\M,\N)$ belong to the same class (I), (II), (III).

Let us now consider (i)--(iii).

If 1 is infinite with respect to $\M$ then it is obviously so for
$\mathrm R(\M,\N)$ too; the same being true for $\N$ and
$\mathrm R(\M,\N)$. This takes care of (i) and of (ii) if $m=\infty$ or
$n=\infty$, (iii) is void. So we need only to consider (i) and (ii) for
$m,n$ finite.

Construct for $\N$ the normalised orthogonal system
$\varphi_1,\varphi_2,\ldots$ from \lemref{5.3.4}, where all $\varphi_i$ are
minimal, $\mathfrak M_{\varphi_1}^{\N'},
\mathfrak M_{\varphi_2}^{\N'},\ldots$ are mutually orthogonal, and
$[\mathfrak M_{\varphi_1}^{\N'},
\mathfrak M_{\varphi_2}^{\N'},\ldots]=\Hspace$. We know by Lemmas
\hyperref[lem:5.3.5]{5.3.5}--\hyperref[lem:5.3.8]{5.3.8} that as $\N$ is in the case
$(\mathrm I_n)$, the number of these $\varphi_i$'s is $n$:
$\varphi_1,\ldots,\varphi_n$ ($n$ is finite!).

\lemref{11.5.1} applies to every $\varphi_i$. So
$(\mathrm R(\M,\N))_{(\mathfrak M_{\varphi_i}^{\N'})}$ is in the same case
as \(\M\), \((\mathrm I_m)\) resp. \((\mathrm{II}_1)\) (for (i) resp.
(ii)). By this, \(\mathfrak M_{\varphi_i}^{\N'}\) has the (relative)
dimension \(m\) in the standard normalisation of
\((\mathrm R(\M,\N))_{(\mathfrak M_{\varphi_i}^{\N'})}\) resp. a finite
relative dimension in any normalisation of
\((\mathrm R(\M,\N))_{(\mathfrak M_{\varphi_i}^{\N'})}\). The same is
true by \lemref{11.4.3}, (i) resp. (ii), for
\(\mathfrak M_{\varphi_i}^{\N'}\) and \(\mathrm R(\M,\N)\). As the
\(\mathfrak M_{\varphi_i}^{\N'}\) are mutually orthogonal and
\[
 [\mathfrak M_{\varphi_1}^{\N'},\ldots,
   \mathfrak M_{\varphi_n}^{\N'}]=\Hspace,
\]
the additivity of the relative dimension implies that it is \(mn\)
(standard normalisation) resp. finite (any normalisation) for \(\Hspace\) in
\(\mathrm R(\M,\N)\). Thus we have case \((\mathrm I_{mn})\) resp.
\((\mathrm{II}_1)\) for \(\mathrm R(\M,\N)\).

This completes the proof.

\lemref{11.5.2} shows that if case (II) exists at all, then
\((\mathrm{II}_\infty)\) exists too. (\((\mathrm{II}_1)\) has been
discussed in \secref{11.4}.) It suffices to take an \(\M_0\) of case (II)
in an \(\Hspace_1\), any \(\N_0\) of case \((\mathrm I_\infty)\) in an
\(\Hspace_2\) (for instance: \(\Hspace_2\) a Hilbert space,
\(\N_0=\B_2\)), and then form
\(\Hspace=\Hspace_1\otimes\Hspace_2\), \(\M=\M_0^{(1)}\),
\(\N=\N_0^{(2)}\). Then \(\mathrm R(\M,\N)\) will certainly be of class
\((\mathrm{II}_\infty)\). But all this will be settled exhaustively in
\hyperref[chap:13]{Chapter XIII}.
% --- End of parts/04-part-iii.tex ---
\Needspace{10\baselineskip}
\PartHeadingTOC{IV}{Examples of case (II)}{Examples of Case II}

\ChapterHeading{12}{XII}{\texorpdfstring{Construction of $\M,\M'$}
{Construction of M and M-prime}}

\SectionStart{12.1}{\texorpdfstring{$\mathfrak G$ and $S$, the spaces
$\Hspace_{\mathfrak G}$, $\Hspace_S$, $\Hspace_{S\mathfrak G}$}
{G and S, the three associated spaces}}
The considerations at the end of \secref{8.6} have established the existence
of all cases $(\mathrm I_n)$, $(\mathrm I_\infty)$ and of all their
combinations for factors $\M$ resp. coupled factorisation $\M,\M'$. We are
going to do now the same for the cases $(\mathrm{II}_1)$,
$(\mathrm{II}_\infty)$. We will thus answer a great part of the questions
raised by Problems \hyperref[prob:3]{3} and \hyperref[prob:4]{4}.

Our first objective is to construct certain coupled factorisations $\M,\M'$
which belong to the classes $(\mathrm{II}_1),(\mathrm{II}_1)$ (the $C$ of
\thmref{X} being $=1$) or
$(\mathrm{II}_\infty),(\mathrm{II}_\infty)$. These factorisations will
contain many arbitrary parameters, and therefore represent a rather wide
variety of examples. In fact, further constructions based on them will answer
\probref{1} in the negative.

\Definition{12.1.1}
We will consider groups $\mathfrak G$ of the following kind:
\begin{enumerate}[label=(\roman*)]
\item $\mathfrak G$ is a group, that is a composition rule $ab$, an inverse
  $a^{-1}$ and a unity $1$ defined in $\mathfrak G$, with the properties
  \[
    (ab)c=a(bc),\qquad aa^{-1}=a^{-1}a=1,\qquad a1=1a=a.
  \]
  (So $ab$ is associative but not necessarily commutative).
\item $\mathfrak G$ is finite or countably infinite.
\end{enumerate}

\Definition{12.1.2}
We will consider spaces $S$ of the following kind: A Lebesgue outer measure is
defined in $S$, that is, for every subset $T$ of $S$ a real number
$\mu^*(T)\geq0$, $\leq\infty$ is given, with the following properties:
\begin{enumerate}[label=(\roman*)]
\item $T_1\subset T_2$ implies $\mu^*(T_1)\leq\mu^*(T_2)$.
\item For every (finite or infinite) sequence $T_1,T_2,\ldots$
  \[
    \mu_1^*(T_1\mathbin{\dot+}T_2\mathbin{\dot+}\cdots)
    \leq\mu_1^*(T_1)+\mu_2^*(T_2)+\cdots.
  \]
\end{enumerate}
We now define measurability following Carath\'eodory
(cf. \cite[p.~246]{ref4}):

\noindent$(*)$ A subset $T$ of $S$ is \emph{measurable}, if and only if, for
every subset $T'$ of $S$,
\[
 \mu^*(T')=\mu^*(TT')+\mu^*(T'-TT').
\]
We then write $\mu(T)$ for $\mu^*(T)$ and call it the \emph{measure} of $T$.

We now continue the enumeration of our postulates.
\begin{enumerate}[label=(\roman*),start=3]
\item If $\mu^*(T)<\alpha$ then a measurable $T_0$ with $T_0\supset T$,
  $\mu(T_0)<\alpha$ exists.
\item A (finite or infinite) sequence $T^{(1)},T^{(2)},\ldots$ of measurable
  sets with finite measure exists, with the following property: If $x,y\in S$,
  and if $x\in T^{(i)}$ is equivalent to $y\in T^{(i)}$ (for all
  $i=1,2,\ldots$), then $x=y$.
\end{enumerate}

Observe the following consequences of \defref{12.1.2}: $\mu^*(T)\geq0$,
$\leq\infty$ and conditions (i)--(ii) correspond to Carath\'eodory's
conditions (I)--(III), without (IV), (Cf. \cite[pp.~238--239]{ref4}). This
omission however is natural, because (IV) makes explicit use of the notion of
distance, whereas we did not require that a distance (or any sort of a topology)
be defined in $S$. (iii) is, (remembering (i)) equivalent to Carath\'eodory's
condition (V), (cf. \cite[p.~258]{ref4}), therefore our outer measure is a
regular measure function (cf. loc. cit.) if we disregard the topological
condition (IV). Closer inspection of Carath\'eodory's deductions shows that in
consequence, all those results on outer measure, measurability, and measure loc.
cit. remain true which are not explicitly topological in their statements. This
applies to the considerations on pp.~246--250 loc. cit. (cf. the remark on
p.~257 there). Thus $0$ and $S$ are measurable; if $T',T''$ are measurable,
then $T'-T'T''$ is too; if $T',T'',\ldots$ are measurable, then
$T'\mathbin{\dot+}T''\mathbin{\dot+}\cdots$ and
$T'\cdot T''\cdots$ are too. Besides, the well known convergence theorems on
measure are valid.

As we see, the chief difference between the characteristics of the situation in
$S$ and those of the one which exists loc. cit., is this: The sets which are
known to be measurable, and which serve as a basis for all constructions of
measurable functions (cf. below \defref{12.1.3}) are there the Borel-sets (a
topological notion) while they are now the sets
$T^{(1)},T^{(2)},\ldots$ from (iv).

In fact, this is essentially the meaning of (iv), which is the only one of our
postulates without an analogue on Carath\'eodory's list.: It serves to replace
the (absent or ignored) topology in $S$, and in particular the ``topological
separability'' of $S$.

Let us mention the following obvious examples of spaces $S$:

\textbf{Example a).} Let $S$ be a (finite or infinite) sequence
$S=(x_1,x_2,\ldots)$. Let a corresponding sequence of real numbers
$\bar\alpha_n\geq0$, $<\infty$ be given. Define
\[
 \mu^*(T)\equiv\sum_{x_n\in T}\bar\alpha_n.
\]
Then (i), (ii) are obviously fulfilled; (iii) is fulfilled because all sets
$T\subset S$ are measurable; (iv) is fulfilled, because we may choose for
$T^{(1)},T^{(2)},\ldots$ the set of all one-element $T\subset S$.

\textbf{Example b).} Let $S$ be any measurable subset of a finite dimensional
Euclidean space. Let $\mu^*(T)$ be the common Lebesgue measure. Then
(i)--(iii) are obviously fulfilled; (iv) is fulfilled because we may choose for
$T^{(1)},T^{(2)},\ldots$ the sequence $Ss_1,Ss_2,\ldots$, where
$s_1,s_2,\ldots$ are all spheres with rational coordinates of the center and a
rational radius.

We now introduce Lebesgue integration in the customary way:

\Definition{12.1.3}
A complex valued function $f(x)$, defined for all $x\in S$ is called
\emph{measurable} if the sets
$(x;\Re(f(x))<\alpha)$, $(x;\Im(f(x))<\alpha)$
($\Re=$ real part, $\Im=$ imaginary part) are measurable for every real
$\alpha$. The notions of summability and of the (Legesgue) integral
$\int_Sf(x)\,dx$ of a measurable function $f(x)$ are then defined in any of the
several equivalent customary ways. (The procedure of Carath\'eodory,
\cite[pp.~418--427]{ref4}, which coincides with Lebesgue's ``geometrical''
definition, \cite[pp.~116--120]{ref10}, would necessitate introducing first the
notion of outer measure in a product space. Lebesgue's ``analytical''
definition, \cite[pp.~112--116]{ref10}, however, is directly applicable;
another very convenient definition is due to Bochner,
\cite[pp.~265--267]{ref2}).

We now form three functional spaces:

\Definition{12.1.4}
Let $\Hspace_{\mathfrak G}$ be the set of all complex valued functions $f(a)$,
defined for all $a\in\mathfrak G$ with a finite
$\sum_{a\in\mathfrak G}|f(a)|^2$. Define
\[
 (f,g)_{\mathfrak G}=\sum_{a\in\mathfrak G}f(a)\overline{g(a)}.
\]

Let $\Hspace_S$ be the set of all complex valued functions $f(x)$, defined for
all $x\in S$ which are measurable in $x$, and with a finite
$\int_S|f(x)|^2\,dx$. Consider $f,g$ as identical if $f(x)=g(x)$ except for an
$x$-set of measure $0$. Define
\[
 (f,g)_S=\int_S f(x)\overline{g(x)}\,dx.
\]

Let $\Hspace_{S\mathfrak G}$ be the set of all complex valued functions
$F(x,a)$ defined for all $x\in S$, $a\in\mathfrak G$, which are measurable in
$x$ (for every fixed $a\in\mathfrak G$) and with a finite
$\sum_{a\in\mathfrak G}\int_S|F(x,a)|^2\,dx$. Consider $F,G$ as identical,
if $F(x,a)=G(x,a)$ except for an $x$-set of measure $0$ (depending on $a$).
Define
\[
 (F,G)_{S\mathfrak G}
 =\sum_{a\in\mathfrak G}\int_SF(x,a)\overline{G(x,a)}\,dx.
\]
$\alpha\cdot$ and $+$ are defined in the obvious way in all three spaces
$\Hspace_{\mathfrak G}$, $\Hspace_S$, $\Hspace_{S\mathfrak G}$.

By \defref{12.1.2} (iv), no two points $x\neq y$ can have the property
\[
 x,y\notin T^{(1)}\mathbin{\dot+}T^{(2)}\mathbin{\dot+}\cdots.
\]
$S-(T^{(1)}\mathbin{\dot+}T^{(2)}\mathbin{\dot+}\cdots)$ is measurable; and
if its measure is finite, we can add it to the sequence
$T^{(1)},T^{(2)},\ldots$ thus forcing
$T^{(1)}\mathbin{\dot+}T^{(2)}\mathbin{\dot+}\cdots=S$. So the only case
where this is not possible is when
\[
 S-(T^{(1)}\mathbin{\dot+}T^{(2)}\mathbin{\dot+}\cdots)=(x_0),
 \qquad \mu((x_0))=\infty.
\]
Then we must have $f(x_0)=0$ and $F(x_0,a)=0$ for all
$f\in\Hspace_S$ resp. $F\in\Hspace_{S\mathfrak G}$, so we may omit $x_0$
from $S$. (If $x\in S$, $x\neq x_0$ then $x\in T^{(i)}$ for some
$i=1,2,\ldots$, and so $\mu^*((x))\leq\mu(T^{(i)})<\infty$. Therefore
\defref{12.1.5}, cf. below, excludes $x_0a\neq x_0$ that is, we have
$x_0a=x_0$. So the omission of $x_0$ does not affect the operation $xa$
either). In what follows we can thus assume
\[
 T^{(1)}\mathbin{\dot+}T^{(2)}\mathbin{\dot+}\cdots=S.
\]

If $\mu(S)=0$, then $\Hspace_S$ and $\Hspace_{S\mathfrak G}$ would consist
of $0$ only. We therefore assume explicitly that $\mu(S)>0$.

\Lemma{12.1.1}
\emph{All three spaces $\Hspace_{\mathfrak G}$, $\Hspace_S$,
$\Hspace_{S\mathfrak G}$ fulfill the postulates A, B, C, E of
\cite[pp.~64--66]{ref16}; thus they are finite dimensional Euclidean or
Hilbert spaces (and $\neq0$).}

\emph{Using the complete normalised orthogonal set of all functions}
\[
 \varphi_{a_0}(a)=
 \begin{cases}
  1&\text{for }a=a_0,\\
  0&\text{for }a\neq a_0,
 \end{cases}
 \qquad a_0\in\mathfrak G\text{ in }\Hspace_{\mathfrak G},
\]
\emph{and setting up the correspondence}
\[
 F(x,a)\sim\langle F(x,\bar a_1),F(x,\bar a_2),\ldots\rangle
\]
\emph{($\bar a_1,\bar a_2,\ldots$ is some enumeration of $\mathfrak G$, so
that the above $a_0$ runs over $\bar a_1,\bar a_2,\ldots$),
$\Hspace_{S\mathfrak G}$ is isomorphic with
$\Hspace_{\mathfrak G}\otimes\Hspace_S$ in the sense of
\defref{2.4.1} and \lemref{2.4.1}.}

We will use in what follows the more symmetric notation
\[
 F(x,a)\sim\langle F(x,a_0);a_0\in\mathfrak G\rangle.
\]

\Proof If we use an enumeration $\bar a_1,\bar a_2,\ldots$ of $\mathfrak G$,
and put $f(\bar a_n)=x_n$, $n=1,2,\ldots$ then
$f(a)\sim(x_1,x_2,\ldots)$ and our definition coincides with the original
definition of finite dimensional Euclidean or of Hilbert space
(cf. \cite[p.~69]{ref16}). Concerning $\Hspace_S$ it suffices to observe that
the considerations of \cite[pp.~108--111]{ref16}, on A, B, E apply immediately
to it; while those on C apply, if the system of neighborhoods occurring there
is replaced by $T^{(1)},T^{(2)},\ldots$.

As
$\mu(S)=\mu(T^{(1)}\mathbin{\dot+}T^{(2)}\mathbin{\dot+}\cdots)>0$ some
$\mu(T^{(i)})>0$. At the same time it is $<\infty$. Then
\[
 f(x)=\begin{cases}
  1,&x\in T^{(i)},\\
  0,&x\notin T^{(i)}
 \end{cases}
\]
belongs to $\Hspace_S$ and is $\neq0$. So $\Hspace_S\neq(0)$. The last part
of our Lemma results immediately by comparing our definitions with
\defref{2.4.1} and \lemref{2.4.1}. As $\Hspace_{S\mathfrak G}$ is isomorphic
with $\Hspace_{\mathfrak G}\otimes\Hspace_S$ it must fulfill A, B, C, E, too.
As $\Hspace_S\neq(0)$, therefore $\Hspace_{S\mathfrak G}\neq(0)$.

Heretofore the connection between $S$ and $\mathfrak G$ was merely formal,
due to our forming the space
$\Hspace_{S\mathfrak G}=\Hspace_{\mathfrak G}\otimes\Hspace_S$. An intrinsic
connection is established by the following assumptions:

\Definition{12.1.5}
$\mathfrak G$ is an $m$-group in $S$, if for every $a\in\mathfrak G$ there is
defined a one to one mapping $x\rightleftarrows xa$ of $S$ on itself, with the
following properties:
\begin{enumerate}[label=(\roman*)]
\item If $T\subset S$ and $T_a$ is the $x\rightleftarrows xa$ image of $T$
  then $\mu^*(T)=\mu^*(T_a)$. (This implies that the measurability of $T$ is
  equivalent to that of $T_a$.)
\item $(xa)b=x(ab)$. (This implies that $x\cdot1=x$ and that
  $x\rightleftarrows xa^{-1}$ is inverse to $x\rightleftarrows xa$.)
\item If $a\neq1$ then $xa=x$ holds only for $x$-sets of measure $0$.
\end{enumerate}
$\mathfrak G$ is \emph{ergodic} in $S$, if besides (i)--(iii) we have further:
\begin{enumerate}[label=(\roman*),start=4]
\item If a measurable $T\subset S$ differs from each $T_a$,
  $a\in\mathfrak G$ only by a set of measure $0$ (but depending on $a$, this
  means $\mu((T\mathbin{\dot+}T_a)-(T\cdot T_a))=0$) then either
  $\mu(T)=0$ or $\mu(S-T)=0$.
\end{enumerate}

Observe that (iv) differs from the customary definitions of ergodicity in two
ways: First $\mathfrak G$ is not necessarily the simple translation group, not
even necessarily Abelian, second (which is essential if $\mu(S)=\infty$) we
did not restrict the validity of (iv) to $T$'s with a finite $\mu(T)$ only.

\SectionStart{12.2}{\texorpdfstring{Discussion of $L=(L_{\varphi(x)})$ and of
$U=(U_a)$ in $\Hspace_S$}{Discussion of L and U in H-sub-S}}
We introduce and discuss first some operators in $\Hspace_S$.

\Lemma{12.2.1}
\emph{Let $\mathfrak G$ be an $m$-group in $S$. Define the operators:}
\begin{enumerate}[label=(\roman*)]
\item \emph{$U_af(x)=f(xa)$ for any $a\in\mathfrak G$.}
\item \emph{$L_{\varphi(x)}f(x)=\varphi(x)f(x)$ for any bounded and
  measurable complex valued function $\varphi(x)$ defined for all $x\in S$.
  These operators are bounded, $U_a$ is even unitary.}
\end{enumerate}

\Proof We have
\[
 \int_S|U_a(f(x))|^2\,dx=\int_S|f(xa)|^2\,dx
 =\int_S|f(x)|^2\,dx
\]
so $U_a$ is bounded, and $\lVert U_af\rVert=\lVert f\rVert$. Now clearly
$U_{a^{-1}}U_a=U_aU_{a^{-1}}=1$ so $U_a$ has an inverse, thus it is unitary
(cf. \cite[p.~71]{ref16}).

If $|\varphi(x)|\leq C$ for all $x\in S$, then
\begin{align*}
 \int_S|L_{\varphi(x)}f(x)|^2\,dx
 &=\int_S|\varphi(x)|^2|f(x)|^2\,dx\\
 &\leq C^2\int_S|f(x)|^2\,dx,
 \qquad \lVert L_{\varphi(x)}f\rVert\leq C\lVert f\rVert.
\end{align*}
So $L_{\varphi(x)}$ too is bounded.

\Lemma{12.2.2}
\emph{Let $\mathfrak G$ be an $m$-group in $S$. Let $L$ be the set of all
operators $L_{\varphi(x)}$ from \lemref{12.2.1}, (iii). Then $L$ is a ring and
$L=L'$.}

\Proof As $L'$ is necessarily a ring, it suffices to prove $L=L'$. As
\[
 (L_{\varphi(x)})^*=L_{\overline{\varphi(x)}},
 \qquad
 L_{\varphi(x)}\mathbin{\cdot}L_{\psi(x)}=L_{\varphi(x)\psi(x)}
\]
therefore every $L_{\psi(x)}$ commutes with every $L_{\varphi(x)}$ and
$(L_{\varphi(x)})^*$ thus $L_{\psi(x)}\in L'$ and so $L\subset L'$. So we
must only prove $L'\subset L$.

Assume therefore $A\in L'$. Then $A$ commutes with every $L_{\varphi(x)}$,
that is $AL_{\varphi(x)}f(x)=L_{\varphi(x)}Af(x)$. That is:
\begin{equation*}
 A(\varphi(x)f(x))=\varphi(x)Af(x).
 \tag{$*$}\label{eq:12.2-star}
\end{equation*}
The assumptions are: $\varphi(x),f(x)$ are measurable, $\varphi(x)$ is
bounded, $\int_S|f(x)|^2\,dx$ is finite, the equation holds with the exception
of an $x$-set of measure $0$. Such exceptions, by the way, are admitted in all
the equations we are going to derive.

Let $T^{(1)},T^{(2)},\ldots$ be the sets from \defref{12.1.2} (iv); assume
(following the remark before \lemref{12.1.1}), that
$T^{(1)}\mathbin{\dot+}T^{(2)}\mathbin{\dot+}\cdots=S$. Define
\[
 e_i(x)=\begin{cases}
  1,&x\in T^{(1)}\mathbin{\dot+}\cdots\mathbin{\dot+}T^{(i)},\\
  0,&x\notin T^{(1)}\mathbin{\dot+}\cdots\mathbin{\dot+}T^{(i)}.
 \end{cases}
\]
Apply now \eqref{eq:12.2-star} to $f=e_i$. Then
$A(e_i(x)\varphi(x))=(Ae_i(x))\varphi(x)$ obtains. For $i\geq j$,
$e_i(x)e_j(x)=e_j(x)$, so $\varphi=e_j$ gives
$Ae_j(x)=(Ae_i(x))e_j(x)$. Thus if
$x\in T^{(1)}\mathbin{\dot+}\cdots\mathbin{\dot+}T^{(j)}$,
$Ae_j(x)=Ae_i(x)$; in other words: for
$x\in T^{(1)}\mathbin{\dot+}\cdots\mathbin{\dot+}T^{(j)}$ all $Ae_i(x)$,
$i\geq j$, agree. Thus for every $x$ all $Ae_i(x)$ with a sufficiently great
$i$ have the same value. Call this value $\psi(x)$ as
$\psi(x)=\lim_{i\to\infty}Ae_i(x)$ this is a measurable function. Now we have
$Ae_i(x)=\psi(x)e_i(x)$. And our original equation becomes:
$Ae_i(x)\varphi(x)=\psi(x)e_i(x)\varphi(x)$. That is:
$Af(x)=\psi(x)f(x)$ where $f(x)=e_i(x)\varphi(x)$. Thus this equation is
proved if $f(x)$ is bounded and $0$ for all
$x\notin T^{(1)}\mathbin{\dot+}\cdots\mathbin{\dot+}T^{(i)}$ for some
$i=1,2,\ldots$.

$A$ is bounded, say $\lVert Af\rVert\leq C\lVert f\rVert$; therefore
\[
 \int_S|\psi(x)|^2|f(x)|^2\,dx
 \leq C^2\int_S|f(x)|^2\,dx
\]
for these $f$'s. Now put
\[
 f(x)=\begin{cases}
  1,&|\psi(x)|\geq C+1,\quad
    x\in T^{(1)}\mathbin{\dot+}\cdots\mathbin{\dot+}T^{(i)},\\
  0,&\text{otherwise}.
 \end{cases}
\]
Denote the set of all $x$ with $|\psi(x)|\geq C+1$ by $T_0$. Then the above
inequality gives:
\[
 (C+1)^2\mu\bigl(T_0(T^{(1)}\mathbin{\dot+}\cdots
 \mathbin{\dot+}T^{(i)})\bigr)
 \leq C^2\mu\bigl(T_0(T^{(1)}\mathbin{\dot+}\cdots
 \mathbin{\dot+}T^{(i)})\bigr),
\]
\[
 \mu\bigl(T_0(T^{(1)}\mathbin{\dot+}\cdots
 \mathbin{\dot+}T^{(i)})\bigr)=0
\]
and passing to the limit as $i\to\infty$, $\mu(T_0)=0$. So we have everywhere
$|\psi(x)|\leq C+1$ (as $T_0$ does not matter), that is: $\psi(x)$ is
bounded. We can therefore write: $Af=L_\psi f$ if $f(x)$ is as described
above.

Now these $f$ form an everywhere dense set in $\Hspace_S$. If any
$f\in\Hspace_S$ is given, define
\[
 f_i(x)=\begin{cases}
  f(x),&|f(x)|<i,\quad
  x\in T^{(1)}\mathbin{\dot+}\cdots\mathbin{\dot+}T^{(i)},\\
  0,&\text{otherwise}
 \end{cases}
\]
for $i=1,2,\ldots$; then every $f_i$ meets these requirements, and the $f_i$
converge for $i\to\infty$ to $f$ (in $\Hspace_S$). As $A,L_\psi$ are both
bounded, we have therefore $Af=L_\psi f$ for every $f$; that is
$A=L_\psi\in L$. So we have established $L'\subset L$ too, thus completing
the proof.

\Lemma{12.2.3}
\emph{Let $\mathfrak G$ be an $m$-group in $S$. The equation
$L_{\varphi(x)}=L_{\psi(x)}U_a$ holds if and only if either $a=1$,
$\varphi(x)=\psi(x)$ (except for an $x$-set of measure $0$) or $a\neq1$,
$\varphi(x)=\psi(x)=0$ (except for an $x$-set of measure $0$).}

\Proof The sufficiency of these conditions is obvious, so we need only to prove
their necessity. Assume therefore $L_{\varphi(x)}=L_{\psi(x)}U_a$. So we
have
\[
 \varphi(x)f(x)=\psi(x)f(xa)
\]
whenever $\int_S|f(x)|^2\,dx$ is finite. This equation, and similarly all
those we are going to derive, holds with the exception of an $x$-set of measure
$0$. Use the $e_i$ from the proof of \lemref{12.2.2}, then we have:
$e_i(x)\varphi(x)=e_i(xa)\psi(x)$, that is $\varphi(x)=\psi(x)$ if $x$ and
$xa\in T^{(1)}\mathbin{\dot+}\cdots\mathbin{\dot+}T^{(i)}$. Summing over all
$i=1,2,\ldots$ we obtain $\varphi(x)=\psi(x)$ for $x\in S$. This settles the
case $a=1$. Assume now $a\neq1$. Substituting $\varphi(x)=\psi(x)$ in our
original equation gives $\varphi(x)f(x)=\varphi(x)f(xa)$ and so
\[
 \int_S|\varphi(x)|^2|f(x)-f(xa)|^2\,dx=0.
\]
Define now
\[
 e_i'(x)=\begin{cases}
  1,&x\in T^{(i)},\\
  0,&x\notin T^{(i)},
 \end{cases}
\]
then $e_i'\in\Hspace_S$ and we can put $f=e_i'$. Then
$|f(x)-f(xa)|^2=1$ for
\[
 x\in(T^{(i)}\mathbin{\dot+}T_{a^{-1}}^{(i)})
       -T^{(i)}T_{a^{-1}}^{(i)}
\]
and so our integral formula becomes
\[
 \int_{S_i}|\varphi(x)|^2\,dx=0,
 \qquad
 S_i=(T^{(i)}\mathbin{\dot+}T_{a^{-1}}^{(i)})
       -T^{(i)}T_{a^{-1}}^{(i)}.
\]

Put now
\[
 S'=\sum_{i=1}^\infty
 \bigl((T^{(i)}\mathbin{\dot+}T_{a^{-1}}^{(i)})
       -T^{(i)}T_{a^{-1}}^{(i)}\bigr).
\]
As the integrand above is $|\varphi(x)|^2\geq0$, adding of the above equations
over all $i=1,2,\ldots$ gives
\[
 \int_{S'}|\varphi(x)|^2\,dx=0.
\]
Now $x\in S-S'$ means
$x\notin(T^{(i)}\mathbin{\dot+}T_{a^{-1}}^{(i)})
-T^{(i)}T_{a^{-1}}^{(i)}$ for all $i=1,2,\ldots$; thus $x\in$ or $\notin$
both $T^{(i)}$ and $T_{a^{-1}}^{(i)}$, that is both $x$ and $xa$ are
$\in T^{(i)}$ or $\notin T^{(i)}$. By \defref{12.1.2}, (iv), this implies
$x=xa$. As $a\neq1$, \defref{12.1.5}, (iii) now necessitates
$\mu(S-S')=0$. Thus $\int_S|\varphi(x)|^2\,dx=0$ too, and therefore
$\varphi(x)=0$ (except for an $x$-set of measure $0$). This completes the
proof.

\Lemma{12.2.4}
\emph{Let $\mathfrak G$ be ergodic in $S$. Let $L$ be as in
\lemref{12.2.2} and let $U$ be the set of all $U_a$, $a\in\mathfrak G$. Then
$L\cdot U'=(\alpha1)$.}

\Proof $\alpha1$ belongs obviously to $U'$, and as
$\alpha1=L_\alpha$, $\alpha1\in L$ so $\alpha1\in L\cdot U'$,
$L\cdot U'\supset(\alpha1)$. So we need only prove
$L\cdot U'\subset(\alpha1)$. Assume therefore $A\in L\cdot U'$. Then
$A=L_{\varphi(x)}$ and as $A\in U'$ therefore
$U_aA=AU_a$, $A=U_a^{-1}AU_a$,
$L_{\varphi(x)}=L_{\varphi(xa^{-1})}$ for every $a\in\mathfrak G$. By
\lemref{12.2.3} this means $\varphi(x)=\varphi(xa^{-1})$ except for an
$x$-set of measure $0$ (depending on $a\in\mathfrak G$).

Put $T_\alpha=(x;\Re(\varphi(x))<\alpha)$. The above relation proves that
$x\in T_\alpha$ and $xa^{-1}\in T_\alpha$ are equivalent, except for an
$x$-set of measure $0$; so
\[
 \mu((T_\alpha\mathbin{\dot+}T_{\alpha a})
       -(T_\alpha T_{\alpha a}))=0.
\]
Therefore \defref{12.1.5}, (iv) (the ergodicity) implies
\[
 \mu(T_\alpha)=0\quad\text{or}\quad\mu(S-T_\alpha)=0.
\]

Denote the set of all (real) $\alpha$'s with $\mu(T_\alpha)=0$ by $N$. If
$\alpha\leq\beta$ then $T_\alpha\subset T_\beta$ so $\beta\in N$ implies
$\alpha\in N$. So if $\alpha_0$ is the least upper bound of $N$,
$-\infty\leq\alpha_0\leq\infty$, then $\alpha<\alpha_0$ implies
$\alpha\in N$, that is $\mu(T_\alpha)=0$, and $\alpha>\alpha_0$ implies
$\alpha\notin N$, that is $\mu(S-T_\alpha)=0$.

So for every rational $\alpha<\alpha_0$,
$\mu((x;\Re(\varphi(x))<\alpha))=0$ and adding over them gives
$\mu((x;\Re(\varphi(x))<\alpha_0))=0$; for every rational
$\alpha>\alpha_0$, $\mu((x;\Re(\varphi(x))\geq\alpha))=0$ and adding over
them gives: $\mu((x;\Re(\varphi(x))>\alpha_0))=0$. Adding again:
\[
 \mu((x;\Re(\varphi(x))\neq\alpha_0))=0.
\]
Similarly we obtain:
$\mu((x;\Im(\varphi(x))\neq\beta_0))=0$ for a suitable $\beta_0$, and if we
put $\alpha_1=\alpha_0+i\beta_0$ then adding these two last equations gives:
\[
 \mu((x;\varphi(x)\neq\alpha_1))=0.
\]
So we have $\varphi(x)=\alpha_1$ except for an $x$-set of measure $0$; and so
$A=L_{\varphi(x)}=\alpha_11$, $A\in(\alpha1)$. Thus
$L\cdot U'\subset(\alpha1)$, completing the proof.

\SectionStart{12.3}{\texorpdfstring{Construction of $\M,\M'$ in
$\Hspace_{S\mathfrak G}$}{Construction of M and M-prime in the product space}}
We pass now to constructions in $\Hspace_{S\mathfrak G}$.

\Lemma{12.3.1}
\emph{Let $\mathfrak G$ be an $m$-group in $S$. Define the operators}
\begin{enumerate}[label=(\roman*)]
\item $\overline U_{a_0}F(x,a)=F(xa_0,aa_0)$,
\item $\overline V_{a_0}F(x,a)=F(x,a_0^{-1}a)$,
\item $\overline WF(x,a)=F(xa^{-1},a^{-1})$,
\item $\overline L_{\varphi(x)}F(x,a)=\varphi(x)F(x,a)$,
\item $\overline M_{\varphi(x)}F(x,a)=\varphi(xa^{-1})F(x,a)$,
\end{enumerate}
\emph{for any $a_0\in\mathfrak G$ and any bounded and measurable complex
valued $\varphi(x)$ defined for all $x\in S$. These operators are bounded,
$\overline U_{a_0}$, $\overline V_{a_0}$, $\overline W$ are even unitary.
Furthermore $\overline W=\overline W^{-1}$ and}
\[
 \overline W\overline U_{a_0}\overline W=\overline V_{a_0};
 \qquad
 \overline W\overline L_{\varphi(x)}\overline W
 =\overline M_{\varphi(x)};
\]
\emph{in other words: $F\mapsto\overline WF$ is an involutory spatial
isomorphism of $\Hspace_{S\mathfrak G}$ which interchanges
$\overline U_{a_0},\overline L_{\varphi(x)}$ with
$\overline V_{a_0},\overline M_{\varphi(x)}$.}

\Proof The boundedness and unitarity are proved in the same way as in the
proof of \lemref{12.2.1}. $\overline W^2=1$ results from direct computation,
and implies $\overline W=\overline W^{-1}$; the two last equations result from
direct computation.

\Lemma{12.3.2}
\emph{Let $\mathfrak G$ be an $m$-group in $S$. Let $I$ be the set of all
$\overline U_{a_0}$ and all $\overline L_{\varphi(x)}$ and $J$ the set of all
$\overline V_{a_0}$ and all $\overline M_{\varphi(x)}$
(cf. \lemref{12.3.1}). Form for each bounded operator $\overline A$ in
$\Hspace_{S\mathfrak G}$ ($\overline A\in\B_{S\mathfrak G}$) the decomposition from
\defref{2.4.2},}
\[
 \overline A\sim\langle A_{a,b}\rangle_{a,b\in\mathfrak G}
\]
\emph{where every $A_{a,b}\in\B_S$. (Cf. the decomposition
$\Hspace_{S\mathfrak G}=\Hspace_{\mathfrak G}\otimes\Hspace_S$ from
\lemref{12.1.1}. As described there, we replace the indices
$t,s=1,2,\ldots$ by indices $a,b\in\mathfrak G$ as already the complete
normalised orthogonal set $\varphi_{a_0}$, $a_0\in\mathfrak G$ in
$\Hspace_{\mathfrak G}$ has been indexed in this way.)}

\emph{Then $\overline A\in I'$ if and only if $A_{a,b}$ has the form}
\[
 A_{a,b}=L_{\chi_{ab^{-1}}(xb^{-1})}
\]
\emph{and $\overline A\in J'$ if and only if $A_{a,b}$ has the form}
\[
 A_{a,b}=L_{\chi_{a^{-1}b}(x)}U_{b^{-1}a}.
\]
\emph{Here $\chi_c(x)$ must be a bounded and measurable function of $x$ for
every $c\in\mathfrak G$. (We do not determine for which systems
$\chi_c(x)$, $c\in\mathfrak G$, a bounded $\overline A$, to which the given
$\chi_c(x)$ corresponds, actually does exist.)}

\Proof Remember the definition of the $A_{a,b}$ (\defref{2.4.2}, with our
present notations):
\[
 \overline A\langle f(x)\varphi_{a'}(a);a\in\mathfrak G\rangle
 =\langle A_{a',b}f(x);b\in\mathfrak G\rangle;
\]
that is (if we replace $a',a$ by $a,b$ so that now $x\in S$ and
$b\in\mathfrak G$ are the variables, and $a\in\mathfrak G$ is a parameter):
\[
 \overline Af(x)\varphi_a(b)=A_{a,b}f(x).
\]
Then the definitions of $\overline U_{a_0}$,
$\overline L_{\varphi(x)}$ and $\overline W$ give: If
$\overline A=\langle A_{a,b}\rangle_{a,b\in\mathfrak G}$, then
\begin{align*}
 \overline U_{a_0}^{-1}\mathbin{\cdot}\overline A
 \mathbin{\cdot}\overline U_{a_0}
 &=\left\langle
 U_{a_0}^{-1}A_{aa_0^{-1},ba_0^{-1}}U_{a_0}
 \right\rangle_{a,b\in\mathfrak G},\\
 \overline W\overline A\overline W
 &=\left\langle U_b^{-1}A_{a^{-1},b^{-1}}U_a
 \right\rangle_{a,b\in\mathfrak G},\\
 \overline L_{\varphi(x)}\overline A
 &=\langle L_{\varphi(x)}A_{a,b}\rangle_{a,b\in\mathfrak G},\\
 \overline A\overline L_{\varphi(x)}
 &=\langle A_{a,b}L_{\varphi(x)}\rangle_{a,b\in\mathfrak G}.
\end{align*}

Now $\overline A\in I'$ means that $\overline A$ commutes with all
$\overline L_{\varphi(x)}$,
$((L_{\varphi(x)})^*=L_{\overline{\varphi(x)}})$ and all
$\overline U_{a_0}$,
$((\overline U_{a_0})^*=(\overline U_{a_0})^{-1}
=\overline U_{a_0^{-1}})$. This means $A_{a,b}\in L'$ that is
$A_{a,b}\in L$ (by \lemref{12.2.1}), and
\[
 U_{a_0}^{-1}A_{aa_0^{-1},ba_0^{-1}}U_{a_0}=A_{a,b}.
\]
The first statement means $A_{a,b}=L_{\omega_{a,b}(x)}$ where
$\omega_{a,b}(x)$ is a bounded and measurable function of $x$ for every
choice of $a,b\in\mathfrak G$. Now
\begin{align*}
 U_{a_0}^{-1}A_{aa_0^{-1},ba_0^{-1}}U_{a_0}
 &=U_{a_0}^{-1}L_{\omega_{aa_0^{-1},ba_0^{-1}}(x)}U_{a_0}\\
 &=L_{\omega_{aa_0^{-1},ba_0^{-1}}(xa_0^{-1})}
\end{align*}
so the remaining requirement is:
\[
 \omega_{aa_0^{-1},ba_0^{-1}}(xa_0^{-1})=\omega_{a,b}(x).
\]
This equation, as the following ones, is valid with the exception of $x$-sets
of measure $0$.

Put $a_0=b$ and $\chi_c(x)=\omega_{c,1}(x)$ then
$\omega_{a,b}(x)=\chi_{ab^{-1}}(xb^{-1})$ obtains. Conversely: If this
equation holds for a system of bounded measurable functions $\chi_c(x)$ then
the $\omega_{a,b}(x)$ are such too, and
\[
 \omega_{aa_0^{-1},ba_0^{-1}}(xa_0^{-1})=\omega_{a,b}(x).
\]
So the characterisation of the $\overline A\in I'$ is given by
$A_{a,b}=L_{\chi_{ab^{-1}}(xb^{-1})}$. This proves our statement concerning
$I'$.

As the spatial isomorphism $F\mapsto\overline WF$ of
$\Hspace_{S\mathfrak G}$ carries $I$ into $J$ it carries $I'$ into $J'$. So
$\overline B\in J'$ means $\overline B=\overline W\overline A\overline W$
for some $\overline A\in I'$. The above result concerning the
$\overline A\in I'$, together with our formula
$\overline W\overline A\overline W$ gives therefore:
\begin{align*}
 B_{a,b}
 &=U_b^{-1}A_{a^{-1},b^{-1}}U_a
 =U_b^{-1}L_{\chi_{a^{-1}b}(xb)}U_a\\
 &=L_{\chi_{a^{-1}b}(x)}U_b^{-1}U_a
 =L_{\chi_{a^{-1}b}(x)}U_{b^{-1}a}.
\end{align*}
This proves our statement concerning $J'$.

\Lemma{12.3.3}
\emph{Let $\mathfrak G$ be an $m$-group in $S$, and $I,J$ as above. Then
$R(I)=J'$, $R(J)=I'$.}

\Proof For $\overline A=\overline V_{a_0}$,
$A_{a,b}=\delta_{a_0ab^{-1}}\cdot1$ (define
$\delta_c=1$ for $c=1$, $=0$ for $c\neq1$), and for
$\overline A=\overline M_{\varphi(x)}$,
$A_{a,b}=\delta_{ab^{-1}}L_{\varphi(xb^{-1})}$. So
$\overline V_{a_0}$ and $\overline M_{\varphi(x)}\in I'$ (with
$\chi_c(x)=\delta_{a_0c}$ resp. $\delta_c\varphi(x)$), that is $J\subset I'$.
As $I'$ is a ring, this implies $R(J)\subset I'$.

If $\overline A\in I'$, $\overline B\in J'$ then we have
\[
 A_{a,b}=L_{\chi_{ab^{-1}}(xb^{-1})},
 \qquad
 B_{a,b}=L_{\xi_{a^{-1}b}(x)}U_{b^{-1}a}.
\]
Now put $\overline C=\overline A\overline B$,
$\overline D=\overline B\overline A$ then the formula (iv) of
\lemref{2.4.3} gives:
\begin{align*}
 C_{a,b}
 &=\sum_cA_{c,b}B_{a,c}\\
 &=\sum_cL_{\chi_{cb^{-1}}(xb^{-1})}
          L_{\xi_{a^{-1}c}(x)}U_{c^{-1}a}\\
 &=\sum_cL_{\chi_{ac^{-1}b^{-1}}(xb^{-1})\xi_{c^{-1}}(x)}U_c,
\end{align*}
(we replaced the summation index $c$ by $ac^{-1}$) and
\begin{align*}
 D_{a,b}
 &=\sum_cB_{c,b}A_{a,c}\\
 &=\sum_cL_{\xi_{c^{-1}b}(x)}U_{b^{-1}c}
          L_{\chi_{ac^{-1}}(xc^{-1})}\\
 &=\sum_cL_{\xi_{c^{-1}b}(x)\chi_{ac^{-1}}(xb^{-1})}U_{b^{-1}c}\\
 &=\sum_cL_{\chi_{ac^{-1}b^{-1}}(xb^{-1})\xi_{c^{-1}}(x)}U_c,
\end{align*}
(we replaced the summation index $c$ by $bc$, remember that the order of
summation in $\sum_c$ does not matter by \lemref{2.4.3} (iv)). So we have
$\overline C=\overline D$, $\overline A\overline B=\overline B\overline A$.
As $\overline B\in J'$ implies $\overline B^*\in J'$, this means
$\overline A\in J''$; as $\overline A$ was an arbitrary element of $I'$ we
have $I'\subset J''$. As $1\in J$, therefore $J''=R(J)$
(cf. \cite[p.~397]{ref18}), so $I'\subset R(J)$.

Thus we have proved $R(J)=I'$. The spatial isomorphism
$F\mapsto\overline WF$ of $\Hspace_{S\mathfrak G}$ interchanges $I$ with
$J$, therefore we have proved $R(I)=J'$ too. Thus the proof is complete.

\Lemma{12.3.4}
\emph{Let $\mathfrak G$ be an $m$-group in $S$. Put $\M=R(I)=J'$ then
$\M'=R(J)=I'$. If $\mathfrak G$ is ergodic in $S$, then $\M$ is a factor.}

\Proof Clearly $\M'=(R(I))'=I'$. If
$\overline A\in\M\cdot\M'$ then we have by \lemref{12.3.2}
\[
 A_{a,b}=L_{\chi_{ab^{-1}}(xb^{-1})}
 =L_{\xi_{a^{-1}b}(x)}U_{b^{-1}a}.
\]
So if $a\neq b$ that is $b^{-1}a\neq1$ then \lemref{12.2.3} gives
$\chi_{ab^{-1}}(xb^{-1})=\xi_{a^{-1}b}(x)=0$. That is:
$\chi_c(x)=\xi_c(x)=0$ if $c\neq1$. For $a=b$, that is, $b^{-1}a=1$
similarly $\chi_1(xb^{-1})=\xi_1(x)$ obtains. This gives for $b=1$,
$\chi_1(x)=\xi_1(x)$ and then in general
$\chi_1(xb^{-1})=\chi_1(x)$. In other words:
\[
 U_b^{-1}L_{\chi_1(x)}U_b=L_{\chi_1(x)}.
\]
So we found the following conditions: $\chi_c(x)=\xi_c(x)$ for all
$c\in\mathfrak G$, if $c\neq1$ then even $=0$; if $c=1$ then
$L_{\chi_1(x)}\in L\cdot U'$.

Now if $\mathfrak G$ is ergodic in $S$, then \lemref{12.2.4} gives
$L\cdot U'=(\alpha1)$ so we have $L_{\chi_1(x)}=\alpha1$,
$\chi_1(x)=\alpha$. This means $\chi_c(x)=\xi_c(x)=\alpha\delta_c$, that is
$\overline A=\alpha1$. Thus we have proved
$\M\cdot\M'=(\alpha1)$ that is: $\M$ is a factor.

\SectionStart{12.4}{\texorpdfstring{Derivation of $t(E)$, that is
$D_\M(\overline E)$ and $D_{\M'}(\overline E)$. Theorem XI}{Derivation of
t(E), that is the two dimensions of E-bar. Theorem XI}}
From now on we assume that $\mathfrak G$ is ergodic in $S$.

\Lemma{12.4.1}
\emph{Write for an $\overline A\in\M$ or $\in\M'$ that}
\[
 \overline A\cong[[\chi_c(x)]]_{x\in S,c\in\mathfrak G}
\]
\emph{if $\overline A\sim\langle A_{a,b}\rangle_{a,b\in\mathfrak G}$ and}
\[
 A_{a,b}=L_{\chi_{ab^{-1}}(x)}U_{b^{-1}a}\quad\text{resp.}\quad
 A_{a,b}=L_{\chi_{ab^{-1}}(xb^{-1})}
\]
\emph{(cf. \lemref{12.3.2}). If now}
\[
 \overline A\cong[[\chi_c(x)]],\qquad
 \overline B\cong[[\xi_c(x)]],\qquad
 \overline C\cong[[\eta_c(x)]]
\]
\emph{then the following rules of computation hold:}
\begin{enumerate}[label=(\roman*)]
\item \emph{If $\overline C=\alpha\overline A$ then
  $\eta_c(x)=\alpha\chi_c(x)$.}
\item \emph{If $\overline C=\overline A^*$ then
  $\eta_c(x)=\overline{\chi_{c^{-1}}(xc^{-1})}$.}
\item \emph{If $\overline C=\overline A+\overline B$ then
  $\eta_c(x)=\chi_c(x)+\xi_c(x)$.}
\item \emph{If $\overline C=\overline A\overline B$ then
  \[
    \eta_c(x)=\sum_a\chi_a(x)\xi_{ca^{-1}}(xa^{-1}).
  \]
  The $\sum_a$ converges ``en mesure'' (cf. the precise description of this
  notion at the end of the proof), irrespective of the order in which the
  $a\in\mathfrak G$ are gone through.}
\end{enumerate}

\Proof Consider first $\overline A\in\M'$. (i), (iii) follow immediately
from \lemref{2.4.3}, (i), (iii); (ii) follows from \lemref{2.4.3}, (ii), by
the following consideration:
\[
 C_{a,b}=L_{\eta_{ab^{-1}}(xb^{-1})},
 \qquad C_{a,b}=A_{b,a}^*
 =L_{\overline{\chi_{ba^{-1}}(xa^{-1})}},
\]
\[
 \eta_{ab^{-1}}(xb^{-1})
 =\overline{\chi_{ba^{-1}}(xa^{-1})},
 \qquad
 \eta_c(x)=\overline{\chi_{c^{-1}}(xc^{-1})}.
\]
Finally (iv) results from \lemref{2.4.3}, (iv):
\begin{align*}
 C_{a,b}&=L_{\eta_{ab^{-1}}(xb^{-1})},\\
 C_{a,b}&=\sum_cA_{c,b}B_{a,c}\\
 &=\sum_cL_{\chi_{cb^{-1}}(xb^{-1})}
           L_{\xi_{ac^{-1}}(xc^{-1})}\\
 &=\sum_cL_{\chi_{cb^{-1}}(xb^{-1})\xi_{ac^{-1}}(xc^{-1})},
\end{align*}
and putting $b=1$, and writing $c,a$ for $a,c$,
\[
 L_{\eta_c(x)}
 =\sum_aL_{\chi_a(x)\xi_{ca^{-1}}(xa^{-1})}
\]
is strongly convergent, irrespective of the order in which the
$a\in\mathfrak G$ are gone through.

Let now $T\subset S$ be a measurable set of finite measure, put
\[
 e_T(x)=\begin{cases}
  1,&x\in T,\\
  0,&x\notin T
 \end{cases}
\]
so that $e_T\in\Hspace_S$. Then if $\bar a_1,\bar a_2,\ldots$ is an
enumeration of $\mathfrak G$ we have
\[
 \left\lVert L_{\eta_c(x)}e_T
 -\sum_{j=1}^i
 L_{\chi_{\bar a_j}(x)
       \xi_{c\bar a_j^{-1}}(x\bar a_j^{-1})}e_T
 \right\rVert\longrightarrow0
 \quad\text{for }i\to\infty
\]
so
\[
 \int_S\left|\eta_c(x)e_T(x)
 -\sum_{j=1}^i\chi_{\bar a_j}(x)
 \xi_{c\bar a_j^{-1}}(x\bar a_j^{-1})e_T(x)\right|^2dx\longrightarrow0,
\]
\[
 \int_T\left|\eta_c(x)
 -\sum_{j=1}^i\chi_{\bar a_j}(x)
 \xi_{c\bar a_j^{-1}}(x\bar a_j^{-1})\right|^2dx\longrightarrow0.
\]
This implies for every $\epsilon>0$ that
\[
 \mu\left((x;\left|\eta_c(x)-
 \sum_{j=1}^i\chi_{\bar a_j}(x)
 \xi_{c\bar a_j^{-1}}(x\bar a_j^{-1})\right|\geq\epsilon)\cdot T\right)
 \longrightarrow0
\]
for $i\to\infty$. That is:
$\sum_{j=1}^\infty\chi_{\bar a_j}(x)
\xi_{c\bar a_j^{-1}}(x\bar a_j^{-1})$ converges ``en mesure'' to
$\eta_c(x)$. As $\bar a_1,\bar a_2,\ldots$ was an arbitrary enumeration of
$\mathfrak G$ this proves (iv). (Of course, if $\mathfrak G$ is finite, the
convergence considerations are unnecessary, the sum $\sum_a$ being simply
equal to $\eta_c(x)$.)

Consider now $\overline A\in\M$. The spatial isomorphism
$F\mapsto\overline WF$ in $\Hspace_{S\mathfrak G}$ carries $I$ into $J$ and
therefore $\M'$ into $\M$. So we have
$\overline A=\overline W\overline B\overline W$, $\overline B\in\M'$. This
$\overline B\cong[[\chi_c(x)]]_{x\in S,c\in\mathfrak G}$. Now the computation
at the end of the proof of \lemref{12.3.2} shows that then
$\overline A\cong[[\chi_c(x)]]_{x\in S,c\in\mathfrak G}$ with the same
$\chi_c(x)$. Therefore, the rules of computation established in $\M'$ carry
over immediately to $\M$.

After these preparations we are in the position to take the decisive step:

\Definition{12.4.1}
Assume $\overline A\in\M$ or $\M'$. Form the system of functions $\chi_c(x)$
with $\overline A\cong[[\chi_c(x)]]_{x\in S,c\in\mathfrak G}$
(cf. \lemref{12.4.1}). Consider the two following cases:
\begin{enumerate}[leftmargin=2.4em]
\item[$(\alpha)$] $\mu(S)$ is finite. Then, as $\chi_1(x)$ is bounded and measurable,
  $\int_S\chi_1(x)\,dx$ exists, and it is a finite complex number.
\item[$(\beta)$] $\chi_1(x)\geq0$ for all $x\in S$. Then
  $\int_S\chi_1(x)\,dx$ exists again, and it is a real number $\geq0$,
  $\leq\infty$.
\end{enumerate}
In both cases define
\[
 t(\overline A)=\int_S\chi_1(x)\,dx.
\]
In what follows immediately, we will really make use of case $(\beta)$ only;
case $(\alpha)$ is needed for later applications (\secref{15.4}).

\Lemma{12.4.2}
\emph{Assume $\overline A,\overline B\in\M$ or $\M'$. Then we have:}
\begin{enumerate}[label=(\roman*)]
\item $t(\alpha\overline A)=\alpha t(\overline A)$.
\item $t(\overline A^*)=\overline{t(\overline A)}$.
\item $t(\overline A+\overline B)=t(\overline A)+t(\overline B)$.
\end{enumerate}
\emph{These equations are to be so understood, that whenever case
$(\alpha)$ or $(\beta)$
holds for the right sides, it holds for the left sides too (for (i) with case
$(\beta)$ however, $\alpha\geq0$ is needed), and both are equal. Finally we
have (automatically with case $(\beta)$ on both sides):}
\begin{enumerate}[label=(\roman*),start=4]
\item $t(\overline A^*\overline A)
  =t(\overline A\overline A^*)\geq0$.
\end{enumerate}

\Proof (i)--(iii) follow directly from \lemref{12.4.1}, (i)--(iii) resp. In
order to prove (iv), put
\[
 \overline A^*\overline A
 \cong[[\omega_c(x)]]_{x\in S,c\in\mathfrak G},
 \qquad
 \overline A\overline A^*
 \cong[[\nu_c(x)]]_{x\in S,c\in\mathfrak G}.
\]
As $\overline A\cong[[\chi_c(x)]]_{x\in S,c\in\mathfrak G}$, we have
\[
 \overline A^*\cong
 [[\overline{\chi_{c^{-1}}(xc^{-1})}]]_{x\in S,c\in\mathfrak G}
\]
and
\begin{align*}
 \omega_c(x)
 &=\sum_a\chi_{ca^{-1}}(xa^{-1})
       \overline{\chi_{a^{-1}}(xa^{-1})}\\
 &=\sum_a\chi_{ca}(xa)\overline{\chi_a(xa)},\\
 \nu_c(x)
 &=\sum_a\chi_a(x)\overline{\chi_{ac^{-1}}(xc^{-1})}.
\end{align*}
(Use \lemref{12.4.1}, (ii) and (iv). The sums converge ``en mesure,''
irrespectively of the order.) Thus
\[
 \omega_1(x)=\sum_a|\chi_a(xa)|^2,
 \qquad
 \nu_1(x)=\sum_a|\chi_a(x)|^2.
\]
We have case $(\beta)$ for both, and
\begin{align*}
 t(\overline A^*\overline A)
 &=\int_S\omega_1(x)\,dx
 =\sum_a\int_S|\chi_a(xa)|^2\,dx\\
 &=\sum_a\int_S|\chi_a(x)|^2\,dx,\\
 t(\overline A\overline A^*)
 &=\int_S\nu_1(x)\,dx
 =\sum_a\int_S|\chi_a(x)|^2\,dx,
\end{align*}
proving our statements.

\Lemma{12.4.3}
\emph{If $\overline E$ runs over all projections in $\M$ or in $\M'$ then
$t(\overline E)$ is always defined (case $(\beta)$ holds), and it is a relative
dimension function for $\M$ resp. $\M'$. In this normalisation the $C$ of
\thmref{X} is $=1$.}

\emph{If $T\subset S$ is measurable, define}
\[
 e_T(x)=\begin{cases}
  1,&x\in T,\\
  0,&x\notin T.
 \end{cases}
\]
\emph{Then an $\overline E\cong
[[\delta_ce_T(x)]]_{x\in S,c\in\mathfrak G}$ exists, in $\M$ as well as in
$\M'$, it is a projection, and $t(\overline E)=\mu(T)$.}

\Proof As $\overline E^*\overline E=\overline E^2=\overline E$ for every
projection $\overline E$, therefore \lemref{12.4.2} (iv) secures case
$(\beta)$ for
$\overline E$ and $t(\overline E)\geq0$, $\leq\infty$. If $T\subset S$ is
measurable, then
\begin{align*}
 \overline L_{e_T(x)}
 &=\langle\delta_{ab^{-1}}L_{e_T(x)}\rangle_{a,b\in\mathfrak G}\\
 &=\langle L_{\delta_{ab^{-1}}e_T(x)}U_{b^{-1}a}
\rangle_{a,b\in\mathfrak G}
 \cong[[\delta_ce_T(x)]]_{x\in S,c\in\mathfrak G}
\end{align*}
in $\M$ and
\begin{align*}
 \overline M_{e_T(x)}
 &=\langle\delta_{ab^{-1}}L_{e_T(xa^{-1})}\rangle_{a,b\in\mathfrak G}\\
 &=\langle L_{\delta_{ab^{-1}}e_T(xb^{-1})}\rangle_{a,b\in\mathfrak G}
 \cong[[\delta_ce_T(x)]]_{x\in S,c\in\mathfrak G}
\end{align*}
in $\M'$. So the $\overline E$'s required at the end of our Lemma exist;
clearly $\overline E^*=\overline E$, $\overline E^2=\overline E$ so that
$\overline E$ is a projection; and
\[
 t(\overline E)=\int_Se_T(x)\,dx=\int_T1\,dx=\mu(T).
\]
Thus the last statement of our Lemma is true.

As we observed in the proof of \lemref{12.1.1}, there exists at least one
element of the sequence $T^{(1)},T^{(2)},\ldots$ from \defref{12.1.2}, (iv),
for which $\mu(T^{(i)})>0$, $<\infty$. So we have for its $\overline E$
(cf. above) $t(\overline E)>0$, $<\infty$.

Now \lemref{8.3.5} implies that $t(\overline E)$ is a relative dimension
function for $\M$ resp. $\M'$ if we can only establish for it the conditions
(ii), (iii) in \defref{8.2.1}. (iii) follows immediately from
\lemref{12.4.2}, (iii), because under those assumptions
$P_{[\mathfrak M,\mathfrak N]}=P_{\mathfrak M}+P_{\mathfrak N}$. (ii)
follows from \lemref{12.4.2}, (iv), because
$\overline E\sim\overline F$ $(\cdots\M)$ resp. $(\cdots\M')$ implies the
existence of a $U$ ($\in\M$ resp. $\in\M'$) with
$U^*U=\overline E$, $UU^*=\overline F$ (cf. \lemref{4.3.1}).

Our last task is to determine the $C$ of \thmref{X}. If
$\overline E\in\M$, $\overline E\cong
[[\chi_c(x)]]_{x\in S,c\in\mathfrak G}$ then we saw at the end of the proof
of \lemref{12.4.1} that $\overline W\overline E\overline W\in\M'$,
$\overline W\overline E\overline W\cong
[[\chi_c(x)]]_{x\in S,c\in\mathfrak G}$. Thus
$t(\overline E)=t(\overline W\overline E\overline W)$. Now, owing to the
character of the spatial isomorphism $F\mapsto\overline WF$ it carries every
$\mathfrak M_F^{\M'}$ into the corresponding $\mathfrak M_{\overline WF}^{\M}$
therefore
\[
 \overline W\overline E_F^{\M'}\overline W
 =\overline E_{\overline WF}^{\M}.
\]
This proves $t(\overline E_F^{\M'})=t(\overline E_{\overline WF}^{\M})$
and thus, if $\overline WF=\pm F$,
$t(\overline E_F^{\M'})=t(\overline E_F^{\M})$.

If besides $F\neq0$, then $\overline E_F^{\M'}\neq0$ and (as
$t(\overline E)$ is a relative weight function for $\M$),
$t(\overline E_F^{\M'})>0$. So the above relations imply, considering
$t(\overline E_F^{\M})=Ct(\overline E_F^{\M'})$, that $C=1$. So we need only
to find an $F\neq0$ with $\overline WF=\pm F$. Choose any $G\neq0$; as
$\overline W^2=1$ we have
$\overline W(G\pm\overline WG)=\pm(G\pm\overline WG)$ and as
\[
 G=\frac12\bigl((G+\overline WG)+(G-\overline WG)\bigr)
\]
therefore both $G\pm\overline WG$ cannot be $=0$. So either
$F=G+\overline WG$ or $F=G-\overline WG$ meets our requirements.

Thus all parts of our Lemma are proved.

We restate the results obtained thus far:

\MainTheorem{XI}
\emph{Let $\mathfrak G$ be a finite or countably infinite group
(\defref{12.1.1}), $S$ a space with a Lebesgue outer measure
(\hyperref[def:12.1.2]{Definitions 12.1.2} and
\hyperref[def:12.1.3]{12.1.3}). Let
$\Hspace_{\mathfrak G}$, $\Hspace_S$, $\Hspace_{S\mathfrak G}$ be the spaces
derived from them (\defref{12.1.4}). Assume that $\mathfrak G$ is ergodic in
$S$ (\defref{12.1.5}).}

\emph{Form the operators $\overline U_{a_0}$, $\overline V_{a_0}$,
$\overline W$, $\overline L_{\varphi(x)}$, $\overline M_{\varphi(x)}$ in
$\Hspace_{S\mathfrak G}$ (\lemref{12.3.1}) and with their help the rings
$\M,\M'$ (\lemref{12.3.4}).}

\emph{$F\mapsto\overline WF$ is an involutory isomorphism of
$\Hspace_{S\mathfrak G}$ which interchanges $\M$ with $\M'$.}

\emph{Certain systems of bounded and measurable $x$-functions $\chi_c(x)$,
$c\in\mathfrak G$ are in a one-to-one correspondence with all
$\overline A\in\M$ and at the same time with all $\overline A\in\M'$. We
denote both correspondences by}
\[
 \overline A\cong[[\chi_c(x)]]_{x\in S,c\in\mathfrak G}
\]
\emph{(\lemref{12.4.1}). Using these correspondences, the computation rules in
$\M$ as well as those in $\M'$ can be expressed in terms of the $\chi_c(x)$;
we get in both cases the same rules (\lemref{12.4.1}, (i)--(iv)).}

\emph{$\M,\M'$ are coupled factors. Relative dimension functions
$D_\M(\overline E)$, $D_{\M'}(\overline E)$ for $\M$ resp. $\M'$ ($\overline
E$ a projection $\in\M$ resp. $\in\M'$) can be defined as follows: Write}
\[
 \overline E\cong[[\chi_c(x)]]_{x\in S,c\in\mathfrak G}
\]
\emph{then}
\[
 D_\M(\overline E)\quad\text{resp.}\quad D_{\M'}(\overline E)
 =\int_S\chi_1(x)\,dx.
\]
\emph{(We have $\chi_1(x)\geq0$ for all $x\in S$, therefore the integral on
the right hand is always defined; it represents a real number
$\geq0$, $\leq\infty$).}

\emph{In this normalisation the $C$ of \thmref{X} is $=1$.}

\emph{If $T\subset S$ is measurable, and}
\[
 e_T(x)=\begin{cases}
  1,&x\in T,\\
  0,&x\notin T,
 \end{cases}
\]
\emph{then we have for both projections
$\overline E=\overline L_{e_T(x)}\in\M$ and
$\overline E'=\overline M_{e_T(x)}\in\M'$ this:
$D_\M(\overline E)$ resp. $D_{\M'}(\overline E')=\mu(T)$.}

\Needspace{8\baselineskip}
\ChapterHeading{13}{XIII}{\texorpdfstring{Properties of $\M,\M'$}
{Properties of M and M-prime}}

\SectionStart{13.1}{\texorpdfstring{Determination of the classes of $\M,\M'$}
{Determination of the classes of M and M-prime}}
The characterisation of $\M,\M'$ by \thmref{XI} makes the determination of
their classes easy. This determination is the object of the discussions which
follow. If we wanted to use these $\M,\M'$ solely to obtain an example of
case $(\mathrm{II})$, then it could be shortened considerably: the examples
\((\alpha)\)--\((\gamma)\) after \lemref{13.2.1} could be obtained in a few
lines, together with
\lemref{13.1.2} for that special case, while \lemref{13.1.1} would be entirely
unnecessary. But owing to the general interest of these $\M,\M'$ we prefer to
give an exhaustive discussion.

In what follows we use throughout the notations and the assumptions of
\thmref{XI}.

\Lemma{13.1.1}
\emph{Case $(\mathrm I)$ holds for $\M,\M'$ if and only if a point $x\in S$
with $\mu^*((x))>0$ exists. If this is the case, $S$ can be characterised,
after the omission of a subset of measure $0$, (which therefore is unessential),
as follows:}

\emph{Every one-point set $(x)$ in $S$ ($x\in S$) is measurable, and
$\mu((x))$ has a value independent of $x$: $\mu((x))=\epsilon>0$,
$<\infty$, $\mathfrak G$ is simply transitive in $S$, that is if $x_0$ is any
fixed element of $S$, then $a\rightleftarrows x_0a$ is a one to one mapping of
$\mathfrak G$ on $S$.}

\emph{If $\mathfrak G$ and $S$ have $n=1,2,\ldots,\infty$ elements, then
$\M,\M'$ are both in case $(\mathrm I_n)$. The $D_\M(E)$, $D_{\M'}(E)$ of
\thmref{XI} go over into the standard normalisation, if both are multiplied by
$1/\epsilon$.}

\Proof Assume first that case $(\mathrm I)$ holds. Then all values of
$D_\M(\overline E)$ are integer multiples of an $\epsilon_0>0$ and so all
$\mu(T)$, $T\subset S$ and measurable, are so. Now a $T$ with
$\mu(T)>0$, $<\infty$ exists (cf. the beginning of the proof of
\lemref{12.4.3}), therefore for one of these $T\subset S$, $\mu(T)$ assumes a
minimal value. Denote it by $T_0$. Thus for any measurable $T\subset T_0$ we
have, owing to $0\leq\mu(T)\leq\mu(T_0)$ either
$\mu(T_0)=\mu(T)$ or $\mu(T)=0$.

Now consider the $T^{(1)},T^{(2)},\ldots$ of \defref{12.1.2}, (iv), assuming
that $T^{(1)}\mathbin{\dot+}T^{(2)}\mathbin{\dot+}\cdots=S$ (cf. the remark
before \lemref{12.1.1}), and that $0$ is one of them (it can be added to their
system without altering its character). Consider an $x\in S$. Denote those
$n$'s for which $x\in T^{(n)}$ by $m_1,m_2,\ldots$ and those for which
$x\notin T^{(n)}$ by $n_1,n_2,\ldots$. Owing to \defref{12.1.2} (iv), we
have
\[
 (x)=T^{(m_1)}\cdot T^{(m_2)}\cdots
 (S-T^{(n_1)})\cdot(S-T^{(n_2)})\cdots.
\]
As $x\in T^{(1)}\mathbin{\dot+}T^{(2)}\mathbin{\dot+}\cdots$ the sequence
$m_1,m_2,\ldots$ is not empty; as $0$ is a $T^{(n)}$ the sequence
$n_1,n_2,\ldots$ is not empty either; by repeating $m$ or $n$ infinitely many
times we can secure that both sequences are infinite. So we see: For every
$x\in S$ there exist two infinite sequences $m_1,m_2,\ldots$ and
$n_1,n_2,\ldots$ so that
\[
 (x)=\prod_{i=1}^\infty T^{(m_i)}\cdot(S-T^{(n_i)}).
\]

Now assume $x\in T_0$ and $\mu^*((x))=0$. Put
\[
 U_j=T_0\cdot\prod_{i=1}^jT^{(m_i)}\cdot(S-T^{(n_i)}).
\]
Then we have $T_0\supset U_1\supset U_2\supset\cdots$,
$(x)=U_1U_2\cdots$. So $\lim_{j\to\infty}\mu(U_j)=\mu((x))=0$ (all $U_j$
are measurable along with $T_0$ and the $T^{(n)}$), and so a $j$ with
$\mu(U_j)<\mu(T_0)$ exists, (remember that $\mu(T_0)>0$). By our assumptions
on $T_0$ this implies $\mu(U_j)=0$.

Now consider all sets $U$ of this form:
\[
 U=T_0\cdot\prod_{i=1}^jT^{(m_i)}\cdot(S-T^{(n_i)}),
\]
$j=1,2,\ldots$; $m_1,\ldots,m_j,n_1,\ldots,n_j=1,2,\ldots$ for which
$\mu(U)=0$. Their set is countable: $U^{(1)},U^{(2)},\ldots$. So
$\mu(U^{(1)}\mathbin{\dot+}U^{(2)}\mathbin{\dot+}\cdots)=0$. Thus an
$x\in T_0$ with
$x\notin U^{(1)}\mathbin{\dot+}U^{(2)}\mathbin{\dot+}\cdots$ exists. By what
we proved above, this excludes $\mu^*((x))=0$. So an $x\in S$ with
$\mu^*((x))>0$ exists.

Assume now conversely the existence of an $x\in S$ with $\mu^*((x))>0$. As
\[
 (x)=\prod_{i=1}^\infty T^{(m_i)}(S-T^{(n_i)})
\]
(cf. above), $(x)$ is measurable, $\mu((x))>0$. Besides
$(x)\subset T^{(m_1)}$, $\mu((x))\leq\mu(T^{(m_1)})<\infty$. Put
$\epsilon=\mu((x))>0$, $<\infty$ then
$\mu((xa))=\mu((x))=\epsilon$ for every $a\in\mathfrak G$. Form now
$S_1=(xa;a\in\mathfrak G)$, $S_1$ is exactly invariant under every
transformation $x\mapsto xb$, $b\in\mathfrak G$. Thus the ergodicity of
$\mathfrak G$ in $S$ (cf. \defref{12.1.5}, (iv)) requires
$\mu(S_1)=0$ or $\mu(S-S_1)=0$. But
$\mu(S_1)\geq\mu((x))=\epsilon>0$ so $\mu(S-S_1)=0$. Omit the set $S-S_1$
from $S$, then we have $S_1=S$, that is $S=(xa;a\in\mathfrak G)$. So we
 have for every $x_0\in S$, $\mu((x_0))=\epsilon$.
 By \hyperref[def:12.1.5]{Def.~12.1.5} (iii),
$a\neq1$ implies $x_0\neq x_0a$, so $a\neq b$ implies $x_0a\neq x_0b$.
Thus all statements of the second section of our Lemma are established for
this case.

Choose an $x_0\in S$ and define
\[
 e_{(x_0)}(x)=\begin{cases}
  1,&x=x_0,\\
  0,&x\neq x_0.
 \end{cases}
\]
Then $\overline L_{e_{(x_0)}(x)}\in\M$, it is a projection, and
$D_{\M'}(\overline L_{e_{(x_0)}(x)})=\mu((x_0))=\epsilon$. We claim that
$\overline L_{e_{(x_0)}(x)}$ is minimal in $\M$. This is the case, if for every
$F$ with $\overline L_{e_{(x_0)}(x)}F=F\neq0$,
$E_F^{\M'}=\overline L_{e_{(x_0)}(x)}$ (cf. \lemref{5.1.2}, obviously
$\overline L_{e_{(x_0)}(x)}\neq0$). Now
$\overline L_{e_{(x_0)}(x)}F=F$ means
$e_{x_0}(x)F(x,a)=F(x,a)$ that is $F(x,a)=0$ if $x\neq x_0$. We must show:
One such $F$, say $F_0$, gives if it is $\neq0$ by application of operators
$\overline A\in\M'$ all others. As $F_0\neq0$ but $F_0(x,a)=0$ for
$x\neq x_0$ an $a_0\in\mathfrak G$ with $F_0(x_0,a_0)\neq0$ exists. Now
\[
 \frac1{\overline{F_0(x_0,a_0)}}
 \overline M_{e_{(x_0a_1^{-1})}(x)}\overline V_{a_1a_0^{-1}}\in\M'
\]
and it transforms $F_0(x,a)$ into
\[
 F_{a_1}'(x,a)=\begin{cases}
  1,&x=x_0, a=a_1,\\
  0,&\text{otherwise}
 \end{cases}
\]
and all desired $F(x,a)$ are linear aggregates of these.

Thus $\overline L_{e_{(x_0)}(x)}$ is minimal in $\M$ and therefore case
$(\mathrm I)$ holds for $\M,\M'$ (cf. \thmref{IV}). So the necessary and
sufficient character of our condition is established.

As the second section of our Lemma has already been verified, all we need to
consider is the last one. As $\overline L_{e_{(x_0)}(x)}$ is minimal in $\M$
and $D_{\M'}(\overline L_{e_{(x_0)}(x)})=\epsilon$ the standard
normalisation of $D_{\M'}(\overline E')$ obtains by multiplying it by
$1/\epsilon$. As the spatial isomorphism $F\mapsto\overline WF$ interchanges
$\M,D_\M(\overline E)$ with $\M',D_{\M'}(\overline E')$ the same is true for
$D_\M(\overline E)$.

Finally, if $S$ (that is $\mathfrak G$) has $n=1,2,\ldots,\infty$ elements,
say $x_1,\ldots,x_n$ (if $n=\infty$ omit $x_n$), then
\begin{align*}
 D_{\M'}(1)&=D_{\M'}(\overline L_1)
 =D_{\M'}(\overline L_{e_{(x_1)}(x)+\cdots+e_{(x_n)}(x)})\\
 &=D_{\M'}(\overline L_{e_{(x_1)}(x)}+\cdots+
              \overline L_{e_{(x_n)}(x)})\\
 &=D_{\M'}(\overline L_{e_{(x_1)}(x)})+\cdots+
   D_{\M'}(\overline L_{e_{(x_n)}(x)})=n\epsilon.
\end{align*}
In the standard normalisation this gives $n$, so we have case
$(\mathrm I_n)$ for $\M'$. The spatial isomorphism gives the same for $\M$.

\Lemma{13.1.2}
\emph{Case $(\mathrm{II})$ holds for $\M,\M'$ if and only if there is for
every $x\in S$, $\mu^*((x))=0$. If $\mu(S)$ is finite, then case
$(\mathrm{II}_1)$ holds for both $\M,\M'$ and the
$D_\M(\overline E)$, $D_{\M'}(\overline E')$ of \thmref{XI} go over into the
standard normalisation, if both are multiplied by $1/\mu(S)$. If $\mu(S)$ is
infinite, then case $(\mathrm{II}_\infty)$ holds for both $\M,\M'$.}

\Proof As an $E'\in\M'$ with $D_{\M'}(\overline E')>0$, $<\infty$ exists
(cf. the beginning of the proof of \lemref{12.4.3}), case $(\mathrm{III})$
cannot hold. So either case $(\mathrm I)$ or case $(\mathrm{II})$ holds, and
therefore \lemref{13.1.1} implies that our present condition is necessary and
sufficient for case $(\mathrm{II})$.

As
\[
 D_\M(1)=D_\M(\overline M_1)=D_\M(\overline M_{e_S(x)})=\mu(S)
\]
and
\[
 D_{\M'}(1)=D_{\M'}(\overline L_1)
 =D_{\M'}(\overline L_{e_S(x)})=\mu(S),
\]
the other statements of our Lemma follow immediately.

The two preceding Lemmas characterise $\M,\M'$ completely. We see (using
the terminology of \thmref{VIII}): The $\M,\M'$ as constructed in
\thmref{XI} are always of the same class; this class is finite or infinite if
$\mu(S)$ is finite resp. infinite; it is discrete or continuous if the measure
$\mu(T)$, $T\subset S$ is discrete resp. continuous (that is, if an $x\in S$
with $\mu^*((x))>0$ does resp. does not exist); it is never purely infinite.
And the difference between the discrete and the continuous cases corresponds
to that one between effective transitivity of $\mathfrak G$ in $S$, and
ergodicity without such transitivity.

\SectionStart{13.2}{Examples of case II}
Examples of case $(\mathrm I)$, based on
\lemref{13.1.1}, are so obvious that we need not be concerned with them.
(Owing to the one-to-one correspondence between $S$ and $\mathfrak G$ we may
even put, without any loss of generality, $S=\mathfrak G$. Then we can make
$\epsilon=1$ and then $\mu^*(T)$, $T\subset S$ becomes
$\mu^*(T)={}$ (number of elements in $T$).) We wish however to give effective
examples of case $(\mathrm{II})$, based on \lemref{13.1.2}. These examples
will be of a very special type, as we will only consider Abelian groups
$\mathfrak G$, and even those highly specialised.

\Lemma{13.2.1}
\emph{Let $S$ be either the set of all real numbers, $S_\infty$, or the set
of all real numbers $\geq0$, $<1$, which can be (and in what follows will be)
looked at as the set of all real numbers $(\bmod\,1)$, $S_1$. Use the common
Lebesgue measure in $S$. Let $\mathfrak G$ be a module of real numbers, that
is a finite or countably infinite set with these properties: $0\in\mathfrak G$,
$a,b\in\mathfrak G$ imply $-a\in\mathfrak G$, $a+b\in\mathfrak G$.
$\mathfrak G$ is a group, if we define unit${}=0$, $a^{-1}=-a$, $ab=a+b$.}

\emph{$\mathfrak G$ is an $m$-group in $S$, if we define $xa=x+a$. For
$S=S_1$ we insist that $1\in\mathfrak G$.}

\emph{Excepting $\mathfrak G=0$ and
$\mathfrak G=(\ldots,-2a_0,-a_0,0,a_0,2a_0,\ldots)$ (for a fixed $a_0>0$;
if $S=S_1$, then necessarily $a_0=1/n$, $n=1,2,\ldots$), $\mathfrak G$ is
always ergodic in $S$.}

\Proof $\mathfrak G$ satisfies obviously \defref{12.1.1}. $S$ satisfies
\defref{12.1.2} because it is a special case of Example b) given after it;
it is easy to verify \defref{12.1.5}, (i)--(iii), for $\mathfrak G$ and $S$,
that is the $m$-group character. It remains for us to decide about
\defref{12.1.5} (iv), that is ergodicity.

Every $T\subset S$ is invariant under $\mathfrak G=(0)$ and the set
\[
 T=(x;\ pa_0\leq x<(p+\tfrac12)a_0,
          \quad p=0,\pm1,\pm2,\ldots)\subset S
\]
(for $S=S_1$ take it $(\bmod\,1)$) is invariant under
\[
 \mathfrak G=(\ldots,-2a_0,-a_0,0,a_0,2a_0,\ldots).
\]
Thus these $\mathfrak G$ are not ergodic.

Assume now that $\mathfrak G$ is not of this form. If $S=S_\infty$ and some
$a_0\neq0$ is in $\mathfrak G$ we can map $S_\infty$ and $\mathfrak G$ by
the one-to-one transformation
$x\longleftrightarrow(1/a_0)x$, $a\longleftrightarrow(1/a_0)a$. This is an
isomorphism for everything that is important now and it carries $a_0$ into
$1$. For this reason we may assume $1\in\mathfrak G$ originally. If
$1\in\mathfrak G$, then all sets $T$ which are invariant under
$\mathfrak G$ can be looked at $(\bmod\,1)$; then we can do this for
$S=S_\infty$, that is we obtain $S=S_1$. So we need to discuss $S=S_1$ only.

Now $\mu(S)=\mu(S_1)=1$ is finite, so every
\[
 e_T(x)=\begin{cases}
  1,&x\in T,\\
  0,&x\notin T
 \end{cases}
\]
is $e_T\in\Hspace_S$. Ergodicity means: $U_a(e_T)=e_T$ for all
$a\in\mathfrak G$ means $e_T=0$ or $1$ or more generally:
$f\in\Hspace_S$, $U_af=f$ for all $a\in\mathfrak G$ means $f={}$ constant.

The $\varphi_n(x)=e^{2\pi i n x}$, $n=0,\pm1,\pm2,\ldots$ form a complete
normalised orthogonal set in $\Hspace_S$,
$U_a\varphi_n=e^{2\pi i n a}\varphi_n$.
Put
\[
 f=\sum_{n=-\infty}^{\infty}\alpha_n\varphi_n,
 \qquad \sum_{n=-\infty}^{\infty}|\alpha_n|^2\ \text{finite};
\]
this orthogonal expansion in the space $\Hspace_S$ is numerically the ``in the
mean'' convergent Fourier expansion of $f$, then
\[
 U_af=\sum_{n=-\infty}^{\infty}e^{2\pi i n a}\alpha_n\varphi_n.
\]
$U_af=f$ means $\alpha_n=e^{2\pi i n a}\alpha_n$, that is $\alpha_n=0$,
except if $na$ is an integer for every $a\in\mathfrak G$. If this happens
for any $n\neq0$ then it is the case for $|n|>0$ too, so we may assume
$n=1,2,\ldots$.

Denote the set of all $na$, $a\in\mathfrak G$ by $\mathfrak G'$.
$\mathfrak G'$ then consists of integers. $0\in\mathfrak G'$;
$p,q\in\mathfrak G'$ imply $-p,p+q\in\mathfrak G'$. Thus $\mathfrak G'$
consists of all multiples of a fixed $q=1,2,\ldots$. As
$1\in\mathfrak G$, $n\in\mathfrak G'$, $q$ is a divisor of $n$, say
$n=qm$. Thus
\[
 \mathfrak G=
 \left(\ldots,-\frac{2q}{n},-\frac qn,0,\frac qn,\frac{2q}{n},\ldots\right)
 =\left(\ldots,-\frac2m,-\frac1m,0,\frac1m,\frac2m,\ldots\right),
\]
contradicting our assumption concerning $\mathfrak G$.

So $\alpha_n=0$ whenever $n\neq0$ and therefore
$f=\alpha_0\varphi_0=\alpha_0={}$ constant. This completes the proof.

Now as $\mu(S_\infty)=\infty$ and $\mu(S_1)=1$ we see by
\lemref{13.1.2} that every ergodic $\mathfrak G$ gives an example of
$\M,\M'$ in class $(\mathrm{II}_\infty)$ resp. $(\mathrm{II}_1)$. Some
simple $\mathfrak G$' which are obviously ergodic by \lemref{13.2.1} are
these:
\phantomsection\label{examples:13.2}
\begin{enumerate}[leftmargin=2.4em]
 \item[$(\alpha)$] $\mathfrak G$ is the set $\mathfrak G_\theta$ of all $m+n\theta$,
 $m,n=0,\pm1,\pm2,\ldots$ where $\theta$ is any given irrational number.
 \item[$(\beta)$] $\mathfrak G$ is the set $\mathfrak G_{\rm rat.}$ of all rational
 numbers.
 \item[$(\gamma)$] $\mathfrak G$ is the set $\mathfrak G_{{\rm rat.},p}$ of all
 rational numbers of the form
 \[
  \frac m{p^n},\qquad m=0,\pm1,\pm2,\ldots,
  \qquad n=0,1,2,\ldots
 \]
 where $p$ is any given number $2,3,\ldots$ (not necessarily a prime!).
\end{enumerate}

\SectionStart{13.3}{Enumerations of the cases which exist. Theorem XII}
The above examples (with $S=S_\infty$) show that
factorisations $\M,\M'$ of the classes
$(\mathrm{II}_\infty),(\mathrm{II}_\infty)$ exist. Then
\lemref{11.4.3}, (ii) and the remarks after this Lemma show that the
combinations of classes
$(\mathrm{II}_1),(\mathrm{II}_\infty)$;
$(\mathrm{II}_\infty),(\mathrm{II}_1)$;
$(\mathrm{II}_1),(\mathrm{II}_1)$ exist too, and that in the last case the
$C$ of \thmref{X} can be prescribed arbitrarily. That any combination of
cases $(\mathrm I_m),(\mathrm I_n)$, $m,n=1,2,\ldots$, exists, has been
remarked at the end of \secref{8.6}. All other combinations, except
$(\mathrm{III}_\infty),(\mathrm{III}_\infty)$, have been excluded by
\thmref{X}. By \thmref{X} again
$(\mathrm{III}_\infty),(\mathrm{III}_\infty)$ exist if and only if factors
in case $(\mathrm{III})$ exist at all. So we have:

\MainTheorem{XII}
\emph{Factorisations $\M,\M'$ belonging to the following classes do exist:
$(\mathrm I_m),(\mathrm I_n)$, $m,n=1,2,\ldots,\infty$;
$(\mathrm{II}_m),(\mathrm{II}_n)$, $m,n=1,\infty$, and in case $m,n=1$
even the $C$ of \thmref{X} can be prescribed arbitrarily.
$(\mathrm{III}_\infty),(\mathrm{III}_\infty)$ exists if and only if factors
in case $(\mathrm{III})$ exist, a problem as yet unsolved. All other
combinations of classes do not exist.}

This answers \probref{2}, negatively (considering \lemref{8.6.1}), and
answers Problems \hyperref[prob:3]{3}, \hyperref[prob:4]{4} as completely
as it is possible at the
present state of our knowledge.

\SectionStart{13.4}{Non-coupled factors}
We can use our $\M,\M'$ to construct a factorisation in
which the factors $\M$ and $\N$ are not coupled ($\M\neq\N'$), and thus
answer \probref{1} negatively. This leads to the partial answer to
\probref{10} discussed in \secref{11.2}.

In what follows we use throughout the notations and the assumptions of
\thmref{XI}.

\Lemma{13.4.1}
\emph{Let $\mathfrak G^0$ be a proper subgroup of $\mathfrak G$ which is
ergodic in $S$. Every $\overline A\in\M'$ is
$\overline A\cong[[\chi_c(x)]]_{x\in S,\,c\in\mathfrak G}$, consider those
for which $\chi_c(x)=0$ whenever $c\notin\mathfrak G^0$. Denote their set by
$\N$. Then $\N$ is a ring, $\N\varsubsetneqq\M'$,
$\M'\mathbin{\cdot}\N'=(\alpha1)$.}

\Proof $\N\subset\M'$ is obvious. As $\overline V_{a_0}\in\M'$,
\[
 \overline V_{a_0}\cong
 [[\delta_{ca_0^{-1}}]]_{x\in S,\,c\in\mathfrak G},
\]
therefore $\overline V_{a_0}\in\N$ if and only if
$a_0\in\mathfrak G^0$, so $\mathfrak G^0\varsubsetneqq\mathfrak G$ implies
$\N\varsubsetneqq\M'$.

The computation rules of \lemref{12.4.1} show that $\overline A,\overline B\in\N$ imply
$\alpha\overline A,\overline A^*,\overline A+\overline B,
\overline A\overline B\in\N$. So we must only prove that $\N$ is weakly
closed in order to show that it is a ring. As $\M'$ is weakly closed (being
a ring) even relative weak closure of $\N$ in $\M'$ suffices.

Now an $\overline A\in\M'$ belongs to $\N$ if and only if
$a_0\notin\mathfrak G^0$ implies $\chi_{a_0}(x)=0$ for
$\overline A\cong[[\chi_c(x)]]_{x\in S,\,c\in\mathfrak G}$; that is if
$ab^{-1}\notin\mathfrak G^0$ implies $A_{a,b}=0$ for
$\overline A\sim\langle A_{a,b}\rangle_{a,b\in\mathfrak G}$. So we need
only to prove: If two $a,b\in\mathfrak G$ are given, then the set of all
$\overline A\sim\langle A_{a',b'}\rangle_{a',b'\in\mathfrak G}$ with
$A_{a,b}=0$ is weakly closed. Now
\[
 (A_{a,b}f,g)=(\overline A f^{[a]},g^{[b]}),
\]
where $h^{[c]}(x,d)=\delta_{c,d}h(x)$ ($h\in\Hspace_S$,
$h^{[c]}\in\Hspace_{S\mathfrak G}$) and so our condition is
$(\overline A f^{[a]},g^{[b]})=0$ for every pair $f,g\in\Hspace_S$. This
describes obviously a weakly closed set. Thus $\N$ is a ring and
$\N\varsubsetneqq\M'$.

Assume now $\overline A\in\M'\mathbin{\cdot}\N'$. As
$\overline M_{\chi(x)}\in\M'$,
\[
 \overline M_{\chi(x)}\cong
 [[\delta_c\chi(x)]]_{x\in S,\,c\in\mathfrak G}
\]
and as $c\notin\mathfrak G^0$ implies $c\neq1$, $\delta_c=0$ so we have
$\overline M_{\chi(x)}\in\N$. So $\overline A$ must commute with every
$\overline M_{\chi(x)}$ and (cf. above) with every
$\overline V_{a_0}$, $a_0\in\mathfrak G^0$. Now
$\overline A\in\M'$, $\overline A\cong
[[\xi_c(x)]]_{x\in S,\,c\in\mathfrak G}$ so our requirements are (by
\lemref{2.4.3}, (iv))
\[
 \xi_c(x)\chi(xc^{-1})=\xi_c(x)\chi(x),
 \qquad
 \xi_{a_0^{-1}ca_0}(xa_0)=\xi_c(x).
\]

The first equation reads
\[
 L_{\xi_c(x)}\chi=L_{\xi_c(x)}U_{c^{-1}}\chi
\]
for all bounded $\chi\in\Hspace_S$. Thus it is for all
\[
 \chi(x)=e_T(x)=\begin{cases}
  1,&x\in T,\\
  0,&x\notin T,
 \end{cases}
 \qquad T\subset S,\quad \mu(T)\ \text{finite};
\]
and for their linear aggregates and their condensation points; that is for
all $\chi\in\Hspace_S$ (cf. the second section of the proof of
\lemref{12.1.1}). So
$L_{\xi_c(x)}=L_{\xi_c(x)}U_{c^{-1}}$ which implies by
\lemref{12.2.3}, that $\xi_c(x)=0$ if $c\neq1$ (and so $c^{-1}\neq1$).

The second equation reads:
\[
 U_{a_0}^{-1}L_{\xi_1(x)}U_{a_0}=L_{\xi_1(x)}.
\]
Denote the set of all $U_{a_0}$, $a_0\in\mathfrak G^0$ by $\mathfrak U^0$.
Then we have $L_{\xi_1(x)}\in\mathfrak L\mathbin{\cdot}(\mathfrak U^0)'$.
As $\mathfrak G^0$ is ergodic in $S$, \lemref{12.2.4} applies:
$L_{\xi_1(x)}=\alpha1=L_\alpha$, $\xi_1(x)=\alpha$. So we have
$\xi_c(x)=\delta_c\alpha$, $\overline A=\alpha1$. This completes the proof
of $\M'\mathbin{\cdot}\N'=(\alpha1)$.

\Lemma{13.4.2}
\emph{Let $\mathfrak G^0,\N$ be as above. Then $\M,\N$ is a factorisation,
but $\M,\N$ are not coupled factors.}

\Proof As $\N\subset\M'$ and
\[
 \R(\M,\N)=(\R(\M,\N))''=(\M'\mathbin{\cdot}\N')'
 =(\alpha1)'=\B,
\]
therefore $\M,\N$ is a factorisation. As $\N\varsubsetneqq\M'$ they are
not coupled.

It is easy to give examples of groups $\mathfrak G,\mathfrak G^0$ as
required by Lemmas \hyperref[lem:13.4.1]{13.4.1} and
\hyperref[lem:13.4.2]{13.4.2}: We use the examples
$(\alpha)$--$(\gamma)$
at the end of \secref{13.1}. So
$\mathfrak G=\mathfrak G_\theta$,
$\mathfrak G^0=\mathfrak G_{k\theta}$, $\theta$ irrational,
$k=2,3,\ldots$, will do; or
$\mathfrak G=\mathfrak G_{\rm rat.}$,
$\mathfrak G^0=\mathfrak G^0_{{\rm rat.},q}$, $q=2,3,\ldots$ or
$\mathfrak G=\mathfrak G_{{\rm rat.},p}$,
$\mathfrak G^0=\mathfrak G_{{\rm rat.},q}$, $p,q=2,3,\ldots$, $q$ such a
divisor of $p$ no power of which is divisible by $p$ (that is: $q$ does not
contain all prime factors of $p$, for instance $p=6$, $q=2$).

Thus \lemref{13.4.2} implies that the $\M,\M'$ of \thmref{XI} are not
normal for $\mathfrak G=\mathfrak G_\theta,\mathfrak G_{\rm rat.},
\mathfrak G_{{\rm rat.},p}$ if $p$ is not the power of a prime. (For $\M'$
this results directly, for $\M$ it follows then from the spatial isomorphism
$F\mapsto\overline WF$ in $\Hspace_{S\mathfrak G}$ which interchanges $\M$ with
$\M'$.) So we have now the examples of not normal factors of class
$(\mathrm{II})$ we wanted. The question whether all factors of class
$(\mathrm{II})$ are not normal, remains nevertheless open; we can even not
decide whether $\M,\M'$ for
$\mathfrak G=\mathfrak G_{{\rm rat.},p}$, $p$ power of a prime, are normal.
(Cf. note at end.)

This is the extent to which we can answer \probref{10}.
% --- End of parts/05-part-iv.tex ---
\Needspace{10\baselineskip}
\PartHeadingTOC{V}{The finite cases}{The Finite Cases}

\ChapterHeading{14}{XIV}{\texorpdfstring{The Generalised additivity of
$D_{\M}(\mathfrak M)$}{The Generalised additivity of the dimension function}}

\SectionStart{14.1}{\texorpdfstring{Analogy between $(\mathrm I_n)$ and
$(\mathrm{II}_1)$}{Analogy between cases I-n and II-1}}
The object of the chapters which follow is the investigation of the finite
cases (cf. \thmref{VIII}), that is, of all factors \(\M\) which belong to
cases \((\mathrm I_n)\), \(n=1,2,\ldots\), or \((\mathrm{II}_1)\). The
outstanding fact is that they all behave very similarly. This is remarkable,
because an \(\M\) in a case \((\mathrm I_n)\) is ring-isomorphic to the ring
of all (bounded) operators of an \(n\)-dimensional Euclidean space and
therefore completely elementary; while an \(\M\) in case
\((\mathrm{II}_1)\) lies in a Hilbert space \(\Hspace\), unbounded operators
can be \(\mathbin\eta\M\), operators in \(\M\) can have continuous spectra,
etc.

All our results will be valid for all finite cases, that is, both for
\((\mathrm I_n)\) and \((\mathrm{II}_1)\). The really interesting
application is of course \((\mathrm{II}_1)\), while most of our results are
void or trivial or well known facts from linear geometry, when applied to
\((\mathrm I_n)\). Our reason for including the \((\mathrm I_n)\) too in all
our statements is only the desire to stress the analogy between
\((\mathrm I_n)\) and \((\mathrm{II}_1)\), and to put every result concerning
\((\mathrm{II}_1)\) at once in the right light by showing what it would mean
for \((\mathrm I_n)\).

Our notations will be these: \(\Hspace\) is a space; \(\M\) a factor in
\(\Hspace\) belonging to a finite case, \((\mathrm I_n)\),
\(n=1,2,\ldots\), or \((\mathrm{II}_1)\); \(D_{\M}(\mathfrak M)\) a relative
dimension function for \(\M\) with the normalisation
\(D_{\M}(\Hspace)=1\). For \((\mathrm I_n)\) this is the standard
normalisation multiplied by \(1/n\); for \((\mathrm{II}_1)\) it is the
standard normalisation itself. In the next \S\ (and only there) however these
assumptions will not be used, the class of \(\M\) and the normalisation of
\(D_{\M}(\Hspace)\) being arbitrary.

\SectionStart{14.2}{\texorpdfstring{Generalisation of the additivity formula for
$D_{\M}(\mathfrak M)$}{Generalisation of the additivity formula for the
dimension function}}
We discuss the behavior of \(D_{\M}(\mathfrak M)\) for such
\(\mathfrak M\)'s which need not to be orthogonal, nor comparable
\((\mathfrak M\subset\mathfrak N)\), thus generalising the statements of
\defref{8.2.1}.

\Lemma{14.2.1}
{\itshape Assume \(\mathfrak M,\mathfrak N\mathbin\eta\M\). Then
\[
 D_{\M}([\mathfrak M,\mathfrak N]-\mathfrak N)
 \leq D_{\M}(\mathfrak M).
\]}

\Proof This follows immediately from \lemref{7.3.4}.

\Lemma{14.2.2}
{\itshape Assume \(\mathfrak M,\mathfrak N\mathbin\eta\M\) and
\(D_{\M}(\mathfrak M)>D_{\M}(\mathfrak N)\). Then
\(\mathfrak M\cdot(\Hspace-\mathfrak N)\ne(0)\).}

\Proof We have, using \lemref{14.2.1},
\[
\begin{aligned}
D_{\M}([\mathfrak N,\Hspace-\mathfrak M]\cdot\mathfrak M)
 &=D_{\M}([\mathfrak N,\Hspace-\mathfrak M]
              -(\Hspace-\mathfrak M))\\
 &\leq D_{\M}(\mathfrak N)<D_{\M}(\mathfrak M),
\end{aligned}
\]
so \([\mathfrak N,\Hspace-\mathfrak M]\cdot\mathfrak M\ne\mathfrak M\).
But \([\mathfrak N,\Hspace-\mathfrak M]\cdot\mathfrak M\subset\mathfrak M\),
therefore
\(\mathfrak M-[\mathfrak N,\Hspace-\mathfrak M]\cdot\mathfrak M\ne(0)\).
Now
\[
\begin{aligned}
\mathfrak M-[\mathfrak N,\Hspace-\mathfrak M]\cdot\mathfrak M
 &=\mathfrak M\cdot(\Hspace-[\mathfrak N,\Hspace-\mathfrak M])\\
 &=\mathfrak M(\Hspace-\mathfrak N)
       (\Hspace-(\Hspace-\mathfrak M))\\
 &=\mathfrak M\cdot(\Hspace-\mathfrak N)\cdot\mathfrak M
 =\mathfrak M\cdot(\Hspace-\mathfrak N).
\end{aligned}
\]
So we proved \(\mathfrak M\cdot(\Hspace-\mathfrak N)\ne(0)\).

\Lemma{14.2.3}
{\itshape Assume \(\mathfrak M,\mathfrak N\mathbin\eta\M\). Then
\[
 D_{\M}([\mathfrak M,\mathfrak N])
 \leq D_{\M}(\mathfrak M)+D_{\M}(\mathfrak N),
\]
and equality holds if \(\mathfrak M\cdot\mathfrak N=(0)\).}

\Proof Put \(\mathfrak M'=[\mathfrak M,\mathfrak N]-\mathfrak N\). Then
\(\mathfrak M'\mathbin\eta\M\), \(\mathfrak M',\mathfrak N\) are
orthogonal, \([\mathfrak M',\mathfrak N]=[\mathfrak M,\mathfrak N]\), and
so \defref{8.2.1}(iii) applies. Thus
\[
\begin{aligned}
D_{\M}([\mathfrak M,\mathfrak N])
 &=D_{\M}([\mathfrak M',\mathfrak N])
  =D_{\M}(\mathfrak M')+D_{\M}(\mathfrak N)\\
 &\leq D_{\M}(\mathfrak M)+D_{\M}(\mathfrak N),
\end{aligned}
\]
considering \(D_{\M}(\mathfrak M')\leq D_{\M}(\mathfrak M)\) by
\lemref{14.2.1}.

We see that we have equality if
\(D_{\M}(\mathfrak M')=D_{\M}(\mathfrak M)\). If this is not the case, then
\(D_{\M}(\mathfrak M')<D_{\M}(\mathfrak M)\), and so \lemref{14.2.2}
applies: \(\mathfrak M\cdot(\Hspace-\mathfrak M')\ne(0)\). Now
\(f\in\mathfrak M\cdot(\Hspace-\mathfrak M')\) means
\(f\in\mathfrak M\) and the orthogonality of \(f\) to
\(\mathfrak M'=[\mathfrak M,\mathfrak N]-\mathfrak N\). But as
\(f\in\mathfrak M\subset[\mathfrak M,\mathfrak N]\), the second statement
means
\[
 f\in[\mathfrak M,\mathfrak N]
  -([\mathfrak M,\mathfrak N]-\mathfrak N)=\mathfrak N.
\]
So we have \(f\in\mathfrak M\cdot\mathfrak N\), and therefore the above
result states \(\mathfrak M\cdot\mathfrak N\ne(0)\). This completes the proof.

\Lemma{14.2.4}
{\itshape Assume \(\mathfrak M,\mathfrak N\mathbin\eta\M\). Then
\[
 D_{\M}([\mathfrak M,\mathfrak N])
 +D_{\M}(\mathfrak M\cdot\mathfrak N)
 =D_{\M}(\mathfrak M)+D_{\M}(\mathfrak N).
\]}

\Proof Put \(\mathfrak M_1=\mathfrak M-\mathfrak M\cdot\mathfrak N\).
Then \(\mathfrak M_1,\mathfrak M\cdot\mathfrak N\) are orthogonal,
\([\mathfrak M_1,\mathfrak M\cdot\mathfrak N]=\mathfrak M\), so by
\defref{8.2.1}(iii),
\[
D_{\M}(\mathfrak M)=D_{\M}(\mathfrak M_1)
 +D_{\M}(\mathfrak M\cdot\mathfrak N).
\]
Furthermore \([\mathfrak M_1,\mathfrak M\cdot\mathfrak N,\mathfrak N]
=[\mathfrak M,\mathfrak N]\), and if
\(f\in\mathfrak M_1\cdot\mathfrak N\), then
\(f\in\mathfrak M_1\cdot\mathfrak M\cdot\mathfrak N\), but
\(\mathfrak M_1\) is orthogonal to
\(\mathfrak M\cdot\mathfrak N\), so \(f=0\), which proves
\(\mathfrak M_1\cdot\mathfrak N=(0)\). So \lemref{14.2.3} applies:
\[
D_{\M}([\mathfrak M,\mathfrak N])
=D_{\M}([\mathfrak M_1,\mathfrak N])
=D_{\M}(\mathfrak M_1)+D_{\M}(\mathfrak N).
\]
Adding \(D(\mathfrak M\cdot\mathfrak N)\) to both sides gives,
considering our previous equation, the desired formula.

\Lemma{14.2.5}
{\itshape Let \(\mathfrak M_1,\mathfrak M_2,\ldots\) be a (finite or
infinite) sequence, all \(\mathfrak M_i\mathbin\eta\M\). Then
\[
D_{\M}([\mathfrak M_1,\mathfrak M_2,\ldots])
 \leq\sum_iD_{\M}(\mathfrak M_i).
\]
The equality holds if and only if one of these two conditions is satisfied:
\begin{enumerate}[label=(\roman*)]
\item \(D_{\M}([\mathfrak M_1,\mathfrak M_2,\ldots])\) is infinite.
\item \(D_{\M}([\mathfrak M_1,\mathfrak M_2,\ldots])\) is finite, and
\([\mathfrak M_1,\ldots,\mathfrak M_{i-1}]\cdot\mathfrak M_i=(0)\) for
\(i=2,3,\ldots\).
\end{enumerate}
In case (ii) there is even
\[
[\mathfrak M_{m_1},\mathfrak M_{m_2},\ldots]
 \cdot[\mathfrak M_{n_1},\mathfrak M_{n_2},\ldots]=(0)
\]
if \(m_1,m_2,\ldots\) and \(n_1,n_2,\ldots\) are any two sequences without
common elements.}

\Proof We can assume that the sequence
\(\mathfrak M_1,\mathfrak M_2,\ldots\) is infinite, as we could otherwise
insert infinitely many terms \(=(0)\).

\lemref{8.3.3} can be applied to
\[
\mathfrak M_1\subset[\mathfrak M_1,\mathfrak M_2]
\subset[\mathfrak M_1,\mathfrak M_2,\mathfrak M_3]\subset\cdots
\]
and it gives
\[
D_{\M}([\mathfrak M_1,\mathfrak M_2,\ldots])
=\lim_{i\to\infty}D_{\M}([\mathfrak M_1,\ldots,\mathfrak M_i]).
\]
Now \lemref{14.2.3} can be applied to
\([\mathfrak M_1,\ldots,\mathfrak M_{i-1}],\mathfrak M_i\), and it gives
\[
D_{\M}([\mathfrak M_1,\ldots,\mathfrak M_i])
\leq D_{\M}([\mathfrak M_1,\ldots,\mathfrak M_{i-1}])
 +D_{\M}(\mathfrak M_i),
\]
with an equality if (ii) holds; therefore
\[
D_{\M}([\mathfrak M_1,\ldots,\mathfrak M_i])
 \leq\sum_{j=1}^{i}D_{\M}(\mathfrak M_j),
\]
and passing to the limit
\[
D_{\M}([\mathfrak M_1,\mathfrak M_2,\ldots])
 \leq\sum_{i=1}^{\infty}D_{\M}(\mathfrak M_i),
\]
again with an equality if (ii) holds. So we have proved the relation with
 \(\leq\) in general, and the sufficience of (ii) for the equality. (i) too
implies equality, because if
\(D_{\M}([\mathfrak M_1,\mathfrak M_2,\ldots])\) is infinite, then
\[
\sum_{i=1}^{\infty}D_{\M}(\mathfrak M_i)
 \geq D([\mathfrak M_1,\mathfrak M_2,\ldots])
\]
is infinite too.

The necessity of (i) or (ii) means that if
\[
D_{\M}([\mathfrak M_1,\mathfrak M_2,\ldots])
=\sum_1^\infty D_{\M}(\mathfrak M_i)
\]
and is finite, then (ii) holds. We prove that in this case
\[
[\mathfrak M_{m_1},\mathfrak M_{m_2},\ldots]
 \cdot[\mathfrak M_{n_1},\mathfrak M_{n_2},\ldots]=(0)
\]
if \(m_1,m_2,\ldots\) and \(n_1,n_2,\ldots\) have no common elements.
This proves (ii) (put \(m_1=1,\ldots,m_{i-1}=i-1\) and \(n_1=i\)), and at
the same time it justifies the last statement of our Lemma too. As we can add
to the sequence \(n_1,n_2,\ldots\) all \(i=1,2,\ldots\) which differ from
\(m_1,m_2,\ldots\) and \(n_1,n_2,\ldots\), we may assume that
\(n_1,n_2,\ldots\) is the complementary set to \(m_1,m_2,\ldots\) in
\(1,2,\ldots\).

Now we have
\[
D_{\M}([\mathfrak M_{m_1},\mathfrak M_{m_2},\ldots])
 \leq\sum_iD_{\M}(\mathfrak M_{m_i}),\qquad
D_{\M}([\mathfrak M_{n_1},\mathfrak M_{n_2},\ldots])
 \leq\sum_iD_{\M}(\mathfrak M_{n_i}),
\]
and so by \lemref{14.2.4}
\[
\begin{aligned}
&D_{\M}([\mathfrak M_1,\mathfrak M_2,\ldots])
 +D_{\M}([\mathfrak M_{m_1},\mathfrak M_{m_2},\ldots]
   \cdot[\mathfrak M_{n_1},\mathfrak M_{n_2},\ldots])\\
&=D_{\M}([[\mathfrak M_{m_1},\mathfrak M_{m_2},\ldots],
            [\mathfrak M_{n_1},\mathfrak M_{n_2},\ldots]])
 +D_{\M}([\mathfrak M_{m_1},\mathfrak M_{m_2},\ldots]
   \cdot[\mathfrak M_{n_1},\mathfrak M_{n_2},\ldots])\\
&=D_{\M}([\mathfrak M_{m_1},\mathfrak M_{m_2},\ldots])
 +D_{\M}([\mathfrak M_{n_1},\mathfrak M_{n_2},\ldots])\\
&\leq\sum_iD_{\M}(\mathfrak M_{m_i})
 +\sum_iD_{\M}(\mathfrak M_{n_i})
 =\sum_iD_{\M}(\mathfrak M_i)
 =D_{\M}([\mathfrak M_1,\mathfrak M_2,\ldots]).
\end{aligned}
\]
As \(D_{\M}([\mathfrak M_1,\mathfrak M_2,\ldots])\) is finite by
assumption, this implies
\[
D_{\M}([\mathfrak M_{m_1},\mathfrak M_{m_2},\ldots]
 \cdot[\mathfrak M_{n_1},\mathfrak M_{n_2},\ldots])\leq0,
\]
so \(=0\), and therefore
\[
[\mathfrak M_{m_1},\mathfrak M_{m_2},\ldots]
 \cdot[\mathfrak M_{n_1},\mathfrak M_{n_2},\ldots]=(0).
\]
This completes the proof.

\Needspace{8\baselineskip}
\ChapterHeading{15}{XV}{The Relative Trace}

\SectionStart{15.1}{Definition of the relative trace}
From now on we make the assumption of \secref{14.1}: \(\M\) is a factor in
a finite case (\((\mathrm I_n)\), \(n=1,2,\ldots\), or
\((\mathrm{II}_1)\)), and \(D_{\M}(\mathfrak M)\) is normalised by
\(D_{\M}(\Hspace)=1\).

We define an analogue of the trace of common (finite dimensional) matrix
theory, more precisely: We define a notion which shares most essential
properties of the trace, and coincides with it in the discrete cases
\((\mathrm I_n)\), \(n=1,2,\ldots\). The importance of this notion lies of
course in its validity in the continuous case \((\mathrm{II}_1)\).

\Definition{15.1.1}
Let \(A\in\M\) be Hermitian. Form its resolution of unity \(E(\lambda)\),
\(-\infty<\lambda<\infty\). (Cf. \cite[p.~92]{ref16}. As \(A\) is bounded
this coincides with Hilbert's spectral form. The description which we will
use is discussed in \cite[pp.~389--390, footnote~42, and p.~418]{ref18}.) It
is characterised and uniquely determined, as discussed loc. cit., by the
following properties:
\begin{enumerate}[leftmargin=2.4em]
\item[$(\alpha)$] \(E(\lambda)\) is a projection, defined for all
  \(-\infty<\lambda<\infty\).
\item[$(\beta)$] \(\lambda\leq\mu\) implies \(E(\lambda)\) part of
  \(E(\mu)\).
\item[$(\gamma)$] \(\lambda\geq\lambda_0\), \(\lambda\to\lambda_0\) implies
  \(E(\lambda)\to E(\lambda_0)\) in the strong topology.
\item[$(\delta)$] There exists a \(c\) so that \(E(\lambda)=0\) if
  \(\lambda\leq-c\),
  and \(E(\lambda)=1\) if \(\lambda\geq c\).
\item[$(\epsilon)$] For all \(f,g\in\Hspace\),
  \[
  (Af,g)=\int_{-\infty}^{\infty}\lambda\,d(E(\lambda)f,g).
  \]
\end{enumerate}
(This is a numerical Stieltjes integral, certainly convergent, because we can
replace \(\int_{-\infty}^{\infty}\) by \(\int_{-c}^{c}\) owing to
\((\delta)\); and \((E(\lambda)f,g)\) is of bounded variation owing to
\((\beta)\), cf.
\cite[p.~93]{ref16}.)

This equation may be written symbolically as
\begin{equation*}
A=\int_{-\infty}^{\infty}\lambda\,dE(\lambda).
\tag{*}\label{eq:15.1-star}
\end{equation*}
As \(1\in\M\), therefore \(A\in\M\) has the consequence that all
\(E(\lambda)\in\M\) (cf. \cite[p.~390]{ref18}). Therefore all
\(D_{\M}(E(\lambda))\) are defined, and we can form the numerical Stieltjes
integral
\begin{equation*}
T_{\M}(A)=\int_{-\infty}^{\infty}\lambda\,dD_{\M}(E(\lambda)).
\tag{**}\label{eq:15.1-double-star}
\end{equation*}
(This integral is convergent, because we can replace
\(\int_{-\infty}^{\infty}\) by \(\int_{-c}^{c}\) owing to \((\delta)\), as
\(D_{\M}(E(\lambda))=0\) resp. \(=1\) if \(\lambda<-c\) resp.
\(\lambda\geq c\), and \(D_{\M}(E(\lambda))\) is monotone owing to
\((\beta)\).)

We call this \(T_{\M}(A)\), which is thus defined for all \(A\in\M\), the
relative trace of \(A\). It is clear how the symbolic equation
\eqref{eq:15.1-star} suggests the definition \eqref{eq:15.1-double-star}.
We shall see in \secref{15.2}, after \lemref{15.2.2}, that in the discrete
cases \((\mathrm I_n)\), \(n=1,2,\ldots\), our \(T_{\M}(A)\) is \(1/n\)
times the trace of \(A\) taken in the usual sense for matrices of degree
\(n\).

\SectionStart{15.2}{The ``Minimax principle.'' New expression for the trace}
An alternative definition of the trace can be obtained by using an extension
of the ``Minimax principle'' (cf. \cite[pp.~26--29]{ref5}). This extension is
as follows:

\Lemma{15.2.1}
{\itshape Let \(A\in\M\) be Hermitian, and \(E(\lambda)\),
\(-\infty<\lambda<\infty\), its resolution of unity. Form the quantity
\begin{equation*}
e(\alpha)=
\underset{\substack{\mathfrak M\mathbin\eta\M\\
                     D_{\M}(\mathfrak M)\geq\alpha}}{\mathrm{g.l.b.}}
\left\{
\underset{\substack{f\in\mathfrak M\\\lVert f\rVert=1}}
{\mathrm{l.u.b.}}(Af,f)\right\}
\tag{$\ddagger$}\label{eq:15.2-ddagger}
\end{equation*}
for every \(\alpha>0\), \(\leq1\). Then at the same time
\begin{equation*}
e(\alpha)=
\underset{[D_{\M}(E(\lambda))\geq\alpha]}{\mathrm{g.l.b.}}\lambda.
\tag{$\ddagger\ddagger$}\label{eq:15.2-double-ddagger}
\end{equation*}}

\Proof Consider the set of all \(\lambda\) with
\(D_{\M}(E(\lambda))\geq\alpha\). As it contains \(\lambda=c\) it is not
empty; as it does not contain \(\lambda=-c\) it is not the set of all real
numbers; it contains with \(\lambda\) every \(\lambda'>\lambda\) by
\defref{15.1.1} \((\beta)\); it contains its g.l.b. \(\lambda_0\) by
\defref{15.1.1} \((\gamma)\). Thus it is the set of all
\(\lambda\geq\lambda_0\). So we must prove \(e(\alpha)=\lambda_0\).

Define \(\mathfrak M_0\) by \(P_{\mathfrak M_0}=E(\lambda_0)\),
\(\mathfrak M_0\mathbin\eta\M\) because \(E(\lambda_0)\in\M\), and
\(D_{\M}(\mathfrak M_0)=D_{\M}(E(\lambda_0))\geq\alpha\). For
\(f\in\mathfrak M_0\) we have \(E(\lambda_0)f=f\), and so
\(E(\lambda)f=f\) if \(\lambda\geq\lambda_0\). Thus
\[
\begin{aligned}
(Af,f)&=\int_{-\infty}^{\infty}\lambda\,d(E(\lambda)f,f)
=\int_{-c}^{\lambda_0}\lambda\,d(E(\lambda)f,f)\\
&\leq\int_{-c}^{\lambda_0}\lambda_0\,d(E(\lambda)f,f)\\
&=\lambda_0\bigl((E(\lambda_0)f,f)-(E(-c)f,f)\bigr)
=\lambda_0(f,f)=\lambda_0\lVert f\rVert^2,
\end{aligned}
\]
and so
\(\mathrm{l.u.b.}_{[f\in\mathfrak M_0,\ \lVert f\rVert=1]}(Af,f)
\leq\lambda_0\). As \(\mathfrak M_0\mathbin\eta\M\),
\(D_{\M}(\mathfrak M_0)\geq\alpha\), this proves
\(e(\alpha)\leq\lambda_0\).

Assume now \(e(\alpha)<\lambda_0\). Then by
\eqref{eq:15.2-ddagger} a \(\mathfrak M\mathbin\eta\M\),
\(D_{\M}(\mathfrak M)\geq\alpha\), with
\[
\underset{\substack{f\in\mathfrak M\\\lVert f\rVert=1}}
{\mathrm{l.u.b.}}(Af,f)<\lambda_0
\]
must exist. Choose a
\[
\lambda'>
\underset{\substack{f\in\mathfrak M\\\lVert f\rVert=1}}
{\mathrm{l.u.b.}}(Af,f),\qquad \lambda'<\lambda_0.
\]
By the definition of \(\lambda_0\), the relation \(\lambda'<\lambda_0\)
has the consequence \(D_{\M}(E(\lambda'))<\alpha\). Define
\(\mathfrak M'\) by \(P_{\mathfrak M'}=E(\lambda')\); then
\(\mathfrak M'\mathbin\eta\M\) because \(E(\lambda')\in\M\). We have
\[
D_{\M}(\mathfrak M')=D_{\M}(E(\lambda'))<\alpha
\leq D_{\M}(\mathfrak M),
\]
so by \lemref{14.2.2}
\(\mathfrak M\cdot(\Hspace-\mathfrak M')\ne(0)\). Therefore an
\(f_0\in\mathfrak M\cdot(\Hspace-\mathfrak M')\), \(f_0\ne0\), exists.

Now as \(f_0\in\Hspace-\mathfrak M'\), we have \(E(\lambda')f_0=0\), and
so \(E(\lambda)f_0=0\) if \(\lambda\leq\lambda'\). Thus
\[
\begin{aligned}
(Af_0,f_0)
 &=\int_{-\infty}^{\infty}\lambda\,d(E(\lambda)f_0,f_0)
  =\int_{\lambda'}^{c}\lambda\,d(E(\lambda)f_0,f_0)\\
 &\geq\int_{\lambda'}^{c}\lambda'\,d(E(\lambda)f_0,f_0)\\
 &=\lambda'\bigl((E(c)f_0,f_0)-(E(\lambda')f_0,f_0)\bigr)
  =\lambda'(f_0,f_0)=\lambda'\lVert f_0\rVert^2.
\end{aligned}
\]
Multiplying \(f_0\) with \(1/\lVert f_0\rVert\), we can obtain
\(\lVert f_0\rVert=1\). But at the same time \(f_0\in\mathfrak M\) too,
therefore we have
\[
\underset{\substack{f\in\mathfrak M\\\lVert f\rVert=1}}
{\mathrm{l.u.b.}}(Af,f)\geq\lambda',
\]
contradicting our original assumptions on \(\lambda'\). Thus there must be
\(e(\alpha)=\lambda_0\), and the proof is completed.

\Lemma{15.2.2}
{\itshape Let \(A\in\M\) be Hermitian, and \(e(\alpha)\) as in
\lemref{15.2.1}. Then
\[
T_{\M}(A)=\int_0^1 e(\alpha)\,d\alpha.
\]}

\Proof We have
\(T_{\M}(A)=\int_{-\infty}^{\infty}\lambda\,
d(D_{\M}(E(\lambda)))\). It suffices to remember the definition of the
Riemann--Stieltjes integral in order to see that this coincides with the
Riemann integral
\[
\int_0^1
\left(\underset{[D_{\M}(E(\lambda))\geq\alpha]}
{\mathrm{g.l.b.}}\lambda\right)d\alpha
\qquad\text{(cf. \cite[p.~198]{ref19}).}
\]
But this is, by formula \eqref{eq:15.2-double-ddagger} of \lemref{15.2.1},
equal to \(\int_0^1e(\alpha)\,d\alpha\).

Observe that in the case \((\mathrm I_n)\), \(D_{\M}(\mathfrak M)\)
assumes no other values than
\(0,1/n,2/n,\allowbreak\ldots,(n-1)/n,1\); therefore \(e(\alpha)\) is constant in each
interval \((p-1)/n<\alpha\leq p/n\), \(p=1,\ldots,n\). Thus
\lemref{15.2.2} gives
\[
T_{\M}(A)=\int_0^1e(\alpha)\,d\alpha
=\frac1n\sum_{p=1}^n e\left(\frac pn\right),
\]
and one verifies easily that \(e(p/n)\) is the proper value No. \(p\) (from
below) of \(A\). Thus \(T_{\M}(A)\) is \(1/n\) times the trace of \(A\). A
more direct proof is given in \secref{15.4}, after \thmref{XIII}. In the case
\((\mathrm{II}_1)\), \(\int_0^1e(\alpha)\,d\alpha\) becomes an integral with
a really continuous domain. Owing to the analogy with the case
\((\mathrm I_n)\), it seems natural to introduce the following terminology:

\Definition{15.2.1}
Let \(A\in\M\) be Hermitian, and \(e(\alpha)\) as in \lemref{15.2.1}.
Then we call \(e(\alpha)\) the proper value No. \(\alpha\) (from below) of
\(A\), on the continuous scale \(0<\alpha\leq1\).

The immediate application of our new way to express \(T_{\M}(A)\) is this:

\Lemma{15.2.3}
{\itshape Let \(A,B\in\M\) be Hermitian, and \(A-B\) (semi-)definite. Then
\[
T_{\M}(A)\geq T_{\M}(B).
\]}

\Proof Form the \(e(\alpha)\) of \lemref{15.2.1} for \(A\) and for \(B\),
denoting it by \(e(\alpha)\) resp. \(\eta(\alpha)\). As the definiteness of
\(A-B\) means \(((A-B)f,f)\geq0\), \((Af,f)\geq(Bf,f)\) for every \(f\),
therefore \eqref{eq:15.2-ddagger} of \lemref{15.2.1} gives
\(e(\alpha)\geq\eta(\alpha)\). Now \lemref{15.2.2} gives
\(T_{\M}(A)\geq T_{\M}(B)\).

Note that \(T_{\M}(A)\geq0\) for definite \(A\) is obvious, and that this
would imply directly our Lemma, if we knew that
\(T_{\M}(A-B)=T_{\M}(A)-T_{\M}(B)\). But this relation is only established
at present if \(A,B\) commute (cf. \lemref{15.3.4} and the remarks after it),
and this makes the above simplification impracticable.

\SectionStart{15.3}{Formal properties of the trace}
We derive now the main formal properties of \(T_{\M}(A)\).

\Lemma{15.3.1}
{\itshape Let \(E\in\M\) be a projection. Then
\(T_{\M}(E)=D_{\M}(E)\).}

\Proof
\[
E(\lambda)=
\begin{cases}
E,&\lambda\geq1,\\
0,&\lambda<1
\end{cases}
\]
is \(E\)'s resolution of unity, as conditions
\((\alpha)\)--\((\epsilon)\) from
\defref{15.1.1} are easily verified for it. So
\[
D_{\M}(E(\lambda))=
\begin{cases}
D_{\M}(E),&\lambda\geq1,\\
0,&\lambda<1,
\end{cases}
\]
and so
\[
T_{\M}(E)=\int_{-\infty}^{\infty}\lambda\,
dD_{\M}(E(\lambda))=1\cdot D_{\M}(E)=D_{\M}(E).
\]

\Lemma{15.3.2}
{\itshape Let \(A\in\M\) be Hermitian, and \(a,b\) two real numbers. Then
\[
T_{\M}(aA+b1)=aT_{\M}(A)+b.
\]}

\Proof We first prove \(T_{\M}(-A)=-T_{\M}(A)\), because this will permit
us to restrict ourselves afterwards to the case \(a\geq0\).

Let \(E(\lambda)\) be \(A\)'s resolution of unity. Then
\(1-E(-\lambda)\) fulfills for \(-A\) all conditions
\((\alpha)\)--\((\epsilon)\) of \defref{15.1.1} except \((\gamma)\). The
function
\[
E_1(\lambda)=
\lim_{\substack{\lambda'>\lambda\\\lambda'\to\lambda}}
\bigl(1-E(-\lambda')\bigr)
\]
fulfills, as a simple argument shows, all of them. So \(E_1(\lambda)\) is
\(-A\)'s resolution of unity. Now \(E_1(\lambda)\ne1-E(-\lambda)\) only in
the points of discontinuity of \(E(-\lambda)\), that is in the points
\(-\lambda\) where \(\lambda\) runs over all point proper values of \(A\).
This is a countable set. A fortiori
\(D_{\M}(E_1(\lambda))\ne D_{\M}(1-E(-\lambda))\) holds only in a countable
set. Therefore we can, by the definition of Riemann--Stieltjes integrals,
replace \(d(D_{\M}(E_1(\lambda)))\) by
\(d(D_{\M}(1-E(-\lambda)))\) in every Riemann--Stieltjes integral in which
it occurs. Thus
\[
\begin{aligned}
T_{\M}(-A)
 &=\int_{-\infty}^{\infty}\lambda\,dD_{\M}(E_1(\lambda))
  =\int_{-\infty}^{\infty}\lambda\,dD_{\M}(1-E(-\lambda))\\
 &=\int_{-\infty}^{\infty}(-\lambda)\,dD_{\M}(E(-\lambda))
  =\int_{\infty}^{-\infty}\lambda\,dD_{\M}(E(\lambda))\\
 &=-\int_{-\infty}^{\infty}\lambda\,dD_{\M}(E(\lambda))
  =-T_{\M}(A).
\end{aligned}
\]

Assume now \(a\geq0\). Denote the \(e(\alpha)\) of \(A\) and of
\(aA+b1\) by \(e(\alpha)\) resp. \(\eta(\alpha)\). For
\(\lVert f\rVert=1\),
\[
((aA+b1)f,f)=a(Af,f)+b(f,f)=a(Af,f)+b,
\]
and as \(a\geq0\), this formula can be used when forming g.l.b.'s and
l.u.b.'s. Therefore \eqref{eq:15.2-ddagger} of \lemref{15.2.1} gives
\(\eta(\alpha)=ae(\alpha)+b\), and then \lemref{15.2.2} gives
\(T_{\M}(aA+b1)=aT_{\M}(A)+b\). This completes the proof.

\Lemma{15.3.3}
{\itshape \(T_{\M}(A)\) is a continuous function of \(A\) if the uniform
topology for operators (cf. \cite[p.~384]{ref18}) is used.}

\Proof We will show: If \(\lVert(A-B)f\rVert\leq\epsilon\lVert f\rVert\)
for every \(f\in\Hspace\), then
\[
|T_{\M}(A)-T_{\M}(B)|\leq\epsilon.
\]
Under the above assumption we have
\[
|((A-B)f,f)|\leq\lVert(A-B)f\rVert\,\lVert f\rVert
\leq\epsilon\lVert f\rVert^2,
\]
\[
(Af,f)\leq(Bf,f)+\epsilon\lVert f\rVert^2=((B+\epsilon1)f,f).
\]
So \((B+\epsilon1)-A\) is definite, and by \lemref{15.2.3},
\(T_{\M}(B+\epsilon1)\geq T_{\M}(A)\). But by \lemref{15.3.2},
\(T_{\M}(B+\epsilon1)=T_{\M}(B)+\epsilon\), so we have
\(T_{\M}(A)-T_{\M}(B)\leq\epsilon\). Interchanging \(A,B\) (the
assumption was symmetric in \(A,B\)), we obtain
\(T_{\M}(A)-T_{\M}(B)\geq-\epsilon\). So
\(|T_{\M}(A)-T_{\M}(B)|\leq\epsilon\), completing the proof.

We will prove elsewhere considerably more about the continuity of
\(T_{\M}(A)\). (Cf. note at the end of this paper.)

\Lemma{15.3.4}
{\itshape Let \(A,B\in\M\) be Hermitian, and commute with each other. Then
\[
T_{\M}(A+B)=T_{\M}(A)+T_{\M}(B).
\]}

\Proof Let \(E(\lambda)\) and \(F(\lambda)\) be the resolutions of unity
belonging to \(A\), resp. \(B\). By
\cite[pp.~390--391, in particular footnote~43]{ref18}, \(A,B\) are limits of
uniformly convergent sequences of linear aggregates of the \(E(\lambda)\),
resp. the \(F(\lambda)\). Therefore it suffices, by \lemref{15.3.3}, to prove
\[
T_{\M}\left(\sum_{i=1}^{m} a_iE(\lambda_i)
 +\sum_{j=1}^{n}b_jF(\mu_j)\right)
=T_{\M}\left(\sum_{i=1}^{m} a_iE(\lambda_i)\right)
 +T_{\M}\left(\sum_{j=1}^{n}b_jF(\mu_j)\right).
\]
Observe that if \(A,B\) commute, all \(E(\lambda),F(\mu)\) commute with
each other (cf. \cite[p.~115]{ref16}).

So it suffices to prove this:
\[
T_{\M}\left(\sum_{\rho=1}^p(c_\rho+d_\rho)E_\rho\right)
=T_{\M}\left(\sum_{\rho=1}^pc_\rho E_\rho\right)
 +T_{\M}\left(\sum_{\rho=1}^pd_\rho E_\rho\right)
\]
if \(E_1,\ldots,E_p\) are commutative projections. (Substitute
\(p=m+n\), \(E(\lambda_1),\ldots,E(\lambda_m),\allowbreak F(\mu_1),\ldots,F(\mu_n)\),
\(a_1,\ldots,a_m,0,\ldots,0\), and
\(0,\ldots,0,b_1,\ldots,b_n\) for
\(p,\allowbreak E_1,\ldots,E_p,\allowbreak c_1,\ldots,c_p,
\allowbreak d_1,\ldots,d_p\).)

This would follow immediately from
\begin{equation*}
 T_{\M}\left(\sum_{\rho=1}^pc_\rho E_\rho\right)
=\sum_{\rho=1}^pc_\rho D_{\M}(E_\rho)
\tag{\#}\label{eq:15.3-hash}
\end{equation*}
(for all choices of \(c_1,\ldots,c_p\), provided that
\(E_1,\ldots,E_p\) commute).

Consider the \(2^p\) terms arising when computing
\[
(E_1+(1-E_1))\cdots(E_p+(1-E_p)).
\]
They are products of commutative projections \(\in\M\), thus projections
\(\in\M\); their sum is \(1\), and the sum of those which arise from the term
\(E_\rho\) in the factor \(E_\rho+(1-E_\rho)\) is \(E_\rho\). Denote them
by \(E'_1,\ldots,E'_q\), \(q=2^p\); as
\(E'_1+\cdots+E'_q=1\), they are mutually orthogonal. Write
\[
E_\rho=E'_{\tau_{\rho,1}}+\cdots+E'_{\tau_{\rho,s_\rho}}.
\]
We claim: If \eqref{eq:15.3-hash} holds for \(E'_1,\ldots,E'_q\) (and all
\(c'_1,\ldots,c'_q\)), then it holds for \(E_1,\ldots,E_p\) too. Indeed:
\[
\begin{aligned}
T_{\M}\left(\sum_{\rho=1}^pc_\rho E_\rho\right)
 &=T_{\M}\left(\sum_{\rho=1}^pc_\rho
       \sum_{\sigma=1}^{s_\rho}E'_{\tau_{\rho,\sigma}}\right)\\
 &=T_{\M}\left(\sum_{\tau=1}^q
       \left(\sum_{\substack{\rho,\sigma\\
              \tau_{\rho,\sigma}=\tau}}c_\rho\right)E'_\tau\right)\\
 &=\sum_{\tau=1}^q
       \left(\sum_{\substack{\rho,\sigma\\
              \tau_{\rho,\sigma}=\tau}}c_\rho\right)D_{\M}(E'_\tau)\\
 &=\sum_{\rho=1}^pc_\rho
       \left(\sum_{\sigma=1}^{s_\rho}D_{\M}(E'_{\tau_{\rho,\sigma}})\right)
  =\sum_{\rho=1}^pc_\rho D_{\M}(E_\rho).
\end{aligned}
\]
In other words: We may assume in \eqref{eq:15.3-hash} that
\(E_1,\ldots,E_p\) are mutually orthogonal, and
\(E_1+\cdots+E_p=1\).

If now two \(c_\rho\)'s in \eqref{eq:15.3-hash} are equal, say
\(c_\rho=c_\sigma\) for some \(\rho\ne\sigma\), then \(E_\rho,E_\sigma\)
coalesce to \(E_\rho+E_\sigma\) in both sides of \eqref{eq:15.3-hash}.
So we may assume that all \(c_\rho\)'s are different. As a permutation does
not matter, we may even assume that \(c_1<c_2<\cdots<c_p\).

Define now
\[
G(\lambda)=
\begin{cases}
0,&\lambda<c_1,\\
E_1+\cdots+E_\rho,&c_\rho\leq\lambda<c_{\rho+1},
  \quad\rho=1,\ldots,p-1,\\
1,&\lambda\geq c_p.
\end{cases}
\]
One verifies immediately the conditions \((\alpha)\)--\((\epsilon)\) of
\defref{15.1.1} for
these \(G(\lambda)\) and \(\sum_{\rho=1}^pc_\rho E_\rho\); therefore it is
the resolution of unity which belongs to \(\sum_{\rho=1}^pc_\rho E_\rho\).
Now \defref{15.1.1} gives
\[
\begin{aligned}
T_{\M}\left(\sum_{\rho=1}^pc_\rho E_\rho\right)
 &=\int_{-\infty}^{\infty}\lambda\,dD_{\M}(G(\lambda))\\
 &=\sum_{\rho=1}^pc_\rho
 \bigl(D_{\M}(E_1+\cdots+E_\rho)
       -D_{\M}(E_1+\cdots+E_{\rho-1})\bigr)\\
 &=\sum_{\rho=1}^pc_\rho D_{\M}(E_\rho).
\end{aligned}
\]
This completes the proof.

Observe that we should expect from the analogy of the trace of finite
dimensional matrices, that the additivity of \(T_{\M}(A)\) expressed in
\lemref{15.3.4} holds even without assuming the commutativity of \(A,B\).
We will come back to this question at the end of \secref{15.4}, \probref{11}.

\SectionStart{15.4}{\texorpdfstring{Characterisation of the trace by inner
properties. Theorem XIII.\protect\linebreak Problem 11}{Characterisation of
the trace by inner properties. Theorem XIII. Problem 11}}
The relation \(\mathfrak M\sim\mathfrak N\) can be expressed in a new way,
owing to \(D_{\M}(\Hspace)=1\).

\Lemma{15.4.1}
{\itshape \(\mathfrak M\sim\mathfrak N\ (\mathbin{\eta}\M)\) is
equivalent to \(\mathfrak M,\mathfrak N\mathbin\eta\M\) and the existence of
a unitary \(V\in\M\) which maps \(\mathfrak M\) on \(\mathfrak N\).}

\Proof If \(\mathfrak M\mathbin\eta\M\) and such a \(V\in\M\) exists,
then \(U=VP_{\mathfrak M}\) is \(\in\M\), partially isometric, with the
initial and final sets \(\mathfrak M\), resp. \(\mathfrak N\), so
\(\mathfrak M\sim\mathfrak N\ (\mathbin\eta\M)\).

Conversely: If \(\mathfrak M\sim\mathfrak N\ (\mathbin\eta\M)\), then
\[
\begin{aligned}
D_{\M}(\Hspace-\mathfrak M)
 &=D_{\M}(\Hspace)-D_{\M}(\mathfrak M)\\
 &=D_{\M}(\Hspace)-D_{\M}(\mathfrak N)
  =D_{\M}(\Hspace-\mathfrak N)
\end{aligned}
\]
(this step is possible because \(D_{\M}(\Hspace)=1\) is finite), and so
\(\Hspace-\mathfrak M\sim\Hspace-\mathfrak N\). Let \(U_1,U_2\) be
 \(\in\M\), partially isometric, and have the resp. initial sets
 \(\mathfrak M,\Hspace-\mathfrak M\) and the resp. final sets
\(\mathfrak N,\Hspace-\mathfrak N\). One verifies easily that
\(V=U_1+U_2\in\M\) is unitary, and that it maps \(\mathfrak M\) on
\(\mathfrak N\). \(\mathfrak M\mathbin\eta\M\) is obvious.

This proof was based on the finiteness of \(D_{\M}(\Hspace)\), that is, of
\(\Hspace\) (with respect to \(\M\)). In fact the finiteness of \(\Hspace\)
is necessary and sufficient for the Lemma itself: Because if \(\Hspace\) is
infinite, then \(\Hspace\sim\mathfrak M\ (\mathbin\eta\M)\) for some
\(\mathfrak M\subsetneqq\Hspace\), and no unitary \(U\) maps \(\Hspace\) on
an \(\mathfrak M\subsetneqq\Hspace\).

We now state the essential invariance property of \(T_{\M}(A)\):

\Lemma{15.4.2}
{\itshape Let \(A\in\M\) be Hermitian and \(U\in\M\) be unitary. Then
\[
T_{\M}(U^{-1}AU)=T_{\M}(A).
\]}

\Proof If \(A\) has the resolution of unity \(E(\lambda)\), then
\(U^{-1}AU\) has clearly \(U^{-1}E(\lambda)U\). Put
\(E(\lambda)=P_{\mathfrak M(\lambda)}\),
\(U^{-1}E(\lambda)U=P_{\mathfrak N(\lambda)}\); then
\(\mathfrak N(\lambda)\) is \(\mathfrak M(\lambda)\)'s image by
\(U^{-1}\). So
\(\mathfrak M(\lambda)\sim\mathfrak N(\lambda)\ (\mathbin\eta\M)\),
\(D_{\M}(\mathfrak M(\lambda))=D_{\M}(\mathfrak N(\lambda))\),
\(D_{\M}(E(\lambda))=D_{\M}(U^{-1}E(\lambda)U)\). Now
\defref{15.1.1} gives \(T_{\M}(A)=T_{\M}(U^{-1}AU)\).

We are now able to characterise the relative trace by inner properties:

\Lemma{15.4.3}
{\itshape Assume that a (real and finite) numerical function \(T'(A)\) is
defined for all Hermitian \(A\in\M\), and that it has the following properties:
\begin{enumerate}[label=(\roman*)]
\item \(T'(1)=1\).
\item \(T'(aA)=aT'(A)\) if \(a\) is real.
\item \(T'(A+B)=T'(A)+T'(B)\) if \(A,B\) commute.
\item \(T'(A)\geq0\) if \(A\) is (semi-)definite.
\item \(T'(U^{-1}AU)=T'(A)\) if \(U\in\M\) is unitary.
\end{enumerate}
This is the case if and only if \(T'(A)\) is relative trace:
\(T'(A)=T_{\M}(A)\) for all Hermitian \(A\in\M\).}

\Proof \(T_{\M}(A)\) has the properties (i)--(v) by Lemmas
\hyperref[lem:15.3.2]{15.3.2}, \hyperref[lem:15.3.4]{15.3.4},
\hyperref[lem:15.2.3]{15.2.3}, and \hyperref[lem:15.4.2]{15.4.2}.
Therefore only the
converse problem remains: to determine \(T'(A)\) if it fulfills (i)--(v).
Assume therefore that such a \(T'(A)\) is given.

If \(E\) is a projection \(\in\M\), then, as \(E\) is definite, we have
\(T'(E)\geq0\) by (iv). If \(E\sim F\ (\mathbin\eta\M)\), then by
\lemref{15.4.1} a unitary \(U\in\M\) with \(F=U^{-1}EU\) exists, and so
by condition (v) \(T'(E)=T'(F)\). If \(E,F\) are orthogonal, that is
\(EF=FE=0\), then they commute, and therefore
\(T'(E+F)=T'(E)+T'(F)\). So the conditions (ii), (iii) of
\defref{8.2.1} are fulfilled, and as \(T'(1)=1>0\), \(<\infty\), therefore
\lemref{8.3.5} applies: \(T'(E)\) is a relative dimension function (with
respect to \(\M\)). So \(T'(E)=cD_{\M}(E)\) with a suitable constant \(c\);
but as \(T'(1)=D_{\M}(1)=1\) by (i), therefore \(c=1\), that is:
\(T'(E)=D_{\M}(E)\).

Thus if \(E_1,\ldots,E_p\) are mutually orthogonal projections,
\(E_1+\cdots+E_p=1\), and \(c_1<c_2<\cdots<c_p\), then we have: If
\(\rho\ne\sigma\), then \(E_\rho E_\sigma=E_\sigma E_\rho=0\), so
\(E_\rho,E_\sigma\) commute, and so by (ii), (iii)
\[
T'\left(\sum_{\rho=1}^pc_\rho E_\rho\right)
=\sum_{\rho=1}^pc_\rho T'(E_\rho)
=\sum_{\rho=1}^pc_\rho D_{\M}(E_\rho),
\]
and this, by the considerations at the end of the proof of \lemref{15.3.4},
is \(T_{\M}(\sum_{\rho=1}^pc_\rho E_\rho)\). So the desired equation
\(T'(A)=T_{\M}(A)\) holds for all \(A\)'s of the above form
\(\sum_{\rho=1}^pc_\rho E_\rho\).

(i)--(iii) imply in general \(T'(aA+b1)=aT'(A)+b\). (iii), (iv) imply
\(T'(A)\geq T'(B)\) if \(A-B\) is (semi-)definite and if \(A,B\) commute.
Thus the argument of the proof of \lemref{15.3.3} applies to \(T'(A)\) too,
if we restrict ourselves to some commutative set of operators: Then
\(T'(A)\) is a continuous function of \(A\), if the uniform topology for
operators is used. We know from \lemref{15.3.3} that this holds for
\(T_{\M}(A)\) too.

Now form the operators mentioned in
\cite[pp.~390--391, in particular footnote~43]{ref18}: They were denoted by
\[
A^{K\epsilon}=\sum_{\nu=1}^n\lambda_\nu
   \bigl(E(\lambda_\nu)-\bigl(E(\lambda_{\nu-1})\bigr)\bigr).
\]
As we pointed out there, they are uniformly convergent with the limit \(A\);
and as all \(E(\lambda)\) commute with each other and with \(A\), the same is
true for the \(A^{K\epsilon}\). Thus \(T'(A)=T_{\M}(A)\) follows, if we can
prove \(T'(A^{K\epsilon})=T_{\M}(A^{K\epsilon})\). But this equation holds,
because \(A^{K\epsilon}\) has the above discussed form. Put \(p=n\),
\(c_\nu=\lambda_\nu\),
\(E_\nu=E(\lambda_\nu)-E(\lambda_{\nu-1})\). Thus the proof is completed.

We restate the results obtained thus far:

\MainTheorem{XIII}
{\itshape Let \(\M\) be a factor of a finite class
(\((\mathrm I_n)\), or \((\mathrm{II}_1)\)). Then there exists one and only
one (real and finite) numerical function \(T'(A)\), defined for all Hermitian
\(A\in\M\), which has the following properties:
\begin{enumerate}[label=(\roman*)]
\item \(T'(1)=1\).
\item \(T'(aA)=aT'(A)\) if \(a\) is real.
\item \(T'(A+B)=T'(A)+T'(B)\) if \(A,B\) commute.
\item \(T'(A)\geq0\) if \(A\) is (semi-)definite.
\item \(T'(U^{-1}AU)=T'(A)\) if \(U\in\M\) is unitary.
\end{enumerate}
This unique \(T'(A)\) is the relative trace of \(A\), \(T_{\M}(A)\).
Further essential properties of this \(T'(A)\) are given in Lemmas
\hyperref[lem:15.2.3]{15.2.3}, \hyperref[lem:15.3.1]{15.3.1}, and
\hyperref[lem:15.3.3]{15.3.3}. Equivalent
definitions are contained in \defref{15.1.1} and in \lemref{15.2.2}.

In particular we have (this follows from \lemref{15.3.1}, so from (i)--(v)):
\begin{enumerate}[label=(\roman*),start=6]
\item For a projection \(E\in\M\), \(T'(E)=0\) implies \(E=0\).
\end{enumerate}}

In case \((\mathrm I_n)\) we obtain by this unicity theorem an easy
determination of \(T_{\M}(A)\): \(\M\) in case \((\mathrm I_n)\) means
\(\Hspace=\Hspace_1\otimes\Hspace_2\), \(\M=\B_1^{(1)}\),
 \(\Hspace_1\) has dimension \(n\). Now
 \hyperref[lem:2.4.6]{Lemma 2.4.6} show, that using the
symbolism
\[
A^0\sim\langle A_{t,s}\rangle_{t,s=1,\ldots,n}
\]
of \defref{2.4.2}, the \(A^0\in\M\) are characterised by
\[
A^0\sim\langle a_{t,s}1\rangle_{t,s=1,\ldots,n},
\]
and thus correspond to the numerical matrices
\(\{a_{t,s}\}_{t,s=1,\ldots,n}\). Now
\[
T'(A^0)=\frac1n\sum_{t=1}^na_{t,t}
\]
fulfills (i)--(v), as one verifies easily; therefore
\[
T_{\M}(A^0)=\frac1n\sum_{t=1}^na_{t,t}.
\]

The examples of case \((\mathrm{II}_1)\) given in \thmref{XI}, considering
\lemref{13.1.2}, can be discussed too. If \(\bar A\in\M\) or
\(\bar A\in\M'\), write
\(\bar A\simeq[\chi_c(x)]_{x\in S,c\in\mathfrak G}\), and form
\[
t(\bar A)=\int_S\chi_1(x)\,dx
\]
(cf. loc. cit. and \lemref{12.4.2}). Then
\(T'(\bar A)=(1/\mu(S))t(\bar A)\) fulfills (i)--(iii), as one verifies
easily with the help of \lemref{12.4.2}(i)--(iii).

As to (iv) observe, that every definite operator \(\bar A\) has the form
\(\bar B^2\), where \(\bar B=\alpha(\bar A)\),
\(\alpha(x)=|x|^{1/2}\). Thus \(\bar B\) belongs to our ring (\(\M\), resp.
\(\M'\)) and is Hermitian. A fortiori \(\bar A=\bar B^*\bar B\), and now
\(t(\bar A)=t(\bar B^*\bar B)\geq0\), \(T'(\bar A)\geq0\), follow from
\lemref{12.4.2}(iv).

(v) is proved as follows: As
\[
A=\tfrac12(A+A^*)+i\mathbin{\cdot}\tfrac1{2i}(A-A^*),
\]
it suffices to consider Hermitian operators \(A\); and as
\[
A=\left(\tfrac12(A+1)\right)^2
 -\left(\tfrac12(A-1)\right)^2,
\]
we may even assume \(A=B^2\), \(B\) Hermitian. So we must prove
\[
T'(U^{-1}B^2U)=T'(B^2),
\]
that is \(t(U^{-1}B^2U)=t(B^2)\). This coincides with
\lemref{12.4.2}(iv), if we put there \(A=BU\).

In all these cases, (iii), that is
\(T_{\M}(A+B)=T_{\M}(A)+T_{\M}(B)\), holds even without the restriction that
\(A,B\) commute. This raises the following problem:

\Problem{11}
{\itshape Is \(T_{\M}(A+B)=T_{\M}(A)+T_{\M}(B)\) always true, without
further restrictions on \(A,B\)?}

The answer is certainly yes in case \((\mathrm I_n)\), and for all known
examples of case \((\mathrm{II}_1)\). (Cf. also note at the end of this
paper.)

\SectionStart{15.5}{Characterisation of the finite cases by the traces.
Theorem XIV}
The statements of \thmref{XIII} concerning the unique determination of
\(T'(A)\) by the properties enumerated there can be formulated and discussed,
even if \(\M\) is not assumed to be a factor, but only a ring with
\(1\in\M\).

So if \(\M\) is a ring with \(1\in\M\), we may ask: When do the conditions
(i)--(vi) of \thmref{XIII} have a unique solution? This is the case if \(\M\)
is a factor and in a finite case. Let us now consider the converse problem.
Assume that \(\M\) is a ring with \(1\in\M\), and that (i)--(vi) have a
unique solution, say \(T^0(A)\).

Consider \(\M\cdot\M'\). This is a ring, and therefore it is the ring
generated by all projections \(E\in\M\cdot\M'\) (cf.
\cite[p.~392]{ref18}). Consider a projection \(E\in\M\cdot\M'\). If
\(A\in\M\), then \(A\) and \(E\) commute, as \(E\in\M'\), and so
\(AE=AEE=EAE\). Thus if \(A\) is Hermitian, resp. (semi-)definite, \(AE\)
is too. Besides \(A,E\in\M\), so \(AE\in\M\). Now take two constants
\(a_1,a_2>0\), and define
\[
T'(A)=a_1T^0(A)+a_2T^0(AE).
\]
Then \(T'(A)\) behaves with respect to the conditions (i)--(vi) of
\thmref{XIII} as follows: (i) holds if
\(a_1+a_2T^0(E)=1\). (ii) holds at any rate. (iii) holds, because
\(E\in\M'\) commutes with the \(A,B\in\M\), so if these \(A,B\) commute,
then \(AE,BE\) do too. (iv) holds, because \(AE\) is (semi-)definite along
with \(A\). (v) holds, because \(E\in\M'\) commutes with the \(U\in\M\),
so \((U^{-1}AU)E=U^{-1}(AE)U\). (vi) holds, because if \(F\) is a
projection, therefore (semi-)definite, then \(FE\) is it too; thus
\(T^0(F)\geq0\), \(T^0(FE)\geq0\), and \(T'(F)=0\) implies
\(T^0(F)=0\), \(F=0\).

Under these conditions our assumption necessitates \(T'(A)=T^0(A)\) for
each Hermitian \(A\in\M\). So in particular \(T'(E)=T^0(E)\), but the
definition gives \(T'(E)=(a_1+a_2)T^0(E)\). Therefore
\((a_1+a_2-1)T^0(E)=0\). Thus we have proved:
\[
a_1,a_2>0,\quad a_1+a_2T^0(E)=1
\quad\text{must imply}\quad (a_1+a_2-1)T^0(E)=0.
\]
This is clearly not the case if \(T^0(E)\ne0,1\); therefore we have either
\(T^0(E)=0\), \(E=0\) (by (vi)), or \(T^0(E)=1\),
\(T^0(1-E)=1-T^0(E)=0\), \(1-E=0\), \(E=1\) (by (i), (iii), (vi)).

Thus \(\M\cdot\M'=\mathbf R(0,1)=(\alpha1)\), and so \(\M\) is a factor.

Now the argument made at the beginning of the proof of \lemref{15.4.3}
shows, remembering \lemref{8.3.5}, that \(T^0(E)\) is a relative dimension
function with respect to \(\M\) (\(E\) runs over all projections
\(E\in\M\)); \(T^0(1)=1\) (by (i)) being finite, \(1\) is finite with
respect to \(\M\), and so is \(\Hspace\). Therefore \(\M\) is in one of the
finite cases: \((\mathrm I_n)\), \(n=1,2,\ldots\), or
\((\mathrm{II}_1)\).

Consequently we have proved this:

\MainTheorem{XIV}
{\itshape Let \(\M\) be a ring with \(1\in\M\). The conditions (i)--(vi)
of \thmref{XIII} admit a unique solution if and only if \(\M\) is a factor,
and belongs to one of the finite cases: \((\mathrm I_n)\),
\(n=1,2,\ldots\), or \((\mathrm{II}_1)\).}

Note that while (i)--(v) imply (vi) if \(\M\) is a factor, they may not
imply it if \(\M\) is merely a ring with \(1\in\M\), and therefore the
unicity of their solution may not guarantee the factor character of \(\M\).
This is the reason why (i)--(v) were sufficient in \thmref{XIII}, but
(i)--(vi) were needed in \thmref{XIV}. In fact there exist rings \(\M\) with
\(1\in\M\) for which (i)--(v) have a unique solution, and which nevertheless
are not factors. Easy considerations, which we will not detail here, show in
fact that the ring \(\M=(P_{[\varphi]})'\) (\(\Hspace\) any space of
\(\infty\) dimensions, \(\varphi\in\Hspace\), \(\lVert\varphi\rVert=1\)) is
such; the unique solution of (i)--(v) being
\(T^0(A)=(A\varphi,\varphi)\), and \(\M\) is not a factor, because
\(P_{[\varphi]}\in\M\cdot\M'\).

We have not attempted in this chapter to give anything like a complete
theory of the trace. We wanted to give only a brief summary of the most
characteristic features.

Many further problems could be handled with our present methods. For
instance, \lemref{15.3.3} could be strengthened so as to apply even to the
strong topology of operators (cf. \cite[p.~381]{ref18}); connections could
be established between \(T_{\M}(A)\) and the \((Af,f)\), etc. We propose to
come back to this subject, as well as to \probref{11}, in subsequent papers.

\ChapterHeading{16}{XVI}{Unbounded Operators}

\SectionStart{16.1}{\texorpdfstring{Solution numbers of $Xf=0$ and $X^*f=0$}
{Solution numbers of Xf=0 and X-star f=0}}
\(\M\) is assumed throughout this chapter too to be a factor in a finite
case (\((\mathrm I_n)\), \(n=1,2,\ldots\), or \((\mathrm{II}_1)\)), and
that \(D_{\M}(\mathfrak M)\) is normalised by \(D_{\M}(\Hspace)=1\).

The following Lemma is a consequence of \lemref{6.2.1}, valid in the finite
cases only:

\Lemma{16.1.1}
{\itshape If \(X\) is a linear, closed operator with an everywhere dense
domain, and \(X\mathbin\eta\M\), then \((f;Xf=0)\) and \((f;X^*f=0)\) are
linear closed sets \(\mathbin\eta\M\), and
\[
D_{\M}((f;Xf=0))=D_{\M}((f;X^*f=0)).
\]}

\Proof \(X^*f=0\) means \((X^*f,g)=(f,Xg)=0\) for all
\(g\in\operatorname{Domain}X\), that is, that \(f\) is orthogonal to
\(\operatorname{Range}X\). So
\((f;X^*f=0)=\Hspace-[\operatorname{Range}X]\), therefore
\((f;X^*f=0)\mathbin\eta\M\), and
\[
D_{\M}((f;X^*f=0))=D_{\M}(\Hspace)-D_{\M}([\operatorname{Range}X])
=1-D_{\M}([\operatorname{Range}X]).
\]
Replacing \(X\) by \(X^*\), we see that
\((f;Xf=0)\mathbin\eta\M\), and
\[
D_{\M}((f;Xf=0))=1-D_{\M}([\operatorname{Range}X^*]).
\]
Now \lemref{6.2.1} states
\[
D_{\M}([\operatorname{Range}X])=D_{\M}([\operatorname{Range}X^*]),
\]
therefore
\[
D_{\M}((f;Xf=0))=D_{\M}((f;X^*f=0)).
\]

In the cases \((\mathrm I_n)\), \(n=1,2,\ldots\), this is the statement
that the number of linearly independent solutions of a linear equation
\(Xf=0\) coincides with the one for the adjoint equation \(X^*f=0\). We see
now, that owing to our notion of relative dimension, it is true in the case
\((\mathrm{II}_1)\) too. Observe, that these are the only cases where it
holds: In any infinite case \(\Hspace\) is infinite, so
\(\Hspace\sim\mathfrak M\subsetneqq\Hspace\), and if \(U\in\M\) is the
partially isometric operator which maps \(\Hspace\) on \(\mathfrak M\), then
\[
(f;Uf=0)=(0),\qquad(f;U^*f=0)=\Hspace-\mathfrak M\ne(0).
\]

\SectionStart{16.2}{Essential density. Main properties}
The key to the deductions which follow is contained in this definition:

\Definition{16.2.1}
Let \(\mathfrak A\) be an arbitrary linear set in \(\Hspace\)
(\(\mathfrak A\) is not necessarily closed, and not necessarily
\(\mathbin\eta\M\)). If a sequence
\(\mathfrak M_1,\mathfrak M_2,\ldots\) of linear closed sets exists which has
the following properties:
\begin{enumerate}[label=(\roman*)]
\item \(\mathfrak M_i\mathbin\eta\M\).
\item \(\mathfrak M_1\subset\mathfrak M_2\subset\cdots\subset\mathfrak A\).
\item \([\mathfrak M_1,\mathfrak M_2,\ldots]=\Hspace\).
\end{enumerate}
then \(\mathfrak A\) is called essentially dense.

We prove:

\Lemma{16.2.1}
{\itshape Every essentially dense \(\mathfrak A\) is everywhere dense too.}

\Proof As all \(\mathfrak M_i\subset\mathfrak A\), and as \(\mathfrak A\)
is a linear set,
\[
\mathfrak M_1,\mathfrak M_2,\ldots\subset\mathfrak A\subset\Hspace,
\qquad
\Hspace=[\mathfrak M_1,\mathfrak M_2,\ldots]
\subset[\mathfrak A]\subset\Hspace,
\]
\([\mathfrak A]=\Hspace\). As \(\mathfrak A\) is linear,
\([\mathfrak A]\) consists of its condensation points only, therefore
\(\mathfrak A\) is everywhere dense.

\Lemma{16.2.2}
{\itshape If \(\mathfrak A_1,\mathfrak A_2,\ldots\) is a (finite or
infinite) sequence of essentially dense sets, then
\(\mathfrak A_1\cdot\mathfrak A_2\cdots\) is essentially dense too.}

\Proof We can assume that the sequence
\(\mathfrak A_1,\mathfrak A_2,\ldots\) is infinite, as we could otherwise
insert infinitely many terms \(\mathfrak A_i=\Hspace\).

Form the sequences from \defref{16.2.1}:
\(\mathfrak M_{i,1},\mathfrak M_{i,2},\ldots\) for \(\mathfrak A_i\),
\(i=1,2,\ldots\). By \lemref{8.3.3},
\[
\lim_{p\to\infty}D_{\M}(\mathfrak M_{i,p})
=D_{\M}([\mathfrak M_{i,1},\mathfrak M_{i,2},\ldots])
=D_{\M}(\Hspace)=1,
\]
and so we can form a sequence \(\rho_{i,1}<\rho_{i,2}<\cdots\), with
\[
D_{\M}(\mathfrak M_{i,\rho_{i,\nu}})>1-\frac1{2^{i+\nu}},
\qquad
D_{\M}(\Hspace-\mathfrak M_{i,\rho_{i,\nu}})
=D_{\M}(\Hspace)-D_{\M}(\mathfrak M_{i,\rho_{i,\nu}})
<\frac1{2^{i+\nu}}.
\]

Form
\(\overline{\mathfrak M}_\nu
=\mathfrak M_{1,\rho_{1,\nu}}\cdot
 \mathfrak M_{2,\rho_{2,\nu}}\cdots\).
Clearly \(\overline{\mathfrak M}_\nu\) is a closed linear set, and
\(\overline{\mathfrak M}_\nu\mathbin\eta\M\). As
\(\rho_{i,\nu}<\rho_{i,\nu+1}\), therefore
\(\mathfrak M_{i,\rho_{i,\nu}}\subset
\mathfrak M_{i,\rho_{i,\nu+1}}\), and so
\(\overline{\mathfrak M}_\nu\subset
\overline{\mathfrak M}_{\nu+1}\); as
\(\mathfrak M_{i,\rho_{i,\nu}}\subset\mathfrak A_i\), so
\(\overline{\mathfrak M}_\nu\subset
\mathfrak A_1\cdot\mathfrak A_2\cdots\). Finally we have
\[
\begin{aligned}
D_{\M}(\Hspace-\overline{\mathfrak M}_\nu)
 &=D_{\M}([\Hspace-\mathfrak M_{1,\rho_{1,\nu}},
            \Hspace-\mathfrak M_{2,\rho_{2,\nu}},\ldots])\\
 &\leq\sum_{i=1}^{\infty}
      D_{\M}(\Hspace-\mathfrak M_{i,\rho_{i,\nu}})
 <\sum_{i=1}^{\infty}\frac1{2^{i+\nu}}=\frac1{2^\nu}
\end{aligned}
\]
by \lemref{14.2.5}, and so
\[
D_{\M}(\Hspace-[\overline{\mathfrak M}_1,
\overline{\mathfrak M}_2,\ldots])
\leq D_{\M}(\Hspace-\overline{\mathfrak M}_\nu)<\frac1{2^\nu}.
\]
This holds for every \(\nu=1,2,\ldots\), therefore
\[
D_{\M}(\Hspace-[\overline{\mathfrak M}_1,
\overline{\mathfrak M}_2,\ldots])=0,
\quad
\Hspace-[\overline{\mathfrak M}_1,
\overline{\mathfrak M}_2,\ldots]=(0),
\quad
[\overline{\mathfrak M}_1,\overline{\mathfrak M}_2,\ldots]=\Hspace.
\]
Thus the sequence
\(\overline{\mathfrak M}_1,\overline{\mathfrak M}_2,\ldots\) fulfills all
conditions (i)--(iii) of \defref{16.2.1} for the linear set
\(\mathfrak A_1\cdot\mathfrak A_2\cdots\); therefore
\(\mathfrak A_1\cdot\mathfrak A_2\cdots\) is essentially dense.

\Lemma{16.2.3}
{\itshape Let \(\mathfrak A\) be an essentially dense set, and \(X\) a
linear closed operator with an everywhere dense domain, and
\(X\mathbin\eta\M\). Then the set
\[
(f;f\in\operatorname{Domain}X,\ Xf\in\mathfrak A)
\]
is essentially dense.}

\Proof Let \(W,B\) be the canonical decomposition of \(X\): \(B\) self
adjoint and definite, \(W\) partially isometric, \(W\in\M\),
\(B\mathbin\eta\M\), and \(X=WB\) (cf. \defref{4.4.1} and
\lemref{4.4.1}).

Let \(E(\lambda)\), \(-\infty<\lambda<\infty\), be the resolution of unity
which belongs to \(B\). This means, as \(B\) is not necessarily bounded,
that \(E(\lambda)\) fulfills the conditions \((\alpha)\)--\((\gamma)\) of
\defref{15.1.1} unaltered, but instead of
\((\delta)\)--\((\epsilon)\) these modifications:

\noindent\((\delta')\) If \(\lambda\to-\infty\), resp.
\(+\infty\), then \(E(\lambda)\to0\), resp. \(1\), in the sense of strong
operator convergence.

\noindent\((\epsilon')\) \(f\in\operatorname{Domain}B\) if and only if
\(\int_{-\infty}^{\infty}\lambda^2\,d\lVert E(\lambda)f\rVert^2\) is
finite, and then
\[
(Bf,g)=\int_{-\infty}^{\infty}\lambda\,d(E(\lambda)f,g)
\qquad\text{(cf. \cite[p.~92]{ref15}).}
\]
The (semi-)definiteness of \(B\) means that \(E(\lambda)=0\) for
\(\lambda<0\) (cf. \cite[p.~302]{ref21}). Put
\(E(\lambda)=P_{\mathfrak M(\lambda)}\).

As \(B\mathbin\eta\M\), \(B\) is invariant under every unitary
transformation \(U'\in\M'\), and so is \(E(\lambda)\); thus
\(E(\lambda)\in\M\), \(\mathfrak M(\lambda)\mathbin\eta\M\), and
\(E(\lambda)\) and \(\mathfrak M(\lambda)\) obviously reduce \(B\). If
\(f\in\mathfrak M(\lambda)\), then \(E(\lambda)f=f\), so
\[
E(\lambda')f=
\begin{cases}
f,&\lambda'\geq\lambda,\\
0,&\lambda'<0.
\end{cases}
\]
Therefore \((\delta')\) shows that \(Bf\) is defined, and \((\epsilon')\) that
\(0\leq(Bf,f)\leq\lambda\lVert f\rVert^2\). Thus the part of \(B\) in
\(\mathfrak M(\lambda)\) is bounded and everywhere defined (in
\(\mathfrak M(\lambda)\)), or: \(BE(\lambda)\) is bounded and everywhere
defined (in \(\Hspace\)).

We know that \(\mathfrak M(1)\subset\mathfrak M(2)\subset\cdots\), and as
\(\lim_{i\to\infty}E(i)=1\), we have
\([\mathfrak M(1),\mathfrak M(2),\allowbreak\ldots]=\Hspace\). Thus by
\lemref{8.3.3},
\[
\begin{aligned}
\lim_{i\to\infty}D_{\M}(\mathfrak M(i))
 &=D_{\M}(\Hspace)=1,\\
\lim_{i\to\infty}D_{\M}(\Hspace-\mathfrak M(i))
 &=\lim_{i\to\infty}
   \bigl(D_{\M}(\Hspace)-D_{\M}(\mathfrak M(i))\bigr)=0.
\end{aligned}
\]

Choose next the sequence \(\mathfrak M_1,\mathfrak M_2,\ldots\) from
\defref{16.2.1} for \(\mathfrak A\). Then
\(\mathfrak M_1\subset\mathfrak M_2\subset\cdots\subset\mathfrak A\),
\([\mathfrak M_1,\mathfrak M_2,\ldots]=\Hspace\), therefore as above,
\[
\lim_{i\to\infty}D_{\M}(\Hspace-\mathfrak M_i)=0.
\]

Form now the set
\[
\mathfrak M^{(i)}=(f;f\in\mathfrak M(i),\ WBE(i)f\in\mathfrak M_i).
\]
As \(BE(i)\) is everywhere defined and bounded, \(\mathfrak M^{(i)}\) is
linear and closed. As \(f\in\mathfrak M(i)\) implies
\(WBE(i)f=WBf=Xf\), we have
\[
\mathfrak M^{(i)}
\subset(f;f\in\operatorname{Domain}X,\ Xf\in\mathfrak M_i)
\subset(f;f\in\operatorname{Domain}X,\ Xf\in\mathfrak A).
\]
Finally \(f\in\mathfrak M^{(i)}\) implies
\(f\in\mathfrak M(i+1)\), and \(E(i+1)f=f\), and so
\(\mathfrak M^{(i)}\subset\mathfrak M^{(i+1)}\). Thus we have
\[
\mathfrak M^{(1)}\subset\mathfrak M^{(2)}\subset\cdots
\subset(f;f\in\operatorname{Domain}X,\ Xf\in\mathfrak A).
\]
The definition of \(\mathfrak M^{(i)}\) makes it obvious that it is invariant
under all unitary transformations \(U'\in\M'\), therefore
\(\mathfrak M^{(i)}\mathbin\eta\M\).

Now
\[
\begin{aligned}
\mathfrak M^{(i)}
 &=(f;f\in\mathfrak M(i),\ WBE(i)f\in\mathfrak M_i)\\
 &=\mathfrak M(i)\cdot(f;WBE(i)f\in\mathfrak M_i)\\
 &=\mathfrak M(i)\cdot
   (f;P_{\Hspace-\mathfrak M_i}WBE(i)f=0).
\end{aligned}
\]
Now by \lemref{16.1.1},
\[
\begin{aligned}
D_{\M}((f;P_{\Hspace-\mathfrak M_i}WBE(i)f=0))
 &=D_{\M}((f;(P_{\Hspace-\mathfrak M_i}WBE(i))^*f=0))\\
 &=D_{\M}((f;(BE(i))^*W^*P_{\Hspace-\mathfrak M_i}f=0)),
\end{aligned}
\]
and
\[
(f;(BE(i))^*W^*P_{\Hspace-\mathfrak M_i}f=0)
\supset(f;P_{\Hspace-\mathfrak M_i}f=0)=\mathfrak M_i.
\]
So
\[
D_{\M}((f;P_{\Hspace-\mathfrak M_i}WBE(i)f=0))
\geq D_{\M}(\mathfrak M_i),
\]
\[
D_{\M}(\Hspace-(f;P_{\Hspace-\mathfrak M_i}WBE(i)f=0))
\leq D_{\M}(\Hspace-\mathfrak M_i).
\]
Now we obtain, using \lemref{14.2.3}, or
\hyperref[lem:14.2.5]{14.2.5}, this:
\[
\begin{aligned}
D_{\M}(\Hspace-\mathfrak M^{(i)})
 &=D_{\M}([\Hspace-\mathfrak M(i),
   \Hspace-(f;P_{\Hspace-\mathfrak M_i}WBE(i)f=0)])\\
 &\leq D_{\M}(\Hspace-\mathfrak M(i))
 +D_{\M}(\Hspace-(f;P_{\Hspace-\mathfrak M_i}WBE(i)f=0))\\
 &\leq D_{\M}(\Hspace-\mathfrak M(i))
 +D_{\M}(\Hspace-\mathfrak M_i).
\end{aligned}
\]
Thus
\[
\lim_{i\to\infty}D_{\M}(\Hspace-\mathfrak M(i))=0,
\qquad
\lim_{i\to\infty}D_{\M}(\Hspace-\mathfrak M_i)=0
\]
give
\[
\lim_{i\to\infty}D_{\M}(\Hspace-\mathfrak M^{(i)})=0.
\]
By \lemref{8.3.4},
\[
\begin{aligned}
D_{\M}(\Hspace-[\mathfrak M^{(1)},\mathfrak M^{(2)},\ldots])
 &=D_{\M}((\Hspace-\mathfrak M^{(1)})
          \cdot(\Hspace-\mathfrak M^{(2)})
          \cdot(\Hspace-\mathfrak M^{(3)})\cdots)\\
 &=\lim_{i\to\infty}D_{\M}(\Hspace-\mathfrak M^{(i)})=0,
\end{aligned}
\]
\[
\Hspace-[\mathfrak M^{(1)},\mathfrak M^{(2)},\ldots]=(0),
\qquad
[\mathfrak M^{(1)},\mathfrak M^{(2)},\ldots]=\Hspace.
\]
Thus \(\mathfrak M^{(1)},\mathfrak M^{(2)},\ldots\) fulfill the conditions
(i)--(iii) of \defref{16.2.1} for
\((f;f\in\operatorname{Domain}X,\ Xf\in\mathfrak A)\). Therefore this set
is essentially dense, and the proof is completed.

Putting \(\mathfrak A=\Hspace\), we see that \(\operatorname{Domain}X\)
itself is essentially dense. This, of course, could have been easily proved
directly too.

The considerations of this \S\ are special cases of a generalisation of the
notion of relative dimensionality (which was only defined for linear closed
sets \(\mathbin\eta\M\) so far) to all linear sets. A theory very similar to
that of Lebesgue's interior measure results. We do not go into the details of
this subject here, but it will be dealt with in subsequent papers.

\SectionStart{16.3}{Calculus with unbounded operators. Definitions}
The \hyperref[lem:16.2.1]{Lemmas 16.2.1}--\hyperref[lem:16.2.3]{16.2.3}
have an interesting application
to unbounded operators.

\Definition{16.3.1}
Let \(x_1,y_1,x_2,y_2,\ldots\) be a (finite or infinite) sequence of symbolic
variables. By a (non-commutative) monomial we mean an expression of the form
\(z_1\cdots z_n\), where \(n=0,1,2,\ldots\), and each \(z_\nu\) is some
\(x_i\) or \(y_i\). If \(n=0\), we use the symbol \(1\). The order in which
the factors \(z_1,\ldots,z_n\) occur is essential; the number \(n\) is the
degree of the monomial. By a (non-commutative) polynomial we mean a finite
linear aggregate of non-commutative monomials, that is, an expression of the
form
\[
p(x_1,y_1,\ldots)=
\sum_{\rho=1}^{p}a_\rho
 z_1^{(\rho)}\cdots z_{n_\rho}^{(\rho)},
\]
where \(p=0,1,2,\ldots\); \(a_1,\ldots,a_p\) are complex numbers, and each
\(z_\nu^{(\rho)}\) is some \(x_i\) or \(y_i\). If \(p=0\), we use the symbol
\(0\). Monomial terms of the form \(0\cdot z_1\cdots z_n\) cannot be
omitted, but two terms \(a\cdot z_1\cdots z_n\) and
\(b\cdot z_1\cdots z_n\) can be contracted to one, and the order of the
additive terms does not matter. We consider monomials as special cases of
polynomials. If
\[
p(x_1,y_1,\ldots)=
\sum_{\rho=1}^{p}a_\rho
 z_1^{(\rho)}\cdots z_{n_\rho}^{(\rho)}
\]
is a polynomial, then \({}^{(r)}p(x_1,y_1,\ldots)\), the reduced polynomial
of \(p(x_1,y_1,\ldots)\), obtains from it by omitting all terms with
\(a_\rho=0\) (that is, of the form \(0\cdot z_1\cdots z_n\)).

If \(p(x_1,y_1,x_2,y_2,\ldots)\),
\(q(x_1,y_1,x_2,y_2,\ldots)\) are two polynomials, then
\[
ap(x_1,y_1,x_2,y_2,\ldots),\qquad
p(x_1,y_1,x_2,y_2,\ldots)+q(x_1,y_1,x_2,y_2,\ldots),
\]
\[
p(x_1,y_1,x_2,y_2,\ldots)\cdot q(x_1,y_1,x_2,y_2,\ldots)
\]
are those which result by carrying out the indicated operations term by term,
but without reducing (if the process \({}^{(r)}\) is not explicitly
indicated).

The symbol \({}^+\) is defined as follows, in successive steps:
\[
x_i^+=y_i,\qquad y_i^+=x_i;\qquad
(z_1\cdots z_n)^+=z_n^+\cdots z_1^+;
\]
\[
\left(\sum_{\rho=1}^{p}a_\rho
 z_1^{(\rho)}\cdots z_{n_\rho}^{(\rho)}\right)^+
=\sum_{\rho=1}^{p}\overline{a_\rho}
 \left(z_1^{(\rho)}\cdots z_{n_\rho}^{(\rho)}\right)^+.
\]
Obviously
\[
\bigl({}^{(r)}p(x_1,y_1,\ldots)\bigr)^+
={}^{(r)}\bigl(p(x_1,y_1,x_2,y_2,\ldots)^+\bigr).
\]

\Definition{16.3.2}
Let \(x_1,y_1,x_2,y_2,\ldots\) be a (finite or infinite) sequence of
symbolic variables; \(X_1,X_2,\ldots\) a corresponding sequence of linear,
closed operators, each one having an everywhere dense domain, so that to
every pair of variables \(x_i,y_i\) one operator \(X_i\) corresponds. If
\[
p(x_1,y_1,x_2,y_2,\ldots)
=\sum_{\rho=1}^{p}a_\rho
 z_1^{(\rho)}\cdots z_{n_\rho}^{(\rho)}
\]
is a polynomial, we mean by
\[
p(X_1,X_1^*,X_2,X_2^*,\ldots)
\]
the operator which arises if we replace in
\[
\sum_{\rho=1}^{p}a_\rho
z_1^{(\rho)}\cdots z_{n_\rho}^{(\rho)}
\]
every \(x_i\) by \(X_i\), and every \(y_i\) by \(X_i^*\). Here the domain
and the values of
\(p(X_1,X_1^*,\allowbreak X_2,X_2^*,\ldots)\) are to be determined with the help of the
following rules (cf. \cite[p.~404]{ref18}):

\(0f\) and \(1f\) are everywhere defined and have the values \(0\), resp.
\(f\). \((aX)f\) is defined if and only if \(Xf\) is defined (even for
\(a=0\)), and its value is \(a(Xf)\). \((X+Y)f\) is defined if and only if
both \(Xf\) and \(Yf\) are defined, and its value is \(Xf+Yf\). \((XY)f\)
is defined if and only if both \(Yf\) and \(X(Yf)\) are defined, and its
value is \(X(Yf)\).

Variables \(x_i,y_i\), resp. the corresponding operators \(X_i,X_i^*\),
which do not occur in the expression
\[
p(x_1,y_1,x_2,y_2,\ldots)
=\sum_{\rho=1}^{p}a_\rho
z_1^{(\rho)}\cdots z_{n_\rho}^{(\rho)}
\]
can be omitted in the expression \(p(x_1,y_1,x_2,y_2,\ldots)\), resp.
\(p(X_1,X_1^*,\ldots)\). (So we can write \(p(x_1)\) for
\(p(x_1,y_1)=x_1\), but not for \(p(x_1,y_1)=x_1+0\cdot y_1\).)
\(p(X_1,X_1^*,\ldots)\) and
\({}^{(r)}p(X_1,X_1^*,\ldots)\) are not identical in general. Thus if
\(p(x_1)=0\cdot x_1\), \({}^{(r)}p(x_1)=0\), then
\(p(X_1)=0\cdot X_1\) and \({}^{(r)}p(X_1)=0\), and these differ, because
they have different domains, if \(X_1\) is not everywhere defined. In this
particular case, however, we can obtain equality by forming the closure
\(\clos{\ }\) of all operators concerned (cf. \cite[pp.~70--71]{ref16},
where a \(\sim\) is used to denote closure): As \(X_1\) has an everywhere
dense domain, therefore \(\clos{0\cdot X_1}=0\), and so
\(\clos{p(X_1)}=\clos{{}^{(r)}p(X_1)}\). But in general this procedure too
breaks down: Define
\[
p(x_1,x_2)=0\cdot x_1+0\cdot x_2,
\qquad{}^{(r)}p(x_1,x_2)=0,
\]
and let \(X_1,X_2\) be two operators for which
\(\operatorname{Domain}X_1\cdot\operatorname{Domain}X_2=(0)\)
(cf. \cite[p.~230]{ref17}). Then \(p(X_1,X_2)\) has the domain \((0)\), and
there of course the value \(0\). Thus the same is true for
\(\clos{p(X_1,X_2)}\), while \({}^{(r)}p(X_1,X_2)=0\),
\(\clos{{}^{(r)}p(X_1,X_2)}=0\). Thus \(\clos{p(X_1,X_2)}\) has the domain
\((0)\), while \(\clos{{}^{(r)}p(X_1,X_2)}\) has the domain \(\Hspace\), and
so they are different.

Another ``pathological'' possibility is that \((X_1\mathbin{\cdot}X_2)^*\) may not be
\(X_2^*\mathbin{\cdot}X_1^*\). In fact it can happen that \(X_1X_2\) has no closure at all,
and therefore no adjoint with an everywhere dense domain (cf.
\cite[pp.~296 and 300]{ref21}; our statement means that \(X_1X_2\) is not a
\(^{**}\)-operator, cf. loc. cit.); or again \((X_1X_2)^*\) may exist and be
a proper extension of \(X_2^*\mathbin{\cdot}X_1^*\). We will not go presently into the
details of this matter.

For \(X_1,X_2,\ldots\mathbin\eta\M\), where \(\M\) is a factor in one of
the finite cases (\((\mathrm I_n)\), \(n=1,2,\ldots\), or
\((\mathrm{II}_1)\)), as we now always assume, nothing of this sort happens,
and the algebra of the
\(\clos{p(X_1,X_1^*,X_2,X_2^*,\ldots)}\) works perfectly smoothly. This is
the main result of the considerations which follow. For the cases
\((\mathrm I_n)\), \(n=1,2,\ldots\), these things are obvious (then even the
\(\clos{\ }\) is superfluous, because every linear operator is closed), and
the essential content of our results is that it is the same way in case
\((\mathrm{II}_1)\). So while the usually considered case
\((\mathrm I_\infty)\) (including \(\M=\B\) for a Hilbert space
\(\Hspace\)) is highly pathological in this respect,
\((\mathrm{II}_1)\) turns out to be the appropriate generalisation of the
\((\mathrm I_n)\).

\SectionStart{16.4}{Calculus with unbounded operators. Results. Theorem XV}
We prove first two Lemmas which have considerable interest of their own. In
particular the first one proves that the ``spectral problem'' (the existence
of a ``resolution of unity'' for every linear, closed Hermitian operator
\(X\mathbin\eta\M\)) has always a satisfactory solution. (Cf.
\cite[p.~92]{ref16}; observe the contrast with \((\mathrm I_\infty)\).)

\Lemma{16.4.1}
{\itshape Every linear, closed, Hermitian operator \(X\mathbin\eta\M\) is
maximal and self adjoint (cf. \cite[p.~88]{ref16}).}

\Proof Let \(U\) be the Cayley transform of \(X\), \(\mathfrak E,
\mathfrak F\) the connected linear, closed sets (cf. \cite[p.~80]{ref16}).
Then \(U,\mathfrak E,\mathfrak F,E=P_{\mathfrak E}\) are invariant under
every unitary transformation \(U'\in\M'\) along with \(X\), so they are all
\(\mathbin\eta\M\). Thus \(UE\mathbin\eta\M\), and it is obviously
partially isometric, with the initial and final sets \(\mathfrak E\), resp.
\(\mathfrak F\), so \(UE\) is bounded, and therefore \(UE\in\M\). Similarly
\(E\in\M\).

We know that \([\operatorname{Range}(U-1)]=\Hspace\), and as \(Uf\) is only
defined if \(f\in\mathfrak E\), we may as well write
\([\operatorname{Range}(UE-E)]=\Hspace\). Now
\(UE-E=(UE-E)E\), \((UE-E)^*=E(UE-E)^*\), and so
\([\operatorname{Range}(UE-E)^*]\subset[\operatorname{Range}E]
=\mathfrak E\). So by \lemref{6.2.1},
\[
\begin{aligned}
D_{\M}(\mathfrak E)
 &\geq D_{\M}([\operatorname{Range}(UE-E)^*])
  =D_{\M}([\operatorname{Range}(UE-E)])\\
 &=D_{\M}(\Hspace)=1.
\end{aligned}
\]
Considering the properties of the partially isometric operator \(UE\), we
have further \(D_{\M}(\mathfrak E)=D_{\M}(\mathfrak F)\). Therefore
\(D_{\M}(\Hspace-\mathfrak E)=D_{\M}(\Hspace-\mathfrak F)=0\), that is
\(0\), and so \(\Hspace-\mathfrak E=\Hspace-\mathfrak F=(0)\). So \(X\)
is maximal and self adjoint by \cite[p.~88]{ref16}.

\Lemma{16.4.2}
{\itshape Let \(X,Y\) be linear, closed operators, with everywhere dense
domains, \(X,Y\mathbin\eta\M\). If \(Y\) is an extension of \(X\) (that is:
whenever \(Xf\) is defined, \(Yf\) is defined too, and \(Xf=Yf\); cf.
\cite[p.~70]{ref15}), then \(X=Y\).}

\Proof Let \(W,B\) be the canonical decomposition of \(Y\): \(B\) self
adjoint and (semi-)definite, \(W\) partially isometric,
\(W\mathbin\eta\M\), \(B\mathbin\eta\M\), and \(Y=WB\) (cf.
\defref{4.4.1} and \lemref{4.4.1}). At the same time \(B=W^*Y\) (cf.
\cite[p.~306]{ref21}), so \(Y=W W^*Y\), and as \(Y\) is a continuation of
\(X\), so \(X=W W^*X\) too.

\(W^*Y=B\) is self adjoint, so it is Hermitian, and so \(W^*X\) is too;
besides it is linear, closed, and has an everywhere dense domain along with
\(X\). As \(W^*X\mathbin\eta\M\), therefore \lemref{16.4.1} applies:
\(W^*X\) is maximal. But \(W^*Y\) is a Hermitian continuation of \(W^*X\),
so \(W^*X=W^*Y\), \(WW^*X=WW^*Y\), \(X=Y\). Thus the proof is completed.

We pass now to the discussion of the (non-commutative) polynomials of
operators.

\Lemma{16.4.3}
{\itshape Let \(X_1,X_2,\ldots\) be a (finite or infinite) sequence of
linear, closed operators, each one having an everywhere dense domain. Assume
that all \(X_i\mathbin\eta\M\). Let
\(\mathfrak D\mathfrak P=\mathfrak D\mathfrak P(X_1,X_2,\ldots)\) be the
common part of the domains of the operators
\(p(X_1,X_1^*,X_2,X_2^*,\ldots)\), where
\(p(x_1,y_1,x_2,y_2,\ldots)\) runs over all (non-commutative) polynomials of
\(x_1,y_1,x_2,y_2,\ldots\). Then \(\mathfrak D\mathfrak P\) is essentially
dense and everywhere dense.}

\Proof It suffices, by \lemref{16.2.1}, to prove the essential density. It
suffices furthermore, obviously, to let \(p(x_1,y_1,x_2,y_2,\ldots)\) run
over all monomials only. These are countable in number, say
\(p^{(\sigma)}(x_1,y_1,\ldots)\), \(\sigma=1,2,\ldots\); therefore
\[
\mathfrak D\mathfrak P
=\prod_{\sigma=1}^{\infty}
 \operatorname{Domain}p^{(\sigma)}(x_1,y_1,\ldots).
\]
By \lemref{16.2.2}, we need only to prove that every
\(\operatorname{Domain}p(X_1,X_1^*,\ldots)\) is essentially dense. That is:
That for every monomial \(p(x_1,y_1,\ldots)=z_1\cdots z_n\) (cf.
\defref{16.3.1}), \(\operatorname{Domain}p(X_1,X_1^*,\ldots)\) is
essentially dense.

We prove this by induction on the degree \(n\). For \(n=0\), the operator in
question is \(1\), its domain is \(\Hspace\), and so our statement is
obvious. Assume now that it is already established for \(n-1\),
\(n=1,2,\ldots\), and let us consider \(n\). Let
\(p(x_1,y_1,\ldots)=z_1\cdots z_n\) be a monomial of degree \(n\). Put
\(q(x_1,y_1,\ldots)=z_1\cdots z_{n-1}\), and write
\(p(X_1,X_1^*,\ldots)=A\), \(q(X_1,X_1^*,\ldots)=B\). \(z_n\) is a
variable \(x_i\) or \(y_i\); put correspondingly \(Y=X_i\), resp.
\(Y=X_i^*\). Now \(A=BY\), so
\[
\operatorname{Domain}A=(f;f\in\operatorname{Domain}Y,
\ Yf\in\operatorname{Domain}B).
\]
\(\operatorname{Domain}B\) is essentially dense by our induction assumption
(the degree of \(q(x_1,y_1,\ldots)=z_1\cdots z_{n-1}\) is \(n-1\)); \(Y\)
is linear, closed, and has an everywhere dense domain (for \(X_i\) this was
assumed, for \(X_i^*\) it follows from \cite[p.~301]{ref21}). Therefore
\lemref{16.2.3} applies, and \(\operatorname{Domain}A\) is essentially dense
too. Thus the proof is completed.

\Lemma{16.4.4}
{\itshape Let \(X_1,X_2,\ldots\) be a (finite or infinite) sequence of
linear, closed operators, each one having an everywhere dense domain. Assume
that all \(X_i\mathbin\eta\M\). Let
\(p(x_1,y_1,x_2,y_2,\ldots)\), \(q(x_1,y_1,x_2,y_2,\ldots)\),
\(r(x_1,y_1,x_2,y_2,\ldots)\) be (non-commutative) polynomials of the
symbolic variables \(x_i,y_i\), corresponding to \(X_i\),
\(i=1,2,\ldots\). Then we have:
\begin{enumerate}[label=(\roman*)]
\item \(\clos{p(X_1,X_1^*,X_2,X_2^*,\ldots)}\) can be formed; it is linear,
closed, has an everywhere dense domain, and it is \(\mathbin\eta\M\) too.
\item If \({}^{(r)}p(x_1,y_1,x_2,y_2,\ldots)
={}^{(r)}q(x_1,y_1,x_2,y_2,\ldots)\), then
\[
\clos{p(X_1,X_1^*,X_2,X_2^*,\ldots)}
=\clos{q(X_1,X_1^*,X_2,X_2^*,\ldots)};
\]
that is, \(\clos{p(X_1,X_1^*,X_2,X_2^*,\ldots)}\) depends on
\({}^{(r)}p(x_1,y_1,x_2,y_2,\ldots)\) only.
\item If \(p(x_1,y_1,x_2,y_2,\ldots)^+=q(x_1,y_1,x_2,y_2,\ldots)\), then
\[
\clos{p(X_1,X_1^*,X_2,X_2^*,\ldots)}^*
=\clos{q(X_1,X_1^*,X_2,X_2^*,\ldots)}.
\]
\item If \(ap(x_1,y_1,x_2,y_2,\ldots)=q(x_1,y_1,x_2,y_2,\ldots)\), then
\[
\clos{a\clos{p(X_1,X_1^*,X_2,X_2^*,\ldots)}}
=\clos{q(X_1,X_1^*,X_2,X_2^*,\ldots)}.
\]
(If \(a\ne0\), then the first \(\clos{\ }\) is obviously unnecessary.)
\item If \(p(x_1,y_1,x_2,y_2,\ldots)+q(x_1,y_1,x_2,y_2,\ldots)
=r(x_1,y_1,x_2,y_2,\ldots)\), then
\[
\clos{\clos{p(X_1,X_1^*,X_2,X_2^*,\ldots)}
      +\clos{q(X_1,X_1^*,X_2,X_2^*,\ldots)}}
=\clos{r(X_1,X_1^*,X_2,X_2^*,\ldots)}.
\]
\item If \(p(x_1,y_1,x_2,y_2,\ldots)\cdot q(x_1,y_1,x_2,y_2,\ldots)
=r(x_1,y_1,x_2,y_2,\ldots)\), then
\[
\clos{\clos{p(X_1,X_1^*,X_2,X_2^*,\ldots)}
      \cdot\clos{q(X_1,X_1^*,X_2,X_2^*,\ldots)}}
=\clos{r(X_1,X_1^*,X_2,X_2^*,\ldots)}.
\]
\end{enumerate}}

\Proof Ad (i): \(p(X_1,X_1^*,\ldots)\) and
\(p(X_1,X_1^*,\ldots)^+\) have everywhere dense domains by
\lemref{16.4.3}. One verifies immediately that
\[
(p(X_1,X_1^*,\ldots)f,g)
=(f,p(X_1,X_1^*,\ldots)^+g)
\]
for all \(f\in\operatorname{Domain}p(X_1,X_1^*,\ldots)\),
\(g\in\operatorname{Domain}p(X_1,X_1^*,\ldots)^+\), and so
\(p(X_1,X_1^*,\ldots)^*\) is an extension of
\(p(X_1,X_1^*,\ldots)^+\). Therefore
\(p(X_1,X_1^*,\ldots)^*\) has an everywhere dense domain, and so we can form
\(\clos{p(X_1,X_1^*,\ldots)}\) (cf. \cite[p.~300]{ref21}).
\(\clos{p(X_1,X_1^*,\ldots)}\) is invariant under every unitary
transformation \(U'\in\M'\), along with \(X_1,X_2,\ldots\), so it is
\(\mathbin\eta\M\).

Ad (ii):
\[
{}^{(r)}p(X_1,X_1^*,X_2,X_2^*,\ldots)
\]
is obviously an extension of
\[
p(X_1,X_1^*,X_2,X_2^*,\ldots);
\]
therefore
\(\clos{{}^{(r)}p(X_1,X_1^*,\ldots)}\) is one of
\(\clos{p(X_1,X_1^*,\ldots)}\). By (i), both operators are linear, closed,
have everywhere dense domains, and are \(\mathbin\eta\M\). So by
\lemref{16.4.2},
\[
\clos{{}^{(r)}p(X_1,X_1^*,\ldots)}
=\clos{p(X_1,X_1^*,X_2,X_2^*,\ldots)}.
\]
Applying this to both \(p(x_1,y_1,x_2,y_2,\ldots)\) and
\(q(x_1,y_1,x_2,y_2,\ldots)\) gives the statement of (ii).

Ad (iii): As we saw in the proof of (i),
\(p(X_1,X_1^*,\ldots)^*\) is an extension of
\(p(X_1,X_1^*,\ldots)^+\), that is, of \(q(X_1,X_1^*,\ldots)\). The first
operator is obviously equal to \(\clos{p(X_1,X_1^*,\ldots)}^*\), and as it
is linear, closed, it must even be an extension of
\(\clos{q(X_1,X_1^*,\ldots)}\). As these are both linear, closed operators
with everywhere dense domains, \lemref{16.4.2} applies again: They must be
equal.

Ad (iv)--(vi): All these statements are proved in the same way; therefore it
suffices to consider one of them, say (vi):

Clearly
\[
\clos{p(X_1,X_1^*,X_2,X_2^*,\ldots)}\cdot
\clos{q(X_1,X_1^*,X_2,X_2^*,\ldots)}
\]
is an extension of
\[
p(X_1,X_1^*,X_2,X_2^*,\ldots)\cdot
q(X_1,X_1^*,X_2,X_2^*,\ldots)
=r(X_1,X_1^*,X_2,X_2^*,\ldots).
\]
The first mentioned operator is linear closed, and therefore it is an
extension of \(\clos{r(X_1,X_1^*,X_2,X_2^*,\ldots)}\), too. Now we can infer
\[
\clos{\clos{p(X_1,X_1^*,X_2,X_2^*,\ldots)}\cdot
      \clos{q(X_1,X_1^*,X_2,X_2^*,\ldots)}}
=\clos{r(X_1,X_1^*,X_2,X_2^*,\ldots)}
\]
from \lemref{16.4.2}, completing the proof.

Of course (iv) for \(a=0\) is trivial, even without the first
\(\clos{\ }\).

We restate the results of this \S.

\MainTheorem{XV}
{\itshape Let \(\M\) be a factor of a finite class
(\((\mathrm I_n)\), \(n=1,2,\ldots\), or \((\mathrm{II}_1)\)). Denote by
\(U(\M)\) the set of all linear, closed operators with an everywhere dense
domain, which are \(\mathbin\eta\M\). Then for \(X,Y\in U(\M)\) the
operators \(\clos{aX}\) (coinciding with \(aX\) for \(a\ne0\)),
\(\clos{X+Y}\), \(\clos{XY}\) can be formed, and are \(\in U(\M)\) again.
These operations satisfy all customary laws of matrix algebra:
\[
\clos{X+Y}=\clos{Y+X};\qquad
\clos{\clos{X+Y}+Z}=\clos{X+\clos{Y+Z}}.
\]
\[
\clos{a\clos{X+Y}}=\clos{a\clos X+b\clos Y};\qquad
\clos{(a+b)X}=\clos{aX+bX}.
\]
\[
\clos{\clos{XY}Z}=\clos{X\clos{YZ}};\qquad
\clos{\clos{aX}Y}=\clos{a\clos{XY}};\qquad
\clos{a\clos{bX}}=\clos{(ab)X}.
\]
\[
\clos{\clos{X+Y}Z}=\clos{\clos{XZ}+\clos{YZ}};\qquad
\clos{X\clos{Y+Z}}=\clos{\clos{XY}+\clos{XZ}}.
\]
\[
\clos{aX}^{*}=\clos{\overline aX^*};\qquad
\clos{X+Y}^{*}=\clos{X^*+Y^*};\qquad
\clos{XY}^{*}=\clos{Y^*X^*}.
\]

More generally: The (non-commutative) polynomials
\(\clos{p(X_1,X_1^*,X_2,X_2^*,\ldots)}\)
(cf. Definitions \hyperref[def:16.3.1]{16.3.1}, and
\hyperref[def:16.3.2]{16.3.2}) are \(\in U(\M)\) again;
\(\clos{p(X_1,X_1^*,X_2,X_2^*,\ldots)}\)
depends only on the reduced form
\({}^{(r)}p(x_1,y_1,x_2,y_2,\ldots)\) of
\(p(x_1,y_1,x_2,y_2,\ldots)\), and the laws of matrix algebra hold for the
\(\clos{p(X_1,X_1^*,X_2,X_2^*,\ldots)}\)
(cf. \lemref{16.4.4}(iii)--(vi)).

Every Hermitian \(X\in U(\M)\) is maximal and self adjoint, and thus it has
a unique resolution of unity (cf. \cite[p.~92]{ref15}).

For \(X,Y\in U(\M)\), whenever \(Y\) is an extension of \(X\) (cf.
\cite[p.~70]{ref15}), then \(X=Y\); that is: proper extensions do not exist
in \(U(\M)\).}

\medskip
\noindent(Added in proof, Jan. 29, 1936.) The authors have since succeeded
in establishing the following further facts concerning the finite cases:
\begin{enumerate}[label=(\roman*)]
\item \(T_{\M}(A+B)=T_{\M}(A)+T_{\M}(B)\) unrestrictedly.
  (Cf. \probref{11}.)
\item \(T_{\M}(A)\) is weakly continuous in \(A\).
  (Cf. Lemma 15.5.5.)
\item If \(C=1\), then an \(f\) exists, for which identically
  \(T_{\M}(A)=(Af,f)\).
\item In the same case certain isomorphisms between \(\M,\M'\), and
  \(\Hspace\) exist.
\item Several examples of Case II, thus \((\alpha)\)--\((\gamma)\) on page
  \hyperref[examples:13.2]{208}, are spacially
  isomorphic. (This answers the question at the end of \secref{13.4}.) The
  proofs will appear in a subsequent paper.
\end{enumerate}

\begin{center}
\textsc{Princeton University and The Institute for Advanced Study.}
\end{center}
% --- End of parts/06-part-v.tex ---

\end{document}
